\documentclass[11pt,a4paper]{article}
\usepackage[margin=25mm,headheight=14pt]{geometry}
\usepackage{fontspec}
\setmainfont{Latin Modern Roman}
\setsansfont{Latin Modern Sans}
\setmonofont{Latin Modern Mono}
\usepackage{xeCJK}
\setCJKmainfont[ItalicFont={Noto Serif CJK SC}]{Noto Serif CJK SC}
\setCJKsansfont{Noto Sans CJK SC}
\setCJKmonofont{Noto Sans Mono CJK SC}
\usepackage{amsmath,amssymb,amsthm,mathtools}
\usepackage{booktabs,longtable,array,enumitem,microtype}
\newcolumntype{P}[1]{>{\raggedright\arraybackslash}p{#1}}
\usepackage{xcolor,xurl}
\usepackage[colorlinks=true,linkcolor=black,citecolor=black,urlcolor=blue!45!black,bookmarksnumbered=true]{hyperref}
\usepackage{fancyhdr}
\pagestyle{fancy}\fancyhf{}\fancyhead[L]{\small Peilin Liu}\fancyhead[R]{\small Star forests}\fancyfoot[C]{\thepage}
\setlength{\parindent}{1.5em}\setlength{\parskip}{0.15em}
\setlength{\emergencystretch}{2em}
\allowdisplaybreaks[2]
\setlist{topsep=4pt,itemsep=2pt,parsep=0pt}
\setcounter{tocdepth}{1}
\makeatletter
\renewcommand{\@pnumwidth}{2.2em}
\renewcommand{\@tocrmarg}{3.2em}
\renewcommand*\l@section[2]{\vskip0.3em\@dottedtocline{1}{0em}{3.0em}{\bfseries #1}{\bfseries #2}}
\makeatother
\newcommand{\Sr}[1]{S_{#1}}
\newcommand{\rh}{\widehat r}
\newcommand{\nn}{\mathbf n}\newcommand{\mm}{\mathbf m}
\newcommand{\N}{\mathbb N}
\newcommand{\p}{\mathcal P}
\newcommand{\sq}{\bigsqcup}
\newcommand{\pos}[1]{\left(#1\right)_+}
\newcommand{\floor}[1]{\left\lfloor#1\right\rfloor}
\newcommand{\ceil}[1]{\left\lceil#1\right\rceil}
\DeclareMathOperator{\defi}{def}
\DeclareMathOperator{\dist}{dist}
\DeclareMathOperator{\supp}{supp}
\newtheorem{theorem}{Theorem}[section]
\newtheorem{proposition}[theorem]{Proposition}
\newtheorem{lemma}[theorem]{Lemma}
\newtheorem{corollary}[theorem]{Corollary}
\theoremstyle{definition}
\newtheorem{definition}[theorem]{Definition}
\newtheorem{remark}[theorem]{Remark}
\renewcommand{\proofname}{Proof}
\renewcommand{\contentsname}{Contents}
\title{Size Ramsey numbers of arbitrary star forests}
\hypersetup{pdftitle={Size Ramsey numbers of arbitrary star forests},pdfauthor={Peilin Liu},pdfsubject={Erdos problem 561; size Ramsey numbers of star forests}}
\author{Peilin Liu\\\small Shandong University}
\date{2 November 2026}
\begin{document}
\maketitle

\begin{abstract}
Let $S_h=K_{1,h}$, and let $\nn=(n_1,\ldots,n_s)$ and $\mm=(m_1,\ldots,m_t)$ be nonincreasing lists of positive integers. We establish the formula
\[
\rh\left(\sq_{i=1}^{s}S_{n_i},\sq_{j=1}^{t}S_{m_j}\right)
=\sum_{v=2}^{s+t}\max_{i+j=v}(n_i+m_j-1)
\]
for the size Ramsey number of two arbitrary star forests. The induction is formulated for the complete family of suffix pairs on a row. A single minimal-counterexample order allows vertex deletion, erasure of edges in actual matching regions, and changes of the target lists to be compared within one framework. Two extensions of a coloring give complete forcing relations on the same residual graph; an intersection operation on fixed lists retains intermediate stars and unit stars. Matching decompositions and factor deficiency lead to operations whose actual edge savings pay for the full original profile. The main text establishes the general mechanisms and the exhaustive final reduction. The appendices contain the required classifications, source transfers, and exact cost proofs.
\end{abstract}
\noindent\textbf{Keywords:} star forests; size Ramsey number; matchings; factors; Erdős problem 561.\par
\noindent\textbf{2020 Mathematics Subject Classification:} 05C55, 05C70, 05D10.
\tableofcontents
\newpage
\section{Introduction and the main theorem}
All graphs are finite and simple. We write $G\to(F,H)$ if every red--blue coloring of $E(G)$ contains a red copy of $F$ or a blue copy of $H$. Copies are ordinary, not necessarily induced, subgraphs. The stars in a copy of a star forest are pairwise vertex-disjoint. The size Ramsey number is
\[
\rh(F,H)=\min\{e(G):G\to(F,H)\}.
\]
Write $S_h=K_{1,h}$, where $h$ denotes the number of edges; in particular, $S_1=K_2$. For nonincreasing lists of positive integers
\[
n_1\ge\cdots\ge n_s\ge1,\qquad m_1\ge\cdots\ge m_t\ge1,
\]
define
\begin{equation}\label{eq:profile}
\ell_v=\max_{\substack{1\le i\le s,\ 1\le j\le t\\i+j=v}}
(n_i+m_j-1),\qquad 2\le v\le s+t.
\end{equation}
Each maximum is over a nonempty set, and the sequence $(\ell_v)$ is nonincreasing.

\begin{theorem}\label{thm:main}
For every pair of lists as above,
\begin{equation}\label{eq:main}
\rh\left(\sq_{i=1}^{s}S_{n_i},\sq_{j=1}^{t}S_{m_j}\right)
=\sum_{v=2}^{s+t}\ell_v.
\end{equation}
\end{theorem}

Formula \eqref{eq:main} is the star-forest conjecture of Burr, Erdős, Faudree, Rousseau, and Schelp, listed as problem 561 in the Erdős problems database \cite{erdos}. For the classification of size Ramsey minimal graphs for uniform star forests, see Fu--Luo--Ni \cite{fln}; for the case of two components on one side, see You \cite{you}. Neither the lengths of the lists, the multiplicities of their entries, the number of distinct star sizes, nor the lengths of their unit tails are bounded here.

The lower bound is the main issue. Deleting a vertex allows one component to be discarded from a target forest. Erasing a larger piece of a host, however, requires all losses at different star sizes to be extracted from a single complete copy. Separate estimates for the largest star and a matching do not control the entire target. We retain the forcing relation for every legal suffix pair and compare each actual operation with every height column of the original profile.

\paragraph{Organization of the proof.}
The main text develops complete suffix families, common minimality, actual matching regions, payable selectors, and full-profile descent, then exhausts the head-gap and small-degree terminals. The appendices prove the required preliminary exclusions, long classifications, and exact cost inequalities. Section \ref{sec:dependencies} provides a reading map. All problem-specific proof dependencies are included here; external references are to the standard literature listed below.

\subsection{The upper bound}
\begin{proof}[Upper bound in Theorem \ref{thm:main}]
Let $Q=\sq_{v=2}^{s+t}S_{\ell_v}$ and fix a coloring of $Q$. Start with $i=j=1$. While neither target is complete, process the component $S_{\ell_{i+j}}$. Since
\[
\ell_{i+j}\ge n_i+m_j-1,
\]
this component contains a red $S_{n_i}$ or a blue $S_{m_j}$. Select such a star and increase the corresponding index by one. Different steps use distinct components, so all selected stars of each color are vertex-disjoint. If neither side were complete after $s+t-1$ steps, then
\[
(i-1)+(j-1)=s+t-1,\qquad i-1\le s-1,\quad j-1\le t-1,
\]
a contradiction. Thus $Q$ has the required arrow relation and $e(Q)=\sum_v\ell_v$.
\end{proof}


\section{Complete suffix families and the exact profile}\label{sec:suffixes}
Put
\[
F_i=\sq_{h=i}^{s}S_{n_h},\qquad H_j=\sq_{h=j}^{t}S_{m_h}.
\]
For $2\le r\le s+t$, define
\begin{equation}\label{eq:Pr}
P_r(G;\nn,\mm)\quad\Longleftrightarrow\quad
G\to(F_i,H_j)\quad\text{for every legal }i+j=r.
\end{equation}
A pair is legal when $1\le i\le s$ and $1\le j\le t$. Write
\[
E_r=\sum_{v=r}^{s+t}\ell_v,\qquad N_r=s+t-r+1.
\]
We prove the stronger statement $P_r(G)\Rightarrow e(G)\ge E_r$. The original problem is the case $r=2$. This formulation permits a change of row after an operation while retaining every other legal pair on the new row.
A forcing datum is a tuple $(G;\nn,\mm;r)$ satisfying $P_r$; it is called strict if $e(G)<E_r$. Minimality below always ranges over all lists, rows, and host graphs.

\begin{lemma}[Vertex deletion]\label{lem:vertices}
If $P_r(G;\nn,\mm)$, $U\subseteq V(G)$, and $u=|U|<N_r$, then
\[
P_{r+u}(G-U;\nn,\mm).
\]
\end{lemma}
\begin{proof}
Fix a legal new pair $(i',j')$. Choose nonnegative integers $a,b$ such that
\[
a+b=u,\quad a\le i'-1,\quad b\le j'-1,\quad
i=i'-a,\quad j=j'-b.
\]
Then $i+j=r$ and $(i,j)$ is an original legal pair. Partition $U$ into $U_R,U_B$ of sizes $a,b$. Extend an arbitrary coloring of $G-U$ by coloring every edge meeting $U_R$ red, and every remaining edge meeting $U_B$ blue.
The original forcing relation gives a red $F_i$ or a blue $H_j$. At most $a$ components of the red copy use added edges: each meets $U_R$, and the components are vertex-disjoint. After discarding them, the remaining stars contain $F_{i+a}$. Indeed, deleting any $a$ entries from a nonincreasing list leaves its $k$th entry at least as large as the original $(k+a)$th entry. The blue case is identical. This proves every new pair.
\end{proof}

\begin{lemma}[Monochromatic restoration of a bounded-degree piece]\label{lem:erasure}
Suppose $G\to(F,H)$ and $J\subseteq G$ is an edge subgraph of maximum degree at most $c$. Define
\[
F^{[-c]}=\sq_{i:n_i>c}S_{n_i-c}.
\]
Then $G-E(J)\to(F^{[-c]},H)$, with the symmetric statement after interchanging the colors.
\end{lemma}
\begin{proof}
Restore every edge of $J$ in red. A blue $H$ supplied by the original relation survives in the residual graph. Otherwise each center of a single complete red $F$ loses at most $c$ edges. Each component with $n_i>c$ therefore retains $S_{n_i-c}$, and these retained stars are still vertex-disjoint. If no component remains, the resulting relation has an empty target and cannot be used as a positive-rank descent.
\end{proof}

\subsection{Height columns}
Define $f_h=\#\{i:n_i\ge h\}$, $g_h=\#\{j:m_j\ge h\}$, and
\[
L_h=\#\{v\in\{r,\ldots,s+t\}:\ell_v\ge h\}.
\]
Monotonicity of the lists gives
\begin{equation}\label{eq:columns}
L_h=\pos{\max_{\substack{a,b\ge1\\a+b-1\ge h}}(f_a+g_b)-r+1},
\qquad E_r=\sum_{h\ge1}L_h.
\end{equation}
Only candidates with $f_a,g_b>0$ are included; the column is zero if there is no candidate. For fixed thresholds $a,b$, indices $i\le f_a$, $j\le g_b$ with $i+j=v$ exist precisely for $2\le v\le f_a+g_b$. Taking the union over all candidates reaching height $h$ proves the first identity. The second counts the cells of the profile by columns.
If the lists are unchanged and the row is increased by $k$, then
\begin{equation}\label{eq:rowloss}
L_h^{\rm new}=(L_h-k)_+,\qquad
E_r-E_{r+k}=\sum_{h\ge1}\min(k,L_h).
\end{equation}
Unit stars are thus included in the same area at height one.

\subsection{Intersecting two fixed forcing families}
Cumulative notation is also used for other nonincreasing lists. We write $\mathbf u\preceq\mathbf v$ for entrywise domination, equivalently domination of every cumulative count.
\begin{lemma}[Intersection of fixed lists]\label{lem:intersection}
Suppose the same graph $Y$ satisfies
\[
P_r(Y;\mathbf R,\mathbf S),\qquad
P_r(Y;\nn,\mathbf T),\qquad \mathbf R\preceq\nn.
\]
Let $1\le k\le r-1$. Define fixed lists $\mathbf K,\mathbf L$ by their cumulative counts:
\begin{equation}\label{eq:intersection}
K_h=\min(f_h,R_h+k),\qquad
L_h=\min(S_h,\max(r-1,T_h+k)).
\end{equation}
Provided both lists are nonempty and the new row has positive rank,
$P_{r+k}(Y;\mathbf K,\mathbf L)$ holds.
\end{lemma}
\begin{proof}
The expressions in \eqref{eq:intersection} are nonincreasing, eventually zero integer sequences, so they define lists. Fix a new legal pair $i+j=r+k$.
If $i>k$, use pair $(i-k,j)$ in the first forcing family. For every $\ell\ge0$,
\[
K_{i+\ell}^{\rm entry}\le R_{i-k+\ell}^{\rm entry},
\qquad L_{j+\ell}^{\rm entry}\le S_{j+\ell}^{\rm entry}.
\]
The superscript distinguishes entries from cumulative counts. These inequalities follow from cumulative domination. The complete monochromatic suffix supplied by the first family contains the required new suffix.
If $i\le k$, then $j\ge r$. Use pair $(i,r-i)$ in the second family. The red side follows from $\mathbf K\preceq\nn$. On the blue side $j+\ell\ge r$, so the cutoff $r-1$ is irrelevant: $L_h\ge j+\ell$ implies $T_h+k\ge j+\ell$. Consequently
\[
L_{j+\ell}^{\rm entry}\le T_{j-k+\ell}^{\rm entry}
=T_{r-i+\ell}^{\rm entry}.
\]
The old pairs are legal. In the first case this follows from the new list lengths being at most $|\mathbf R|+k$ and $|\mathbf S|$; in the second it follows from $|\mathbf K|\le|\nn|$ and the displayed blue domination. The two cases exhaust the new row.
\end{proof}

\begin{remark}
Lemma \ref{lem:intersection} uses two universal forcing relations on one residual graph. It never joins copies obtained in different restorations. The new lists are fixed before any coloring is chosen. Positivity of the new rank and the full edge-cost comparison remain separate obligations.
\end{remark}

\subsection{A sufficient condition for actual descent}
\begin{lemma}\label{lem:descent}
Among all strict sources $P_r(G)$ with $e(G)<E_r$, choose one minimizing
\begin{equation}\label{eq:order}
(N_r,e(G),|V(G)|)
\end{equation}
lexicographically. Suppose an actual subgraph $Y$ has a fixed complete forcing family $P_{r'}(Y;\nn',\mm')$ and
\[
1\le N'=|\nn'|+|\mm'|-r'+1<N_r,\qquad
e(G)-e(Y)\ge E_r-E_{r'}(\nn',\mm').
\]
Then a contradiction follows.
\end{lemma}
\begin{proof}
The edge-cost inequality and $e(G)<E_r$ imply
$e(Y)<E_{r'}(\nn',\mm')$. The new rank is positive and smaller, contrary to \eqref{eq:order}.
\end{proof}
The new graph need not inherit the original matching decomposition or any minimal-counterexample structure. Its forcing relation, rank, and edge count are established on that graph itself.


\section{Matching potential and the minimal source}\label{sec:potential}
Let $\nu(X)$ denote the matching number. A basic quantity is
\[
\Phi(X)=e(X)-\nu(X).
\]
Deleting an edge or a vertex cannot increase $\Phi$. This potential connects the degree requirement in a single-star arrow with the number of components in a whole forest.

\subsection{A covering matching permits a degree split}
\begin{lemma}\label{lem:split}
Suppose $\Delta(X)\le p+q$, where $p,q\ge1$. If a matching covers every vertex of degree $p+q$, then $E(X)$ can be partitioned into two subgraphs of maximum degrees at most $p$ and $q$.
\end{lemma}
\begin{proof}
First suppose both bounds are even. A graph of maximum degree at most $2a+2b$ can be augmented by auxiliary edges so that every degree is even and at most $2a+2b$. Orient Euler circuits. Every indegree and outdegree is at most $a+b$. Splitting each vertex into a left and a right copy gives a bipartite multigraph, which can be completed to an $(a+b)$-regular bipartite multigraph and decomposed into $a+b$ matchings by Hall's theorem. Each matching corresponds to a subgraph of maximum degree two in the original graph. Group $a$ classes on one side and $b$ on the other.

If exactly one bound is odd, delete the given matching. Each original maximum-degree vertex loses an edge, so the maximum degree falls by at least one. Apply the even case and put the matching back on the originally odd side.

Finally suppose $D=p+q$ is even and $p,q$ are odd. A bound equal to one is handled by the matching and its complement. Otherwise apply the ordinary $(g,f)$-factor theorem \cite{lovasz}. Set
\[
f_x(v)=x,\qquad g_x(v)=\max(0,d_X(v)-D+x).
\]
For disjoint sets $S,T$, the left side of the factor criterion is
\[
A_x(S,T)=x|S|+\sum_{v\in T}\min(d_X(v),D-x)-e_X(S,T).
\]
For odd $x$, the components counted on the right do not depend on $x$: each consists entirely of original degree-$D$ vertices and has $|C|+e(C,T)$ odd. Denote their number by $q(S,T)$. The matching and its complement establish feasibility at $x=1,D-1$, so $A_1,A_{D-1}\ge q(S,T)$. For fixed $S,T$, the function $A_x$ is concave. Hence $A_p\ge q(S,T)$. The factor theorem gives a subgraph with degrees between $g_p$ and $f_p$; its complement has maximum degree at most $q$.
\end{proof}

\begin{lemma}[Local potential for a star pair]\label{lem:local}
If $X\to(S_a,S_b)$, some vertex $v$ satisfies
\[
\Phi(X)-\Phi(X-v)\ge a+b-2.
\]
\end{lemma}
\begin{proof}
Deleting a vertex lowers the matching number by zero or one. If $\Delta(X)\ge a+b-1$, take a vertex of such degree. Otherwise the maximum degree cannot be smaller than $a+b-2$: Vizing's theorem would give at most $a+b-2$ matching classes, which could be grouped into $a-1$ and $b-1$ classes to avoid both stars.
Put $D=a+b-2$. If $a=1$ or $b=1$, the all-blue or all-red coloring forces the larger degree already treated. In the remaining case, Lemma \ref{lem:split} implies that no maximum matching covers all degree-$D$ vertices. Choose an uncovered degree-$D$ vertex $v$ of a maximum matching. Then $\nu(X-v)=\nu(X)$ and the potential decreases by $D$.
\end{proof}

\begin{proposition}\label{prop:potential}
If $P_r(G;\nn,\mm)$, then
\begin{equation}\label{eq:potential}
e(G)-\nu(G)\ge E_r-N_r.
\end{equation}
\end{proposition}
\begin{proof}
Choose a legal pair attaining $\ell_r$. Its complete suffix relation implies the corresponding single-star relation. Lemma \ref{lem:local} supplies a vertex whose deletion lowers the potential by at least $\ell_r-1$, while Lemma \ref{lem:vertices} preserves the next complete row. Repeat $N_r-1$ times. On the last row, the same local lemma gives a decrease of at least $\ell_{s+t}-1$. The remaining potential is nonnegative. Summing gives
$\sum_{v=r}^{s+t}(\ell_v-1)=E_r-N_r$.
\end{proof}

\subsection{Normalization and the two foundational exclusions}
Fix a minimal strict source in the order \eqref{eq:order}, and write $N=N_r$, $E=E_r$. The following foundational results use an equality classification for the potential and a classification of actual matching regions. Their precise technical dependencies are recorded in Section \ref{sec:dependencies}.

\begin{proposition}[Normalization]\label{prop:normal}
The minimal strict source satisfies
\begin{equation}\label{eq:normal}
\begin{gathered}
e(G)=E-1,\qquad \Delta(G)=D=\ell_r-1\ge3,\\
d:=N-\nu(G)\ge2,\qquad r\le\min(s,t).
\end{gathered}
\end{equation}
It has no isolated vertices, and both heads in every pair attaining $\ell_r$ are at least two.
\end{proposition}
\begin{proof}
An isolated vertex contradicts minimality. Rank one is the single-star edge bound, so $N\ge2$. Write $\tau=E-e(G)\ge1$. Lemma \ref{lem:vertices} and minimality at smaller positive rank give, for every $v$,
\[
e(G)-d_G(v)\ge E_{r+1},\qquad d_G(v)\le\ell_r-\tau.
\]
The proof of Lemma \ref{lem:local} also gives $\Delta(G)\ge\ell_r-1$. Thus $\tau=1$ and $\Delta(G)=\ell_r-1$. Inequality \eqref{eq:potential} gives $d\ge1$.

Appendix \ref{app:binary} proves $e(G)\ge E_r$ for complete suffix families on graphs of maximum degree at most two. Thus $D\ge3$. We can now use the universal reserve \eqref{eq:reserve} from the next section. It is a structural theorem for arbitrary graphs and does not depend on this normalization. Column domination and the highest-column exposure bound give
\[
e(G)\ge E_r+\rho_D(G)-d,\qquad \rho_D(G)\ge L_{D+1}\ge1.
\]
Since $e(G)=E_r-1$, it follows that $d\ge\rho_D(G)+1\ge2$.

If $r>s$, the legal pair $(s,r-s)$ in the all-blue coloring forces a blue forest with $N$ components, giving $\nu(G)\ge N$, a contradiction. The other side is symmetric. A head equal to one in an attaining pair similarly forces maximum degree at least $\ell_r$ by a monochromatic coloring.
\end{proof}

For the second foundational exclusion, use the ordinary maximum-degree matching decomposition $A_0,\mathcal R,B_0$ described in the next section, and put $\mu=\nu(\mathcal R)$.
\begin{proposition}[Foundational region exclusion]\label{prop:base}
Every minimal strict source must satisfy
\begin{equation}\label{eq:base}
D\ge6,\qquad \mu<\min(s,t),\qquad H:=\max(n_1,m_1)\le D-1.
\end{equation}
\end{proposition}
\begin{proof}
The emptiness interface identifies an empty residual list with an actual complete forest in the same intrinsic region. Appendices \ref{app:sources}--\ref{app:base-gluing} give the full proofs retaining arbitrary unit tails, roots, and exteriors: first exclude one empty side and original heads equal to the degree cap, then use deletion-resistant embeddings and dominating-set colorings to exclude the double-empty cases $D=3,4,5$. Appendix \ref{app:base-gluing} assembles these complementary cases under the same order \eqref{eq:order}, proving \eqref{eq:base}.
\end{proof}


\section{Actual matching regions and complete length constraints}\label{sec:regions}
\subsection{The decomposition and its weight identity}
For $\Delta(X)\le D$, let $\rho_D(X)$ be the minimum number of degree-$D$ vertices left uncovered by a maximum matching. Define also
\[
\begin{aligned}
\kappa_D(X)=\min\{|U|:\ &E(X-U)\text{ partitions into a matching}\\
&\text{and a graph of maximum degree at most }D-1\}.
\end{aligned}
\]
Always $\kappa_D(X)\le\rho_D(X)$: delete the degree-$D$ vertices missed by the chosen maximum matching. That matching covers every remaining maximum-degree vertex. Equality is neither needed nor assumed.

The ordinary maximum-degree matching decomposition of Appendix \ref{app:matching} gives an actual partition $A_0,\mathcal R,B_0$. It is obtained from the Gallai--Edmonds structure and a matching saturating the roots. The components of $\mathcal R$ are odd factor-critical graphs, every vertex of $\mathcal R$ has degree $D$ in the original graph, and there are no edges from $\mathcal R$ to $B_0$. Put
\[
a=|A_0|,\quad \mu=\nu(\mathcal R),\quad \rho=\rho_D(X).
\]
The selected decomposition satisfies
\begin{equation}\label{eq:region}
|\mathcal R|=2\mu+a+\rho,\qquad
\nu(X)=a+\mu+\nu(B_0).
\end{equation}
A root assignment matches $A_0$ to distinct components of $\mathcal R$. Suitable near-perfect matchings in the components realize \eqref{eq:region}.

For a whole component $Q$ of order $2z_Q+1$, put
\[
w_Q=e(Q)-Dz_Q.
\]
Every vertex of $Q$ has original degree $D$, so
\begin{equation}\label{eq:coreweight}
2w_Q=D-\delta_Q,\qquad
\delta_Q=e_X(V(Q),V(X)\setminus V(Q)).
\end{equation}
Write
\[
\zeta=Da-\sum_{v\in A_0}d_X(v),\quad
\varepsilon=e(X[A_0]),\quad \beta=e_X(A_0,B_0).
\]
Sum degrees first over $\mathcal R$ and then over the roots to eliminate $e(A_0,\mathcal R)$. This gives
\begin{equation}\label{eq:handshake}
2\sum_Qw_Q=D\rho+\zeta+2\varepsilon+\beta.
\end{equation}
Every term belongs to the same decomposition. Whole components are vertex-disjoint, so their edges and matching capacities add. Graph-theoretic blocks inside a component may share cutvertices and require a separate block interface.

\subsection{A reserve on the same graph}
For $h\ge1$, define the uniform star-forest envelope
\[
u_h(X)=\max\{p+q-1:X\to(pS_\alpha,qS_\beta),\
\alpha+\beta-1=h,\ p,q\ge1\},
\]
with value zero if the set is empty. Height one is handled separately by the actual matching number. Appendix \ref{app:reserve} proves that, for $D\ge3$ and $\Delta(X)\le D$,
\begin{equation}\label{eq:reserve}
e(X)-\nu(X)-\sum_{h\ge2}u_h(X)\ge\rho_D(X).
\end{equation}
The assembly of this reserve is as follows. Put
\[
b(X)=e(X)-\nu(X)-\sum_{h\ge2}u_h(X),\qquad
J(X)=e(X)-\nu(X)-\sum_{j\ge1}\kappa_j(X).
\]
Appendix \ref{app:reserve} proves $b\ge J\ge0$ and superadditivity of $b$ on disjoint unions.
Separate all connected $D$-regular factor-critical components $Q_1,\ldots,Q_c$, and denote the remainder by $B'$.
Lemma \ref{lem:regular-reserve} gives $b(Q_i)\ge1$, while Lemma \ref{lem:ported-reserve} gives
$b(B')\ge J(B')\ge\rho_D(B')$. Each $Q_i$ contributes exactly one unit of full-degree exposure.
Maximum matchings can be chosen independently in connected components, so
\[
 \rho_D(X)=c+\rho_D(B'),\qquad
 b(X)\ge\sum_i b(Q_i)+b(B')\ge\rho_D(X).
\]
This proves \eqref{eq:reserve} without identifying the height-one envelope with the matching number.

If $P_R(X)$ has rank $N_X\le\nu(X)$, each original height column is dominated by $u_h(X)$, and \eqref{eq:reserve} gives
\begin{equation}\label{eq:rich}
e(X)\ge E_R+\rho_D(X).
\end{equation}
For $N_X>\nu(X)$, the actual selector in the next section gives a different price. The two prices cover both signs of the residual matching deficit.

\subsection{Strong lengths forced by complete sources}
\begin{lemma}\label{lem:length}
Suppose $X=X_1\sqcup X_2$, $P_\ell(X;\nn,\mm)$, and
\[
N_\ell=s+t-\ell+1,\quad \nu(X)=N_\ell-d,\quad d\ge1,
\quad \nu(X_1)=\mu<\min(s,t).
\]
Then
\begin{equation}\label{eq:stronglength-general}
s,t\ge\mu+\ell+d-1.
\end{equation}
\end{lemma}
\begin{proof}
Put $x=s-\mu\ge1$, $i=\min(\ell-1,x)$, and $j=\ell-i$. This pair is legal. If $i=\ell-1$, then $j=1$; otherwise $i=x$ and $N_\ell=\mu+\nu(X_2)+d\ge\mu+1$ gives $j=\ell-s+\mu\le t$. Color $X_1$ red and $X_2$ blue. The red target has $s-i+1\ge\mu+1$ components and is absent. The excess of the blue target's component count over the blue matching number is
\[
(t-j+1)-\nu(X_2)=i-x+d.
\]
If $x<\ell+d-1$, this is $d>0$ when $x\le\ell-1$, and $\ell-1-x+d>0$ otherwise. The blue target is then absent as well, a contradiction. Thus $s\ge\mu+\ell+d-1$; interchanging colors proves the other inequality.
\end{proof}

Apply this lemma to $G-A_0=\mathcal R\sqcup B_0$. Vertex deletion gives $P_{r+a}$, and \eqref{eq:region} preserves exactly the original deficit $d$. Proposition \ref{prop:base} supplies $\mu<\min(s,t)$. Consequently
\begin{equation}\label{eq:stronglength}
\min(s,t)\ge\mu+a+r+d-1,
\qquad \nu(B_0)\ge\max(s,t).
\end{equation}
The new row $r+a$ is allowed to exceed one list length: legality of the relevant pairs was proved directly.

\subsection{Costs after removing all roots}
In the original ordinary region, let $C$ be the union of all positive-weight whole cores, and write
\[
Z=\nu(C),\quad |C|=n=2Z+c,\quad e(C)=DZ+W.
\]
Equation \eqref{eq:coreweight} implies that every positive core has at least $D+1$ vertices and weight at most $D/2$. Hence
\begin{equation}\label{eq:corecaps}
0<W\le Dc/2\le Z,\qquad 2Z\ge Dc.
\end{equation}
Let $\lambda=L_{D+1}$. The uniform source and $\kappa_D\le\rho_D$ give $\lambda\le\rho$. From \eqref{eq:handshake},
\begin{equation}\label{eq:rootdebt}
\zeta+\varepsilon+\lambda\le2W-(D-1)\rho.
\end{equation}
Delete $A_0$ and then erase the internal edges of $C$, obtaining
\[
Y=(G-A_0)-E(C),\quad r_0=r+a,\quad R_0=r_0-1.
\]
All original ports are included in the root-deletion cost $Da-\zeta-\varepsilon$. The erased internal edges cost $DZ+W$ and are disjoint from the root edges. The highest column contributes at most the additional $\lambda$, giving
\begin{equation}\label{eq:Ycost}
e(Y)\le E_{r_0}-DZ-W-1+\zeta+\varepsilon+\lambda.
\end{equation}
Also
\[
\nu(Y)=\nu(G)-a-Z,\qquad L_h^0=(L_h-a)_+.
\]
These relations concern this particular residual graph. They do not assert that $Y$ automatically satisfies $P_{r_0+Z}$.


\section{Whole-core selection and residual prices}\label{sec:selector}
The next argument explains why regions with ports still permit actual descent. Only internal core edges are erased; root edges and the exterior remain. The saving is an edge count in the host, not merely a formal capacity.

\begin{lemma}[The capacity menu]\label{lem:menu}
Suppose $P_R(X;\nn,\mm)$. Select pieces $Q_1,\ldots,Q_m$ with disjoint internal edge sets such that, for every selected subfamily, the matching number of its union is at most the sum of the associated integers $z_i$. Put
\[
z=\sum_i z_i,\qquad Z=\max_i z_i.
\]
Delete $b$ vertices avoiding the endpoints of all selected internal edges, and erase those internal edges, obtaining $Y$. If $R+b>Z$ and $z+b<N_R$, then
\[
P_{R+b+z}(Y;\nn,\mm).
\]
\end{lemma}
\begin{proof}
First delete the $b$ vertices using Lemma \ref{lem:vertices}. The subset sums of the $z_i$ contain $0,z$, and the gap between consecutive subset sums is at most $Z$. This follows by induction as each summand $z_i\le Z$ is inserted.
Fix a new pair $i'+j'=R+b+z$. We need a subset sum $x$ satisfying
\begin{equation}\label{eq:menuinterval}
\max(0,z-j'+1)\le x\le\min(z,i'-1).
\end{equation}
If the interval reaches $0$ or $z$, that endpoint works. Otherwise it has length $R+b-2\ge Z-1$ and contains at least $Z$ consecutive integers, so cannot lie entirely between consecutive subset sums.
Restore the corresponding pieces in red and the other selected pieces in blue. Apply the old complete relation at
$(i'-x,j'-(z-x))$. From each red component using a restored red edge, choose one such edge. They form a matching, so at most $x$ components are affected. Discarding and trimming them leaves the complete $F_{i'}$ in $Y$. The blue case is identical. Small stars crossing pieces, unit stars, and root attachments are all handled within this single complete copy.
\end{proof}

\begin{proposition}[A universal whole-core price]\label{prop:selector}
Relative to the original minimum rank $N_0$, suppose an actual graph $X$ satisfies
\[
P_R(X;\nn,\mm),\quad 1\le N_X<N_0,\quad
\Delta(X)\le D,\quad D\ge3,
\]
and its highest profile value is at most $D+1$. Let $h$ be its height-$(D+1)$ column width, $\rho=\rho_D(X)$, and $p=\floor{D/2}$. Then
\begin{equation}\label{eq:selector}
e(X)\ge E_R+\max\left(0,\ceil{D\rho/2}-p(p+1)-h\right).
\end{equation}
\end{proposition}
\begin{proof}
Use the ordinary maximum-degree decomposition of $X$ itself. Equation \eqref{eq:handshake} gives total positive whole-core weight at least $\ceil{D\rho/2}$, with each weight at most $p$. Put
\[
T=\ceil{D\rho/2}-p(p+1).
\]
If $T\le h$, minimality already gives $e(X)\ge E_R$. Assume $T>h$.

Order positive cores by capacity, leave the largest $p+1$ unselected, and select all the others. Their total weight $w$ is at least $T>0$. Let their total capacity be $z$ and their largest capacity $Z$. Each unselected core has capacity at least $Z$, so the unselected cores contain at least
\[
(p+1)(2Z+1)>(D+1)Z
\]
actual degree-$D$ vertices, all outside the selected cores. A greedy closed-neighborhood argument supplies at least $Z+1$ independent vertices. Choose
$b=(Z-R+1)_+\le Z-1$ of them.

Let this independent set be $I$, and put
$Y=(X-I)-E(\sq Q_i)$. The actual saving is exactly
\[
e(X)-e(Y)=Db+\sum_i(Dz_i+w_i)=D(z+b)+w.
\]
Root edges may be among the edges deleted with $I$, but none is an erased internal edge of a selected core. Thus no edge is counted twice.
Put $K=z+b\ge1$. If $K<N_X$, the new rank is positive and smaller,
and $R+b>Z$ permits Lemma \ref{lem:menu} at every new pair.
The full profile satisfies $E_R-E_{R+K}\le DK+h$.
Thus $e(X)<E_R+T-h$ would give $e(Y)<E_{R+K}$, contradicting the same
original choice of $N_0$.
If $K\ge N_X$, no zero-rank source is invoked: the actual edge saving gives
\[
e(X)\ge DK+w\ge DN_X+T\ge E_R+T-h.
\]
Take the maximum with the smaller-rank baseline $e(X)\ge E_R$.
This proves the price for either sign of $N_X-\nu(X)$.
At $N_X=N_0$, the same construction still proves the nonzero term when $T>h$,
but the zero baseline is not available without proof.
\end{proof}

For $D=5$, \eqref{eq:selector} gives
\[
e(X)-E_R\ge\max(0,\ceil{5\rho/2}-6-h).
\]
Combining $h\le\rho$ with \eqref{eq:rich} yields, for either sign of the residual matching deficit,
\begin{equation}\label{eq:five-safe}
e(X)\ge E_R+\max(0,\rho_5(X)-4).
\end{equation}
Use \eqref{eq:rich} when the matching number is at least the rank. Otherwise nonnegativity handles $\rho\le4$, while
$\ceil{3\rho/2}-6\ge\rho-4$ handles $\rho\ge4$. This avoids transferring the original deficit to a residual graph without proof.

\subsection{A complete forest versus a matching}
\begin{lemma}\label{lem:endpoint}
Let $F=\sq_{i=1}^{s}S_{n_i}$ with $n_1\ge\cdots\ge n_s\ge1$, and let $q\ge1$. Then
\begin{equation}\label{eq:endpoint}
\rh(F,qK_2)=(q-1)n_1+\sum_i n_i.
\end{equation}
\end{lemma}
\begin{proof}
The explicit star-host construction already proved for the upper bound in Theorem \ref{thm:main}, with the second list consisting of ones, gives the upper bound without using the unproved lower bound.
For the lower bound, induct on $q$. When $q=1$, the all-red coloring forces $F$, so $e(X)\ge e(F)$. For $q\ge2$, the all-red coloring forces $\Delta(X)\ge n_1$. Choose $v$ of degree at least $n_1$. Restoring all edges at $v$ in blue affects at most one edge of a blue matching, so
$X-v\to(F,(q-1)K_2)$. Induction gives
$e(X-v)\ge(q-2)n_1+e(F)$; add $d_X(v)\ge n_1$.
\end{proof}
This price applies to residual graphs retaining one complete target forest. It is one lower bound on that graph. Other prices on the same graph are combined by taking a maximum unless an actual decomposition into disjoint edge sets has been established.


\section{A payable root package from a factor obstruction}\label{sec:fractional}
Choose actual heads $p+1,q+1$ attaining the original highest column, where $p+q=D$ and the associated uniform source has component rank $\lambda=L_{D+1}$. By \eqref{eq:base}, $p,q\ge2$. If both bounds are even, the even part of Lemma \ref{lem:split} gives an avoiding coloring for the complete original pair. Otherwise we use factor deficiency rather than assume a port-free core.

\subsection{The fractional root-assignment interface}
We first state the structural interface, then prove it from a real flow and the factor criterion. Write $\delta_x=D-d_G(x)$. For even $D$ and odd $p<q$, assign root capacity
\[
\alpha_x=1+\frac{(p-1)\delta_x}{q-p}.
\]
For odd $D$, interchange the bounds so that $p$ is odd and $q$ even, and assign
\[
\alpha_x=1+\frac{(p-1)\delta_x}{q}.
\]
Each original whole core has capacity one. The integral root assignment ensures that every root can receive at least one unit of flow. Put
\[
h_f=\max_{\mathcal J}\left(|\mathcal J|-
\sum_{x\in N(\mathcal J)}\alpha_x\right)_+.
\]
The ordinary factor-deficiency formula, together with an integral decomposition of the maximum flow, gives either a complete original avoiding coloring or a residual reachable core family $\mathcal J$ and all its neighboring roots $A_{\mathcal J}$ satisfying the following failure inequalities. Here $c_*$ counts actual isolated odd-order $D$-regular components of the original graph:
\begin{align}
(p-1)\zeta_{\mathcal J}
&\le(q-p)\rho_{\mathcal J}+2(p-1)c_*-(D-2)\lambda,
&& D\text{ even},\ p<q,\label{eq:failure-even}\\
(p-1)\zeta_{\mathcal J}
&\le q\rho_{\mathcal J}-(D-1)\lambda,
&& D\text{ odd}.\label{eq:failure-odd}
\end{align}
Here $\rho_{\mathcal J}=|\mathcal J|-|A_{\mathcal J}|$ and
$\zeta_{\mathcal J}=\sum_{A_{\mathcal J}}\delta_x$. The next subsection derives these inequalities from an actual matching edge subgraph and the boundary count in the factor-deficiency theorem of Lovász \cite{lovasz} and Lu--Yu \cite{luyu}.

The interface retains three structural properties. The family $\mathcal J$ has no ports to roots outside $A_{\mathcal J}$; the roots can be matched simultaneously into distinct cores of $\mathcal J$; and
\begin{equation}\label{eq:packmatching}
\nu(G)=|A_{\mathcal J}|+\sum_{Q\in\mathcal J}\nu(Q)
+\nu\bigl(G-(A_{\mathcal J}\cup V(\mathcal J))\bigr).
\end{equation}
When the deficiency resource is less than $\lambda$, exceptional centers can be assigned to the two colors below their required largest-star counts. This avoids a complete original pair, not merely reduced targets.

\subsection{From a real flow to one coloring of the whole graph}
We prove the interface in each unbalanced parity case. For even $D$ put
\[
\theta=\frac{q-p}{D-2},\qquad t=\frac{p-1}{D-2},\qquad \theta+2t=1;
\]
for odd $D$ put $\theta=q/(D-1)$. In the source--core--root--sink network, cores have capacity one and roots capacity $\alpha_x$. Augment an integral assignment covering every root to a maximum real flow, maintaining $1\le m_x\le\alpha_x$ at every root: a simple augmenting path does not traverse a root--sink edge in reverse.

Now retain only integral constraints: capacity one at cores, lower bound one at roots, and the actual neighbor count as each root's upper bound. Integrality of the bipartite flow polytope expresses the real flow as a convex combination of actual integral assignments. For each assignment $\mathcal M$, a core $Q$ assigned to a root $x$ supplies an actual port edge $xw$ and a perfect matching of $Q-w$. An unassigned core exposes one vertex and matches the others. Include the fixed matching of $B_0$ covering its full-degree vertices. This gives an actual edge subgraph $L_{\mathcal M}$ of the same $G$. Low-degree roots may have several incident edges, but every original full-degree vertex has degree one in this subgraph except the exposed set $U_{\mathcal M}$ from unassigned cores. Moreover,
\[
\mathbb E|U_{\mathcal M}|=h_f,\qquad
\mathbb E\,d_{L_{\mathcal M}}(x)=m_x\le\alpha_x.
\]

Define vertex bounds
\[
f(x)=\min(p,d_G(x)),\qquad
g(x)=\max(0,d_G(x)-q).
\]
The ordinary minimum-deficiency theorem for $(g,f)$-factors \cite[Theorem 2.7 and Corollary 2.9]{luyu} states that the minimum of $\sum_x(g(x)-d_F(x))_+$ over spanning subgraphs with $d_F\le f$ equals
\begin{equation}\label{eq:factor-deficiency}
\max_{S\cap T=\varnothing}
\bigl(Q(S,T)+g(T)-d_{G-S}(T)-f(S)\bigr).
\end{equation}
Here $Q(S,T)$ counts components $C$ of $G-S-T$ on which $g=f$ at every vertex and
$f(V(C))+e(C,T)$ is odd. A nonisolated vertex with $g=f$ must have original degree $D$ and $g=f=p$. Thus all vertices in these odd components have full degree.

Fix the same $S,T$ and the same odd components throughout. For each assignment $\mathcal M$, the matching edge subgraph on an odd component $C$ gives
\[
e_{L_{\mathcal M}}(C,S)+e_{G-L_{\mathcal M}}(C,T)
\equiv1-|U_{\mathcal M}\cap C|\pmod2.
\]
Summing over these components and charging only exposed vertices within them gives
\begin{align}
Q(S,T)\le {}&
d_{L_{\mathcal M}}(S)+(d_G-d_{L_{\mathcal M}})(T)-e(S,T)\notag\\
&+|U_{\mathcal M}\setminus(S\cup T)|. \label{eq:matching-cut}
\end{align}
For even $D$, a positive-boundary odd component has at least two boundary edges; a zero-boundary one is an actual isolated odd-order $D$-regular component. Hence
\[
Q(S,T)\le
\frac{d_G(S)+d_G(T)}2-e(S,T)-e(S)-e(T)+c_*.
\]
Multiply \eqref{eq:matching-cut} by $\theta$ and this inequality by $2t$, and take expectations. After discarding nonpositive internal-edge terms, the pointwise coefficients on $S,T$ are
\[
t\,d+\theta m,\qquad (1-t)d-\theta m.
\]
At covered full-degree vertices these are exactly $p,q$. At low-degree roots, $m\le\alpha$ and $m\ge1$ give the respective bounds $f,d-g$; low-degree exterior vertices obey the same bounds. An exposed full-degree vertex in $S$ costs nothing extra, whereas one in $T$ costs $\theta$. These and the exposed vertices within odd components are disjoint parts of the same $U_{\mathcal M}$. Their total contribution is at most $\theta h_f$, counted once. Substitution in \eqref{eq:factor-deficiency} bounds the deficiency by
\[
B_E=\theta h_f+2t c_*.
\]

For odd $D$, the parity identity for an odd component implies that $e(C,S)$ is positive and odd. Thus
\[
Q(S,T)\le d_G(S)-e(S,T)-2e(S).
\]
Combine this with \eqref{eq:matching-cut} with weights $1-\theta,\theta$. The coefficients become
\[
(1-\theta)d+\theta m,\qquad \theta(d-m).
\]
The same partition of the exposed set bounds the deficiency by $B_O=\theta h_f$. In either case there is an actual red edge subgraph $F$ of maximum degree at most $p$, with total blue excess $\sum_x(d_{\mathrm{blue}}(x)-q)_+$ at most $\floor{B_E}$ or $\floor{B_O}$.

Suppose the original attaining pair requires $u$ red stars with $p+1$ edges or $v$ blue stars with $q+1$ edges, where $u+v-1=\lambda$. If $B<\lambda$, partition the blue excess vertices into $T_R,T_B$ of sizes at most $u-1,v-1$. Recolor all edges meeting $T_R$ red. Every newly created red $(p+1)$-star meets $T_R$, and every remaining blue $(q+1)$-star is centered in $T_B$. Thus the required $u$ or $v$ vertex-disjoint largest stars cannot occur. This avoids both complete original targets.

If $B\ge\lambda$, start from unsaturated cores in the maximum-flow residual network and let $\mathcal J$ be the reachable cores. Their entire root neighborhood consists of reachable roots, whose sink edges are saturated and whose positive inflow comes only from $\mathcal J$. Therefore
\[
h_f=|\mathcal J|-\sum_{x\in A_{\mathcal J}}\alpha_x,\qquad
\rho_{\mathcal J}=h_f+\sum_{x\in A_{\mathcal J}}(\alpha_x-1).
\]
Substituting the capacities in $B_E\ge\lambda$ or $B_O\ge\lambda$ gives precisely \eqref{eq:failure-even} and \eqref{eq:failure-odd}. Each selected root receives at least one unit from $\mathcal J$, so an actual integral matching saturates these roots into distinct cores. Internal near-perfect matchings and the original exterior matching give \eqref{eq:packmatching}, completing the proof of the interface.

\subsection{All actual odd regular components enter the same package}
We have
\begin{equation}\label{eq:cstar}
c_*\le\rho_{\mathcal J}.
\end{equation}
Indeed, within an actual isolated odd regular component all root capacities are one, and its network is separate from the others. If a core were unreachable, its source edge would be saturated. Its flow could not enter a reachable root, since the reverse residual edge would make the core reachable. Every neighbor of an unreachable root is also unreachable. Thus unreachable cores and roots would form a nonempty equal-sized subnetwork, contrary to the strict Hall condition of the original root assignment. All the original cores and roots in this component therefore belong to the package. Its contribution to $|\mathcal J|-|A_{\mathcal J}|$ is its positive odd matching deficiency, at least one. Contributions from other actual components are nonnegative, proving \eqref{eq:cstar}.

\begin{proposition}[A nonpositive root cost on one package]\label{prop:pack}
Unless there is a complete original avoiding coloring, the actual nonempty package can be chosen with $\rho_{\mathcal J}\ge\lambda$. Let $C$ be the union of its positive-weight whole cores and put
\[
Z=\nu(C),\quad W=e(C)-DZ,\quad
\Gamma=\zeta_{\mathcal J}+\varepsilon_{\mathcal J}+\lambda-W.
\]
For $D\ge6$,
\begin{equation}\label{eq:Gamma}
\Gamma\le-\frac{D+1}{4}\rho_{\mathcal J}<0.
\end{equation}
\end{proposition}
\begin{proof}
In the unbalanced even case, \eqref{eq:failure-even} and \eqref{eq:cstar} give
\[
\zeta_{\mathcal J}\le\frac{D-2}{p-1}(\rho_{\mathcal J}-\lambda).
\]
In the odd case, \eqref{eq:failure-odd} gives
$\zeta_{\mathcal J}\le(D-1)(\rho_{\mathcal J}-\lambda)/(p-1)$.
Both force $\rho_{\mathcal J}\ge\lambda$.
The handshake identity on all cores of the same package is
\[
\sum_{Q\in\mathcal J}w_Q
=\frac{D\rho_{\mathcal J}+\zeta_{\mathcal J}+\beta_{\mathcal J}}2
+\varepsilon_{\mathcal J}.
\]
Since $W$ is the positive part of this sum,
\[
\Gamma\le
\frac{\zeta_{\mathcal J}-D\rho_{\mathcal J}-\beta_{\mathcal J}}2+\lambda.
\]
Using $p\ge3$ yields, respectively,
\begin{align*}
\Gamma&\le-\frac{D-2}{2}\lambda
-\frac{D+2}{4}(\rho_{\mathcal J}-\lambda)-\frac{\beta_{\mathcal J}}2,\\
\Gamma&\le-\frac{D-2}{2}\lambda
-\frac{D+1}{4}(\rho_{\mathcal J}-\lambda)-\frac{\beta_{\mathcal J}}2.
\end{align*}
Both imply \eqref{eq:Gamma}. In the odd case, retaining $q\le D-3$ gives the sharper form
\begin{equation}\label{eq:Gamma-odd-sharp}
\Gamma\le-\frac{D+3}{4}\rho_{\mathcal J}
-\frac{D-5}{4}\lambda-\frac{\beta_{\mathcal J}}2.
\end{equation}
Indeed, \eqref{eq:failure-odd} implies
$2\zeta_{\mathcal J}\le(D-3)\rho_{\mathcal J}-(D-1)\lambda$; substitution in the handshake estimate proves the claim. It also gives $\rho_{\mathcal J}\ge\lambda+1$ in the odd case.

The remaining case has even $D$ and equal odd bounds $p=q$. We do not use a capacity with denominator $q-p$. Each positive-boundary bad component in the same factor criterion has at least two boundary edges, while zero-boundary bad components are precisely actual isolated odd regular components. The half-cut count bounds the total blue excess by $c_*$. If $c_*<\lambda$, an avoiding coloring follows. Otherwise take all original cores and roots in these actual components. Their roots have original degree $D$, so $\zeta_{\mathcal J}=0$, and $\rho_{\mathcal J}\ge c_*\ge\lambda$. The handshake identity gives
\[
\Gamma\le\lambda-\frac D2\rho_{\mathcal J}-\frac{\beta_{\mathcal J}}2
\le-\frac{D-2}{2}\rho_{\mathcal J},
\]
which again implies \eqref{eq:Gamma}.
\end{proof}

Put $a=|A_{\mathcal J}|$ and
$Y=(G-A_{\mathcal J})-E(C)$. Equation \eqref{eq:packmatching} gives
\[
\nu(Y)=\nu(G)-a-Z.
\]
The strong lengths \eqref{eq:stronglength} imply $s,t\ge Z+a+r+d-1$. The exact root-and-core edge cost is
\[
Da-\zeta_{\mathcal J}-\varepsilon_{\mathcal J}+DZ+W.
\]
Thus the net root term in every subsequent full-profile comparison is exactly $\Gamma$. It replaces the coarse all-root estimate and is not added to it. Any deletion inventory or highest-column bound used on this subpackage must be proved from the actual forcing relations on this $Y$; additional properties of the all-root decomposition do not transfer automatically.


\section{Profile descent retaining every original pair}\label{sec:profile-descent}
Interchange colors so that the original largest heads are $A\le B$, and put
\[
\gamma=D-B.
\]
Since the original row attains $D+1$,
\begin{equation}\label{eq:gaprange}
A+B\ge D+2,\qquad A,B\ge\gamma+2,\qquad D\ge2\gamma+2.
\end{equation}
These are conditions on the original input, not on a successor with changed lists.

\subsection{One deletion history and actual column losses}
Use the all-root notation of Section \ref{sec:regions}. Fix a history that repeatedly deletes a current maximum-degree vertex of $Y$. Let $C_j(Y)$ count the steps whose current degree is at least $j$. These counts decrease with $j$. If $Y$ has a positive uniform $(p,q)$-star source of rank $w$, where $p+q-1=h$, then
\begin{equation}\label{eq:inventory}
C_{h-1}(Y)\ge w.
\end{equation}
Otherwise fewer than $w$ vertex deletions would lower the maximum degree to $h-2$. Vizing's theorem would split its matching classes into $p-1$ and $q-1$ classes, contradicting the positive source after those deletions.

For the isolated core union $C$, put
\[
M_j=\min\left(Z,\floor{\frac n{j+1}}\right).
\]
At most $M_j$ stars of size at least $j$ in one complete original copy can be consumed entirely inside $C$. A vertex-cover coloring and the two monochromatic restorations yield the following column loss on the same $Y$:
\begin{equation}\label{eq:pointtax}
J(p,q)=\max\left\{M_p,M_q,
\floor{\frac{n(p+q)-R_0pq}{pq+p+q}}\right\}.
\end{equation}
The cover and complete-source derivation is in Appendix \ref{app:prices}: each restoration first produces a fixed residual list, and Lemma \ref{lem:intersection} is then applied. The resulting uniform width is at least $(L_h^0-J(p,q))_+$; copies from different restorations are not combined. For $h=B+j$, $j\ge2$, the coarser summable bound
\begin{equation}\label{eq:eta}
\eta_j=\frac n{j+2}+\pos{\frac n{B+1}-R_0}
\end{equation}
is also available.

\subsection{The low platform at a large head gap}
\begin{proposition}\label{prop:largegap-low}
If $\gamma\ge7$ and $L_{B+2}^0\le Z$, the original strict source does not exist.
\end{proposition}
\begin{proof}
Compare the column price on the same $Y$ with \eqref{eq:Ycost}, using \eqref{eq:pointtax}. This gives
\[
d\ge(\gamma-1)(Z-U)-W+(D-1)\rho+1-\tau-c_0,
\]
where $U=\floor{n/3}$, $0\le\tau\le c/2$, and $c_0\le\rho$ counts original port-free cores. Since $\gamma\ge7$, $W\le Z$, and $c\le2Z/D$,
\[
d\ge(1-5/D)Z+(D-2)\rho+1>Z/2,
\]
using $D\ge16$. Each lower column is charged on the same envelope $u_h(Y)$ and against the original root-core identity.

For $B+2\le h\le D+1$, fixed-list uniform source widths satisfy
\[
0\le w_h\le L_h^0\le Z,\qquad
w_h\ge(L_h^0-\eta_{h-B})_+.
\]
Set
\[
k=\max\left(\max_{B+2\le h\le D+1}w_h,(Z-R_0)_+\right)\le Z.
\]
The complete-forest/matching endpoint retains deficit $d$ and has matching-side width $q_H\ge2d>Z$. Thus the same history actually contains more than $Z$ vertices of current degree at least $B$. Delete its first $k$ vertices, obtaining $X$. They avoid $C$, which is already isolated in $Y$.

Since $r_0+k>Z$, Lemma \ref{lem:menu} gives
$P_{r_0+Z+k}(X;\nn,\mm)$. The new rank is
$N'=\nu(Y)+d-k\ge d>0$ and is smaller than $N$. The original lists have not changed.
On this single history,
\[
\min(Z+k,L_h^0)-\min(k,C_{h-1}(Y))\le\eta_{h-B}
\]
for the tracked high columns, and the first high column $B+1$ loses at most $2Z$. Combining the full row shift with the root and core costs yields
\[
e(X)-E_{r_0+Z+k}\le W-(D-1)\rho-1+T,
\]
\[
\frac TZ\le2-\gamma+
\left(2+\frac1{\gamma+1}\right)
\left(H_{\gamma+3}-H_3+\frac{\gamma}{\gamma+3}\right).
\]
Here $H_m=\sum_{j=1}^{m}1/j$. Write the right side as $-1-G(\gamma)$. Directly,
$G(7)=743/4032>0$. With
$S(\gamma)=H_{\gamma+3}-H_3+\gamma/(\gamma+3)$,
\[
S(\gamma+1)-S(\gamma)
=\frac{\gamma+6}{(\gamma+3)(\gamma+4)}\le\frac{13}{110}.
\]
Hence $G(\gamma+1)-G(\gamma)\ge1-(17/8)(13/110)>0$.
Therefore $T<-Z$, and $W\le Z$ gives strict descent.
\end{proof}

\subsection{An unconditional inventory at head gap six}
\begin{lemma}\label{lem:gap6}
If $\gamma=6$, either a complete original avoiding coloring exists, or the same all-root residual construction satisfies $d>Z/2$ and $C_B(Y)>Z$.
\end{lemma}
\begin{proof}
Here $B\ge8$ and $D=B+6\ge14$. The total loss on columns $2,\ldots,B+1$ is at most $BZ+\tau$, where $0\le\tau\le c/2$; only candidate $(2,2)$ needs this correction. For $h=B+j$, $2\le j\le6$, an attaining pair has $p,q\ge j+1$ and $p+q=B+j+1$. All three terms in \eqref{eq:pointtax} are at most
\[
n\theta_j(B),\qquad
\theta_j(B)=\frac{B+j+1}{B(j+2)+j+1}.
\]
The highest-column loss is at most $c_0\le\rho$. Put $S_6(B)=\sum_{j=2}^{6}\theta_j(B)$. Comparing the same net price $e(Y)-\nu(Y)$ gives
\[
d\ge5Z-(2Z+c)S_6(B)-W+(D-1)\rho+1-\tau-c_0.
\]
Consequently
\[
d\ge\left[4-2S_6(B)-\frac{2S_6(B)+1}{D}\right]Z
+(D-2)\rho+1.
\]
Every $\theta_j$ decreases with $B$, so the bracket increases. At $B=8$,
\[
S_6(8)=\frac{57009193}{44911405},\qquad
4-2S_6(8)-\frac{2S_6(8)+1}{14}
=\frac{151970297}{125751934}>\frac12.
\]
Thus $d>Z/2$. The same complete-forest/matching endpoint gives $q_H\ge2d>Z$, and \eqref{eq:inventory} yields $C_B(Y)>Z$.
\end{proof}

\subsection{Near-top columns retaining earlier large heads}
\begin{lemma}\label{lem:near-top}
In the same root-deleted graph, let $C$ be a union of isolated whole factor-critical cores, with $n$ vertices, of which $c_0$ cores had no original ports. For an actual pair $a+b=D+1$ attaining column $D$, suppose $A>a$ or $B>b$. The same erased graph $Y$ then has a uniform source for that column with loss at most
\[
\floor{\frac n{\min(a,b)+1}}+c_0.
\]
For an attaining pair $a+b=D$ at height $D-1$, the same conclusion holds if $A\ge a+2$ or $B\ge b+2$.
\end{lemma}
\begin{proof}
A core with a port has an internally lower-degree vertex. A near-perfect matching exposing that vertex covers every internal degree-$D$ vertex, so Lemma \ref{lem:split} supplies every strict split. A port-free core is an actual isolated odd-order $D$-regular graph and occurs only for even $D$. Even/even bounds split strictly. For odd/odd bounds the ordinary factor deficiency is at most one, with the exceptional unit assignable to either color. Recoloring one edge at the exceptional vertex moves the exception to the other color.

Suppose $A>a$ and $a+b=D+1$. Use bounds $a/(b-1)$ and put $M_a=\floor{n/(a+1)}$, $L=f_a+g_b-R_0$. Only $L>M_a+c_0$ requires proof.
If $f_a>M_a$, put all exceptions on the red side. There is no internal blue $b$-star, and a complete red copy loses at most $M_a$ components of size at least $a$. Choose
\[
i=\min(f_a-M_a,R_0),\quad j=R_0+1-i.
\]
This legal original pair produces positive uniform targets of rank $L-M_a$.
If $f_a\le M_a$, put the exceptions on the blue side. Internal red degree is at most $a<A$, so the largest star of the complete $F_1$ survives in $Y$. The extreme original pair $(1,R_0)$ gives
\[
Y\to\left(S_a,(g_b-R_0+1-c_0)S_b\right).
\]
The second target is positive because $L>M_a+c_0$, and the rank is at least $L-M_a-c_0$. Strong lengths ensure legality of the extreme pair.

Interchange colors if the earlier larger head is blue. When $a+b=D$, use bounds $(a+1)/(b-1)$. The condition $A\ge a+2$ still preserves the largest original star in the second restoration; the rest of the proof is unchanged.
\end{proof}


\subsection{One fixed parameter protecting every large head}
\begin{lemma}\label{lem:fixed-restoration}
Suppose $C$ is isolated from the exterior in an actual graph $X_0$, with $\nu(C)=Z$, $|C|=n$, and $P_q(X_0;\nn,\mm)$. Put $M_h=\min(Z,\floor{n/(h+1)})$ and define residual lists by
\[
\widetilde f_h=\max_{j\ge h}(f_j-M_j)_+,\qquad
\widetilde g_h=\max_{j\ge h}(g_j-M_j)_+.
\]
Whenever the relevant suffixes are nonempty, $Y=X_0-E(C)$ has the two fixed complete sources
\[
P_q(Y;\widetilde{\nn},\mm),\qquad
P_q(Y;\nn,\widetilde{\mm}).
\]
For $1\le\ell\le q-1$, Lemma \ref{lem:intersection} yields
$P_{q+\ell}(Y;\mathbf K,\mathbf L)$, with cumulative loss at height $h$ at most $(M_h-\ell)_+$ on either side.
\end{lemma}
\begin{proof}
Restore $C$ in red. The blue forest survives entirely. At most $M_j$ stars of size at least $j$ in a single red $F_i$ can lie inside $C$. Hence at least
$(f_j-i+1-M_j)_+$ such stars survive. Maximizing over $j\ge h$ leaves at least
$(\widetilde f_h-i+1)_+$ stars of size at least $h$, exactly the cumulative requirement of the fixed residual suffix. All height requirements hold in the same remaining forest. The all-blue restoration is symmetric, and \eqref{eq:intersection} gives the asserted loss bound.
\end{proof}

For the monochromatic fixed-restoration use of Lemma \ref{lem:allowance-tax},
take the source with the red residual as its first source. Its blue list is
the original full $\mm$, so every original candidate has $S_q=g_q>0$.
The candidate-survival condition is therefore satisfied. Nonemptiness of both residual lists and positivity of the shifted rank are checked in each descent before the complete intersection source is used.


\begin{proposition}\label{prop:largegap-high}
If $\gamma\ge8$ and $L_{B+2}^0>Z$, the original strict source does not exist.
\end{proposition}
\begin{proof}
Put $L=L_{B+2}^0$. An actual attaining pair has
$a_1+b_1=B+3$ and $a_1,b_1\ge3$. Since $B\ge10$ and $D\ge18$,
\[
M_{a_1}+M_{b_1}\le n\left(\frac14+\frac1{B+1}\right)
\le\frac{95}{132}Z<Z.
\]
At least one original count exceeds its capacity inside $C$. Restoring $C$ in that color gives a positive uniform source on the same $Y$ of rank $w\ge L-M_3>Z-M_3$. Hence $C_{B+1}(Y)\ge w$.
Choose
\[
\ell=M_6=\floor{n/7},\qquad h=(\ell-R_0)_+.
\]
As $M_3+\ell\le11n/28\le209Z/252<Z$, the same history supplies the required $h$ vertices of current degree at least $B+1$. Delete them to obtain $X$.

Equivalently, first delete $A_0$ and these $h$ vertices, and then restore the isolated $C$. The single parameter $\ell$ in Lemma \ref{lem:fixed-restoration} is legal because $\ell\le R_0+h$. The resulting fixed lists preserve every original cumulative count at heights at least six. The unit loss is at most $Z-\ell$, and the losses at heights two through five are at most $M_j-\ell$. The new rank satisfies
\[
N'\ge N-a-h-2Z+\ell
\ge2(\mu-Z)+R_0+2d+\ell-h>0,
\quad N'\le N-a-h-\ell<N.
\]
Both original largest heads are at least ten, so the new lists are nonempty.

Since $Z\ge9c$, the following integer bounds hold:
\[
\begin{gathered}
2M_2\le Z+M_3,\quad M_2+M_3\le Z+M_4,\\
M_2+M_4\le Z+M_5,\quad2M_3\le Z+M_5,\\
M_2+M_5-M_6\le Z,\quad M_3+M_4-M_6\le Z.
\end{gathered}
\]
For example, the last two left sides are at most $5n/14+6/7$ and $43n/140+6/7$. Substitute $n=2Z+c$ and $Z\ge7c$. The other bounds follow from the same floor estimates.
Charge every original height candidate. Column one loses $2Z-\ell$; columns two through five lose at most $Z+M_j-\ell$; columns six through $B$ lose at most $Z$ each. The next four columns lose at most $M_2,M_3,M_4,M_5$. In the remaining $\gamma-3$ columns both heads are protected, and only row loss $\ell$ remains. Thus the full loss for this single pair of fixed lists is
\[
\mathrm{Fee}\le(B+1)Z+2\sum_{j=2}^{5}M_j+(\gamma-8)\ell.
\]
Combine the original root shift and the later vertex shift, charging the highest column $\lambda$ only once. The actual later saving is at least $(B+1)h$. With \eqref{eq:rootdebt},
\[
\begin{split}
e(X)-E' \le{}& W-(D-1)\rho-1+(1-\gamma)Z\\
&+2\sum_{j=2}^{5}M_j+(\gamma-8)\ell+(\gamma-1)h.
\end{split}
\]
For fixed $Z,n,\ell,h$, increasing $\gamma$ by one changes this bound by
$-Z+\ell+h\le-Z+2\ell<0$. It therefore suffices to consider $\gamma=8$.
If $h=0$, then $2\sum_{j=2}^{5}M_j\le19n/10$ and $W\le Z$ give
\[
e(X)-E'\le-\frac{11}{5}Z+\frac{19}{10}c-(D-1)\rho-1<0.
\]
If $h>0$, use $h=\ell-R_0$ and $7\ell\le n$ to obtain
\[
e(X)-E'\le-\frac15Z+\frac{29}{10}c-7R_0-(D-1)\rho-1.
\]
Since $Z\ge9c$, $c\le a+\rho$, and $R_0=a+r-1$, this is at most
\[
-\frac{59}{10}a-7(r-1)-\left(D-\frac{21}{10}\right)\rho-1<0.
\]
The positive smaller rank now gives the contradiction.
\end{proof}

Propositions \ref{prop:largegap-low} and \ref{prop:largegap-high} exclude the entire range $\gamma\ge8$, retaining unbounded multiplicities, unit tails, core orders, ports, and exterior graphs. The attaining pair need not contain both original largest heads.


\section{Complete classification near the largest head}\label{sec:near-limit}
The preceding section handles large head gaps. We now state the exact interfaces of the remaining classifications and expand a principal cost argument at gap one. Every classification concerns the same original minimal source; roots and cores belong to the specified all-root graph or fractional package.

\subsection{Head gaps two through seven}
\begin{proposition}\label{prop:gap2to7}
A strict source satisfying \eqref{eq:normal} and \eqref{eq:base} cannot have
$\gamma=3,4,5,6,7$, or $D\ge9$ and $\gamma=2$.
\end{proposition}
\begin{proof}
Every complementary platform is treated by a fixed complete successor on one actual residual graph. The exhaustive splits and proofs are:
\begin{center}\small
\begin{tabular}{P{.14\textwidth}P{.58\textwidth}P{.17\textwidth}}
\toprule
Gap & Exhaustive split & Full proof\\\midrule
$6,7$ & Low and high platforms; same-history sources supply the required current-degree vertices & Appendix \ref{app:gaps67}\\
$4,5$ & High and low platforms within a fractional package and joint low-column losses & Appendix \ref{app:gaps45}\\
$3$ & Low platform; actual $3/B$ or two-heads-at-least-four winners on the high platform & Appendix \ref{app:gap3}\\
$2$ & Low, middle, and high platforms; $2/B$, $3/B$, and two heads at least four; failed selection produces an actual dominant core & Appendices \ref{app:gap2-stages}--\ref{app:gap2-dominant}\\
\bottomrule
\end{tabular}
\end{center}
The final-degree-$B$ states at gap two are included: for $\ell_0=R+C_{B+1}(Y)$ the proofs treat $\ell_0<M_3$, $M_3\le\ell_0\le M_2$, and $\ell_0>M_2$. The following lemma handles package deficiency one and the two-positive-core boundary of a dominant core. Each actual operation supplies every new pair, fixed lists, and $1\le N'<N$, and strictly pays the full original profile loss. This contradicts Lemma \ref{lem:descent}.
\end{proof}

\begin{lemma}\label{lem:rho1}
Let $D\ge8$ be even. A fractional package with $\rho_{\mathcal J}=1$ excludes the original strict source. The two-positive-core boundary in Appendix \ref{app:gap2-dominant} also has a complete exit.
\end{lemma}
\begin{proof}
Since $1\le\lambda\le\rho_{\mathcal J}$, we have $\lambda=1$. For unequal bounds $p<q$, \eqref{eq:failure-even} and $c_*\le1$ give
\[
0\le(p-1)\zeta_{\mathcal J}
\le(q-p)+2(p-1)c_*-(D-2)=2(p-1)(c_*-1)\le0.
\]
Hence $c_*=1$ and $\zeta_{\mathcal J}=0$. With equal bounds, the factor construction itself gives $c_*\ge\lambda$. Thus in either case the original graph has an actual isolated odd-order $D$-regular component. It is not asserted to be factor-critical or near-perfect.
The first part of Proposition \ref{prop:regular-components} gives a complete exit for $\lambda=1$ and any such actual component. Its proof derives strong lengths from a Tutte decomposition inside the component and then descends on a fixed exterior and one deletion history, independently of the final small-degree classification.

If the package of Appendix \ref{app:gap2-dominant} has only two positive cores, then $W\ge D\rho_{\mathcal J}/2$, each core weight is at most $D/2$, and $\rho_{\mathcal J}=1$ has just been excluded. Thus $\rho_{\mathcal J}=2$ and both core weights equal $D/2$. Their boundaries are zero. They are actual isolated odd regular components and $\lambda\le2$. The strong lengths \eqref{eq:stronglength} supply the hypotheses of the second part of Proposition \ref{prop:regular-components}, which also does not use the final classification. Hence Appendix \ref{app:gap2-dominant} may proceed with at least three positive cores.
\end{proof}

\subsection{The automatic entry at head gap one}
Now $B=D-1$. Use the actual fractional package from Section \ref{sec:fractional}, abbreviate its deficiency by $\rho$ and its number of roots by $a$, and put $R=r+a-1$, $M_j=\floor{(2Z+c)/(j+1)}$. On the same maximum-degree deletion history write $b=C_D(Y)$. Appendix \ref{app:gap1-entry} first excludes the low platform by original-list restoration and the complete unit-column comparison. On its complement,
\begin{equation}\label{eq:gamma1gate}
R+b<M_2,\qquad C_B(Y)>M_2-R,\qquad
L:=L_D^0\ge Z+b+1-\Gamma>Z.
\end{equation}
These are conditions on the present actual package, not an inventory borrowed from another residual graph.

\begin{proposition}\label{prop:gamma1largeheads}
Suppose every pair attaining the original height-$D$ column has both heads at least four. Then no gap-one strict source exists for even $D\ge8$ or odd $D\ge11$.
\end{proposition}
\begin{proof}
Choose $\ell=\max(M_4,R+b)<M_2$ and $h=\ell-R\ge b$. Equation \eqref{eq:gamma1gate} supplies enough vertices on the same history: the first $b$ have degree $D$, and the next $h-b$ have degree $B$. The actual saving is $Bh+b$, and the resulting $X$ has maximum degree at most $D-1$.
Use the single parameter $\ell$ in Lemma \ref{lem:fixed-restoration}. All cumulative counts at heights at least four are preserved; losses at heights two and three are at most $(M_2-\ell)_+$ and $(M_3-\ell)_+$. The new rank satisfies
\[
N'\ge N-a-h-2Z+\ell=\nu(Y)+d-Z+\ell-h>0,\qquad N'<N.
\]
Its highest-column width is exactly
\begin{equation}\label{eq:topwidth}
T=L-2\ell+R>0,
\end{equation}
since all original attaining pairs are protected and other candidates only decrease. If height $D+1$ remains positive, a $D$-edge-coloring of $X$ already avoids its corresponding pair.

For the same fixed successor, the complete upper cost is
\[
g:=e(X)-E'\le
\begin{cases}
U=\Gamma+M_2+M_3-\ell-R-b-1,&\ell\le M_3,\\
U=\Gamma+M_2-R-b-1,&\ell\ge M_3.
\end{cases}
\]
If $U<0$, we are done; assume $U\ge0$. The whole-core price on the same graph and the signed uniform exposure bound (Appendices \ref{app:matching} and \ref{app:prices}) give
\begin{equation}\label{eq:topprice}
g\ge\alpha T-K,\quad
(\alpha,K)=
\begin{cases}
\bigl((2m^2-5m+5)/(2m-4),\,m(m-1)\bigr),&D=2m,\\
(k-1,\,k(k+1)),&D=2k+1.
\end{cases}
\end{equation}
This price includes both signs of the actual successor's matching deficit.

If $\ell=M_4$, floor estimates and \eqref{eq:gamma1gate} give
\[
U\le\Gamma+\frac{23}{30}Z+\frac{23}{60}c-R-b,\quad
T\ge\frac15Z-\frac25c+R+b+1-\Gamma.
\]
Use $Z\ge mc$ in the even case or $Z\ge(k+1)c$ in the odd case, $R\ge c-\rho+1$, and the precise $\Gamma$ bounds of Section \ref{sec:fractional}. The result is
\[
\alpha T-K-U\ge
\begin{cases}
\dfrac{12m^3+50m^2-81m+136}{60(m-2)}>0,&m\ge4,\\[4pt]
\dfrac{22k^2-10k-77}{20}>0,&k\ge5.
\end{cases}
\]
In the odd case, the same failure inequality also forces $\rho\ge\lambda+1$ and $c\ge\rho+1$.

If $\ell=R+b\le M_3$, put $x=U\ge0$. Then
\[
T\ge\frac5{12}Z-\frac7{24}c+\frac32-\frac32\Gamma+\frac x2.
\]
Since $\alpha/2>1$, it suffices to bound the part independent of $x$. Positive lower bounds for $\alpha T-K-U$ are
\[
\frac{44m^3-100m^2+169m-35}{48(m-2)}
\quad\text{and}\quad
\frac{20k^2-9k-27}{8},
\]
respectively.
If $\ell=R+b\ge M_3$, then
$T\ge Z-M_2+2-2\Gamma+x$, and $\alpha>1$. The positive lower bounds become
\[
\frac{8m^3-19m^2+28m-5}{6(m-2)}
\quad\text{and}\quad (k+1)(3k-4).
\]
The numerators have positive coefficients after setting $m=4+y$ or $k=5+y$, $y\ge0$. These three locations exhaust $\ell$, contradicting \eqref{eq:topprice}.
\end{proof}

\subsection{Prioritizing an attaining pair at gap one}
Among all pairs attaining height $D$, choose one whose smaller head is smallest. Write it as $a+b=D+1$, $a\le b$. Exactly one of
\[
a=2,\qquad a=3,\qquad a\ge4
\]
holds. The third alternative means that \emph{every} attaining pair has both heads at least four, precisely the hypothesis of Proposition \ref{prop:gamma1largeheads}. For $D\ge13$, Appendix \ref{app:gap1-two} excludes the first alternative by a two-neighbor cover, one fixed successor, and a cap price. For $D\ge8$, Appendix \ref{app:gap1-three} excludes the second. Both complete proofs retain earlier larger original heads and arbitrary unit tails. Cases outside these ranges are exactly the small-degree terminals listed next; estimates for $D\ge13$ are not extrapolated downward.


\section{Small-degree terminals and their interfaces}\label{sec:terminals}
\begin{proposition}\label{prop:small}
There is no original minimal strict source for $D=6,7,8,9,10$, nor for $D=11,12$ with $\gamma=1$.
\end{proposition}
We distinguish general structural interfaces, exact small-degree scopes, and their final uses. A small degree does not mean finitely many hosts: every case below retains arbitrary core orders, exterior graphs, multiplicities, and unit tails.

\subsection{Complete coverage at degree six}
Appendix \ref{app:D6-structure} excludes actual isolated $K_7$ components and uses the original ordinary-region selector to prove
\[
 1\le\lambda\le c\le\rho_6\le4.
\]
Each actual odd six-regular component has a near-perfect matching. Appendix \ref{app:D6-one} excludes $\lambda=1$ by one whole-component argument; Appendix \ref{app:D6-two} excludes $\lambda=2$ by two whole components. Exactly $(c,\lambda,\rho_6)=(3,3,3)$ and $(4,4,4)$ remain. Appendix \ref{app:D6-three-four} excludes both using complete profile taxes, actual root inventory, and the final cap-five price on the same exterior. Lemma \ref{lem:isolated-length} proves the general strong-length interface, and Appendix \ref{app:small-prices} supplies successor prices for both deficit signs. All original rows, exterior graphs, head multiplicities, and unit tails are retained.

\subsection{Degree seven: actual descent with four positive blocks}
We expand a complete low-degree terminal to illustrate how all original parameters are retained.
The original signed exposure bound and positive-weight account (Section \ref{sec:selector} and Appendix \ref{app:matching}) give
\[
\rho_7\le4,\qquad
W_+\ge\ceil{7\rho_7/2},\qquad W_+\le\lambda+11.
\]
A highest-column $3/6$ pair forces $\rho_7\ge5$. Only $4/5$ or its reverse remains, with
$\rho_7\ge\lambda+2\ceil{\lambda/4}$. Hence
\[
(\lambda,\rho_7)=(1,3),\qquad W_+\in\{11,12\}.
\]
Each positive block weight is one, two, or three, so at least four blocks occur. If there are at least five, choose a minimum-capacity block $Q$, with $z=\nu(Q)$ and $w=e(Q)-7z\ge1$. The other capacities imply $\mu\ge5z$. Thus $\mathcal R\setminus V(Q)$ contains more than $8(z-1)$ actual degree-seven vertices. Greedily choose $b=(z-r+1)_+\le z-1$ independent ones. Erase the internal edges of $Q$ and delete these vertices. The saving $7(z+b)+w$ pays the full loss $7(z+b)+\lambda$; the capacity menu gives positive smaller rank, a contradiction.

Exactly four positive blocks remain, with weight multiset
$\{3,3,3,2\}$ or $\{3,3,3,3\}$. Choose a minimum-capacity block $Q$. It has $z\ge4$, $w\ge2$, and boundary $\delta_Q=7-2w\le3$. Although blocks may share cutvertices, the size of the entire ordinary region gives
\[
|\mathcal R\setminus V(Q)|=2(\mu-z)+a+2\ge6z+2.
\]
After erasing the internal edges of $Q$, these exterior vertices still have degree seven, while vertices of $Q$ have degree at most three. Their induced graph is seven-colorable. A $K_8$ component would, by degree saturation, be an original isolated component, contradicting the odd factor-critical components of $\mathcal R$; apply Brooks' theorem to the rest. Hence there is an independent set of size at least $b_0=\ceil{(6z+2)/7}$. Extend it to a maximal independent set among all degree-seven vertices and delete this complete degree-seven stage first. Let its length be $b\ge b_0$.

Put $t=(z-r+1)_+$. If $t=0$ or $b\ge t$, the required vertices all have degree seven and give immediate paid descent. If $0<b<t$, the first $t$ vertices on the same history have degree at least six: after deleting $k\le z-2$ vertices, lowering every remaining original exterior vertex to degree at most five would require
$2(6z+2-k)\le7k$, which is impossible. Thus the first $b$ vertices have degree seven and the next $t-b$ degree six. The final graph $Y$ has maximum degree six, and the saving is
\[
7z+w+7b+6(t-b)=6(z+t)+z+b+w.
\]
The capacity menu gives $P_{r+z+t}(Y)$ with the original lists, and
$z+t\le2z-1<4z\le\mu<N$. Writing $L_7$ for the original seventh-column width,
\begin{equation}\label{eq:D7upper}
e(Y)-E'\le L_7-z-b-w=:\mathcal D.
\end{equation}
A negative right side already contradicts minimality.

Assume $\mathcal D\ge0$ and put $M=\floor{(2z+4)/3}$. Each star of size at least two using a restored internal edge has its center in $Q$, and external leaves use at most three actual port endpoints. Vertex-disjointness bounds their number by $M$. For every $z\ge4$, $b_0\ge M$, whence $L_7\ge z+b+w>2M$. At least one count in a pair attaining height seven exceeds $M$. Restore all internal edges of $Q$ in that color, then delete the same $t$ vertices. This gives a positive uniform source on the same $Y$ of width
\[
L=L_7-M-t>0,\qquad \rho_6(Y)\ge L.
\]
The height-seven width of the fixed original-list successor is
$T=(L_7-z-t)_+=(L-v)_+$, where $v=z-M$.

The condition $b<t$ forces $z\ge16$, and hence
$v\ge4$, $b-M\ge2$, and $L\ge r+w+1\ge5$. If the successor's matching number is at least its rank, \eqref{eq:rich} gives $e(Y)-E'\ge L>\mathcal D$.
Otherwise the actual degree-six price \eqref{eq:price6}, valid for $\rho_6(Y)\ge5$, gives
$e(Y)-E'\ge2\rho_6(Y)-T-4$. Therefore
\[
(2L-T-4)-\mathcal D
=(b-M+r-1)+w-4+\min(L,v)\ge r+3>0.
\]
Both signs contradict \eqref{eq:D7upper}. This excludes the complete original degree-seven domain.

\subsection{Interfaces for degrees eight through twelve}
\begin{center}\small
\begin{tabular}{P{.14\textwidth}P{.49\textwidth}P{.22\textwidth}}
\toprule
Input & Actual exit and retained hypotheses & Full proof\\
\midrule
$D=8$ & Even factor obstruction, strong lengths for actual components, two fixed restorations; isolated $K_9$ has its own analytic cost & \ref{app:even-small}\\
$D=9$ & $4/6$, $5/5$, and the gap-two dominant core; actual degree-$4/5$ stage costs on one history & \ref{app:D9}\\
$D=10$ & $\lambda=1,2$; a real region with total degree deficit plus boundary two, charged jointly with actual exterior components; $K_{11}$ separately & \ref{app:even-small}, \ref{app:D10}\\
$D=11,\gamma=1$ & One original row, unit axes, every intermediate head, and weighted whole-core deletion & \ref{app:D11}\\
$D=12,\gamma=1$ & The original weak entry, joint actual vertex selection, the strong odd-cap price, and every new pair & \ref{app:even-weighted}\\
\bottomrule
\end{tabular}
\end{center}
The isolated regular-component result is proved independently before the final small-degree classification, so Lemma \ref{lem:rho1} is noncircular. Appendix \ref{app:even-small} computes the complete successors for actual $K_9,K_{11}$ column by column. The general even weighted formula is needed here only for $D=12$ and retains the original-row guard and all new pairs. The degree-nine stage fee is always the actually available $\min(b,\tau)$. These full proofs establish Proposition \ref{prop:small}.



\section{Completion for all parameters}\label{sec:completion}
\begin{theorem}\label{thm:family}
For all nonincreasing positive integer lists $\nn,\mm$, all $2\le r\le s+t$, and every finite simple graph $G$,
\[
P_r(G;\nn,\mm)\quad\Longrightarrow\quad e(G)\ge E_r(\nn,\mm).
\]
\end{theorem}
\begin{proof}
Suppose there is a counterexample. Choose one by \eqref{eq:order} over all lists, all legal rows, and all host graphs. Propositions \ref{prop:normal} and \ref{prop:base} give
\[
e(G)=E_r-1,\quad \nu(G)=N-d,\quad d\ge2,\quad
D\ge6,\quad \mu<\min(s,t),\quad H\le D-1.
\]
Order the two original maximum heads as $A\le B=H$, and set $\gamma=D-B$. Then
\[
1\le\gamma,\qquad A,B\ge\gamma+2,\qquad D\ge2\gamma+2.
\]
The following cases exhaust these integer parameters.

\begin{enumerate}
\item If $6\le D\le10$, Proposition \ref{prop:small} excludes the entire original domain.
\item If $\gamma\ge8$, split according to $L_{B+2}^0\le Z$ or $>Z$ on the same all-root residual graph, and apply Propositions \ref{prop:largegap-low} and \ref{prop:largegap-high}, respectively. Both give a paid descent with positive smaller rank.
\item If $2\le\gamma\le7$ and $D\ge11$, apply Proposition \ref{prop:gap2to7}. Its proof includes the late high-platform states and dominant-core boundaries for gap two.
\item It remains that $\gamma=1$. Degrees $11$ and $12$ are excluded by Proposition \ref{prop:small}. Let $D\ge13$. Among all pairs attaining the original height-$D$ column $L_D$, choose one whose smaller head is smallest. Write $a+b=D+1$, $a\le b$. Since $H=D-1$, necessarily $a\ge2$. If $a=2$, use the $2/(D-1)$ exit of Appendix \ref{app:gap1-two}; if $a=3$, use the $3/(D-2)$ exit of Appendix \ref{app:gap1-three}. If $a\ge4$, every pair attaining height $D$ has both heads at least four, exactly the hypothesis of Proposition \ref{prop:gamma1largeheads}. All three degree ranges include $D\ge13$.
The common automatic entry \eqref{eq:gamma1gate} has already been established in \ref{app:gap1-entry}.
\end{enumerate}

Each exit gives either a coloring avoiding a complete original pair, or fixed new lists on one actual residual graph with $1\le N'<N$ and $e(Y)<E_{r'}(\nn',\mm')$. The former contradicts $P_r$ and the latter minimality. No step adds the gain of one coloring to the copy supplied by another. Nor does it require an arbitrary residual graph to inherit the structural properties of the original minimum source. Thus no minimal counterexample exists.
\end{proof}

\begin{proof}[The lower bound in Theorem \ref{thm:main}]
At $r=2$, the only legal pair is $(1,1)$, so $P_2(G)$ is the original relation $G\to(F,H)$. Theorem \ref{thm:family} gives $e(G)\ge E_2=\sum_{v=2}^{s+t}\ell_v$. Together with the construction in Section 1, this proves \eqref{eq:main}.
\end{proof}

\section{Notation and proof dependencies}\label{sec:dependencies}
\subsection{Principal notation}
\begin{longtable}{P{.22\textwidth}P{.7\textwidth}}
\toprule
Symbol & Meaning\\\midrule\endhead
$P_r,E_r,N$ & Forcing for every legal suffix pair, original profile tail sum, and component rank $s+t-r+1$.\\
$\ell_v,L_h$ & Maximum head sum minus one on row $v$, and the number of profile entries at least $h$.\\
$D,H,\gamma$ & Normalized host maximum degree, largest original head, and gap $D-H$.\\
$d$ & Difference $N-\nu(G)$ between original rank and actual matching number.\\
$A_0,\mathcal R,B_0$ & Roots, odd-core union, and exterior in the ordinary full-degree matching decomposition.\\
$\rho_D(X),\kappa_D(X)$ & Minimum number of full-degree vertices exposed by a maximum matching; minimum vertex deletion making a matching plus low-degree decomposition possible. Only $\kappa_D\le\rho_D$ is used.\\
$z_Q,w_Q$ & Matching number of an original whole core, and its weight $e(Q)-Dz_Q$.\\
$\zeta,\varepsilon,\beta$ & Original degree deficit of the chosen roots, edges within them, and edges from them to the exterior.\\
$\mathcal J,A_{\mathcal J}$ & Actual core package selected by one factor obstruction and all its neighboring roots.\\
$C,Z,c,W$ & Union of the chosen positive whole cores, matching number, core count, and total weight; $|C|=2Z+c$.\\
$\Gamma$ & Net root term $\zeta+\varepsilon+\lambda-W$ for the same package.\\
$\lambda,R$ & Original highest-column length; local row offset $r+a-1$. The numerical $R$ is distinct from the region $\mathcal R$.\\
$M_h$ & Erased-core star-capacity bound $\floor{(2Z+c)/(h+1)}$; height one is treated separately by matching capacity $Z$.\\
$C_j(Y)$ & Number of vertices removed through degree at least $j$ in one maximum-degree deletion history.\\
\bottomrule
\end{longtable}

\subsection{Internal proof map}
\begin{longtable}{P{.35\textwidth}P{.55\textwidth}}
\toprule
Content & Proof location\\\midrule\endhead
Matching reserves, ordinary regions, complete endpoints & Appendices \ref{app:reserve}--\ref{app:matching}\\
Complete residual sources and unified prices & Appendices \ref{app:prices}--\ref{app:sources}\\
Base-region, single-empty, and low-degree double-empty exclusions & Appendices \ref{app:binary}--\ref{app:base-gluing}\\
Successive degree reductions and actual regular components & Appendices \ref{app:stage-prices}--\ref{app:regular-components}\\
Head gaps two through seven & Appendices \ref{app:gaps67}--\ref{app:gap2-dominant}\\
Head gap one & Appendices \ref{app:gap1-entry}--\ref{app:gap1-two}\\
Degrees eight through twelve & Appendices \ref{app:even-weighted}--\ref{app:D10}\\
Fixed-port factors and low-cap prices & Appendices \ref{app:factor}--\ref{app:small-prices}\\
The complete degree-six classification & Appendices \ref{app:D6-structure}--\ref{app:D6-three-four}\\
\bottomrule
\end{longtable}
Successors use the inductive lower bound at smaller rank. Structural conclusions about the original minimum source are used only after checking their hypotheses. The actual regular-component descent is proved independently before its use in the low-degree exhaustion; final coverage is not used to justify an earlier lemma.

Corollary \ref{cor:local-root-core-interface} proves the local root--core
interface from the complete residual-source and endpoint arguments and
the short-head calculation. Its hypothesis $H\le D-1$ is explicit;
it does not invoke the subsequent degree-six or regular-component
classification. Proposition \ref{prop:regular-components} and Lemma
\ref{lem:isolated-length} verify this interface on their actual
decompositions. The original degree-nine and degree-ten exposure bounds
in Appendix \ref{app:gap2-high}, the Brooks selection in Appendix
\ref{app:even-small}, and the degree-ten successor price in Appendix
\ref{app:D11} include their positive-rank and full-fee checks at the
point of use.


\appendix\renewcommand{\thesection}{A\arabic{section}}
\section{Canonical envelopes and actual matching reserves}\label{app:reserve}
Only the graph inequalities used in the main proof are established here. A canonical envelope is not a free replacement subgraph of the host. Write $u_h(X)$ for the uniform-arrow width and retain $\rho_D(X)$ for full-degree exposure by a maximum matching.

\subsection{A common envelope and layer losses}
\begin{lemma}\label{lem:envelope}
Let $T(X)$ be the star forest having exactly $u_h(X)$ components of head at least $h$. It preserves every star-forest arrow of $X$. For vertex-disjoint $X_i$,
\[
u_h\left(\sq_iX_i\right)\le\sum_i u_h(X_i).
\]
Set $J(X)=e(X)-\nu(X)-\sum_{j\ge1}\kappa_j(X)$. Then $J\ge0$ and
\[
u_{j+1}(X)\le\kappa_j(X),\qquad
b(X):=e(X)-\nu(X)-\sum_{h\ge2}u_h(X)\ge J(X).
\]
The function $b$ is superadditive on vertex-disjoint unions.
\end{lemma}
\begin{proof}
A star host with heads $c_1\ge c_2\ge\cdots$ arrows $(\nn,\mm)$ precisely when $c_{i+j-1}\ge n_i+m_j-1$ for every $i,j$. For necessity, assign the at most $i+j-2$ large components to at most $i-1$ all-red components and $j-1$ all-blue components. Split each remaining component into at most $n_i-1$ red edges and $m_j-1$ blue edges. For sufficiency, an avoiding coloring has a failed red and a failed blue order-statistic threshold. Each sufficiently large component supplies one threshold, contradicting their combined capacities.

The relation $X\to(F,H)$ implies $X\to(iS_{n_i},jS_{m_j})$, so this criterion gives $T(X)\to(F,H)$. Monotonicity of $u_h$ follows by reducing one head. Suppose a disjoint union of envelopes fails the threshold $(a,b)$, and let $t_i$ count the components of $T(X_i)$ of head at least $a+b-1$. Distribute $\sum t_i$ between the red and blue packing budgets and choose, on each $X_i$, a coloring avoiding the corresponding uniform pair. A connected star cannot cross the components. The combined coloring avoids the original pair, proving subadditivity.

Deleting one vertex changes each $\kappa_j$ by $b_j\in\{0,1\}$ and the matching number by $\delta\in\{0,1\}$. At current maximum degree $k$, if $\kappa_k=0$, delete any degree-$k$ vertex. Then
\[
J(X)-J(X-v)=(1-\delta)+\sum_{j<k}(1-b_j)\ge0.
\]
If $\kappa_k>0$, a maximum matching exposes a full-degree vertex. Delete it; then $\delta=0$ and
\[
J(X)-J(X-v)=\sum_{j\le k}(1-b_j)\ge0.
\]
Iteration to an edgeless graph proves $J\ge0$.
If $\kappa_j<u_{j+1}$, delete an optimal $\kappa_j$ set and choose a uniform pair attaining $u_{j+1}$. The vertex-deletion lemma consumes the corresponding number of target components and leaves a forced star pair whose heads sum to $j+2$. The residual graph, however, splits into a matching and a graph of maximum degree $j-1$. The degree-splitting lemma splits it into caps one below those two heads, a contradiction. Thus $u_{j+1}\le\kappa_j$. Additivity of edges and matching number, together with subadditivity of $u_h$, proves the assertion for $b$.
\end{proof}

\subsection{Reserve outside isolated regular cores}
We use the standard Gallai--Edmonds decomposition: vertices exposable by maximum matchings induce factor-critical components; neighboring roots are matched to distinct such components; the remaining vertices have a perfect matching. Any specified exposure in an unmatched component, or any actual root-edge endpoint in an assigned component, is compatible with the corresponding internal perfect matching.

\begin{lemma}\label{lem:ported-reserve}
Suppose $\Delta(X)=D\ge3$ and no connected component of $X$ is both $D$-regular and factor-critical. There is an independent set $I$ of original full-degree vertices such that
\[
\kappa_D(X-I)=0,\qquad J(X)-J(X-I)\ge\rho_D(X).
\]
\end{lemma}
\begin{proof}
Put $\rho=\rho_D(X)$ and $k=\kappa_D(X)$. If $k=0$, augmenting a matching covering all full-degree vertices to a maximum matching preserves that coverage, so $\rho=0$. Assume $k>0$ and choose a maximum matching with minimum full-degree exposure. From its unassigned full-degree components, follow alternating paths in the component--root incidence graph. Reaching a nonfull component would allow a path exchange and exposure at a low-degree vertex, reducing $\rho$. Thus the reached closed region consists entirely of full components, all their external neighbors are reached roots, and the component count minus root count is $\rho$. Every reached component has a port by the hypothesis.

Call a vertex $y$ of a factor-critical graph $Q$ good if, for every $v\ne y$ and every perfect matching of $Q-v$, the mate of $y$ is adjacent to $v$. If that mate is $z$, a matching-edge exchange gives
\[
d_Q(z)\ge d_Q(v)+1-\mathbf1_{vy\in E(Q)}. \tag{A.1}
\]
If every vertex of $Q$ of internal degree below $D$ is good, then $Q$ is $D$-regular, an odd clique of order at most $D$, or, for even $D$,
\[
Q=K_f\vee\overline K_u,\qquad u\ge2,\quad f\ge u+1,\quad f+u=D+1.
\]
If $\Delta(Q)<D$, all vertices are good. Applying the definition to a near-perfect matching exposing any specified vertex shows that it is adjacent to every other vertex, so $Q$ is complete. The regular case needs no further classification. Otherwise let $F$ be the internal full-degree set and $U$ its complement, with $f=|F|$ and $u=|U|$. Both are nonempty. Equation (A.1) makes $U$ complete to $F$ and matches all of $U$ into $F$ whenever a vertex of $F$ is exposed. Hence $f\ge u+1$ and
\[
m_0=(f-u-1)/2
\]
is a nonnegative integer. Any $m_0$-edge matching $H$ in $F$ has an uncovered set $W$ of size $u+1$. For distinct $v,z\in W$, match $U$ bijectively to $W-\{v\}$, assigning some $y\in U$ to $z$, and adjoin $H$. This perfectly matches $Q-v$. Goodness of $y$ forces $vz\in E(Q)$, so every such $W$ is a clique. A maximum matching of $F$ exposes at most one vertex: otherwise an $m_0$-edge submatching leaves two originally exposed vertices in one uncovered clique, whose edge augments the original matching.

Every vertex of $U$ has degree below $D$ and is adjacent to all $f$ vertices of $F$, giving $D\ge f+1$. If $u=1$, vertices of $F$ have degree at most $f$, a contradiction. If $u=2$, then
\[
D-2=d_{Q[F]}(x)\le f-1\quad(x\in F),
\]
so $D=f+1$ and $F$ is complete. If $u\ge3$, choose a maximum matching $M$ in $F$. It exposes at most one vertex, and the odd parity of $f+u$ gives $|M|-m_0\ge2$. For specified vertices on different matching edges, remove those two edges and enough further edges to leave $m_0$ edges. Both vertices are then in one uncovered clique. If one is already exposed, remove the other vertex's matching edge first. Vertices paired by $M$ are already adjacent. Thus $F$ is complete in every case.

Now $D=f+u-1$ is even, and every vertex of $U$ has a nonneighbor in $U$. If $yz$ were an edge in $U$, choose a nonneighbor $v\in U$ of $z$, retain $yz$, match $U-\{v,y,z\}$ into $F$, and perfectly match the remaining even number of clique vertices. This perfectly matches $Q-v$, but the mate $z$ of the good vertex $y$ is not adjacent to $v$, a contradiction. Hence $U$ is independent, completing the classification.

Call a component good when all its ports are good. Among assignments saturating all roots, maximize the number of assigned good components. The unassigned good components and the alternating paths using only good components form a closed subregion of deficiency $h$. There are $\rho-h$ unassigned bad components. Put $\eta=\rho-k$. If a root set $T$ neighbors $q(T)$ unassigned components, then
\[
q(T)\le |T|+\eta. \tag{A.2}
\]
Otherwise delete $T$ and one exposed vertex from each remaining unassigned component. Fewer than $k$ deletions would leave a matching covering every remaining full-degree vertex. Applying Hall's theorem to the bad-port incidence graph, with $\eta$ auxiliary universal roots, assigns at least $b\ge k-h$ unassigned bad components to distinct bad-port roots.

For each chosen port $ay$, fix an exposure $v$ and an internal matching edge $yz$ with $vz\notin E(Q)$. These roots and components are distinct. In the good subregion, select one vertex from each odd clique and the entire independent side from each split graph; call their union $P$. Choose the original maximum matching $M$ with the unassigned good exposures and the good roots' mates in $P$, and with the prescribed bad certificates. Let $S$ be its $\rho$ exposed full-degree vertices.

For odd $D$, each good component has at least $D$ boundary edges. Root degree capacity in the closed subregion forces $h=0$. For even $D$, write $a_g$ for the good-root count and $t_g$ for the split-component count. Compare actual deletion of $P$ with the competing witness that deletes the good roots and one clique vertex in each split component. At layer $D-1$ this yields
\[
 |P|-\bigl[\kappa_{D-1}(X-(S\setminus P))-
 \kappa_{D-1}(X-(S\cup P))\bigr]
 \ge |P|-a_g-t_g
 =h+\sum_{\rm split}(u-2)\ge h. \tag{A.3}
\]
The witness is valid: the remaining odd cliques have maximum degree at most $D-2$, and a matching covers all degree-$D-1$ vertices of a split graph after the indicated clique deletion.

Also delete the original mates of the $b$ chosen bad roots, omitting those already in $P$, and call the additional set $B$. After deleting $S\cup P\cup B$, replace each $yz$ by $ay$. The exchanges are disjoint and expose vertices $z$ still of original degree $D$. Delete these $z$ as well. The resulting set $I$ is independent and consists of original full-degree vertices. With $t=|P\setminus S|+|B|$,
\[
|I|=\rho+t+b,\quad \nu(X-I)\ge\nu(X)-t,\quad
\kappa_D(X-I)=0.
\]
Each layer loses at most $|I|$, and (A.3) saves at least $h$. Therefore
\[
J(X)-J(X-I)\ge\rho+b-k+h\ge\rho.
\]
Delete $S$ first and the rest of $I$ afterwards. The maximum matching survives the first stage; the second has $\kappa_D=0$. Thus every step is permitted by the preceding nonnegative-loss recurrence.
\end{proof}

\subsection{Regular factor-critical components}
\begin{lemma}\label{lem:regular-reserve}
If $Q$ is connected, $D$-regular and factor-critical, with $D\ge3$, then $b(Q)\ge1$.
\end{lemma}
\begin{proof}
Write $|Q|=2z+1$; $D$ is positive and even, and $\kappa_D(Q)=1$. Deleting any nonempty independent set $I$ one vertex at a time gives
\[
J(Q)=J(Q-I)+|I|-1-z+\nu(Q-I)
 +\sum_{j<D}\bigl(|I|-\kappa_j(Q)+\kappa_j(Q-I)\bigr). \tag{A.4}
\]
Every term on the right is nonnegative. If $J(Q)=0$, apply this to the exposed set of a maximal matching. Every maximal matching must then be maximum.

Suppose $Q$ is noncomplete and $\alpha(Q)\le2$. Its complement is triangle-free, nonbipartite and $c$-regular, for positive even $c$. Every vertex outside a shortest odd cycle has at most two neighbors on it, so $|Q|\ge5c/2$. Together with $D\ge4$ and odd order, this gives $|Q|\ge2c+3$. For a complement edge $vw$, the graph $Q-\{v,w\}$ is connected, has no cut vertex, has minimum degree at least half its order minus one half, and is nonregular. Every connected even-order graph with independence number at most two has a perfect matching: otherwise a maximum matching exposes exactly two vertices; their neighbor classes, and the absence of augmenting paths of lengths three and five, disconnect the graph. Applying this observation after deleting each vertex shows that $Q-\{v,w\}$ is factor-critical. We supply the positive reserve using the deletion recurrence already proved. Let $X$ be nonregular and factor-critical, put $k=\Delta(X)$, and choose $y$ with $d_X(y)<k$. A perfect matching of $X-y$ covers all degree-$k$ vertices, so $\kappa_k(X)=0$. Choose a degree-$k$ vertex $x$. Factor-criticality gives $\nu(X-x)=\nu(X)$. Write $b_j=\kappa_j(X)-\kappa_j(X-x)\in\{0,1\}$; both terms vanish for every $j\ge k$. Consequently
\[
J(X)-J(X-x)=k-\sum_{j<k}b_j
=1+\sum_{j<k}(1-b_j)\ge1.
\]
Lemma \ref{lem:envelope} already gives $J(X-x)\ge0$, hence $J(X)\ge1$. Apply this to $X=Q-\{v,w\}$ and use the independent set $\{v,w\}$ in (A.4) to obtain $J(Q)\ge1$. This argument does not invoke the lemma being proved.

It remains that $\alpha(Q)\ge3$. The regular-equimatchable classification in \cite{zhao}, Theorems 1.9 and 1.7(iii), reduces $D\ge6$ to $F_D$, obtained from $K_{D,D}$ by deleting $D/2$ matching edges and adding a vertex adjacent to their endpoints. At $D=4$ there are also $\overline F_4$ and $\overline C_7$. In $F_D$, an independent set has size $D$. Deleting the endpoints of one replaced matching edge leaves a matching covering every degree-$D-1$ vertex, so $\kappa_{D-1}\le2$, and (A.4) gives $J\ge D-2$. In $\overline F_4$, there is an independent triple; deleting one unreplaced endpoint from each original bipartition class leaves three edges covering all degree-three vertices. Thus $\kappa_3\le2$ and $J\ge1$. The graph $\overline C_7$ belongs to the preceding independence-number-two case.

Consequently $J(Q)=0$ is possible only for odd cliques. For $Q=K_{2z+1}$, $z\ge2$, we have $\kappa_1=2z-1$. The vertex requirements of the all-red and all-blue colorings give
$u_2\le z+\floor{(2z+1)/3}-1\le2z-2$. Hence $b\ge J+\kappa_1-u_2\ge1$. In the noncomplete case, $b\ge J\ge1$ already suffices.
\end{proof}

Remove all isolated $D$-regular factor-critical components and apply Lemma \ref{lem:ported-reserve} to the remaining graph. Superadditivity of $b$ and Lemma \ref{lem:regular-reserve} prove the inequality needed in the text:
\[
\boxed{\ e(X)-\nu(X)-\sum_{h\ge2}u_h(X)\ge\rho_D(X),\quad D\ge3.\ }
\]
If $\Delta(X)<D$, the right side is zero. Neither a clique-replacement theorem nor the unused parts of the zero-reserve classification are required.


\section{Ordinary matching regions and complete endpoints}\label{app:matching}
\subsection{An actual region for full-degree exposure}
\begin{lemma}\label{lem:intrinsic-region}
Let $\Delta(X)\le D$ and $W=\{v:d_X(v)=D\}$. The decomposition in \eqref{eq:region} exists. Its region $\mathcal R$ consists exactly of the vertices of $W$ that some maximum matching with minimum full-degree exposure can miss, and its root set is $A_0=N_X(\mathcal R)\setminus\mathcal R$. Its root--core incidence graph satisfies
\[
 |N_{\mathcal R}(T)|\ge |T|+1\quad(\varnothing\ne T\subseteq A_0).
\]
Each core can be avoided by an assignment covering all roots, and each actual root--core edge can belong to such an assignment. The exterior $B_0$ has a maximum matching covering every vertex of original degree $D$. For $U\subseteq A_0$, one has $\nu(X-U)=\nu(X)-|U|$.
\end{lemma}
\begin{proof}
By the Tutte--Berge formula choose $S$ maximizing $o(X-S)-|S|$, and then maximizing $|S|$. Every component of $X-S$ is odd and factor-critical. An even component would allow one vertex to be moved into $S$ without decreasing the first objective. If an odd component $C$ were not factor-critical, choose $v$ for which $C-v$ has no perfect matching. A Tutte obstruction $T$ with $o(C-v-T)\ge |T|+2$ allows $T\cup\{v\}$ to be moved into $S$, again contradicting maximal cardinality. Thus
\[
 \nu(X)=|S|+\sum_C\frac{|C|-1}{2}.
\]
Every maximum matching assigns $S$ to distinct components. Indeed, if it uses $b$ edges within $S$, misses $u$ roots, and uses $h_C$ edges between $S$ and $C$, its size is at most
\[
 \sum_C\frac{|C|-1}{2}+|S|-b-u-\sum_C\floor{h_C/2}.
\]
Equality forces $b=u=0$ and $h_C\le1$.

Retain as full core nodes only those components whose vertices all belong to $W$, and form their incidence graph with $S$. If there are $c$ full cores and the maximum incidence matching has size $m$, any such matching extends to an assignment covering $S$ by augmentation without losing a covered core. In each unassigned nonfull component a non-$W$ vertex can be exposed. Consequently the minimum full-degree exposure is $\rho=c-m$.

Starting at unmatched full cores in a maximum incidence matching, follow alternating paths and let $\mathcal I,A_0$ be the reached cores and roots. Every reached root is matched, and
$N(\mathcal I)=A_0$, $|\mathcal I|-|A_0|=\rho$. Every core family has deficiency at most $\rho$. A family attaining equality contains each core that a maximum incidence matching can miss. Alternating-path flips show that these cores are exactly $\mathcal I$. Thus $\mathcal I$ is the unique least family of maximum deficiency. Factor-criticality allows any vertex in an unassigned core to be exposed, so its union is precisely the intrinsic set $\mathcal R$ in the statement.

If a nonempty $T\subseteq A_0$ had $|N(T)|\le|T|$, removing $N(T)$ from $\mathcal I$ would leave a proper family with deficiency at least $\rho+|T|-|N(T)|\ge\rho$, a contradiction. This proves strict Hall expansion. Deleting one core preserves root-side Hall; prescribing one root--core edge and deleting its endpoints also preserves Hall. The asserted assignments therefore exist and lift through perfect or near-perfect core matchings and a fixed exterior maximum matching. That exterior matching covers all its original full-degree vertices, since otherwise the exposure would not be optimal. Deleting the entire root set loses exactly $|A_0|$ matching edges; each single deletion loses at most one. Every root subset therefore loses exactly its cardinality.
\end{proof}

\begin{lemma}\label{lem:fc-blocks}
Every block of a factor-critical graph is odd and factor-critical, and
\[
 \nu(Q)=\sum_{K\text{ a block}}\nu(K),\qquad e(Q)=\sum_K e(K).
\]
The sum of separate red or blue matching bounds on the blocks bounds the corresponding matching number of the whole graph, even when blocks share cutvertices.
\end{lemma}
\begin{proof}
If $v$ is a cutvertex, every component of $Q-v$ has even order because it has a perfect matching. Add $v$ to each such component to obtain $Q_i$. Deleting $v$ leaves a perfect matching. For $x\in Q_i-v$, a perfect matching of $Q-x$ matches each other even component internally, so its restriction to $Q_i-x$ is perfect. Recursion proves factor-criticality of the blocks and $|Q|-1=\sum_K(|K|-1)$. Block edge sets partition $E(Q)$. The restriction of any matching to a block is a matching in that block, proving the last assertion.
\end{proof}

\subsection{The full forest--matching endpoint}
\begin{lemma}\label{lem:strong-endpoint}
Let $F=\sq_{i=1}^s S_{n_i}$ and $\omega(F)=\sum_i\floor{(n_i-1)/2}$. If $Z\to(F,qK_2)$, then
\[
 \nu(Z)\ge s+\min\left(q-1,\omega(F)+\floor{(q-1)/2}\right).
\]
In particular, if $\nu(Z)=s+q-1-d$ with $d\ge1$, then $q\ge b+2(d+\omega(F))$, where $b$ is the size of any maximum-cardinality optimal Tutte barrier. In particular, $q\ge2(d+\omega(F))$.
\end{lemma}
\begin{proof}
Monochromatic colorings give $\nu(Z)\ge s,q$. Write $\nu(Z)=s+u$; only $u<q-1$ needs consideration. Choose the optimal barrier $S$ from the preceding proof and put $b=|S|$. If $b\ge u+1$, color all edges incident with any $u+1$ roots blue and all other edges red. The blue matching number is at most $u+1\le q-1$, while the red matching number is at most $s-1$, a contradiction. Hence $b\le u$.

Put $j=q-1-b$. The odd components of $Z-S$ have total matching capacity $\nu(Z)-b\ge j+1$. Choose $\floor{j/2}$ pairs of vertices in these components, and one further vertex if $j$ is odd, to obtain a set $U$ of $j$ vertices. These choices reduce the sum of the half-order capacities by at least $\floor{j/2}$. Color the edges incident with $S\cup U$ blue and all others red. The blue graph has a vertex cover of size $q-1$, so the red graph contains one complete copy of $F$. A star with $n_i$ edges costs at least $\floor{(n_i+1)/2}$ units of the half-order capacity of its original odd component. Thus
\[
 s+\omega(F)\le\nu(Z)-b-\floor{j/2}.
\]
It follows that $u\ge\omega(F)+\floor{(q-1+b)/2}\ge\omega(F)+\floor{(q-1)/2}$. Substitution of $u=q-1-d$ gives $q\ge b+2(d+\omega(F))$. If a specified root set $A_0$ gives an actual decomposition
$Z-A_0=\mathcal R\sqcup B$, where the components of $\mathcal R$ are odd
and factor-critical and
$\nu(Z)=|A_0|+\nu(\mathcal R)+\nu(B)$, choose an optimal Tutte barrier
$S_B$ in $B$. Then
\[
 o\bigl(Z-(A_0\cup S_B)\bigr)-|A_0\cup S_B|
 =|\mathcal R|-2\nu(\mathcal R)-|A_0|+|B|-2\nu(B)
 =|Z|-2\nu(Z).
\]
Thus $A_0\cup S_B$ is an optimal barrier of $Z$. The size $b$ of a
maximum-cardinality optimal barrier is at least $|A_0|$; no containment
property of that maximum-cardinality barrier is required.
\end{proof}

\subsection{Signed exposure under an odd degree cap}
\begin{lemma}\label{lem:signed-exposure}
Suppose $D=a+b-2$ is odd, $a\ge2$ is even, $b\ge3$ is odd, $\Delta(X)\le D$, and $X\to(pS_a,qS_b)$. For $L=p+q-1$,
\[
 \rho_D(X)\ge L+(a-2)\ceil{L/(b-1)}.
\]
\end{lemma}
\begin{proof}
Choose a maximum matching with minimum full-degree exposure. Its exposed full-degree set $U$ is independent and has size $\rho$. The following regular completion is used only to color the original graph, never as original edge income. Let $V_0$ be the nonfull vertices exposed by the chosen matching.
If $|V_0|$ is odd, add an isolated vertex $z_0$ and include it in $V_0$.
Pair the vertices of $V_0$.  For each pair $x,y$, take a fresh copy of
$K_{D+1}-ab$, add $xa,yb$, and extend the matching by these two edges
and a perfect matching of the remaining $K_{D-1}$.
All vertices of $V_0$, including $z_0$ when present, are now matched;
the exposed set is exactly $U$.
Since $D$ is odd, the number of remaining degree-deficit stubs has
the parity of the current vertex number, hence the parity of $\rho$.
If $\rho$ is odd, add another isolated vertex $z_1$, leaving it exposed.
The deficit stubs can now be paired.  For each pair attach a fresh
$K_{D+1}-ab$ through its two deficient ports, and match this fresh
gadget internally by a perfect matching, leaving all previously
chosen matching edges unchanged.  Even if two stubs belong to the
same vertex, the fresh ports are different, so the graph remains simple.
The resulting graph is $D$-regular and has a specified matching
exposing exactly $U$, together with $z_1$ precisely when $\rho$ is odd.
No edge was added at any vertex of $U$.

For each $x_i\in U$ choose an original edge $e_i=x_iw_i$. These edges lie outside the matching; the vertices $w_i$ need not be distinct. For odd exposure choose one further edge at the new exposed vertex. Remove the specified matching and all these edges. The resulting degrees are $D-1$, except for one labeled missing stub at each $w_i$. Pair these stubs by auxiliary edges, allowing loops and parallel edges. The resulting even-regular multigraph splits into $(D-1)/2$ two-factors: orient Euler tours and apply a perfect-matching decomposition to the resulting regular bipartite graph.

Let $c_j$ count the original labeled stubs in factor $j$. Their sum is $\rho$; they are all even if $\rho$ is even, and exactly one is odd otherwise. Color any $(b-1)/2$ factors blue and all remaining factors, the specified matching, and the as-yet-unassigned $e_i$ red; then discard all auxiliary vertices and edges. If the selected factors contain at most $L-1$ labels, assign at most $p-1$ of these labels red and at most $q-1$ blue, recoloring the corresponding $e_i$ blue in the latter case. Red $a$-star centers can occur only at red-labeled endpoints, and blue $b$-star centers only at the blue-labeled $x_i$. This is an avoiding coloring, a contradiction. Every set of $(b-1)/2$ factors therefore contains at least $L$ original labels.

Write $\ell=(b-1)/2$. In the even case, order $c_j=2y_j$. The first $\ell$ terms have sum at least $L/2$, so $y_\ell\ge\ceil{L/(b-1)}$; each of the remaining $(a-2)/2$ terms is at least this large. This proves the bound. In the odd case write the unique odd term as $2y_0+1$. The smallest $\ell$ even terms give the same threshold, while the odd term together with the smallest $\ell-1$ even terms has total at least $L$. Each remaining $(a-2)/2$ even term again meets the threshold. If $a=2$, all factors together already have sum at least $L$. The conclusion follows in every case.
\end{proof}


\section{Unified prices on one residual graph}\label{app:prices}
Fix the rank $N_0$ of the original minimal counterexample. Every complete source of rank $0<N<N_0$ satisfies $e(X)\ge E_r$ by minimality. No structural property of the original minimum source is imposed on $X$.

\subsection{Whole-core prices with the entire highest-column tax}
\begin{proposition}\label{prop:general-price}
Suppose $P_r(X)$, $r\ge2$, $0<N<N_0$, $\Delta(X)\le q$, $\ell_r\le q+1$, and $q\ge3$. Put $h=L_{q+1}$, $\rho=\rho_q(X)$, $m=\floor{q/2}$, $c=\ceil{q/2}$, and
\[
\begin{split}
 C_q&=2c(m-1)+mq(c-2)+\mathbf1_{2\mid q}\,m\max(c-2,0),\\
 T_q(\rho)&=\max\left(0,\min\left\{\ceil{q\rho/2}-mc,
 \floor{\frac{(cq+m)\rho-C_q-1}{2c+m}}\right\}\right).
\end{split}
\]
Then, regardless of the sign of $N-\nu(X)$,
\begin{equation}\label{eq:general-price}
 e(X)-E_r\ge\max(0,T_q(\rho)-h).
\end{equation}
\end{proposition}
\begin{proof}
It suffices to prove the second bound when $T:=T_q(\rho)>0$. Use Lemma \ref{lem:intrinsic-region}. Positive cores have $1\le w\le m$, $z\ge c$, and total positive weight at least $\ceil{q\rho/2}$. Order them by nondecreasing $z$ and stop when the accumulated weight first reaches $T$. The selected union has weight $w\in[T,T+m-1]$, at most $l\le T$ cores, largest capacity $Z$, and total capacity $z$.

There remain at least
\[
 u\ge\ceil{\frac{\ceil{q\rho/2}-T-m+1}{m}}\ge c
\]
positive cores, the last inequality following from $T\le\ceil{q\rho/2}-mc$. Each has capacity at least $Z$. The unselected region contains at least $n_0\ge2uZ+\rho-T$ full-degree vertices: sum the orders of all remaining cores and subtract at most all roots. Zero and negative cores are retained in this count.

If $q$ is odd, or if $q$ is even and $Z>c$, the induced outside graph has no component $K_{q+1}$. For odd $q$ this clique has even order and cannot be an original full-degree odd core. In the even case its capacity would be $c<Z$, contradicting the ordering. Brooks' theorem \cite{brooks} gives an independent set of size at least $n_0/q$; odd cycles do not affect this bound for $q\ge3$. Substitution yields
\[
 n_0-q(Z-2)\ge
 \frac{(cq+m)\rho-(2c+m)T-[2c(m-1)+mq(c-2)]}{m}>0.
\]
Thus $Z-1$ independent full-degree vertices are available. If $q$ is even and $Z=c$, greedy $(q+1)$-coloring and
\[
 n_0-(q+1)(c-2)\ge
 \frac{(cq+m)\rho-(2c+m)T-C_q}{m}>0
\]
give the same conclusion.

Choose $b=(Z-r+1)_+\le Z-1$ such vertices. Erase the selected core edges and delete these outside independent full-degree vertices. The unique actual cost is $q(z+b)+w$. The whole-core subset-sum menu gives the complete source $P_{r+z+b}$; its proof uses the maximum capacity $Z$, not equality of all core capacities. Put $K=z+b$. If $K<N$, the new rank is positive and smaller than $N_0$, so
\[
e(X)\ge E_{r+K}+qK+w\ge E_r+w-h\ge E_r+T-h.
\]
Every column below the highest loses at most $K$, and the highest loses at most $h$. If $K\ge N$, no zero-rank source is invoked: the actual edge cost itself gives $e(X)\ge qK+w\ge qN+T\ge E_r+T-h$. Taking the maximum with $e(X)\ge E_r$ proves the assertion.
\end{proof}

Direct substitution gives the following forms, with negative values replaced by zero:
\[
\begin{array}{c|l}
q&T_q(\rho)\\\hline
3&\min(\ceil{3\rho/2}-2,\floor{(7\rho-1)/5})\\
4&\min(2\rho-4,\floor{(10\rho-5)/6})\\
5&\min(\ceil{5\rho/2}-6,\floor{(17\rho-17)/8})\\
6&\min(3\rho-9,\floor{(7\rho-12)/3})\\
7&\min(\ceil{7\rho/2}-12,\floor{(31\rho-59)/11})\\
8&\min(4\rho-16,3\rho-9).
\end{array}
\]
Since $h\le\rho$, the case $q=5$ implies the safe all-sign bound
$e(X)-E_r\ge\max(0,\rho_5(X)-4)$.

\subsection{Exposure at arbitrary nonnegative surplus}
\begin{lemma}\label{lem:surplus-price}
In the positive smaller-rank setting of Proposition \ref{prop:general-price}, put $v=e(X)-E_r\ge0$. Then
\[
\begin{array}{ll}
q=2p-1,\ p\ge4:&\rho_q(X)\le v+p,\\
q=2s,\ s\ge3:&\rho_q(X)\le v+s+1.
\end{array}
\]
\end{lemma}
\begin{proof}
If $\nu(X)\ge N$, the canonical reserve gives $\rho\le v$. Suppose therefore that matching is deficient. For the odd cap, assume $\rho\ge v+p+1$ and write $\rho=p+1+j$. Select a positive-core prefix whose weight first reaches $T=h+v+1\le2\rho-p$; its actual weight satisfies $w\le2\rho-2$. Total positive weight is at least $(p-\tfrac12)\rho$, so the number of remaining positive cores satisfies
\[
u\ge\frac{(p-\tfrac52)\rho+2}{p-1}>p-1+j/2.
\]
The difference has numerator $p-3+(p-4)j>0$, whence $u\ge p+\floor{j/2}$. For largest selected capacity $Z$, there are at least $2uZ-\rho+p$ outside full-degree vertices, and
\[
n_0-2p(Z-1)\ge2\floor{j/2}Z-j+2p-1>0.
\]
Greedy $2p$-coloring supplies the required $b\le Z-1$ independent full-degree vertices. The actual deletion parameter satisfies $K=z+b\le\nu(X)<N$, while $w\ge h+v+1$. The complete successor therefore has surplus at most $v+h-w\le-1$, a contradiction.

For an even cap assume $\rho\ge v+s+2$ and write $\rho=s+2+j$. Select $T=h+v+1\le2\rho-s-1$, again with $w\le2\rho-2$. Then
\[
u\ge\frac{(s-2)\rho+2}{s}>s-1+j/3,
\]
so $u\ge s+\floor{j/3}$ and $n_0\ge2uZ-j-1\ge2sZ-3$. If $Z=s$, use $2s^2-3>(2s+1)(s-1)$ and greedy coloring. If $Z>s$, there is no outside $K_{2s+1}$, and Brooks' $2s$-coloring applies. Both supply $Z-1$ independent vertices and give the same successor contradiction.
\end{proof}

\subsection{Actual point-cover taxes in the common envelope}
\begin{lemma}\label{lem:point-cover-tax}
After root deletion, suppose $C$ is an isolated core union in $G_0$, with $|C|=n$, $\nu(C)=Z$, and $P_{r_0}(G_0)$. Put $R_0=r_0-1$, $Y=G_0-E(C)$, and $M_h=\min(Z,\floor{n/(h+1)})$. For original heights $p,q\ge2$, the tax on the candidate width $f_p+g_q-R_0$ in the uniform coordinate of this same $Y$ is at most
\[
J(p,q)=\max\left\{M_p,M_q,
\floor{\frac{n(p+q)-R_0pq}{pq+p+q}}\right\}.
\]
For a unit candidate $(1,q)$, the tax is at most $Z$ provided $s>Z$; for $(p,1)$ it is at most $Z$ provided $t>Z$. In particular the bound holds for every unit candidate if $\min(s,t)>Z$.
\end{lemma}
\begin{proof}
Put $L=f_p+g_q-R_0$. If $f_p>M_p$, the tax bound is trivial for $L\le M_p$. Otherwise the integer interval
\[
\max(1,r_0-g_q)\le i\le\min(f_p-M_p,r_0-1)
\]
is nonempty. Set $j=r_0-i$. This is an original legal pair, and both $f_p-i+1-M_p$ and $g_q-j+1$ are positive. Restore all internal edges of $C$ in red. At most $M_p$ red stars of head at least $p$ are lost from one complete copy, while a blue copy is unchanged. The uniform width on the same $Y$ is therefore at least $L-M_p$. Interchange colors when $g_q>M_q$ to obtain tax at most $M_q$. Otherwise $f_p\le M_p$ and $g_q\le M_q$. If $n-pf_p<L$, choose
\[
v=n-(p+1)f_p+1,\qquad1\le v\le g_q-R_0,
\]
actual vertices of $C$, color internal edges incident with them blue, and all other internal edges red. The trigger gives $g_q\ge R_0+1$, so the original pair $(1,r_0-1)$ is legal. If $Y$ has no red $S_p$, all $f_p$ stars of height at least $p$ in the original complete $F_1$ would have to occupy the remaining $n-v$ core vertices, which is impossible. This pair therefore forces a complete blue suffix. At most $v$ blue stars are damaged by the cover, giving a genuine coordinate at least $L-n+pf_p$. Reversing colors gives $L-n+qg_q$; negative bounds are replaced by zero.

If the tax exceeded the integer $J=J(p,q)$, then $L>J$ and both $n-pf_p\ge J+1$ and $n-qg_q\ge J+1$. Hence
\[
L\le(n-J-1)(1/p+1/q)-R_0<J+1.
\]
The last inequality is exactly the assertion $J\ge\floor{[n(p+q)-R_0pq]/(pq+p+q)}$. Integrality gives $L\le J$, a contradiction. For a unit candidate $(1,q)$ put $L=s+g_q-r_0+1$. If $L\le Z$, the required lower bound is zero. Otherwise assume $L>Z,s>Z$ and fix the nonempty list $\nn'=\nn[Z+1:s]$. Its rank satisfies
\[
N_{r_0}-Z=s+t-r_0+1-Z\ge L-Z>0.
\]
Every new legal pair $(i,j)$, with $i+j=r_0$, $1\le i\le s-Z$ and $1\le j\le t$, is also an original legal pair. Extend an arbitrary coloring of $Y$ by restoring $E(C)$ in red. The original source supplies one complete red $F_i$ or blue $H_j$. A blue copy remains intact in $Y$. In a red copy choose one restored edge from each affected component. These edges form a matching in $C$, so at most $Z$ components are affected. Remove them and trim in the original order to obtain the complete $F_{i+Z}$. Hence on this one $Y$,
\[
P_{r_0}\bigl(Y;\nn[Z+1:s],\mm\bigr).
\]
The unit candidate in this fixed source has width $(s-Z)+g_q-r_0+1=L-Z$, proving the tax bound. Interchanging colors treats $(p,1)$ under $t>Z$. Every restoration trims one complete copy; copies from different restorations are never combined.
\end{proof}

\begin{remark}\label{rem:unit-tax-scope}
The unit tax uses the current list lengths and the matching capacity of the current actual piece. In an original normalized source, Proposition \ref{prop:base} gives $\min(s,t)>\mu$, so every ordinary-core subunion with $Z\le\mu$ satisfies the condition; root deletion shifts the row without changing these lists. For separately selected regular components or ported pieces, use the strong lengths $s,t>Z$ proved locally before the tax is invoked. The nonunit tax formula needs no such additional condition.
\end{remark}
\begin{lemma}\label{lem:allowance-tax}
Let $\mathbf n,\mathbf m,\mathbf R,\mathbf S,\mathbf T$ be nonempty
nonincreasing lists of positive integers, and let
$2\le r\le\min(|\mathbf R|+|\mathbf S|,|\mathbf n|+|\mathbf T|)$;
thus the row is legal for both displayed source data.
Suppose the same graph $Y$ satisfies
\[
 P_r(Y;\mathbf R,\mathbf S),\qquad
 P_r(Y;\mathbf n,\mathbf T),\qquad \mathbf R\preceq\mathbf n.
\]
Write $f_j,g_j$ for the cumulative counts of $\mathbf n,\mathbf m$,
and $R_j,S_j,T_j$ for those of the other three lists.
Fix a reference candidate $p,q\ge1$ with $f_p,g_q>0$, and set
\[
 x=f_p-R_p,\qquad y=g_q-S_q,\qquad z=g_q-T_q,
 \qquad R=r-1\ge1.
\]
Then
\[
 \min_{1\le\ell\le R}
 \{\max(x,\ell)+\max(y,z-\ell)\}
 =M:=\max(x+y,z,x+z-R,y+1).
\]
If $S_q+T_q>0$, then
\[
 u_{p+q-1}(Y)\ge(f_p+g_q-R-M)_+.
\]
If $S_q=T_q=0$, this candidate is omitted.
Indices with $f_p=0$ or $g_q=0$ are not reference candidates and are
omitted as well.
Only $x$ is necessarily nonnegative; the identity and bound also allow
negative $y$ or $z$. When $\mathbf S,\mathbf T\preceq\mathbf m$,
all three differences are nonnegative.
All coordinate maxima concern this same graph $Y$; whole-source prices
are not added. In the canonical reserve, only heights $p+q-1\ge2$
are summed; height one is handled by the actual matching number.
\end{lemma}
\begin{proof}
For an integer allowance $1\le\ell\le R$, the expression minimized is
$\max(x+y,z,x+z-\ell,y+\ell)$.
Each of the four terms defining $M$ is a lower bound for its minimum.
Conversely, the integer interval
\[
 \max(1,x+z-M)\le\ell\le\min(R,M-y)
\]
is nonempty: the endpoint conditions follow from
$M\ge x+z-R,y+1$, and the internal comparison follows from
$2M\ge x+y+z$.

Assume first that $S_q>0$. For a minimizing allowance, the fixed
intersection lists of Lemma~\ref{lem:intersection} satisfy
\[
 K_p=f_p-(x-\ell)_+,
 \qquad L_q\ge g_q-\max(y,z-\ell).
\]
Hence their candidate width at row $r+\ell$ is at least
\[
 f_p+g_q-R-\max(x,\ell)-\max(y,z-\ell).
\]
If this lower bound is positive, then $K_p>0$ because $f_p>0$ and
$\ell\ge1$, while $L_q>0$ because $S_q>0$ and $R\ge1$.
The fixed lists are therefore nonempty and their rank is positive.
Lemma~\ref{lem:intersection} applies. Choose a legal pair $i+j=r+\ell$
with $i\le K_p$ and $j\le L_q$; such a pair exists because
$K_p,L_q>0$ and $K_p+L_q\ge r+\ell$.
Choose this pair independently of the coloring. Its complete arrow
relation weakens to
\[
 Y\longrightarrow
 \bigl((K_p-i+1)K_{1,p},(L_q-j+1)K_{1,q}\bigr).
\]
Both multiplicities are positive, and their sum minus one is
$K_p+L_q-(r+\ell)+1$. This is a source included in the definition of
$u_{p+q-1}(Y)$.
If the lower bound is nonpositive, use $u_{p+q-1}(Y)\ge0$.

If $S_q=0<T_q$, then $M\ge z$ gives
$f_p+g_q-R-M\le f_p+T_q-R$.
When the latter is positive, choose a fixed legal pair $i+j=r$ with
$i\le f_p$ and $j\le T_q$ in the second given source.
It yields the arrow to $(f_p-i+1)K_{1,p}$ versus $(T_q-j+1)K_{1,q}$,
of width $f_p+T_q-R$.
Otherwise the required bound is zero.
The case $S_q=T_q=0$ is omitted, as stated.
\end{proof}

In one maximum-degree deletion history, let $C_j$ count vertices removed while their current degree is at least $j$. Then
\begin{equation}\label{eq:inventory-envelope}
C_j\ge u_{j+1}(Y).
\end{equation}
Otherwise, after these $C_j$ deletions some uniform source still forces a positive star pair with head sum $j+2$, whereas the remaining maximum degree is at most $j-1$ and the degree-splitting lemma supplies an avoiding coloring. All coordinate inventories can therefore be used in the same actual deletion history.


\section{Emptiness, complete residuals, and two actual successors}\label{app:sources}
Use one ordinary region $A,\mathcal R,B$, with $a=|A|$, $\mu=\nu(\mathcal R)$, $k=\nu(B)$, $a+\rho$ cores, and $Z=G-E(\mathcal R)$. Here $Z$ denotes a graph; later core-selection arguments redeclare $Z$ as a capacity.

\subsection{Emptiness and an actual embedding of the whole forest}
\begin{lemma}\label{lem:empty-embedding}
Suppose the complete source has deficit $d=N-\nu(G)>0$. For either full original list $Q$, let $l$ be its length, $x(Q)=\sum\floor{(Q_i-1)/2}$, and $q_+$ the number of nontrivial cores. Exactly one of the following alternatives holds:
\[
 \mu\le l-d-1,\qquad\text{or}\qquad\mu\ge l+d+2x(Q)-1.
\]
In the second alternative, $\mathcal R-W$ contains the whole $Q$ for every $|W|\le\mu-l+d$, and
\[
 e(Q)+d\le\mu+q_+.
\]
Thus $l\le\mu$ is equivalent to an embedding of the whole $Q$ in $\mathcal R$, and to feasibility of all its entries in the half-order and vertex budgets. If $t>\mu$, then
\[
 P_r(Z;\nn,\mm[\mu+1:t]),\quad N'=N-\mu=a+k+d,
 \quad\nu(Z)=a+k.
\]
The color-reversed assertions also hold.
\end{lemma}
\begin{proof}
The actual root assignment gives $\nu(G)=a+\mu+k$ and $\nu(Z)=a+k$. Deleting all roots preserves positive deficit $d$, so the complete-source endpoints give $r+a\le\min(s,t)$.

Take $Q=\nn$. If $b=\mu-s+d\ge0$, then $\mathcal R\to(F_1,(b+1)K_2)$. Otherwise color all of $Z$ blue; the entire blue matching number is at most $a+k+b=t-r+1$, while no red $F_1$ occurs, contradicting the original corner $(1,r-1)$. Coloring the core edges incident with $W$ blue proves robust embedding.

Delete an even number of vertices in each original odd core, totaling $2\floor{b/2}$. Every star of head $n_i$ consumes at least $\floor{(n_i+1)/2}$ half-order units, so
\[
s+x(\nn)\le\mu-\floor{(\mu-s+d)/2}.
\]
Integer rearrangement gives $\mu\ge s+d+2x(\nn)-1$. If the trigger fails, $\mu\le s-d-1$. Deleting any $b$ vertices in the nontrivial cores still leaves the whole forest on $s+e(\nn)$ vertices. Thus $s+e(\nn)\le2\mu+q_+-b$, proving the mass bound. Interchange colors for the other side.

An embedding implies $l\le\mu$; the converse follows from robust embedding with $W=\varnothing$. An actual embedding satisfies both budgets. Conversely the half-order cost is at least $l$, so numerical feasibility implies $l\le\mu$.

Finally restore $E(\mathcal R)$ entirely blue. At most $\mu$ components of a complete blue suffix use erased edges, since one such edge from each gives a matching in the cores. Discard them and trim in the original order to the suffix shifted by $\mu$. A red suffix is already in $Z$. This works at every new legal corner, giving the fixed complete successor and the stated rank and matching identities.
\end{proof}

\subsection{One residual list controlling both actual graphs}
For a positive list $Q$, let $T_{u,v}(Q)_i$ be the last entry of the first window starting at $i$ that violates $\sum\ceil{Q_j/2}\le u$ or $\sum(Q_j+1)\le v$; output nothing if no violation occurs. Put
\[
r_h=\sum_C\floor{|C|/(h+1)},\quad
c_h(U(Q))=\max_{j\ge h}(c_j(Q)-r_j)_+,
\quad L(Q)=T_{\mu,|\mathcal R^+|}(Q)\vee U(Q).
\]
Here $\vee$ is the componentwise maximum after zero padding. The operation $C_a(Q)$ changes the largest at most $a$ entries to $1$ and reorders them.

\begin{lemma}\label{lem:two-successors}
If $s>\mu$, set $L=L(\nn)$, $e_0=s-|L|$, $\delta=\mu-e_0$, and $J=C_a(L)\vee\nn[\mu+1:s]$. Then $1\le e_0\le\mu$ and the actual complete sources are
\[
 P_{r+a}(B;L,\mm),\qquad P_r(Z;J,\mm).
\]
Their ranks are $k+d+\delta$ and $a+k+d+\delta$, both in $[1,N-1]$, and $r+a\le\min(|L|,t)$. All original blue heads, units, and retained root edges are preserved.

More generally, if a fixed internal coloring has red matching number at most $K$ and half-order and vertex budgets $u,v$ for its wholly internal blue star forests, then
\[
 P_{r+a}(B;\nn[K+1:s],T_{u,v}(\mm)),\quad
 P_r(Z;\nn[K+1:s],C_a(T_{u,v}(\mm))).
\]
Each assertion is used only when both its lists are nonempty and its rank is positive.
\end{lemma}
\begin{proof}
From one actual copy of $Q$, delete the components wholly contained in $\mathcal R$. Their half-order cost is at most $\mu$, their vertex cost at most $|\mathcal R^+|$, and their count at height $h$ at most $r_h$. If entry $i$ of the surviving list were below the last entry of the first violating window, all the cheapest entries of that window would have been deleted, violating a budget. The same surviving list therefore dominates $T$. The height counts simultaneously make it dominate $U$, and hence $L$. These operations commute with taking suffixes.

At least one core has order greater than $D$: a full-degree core of order at most $D$ has at least $D$ root-cut edges, so there are at most $a$ such cores, whereas there are $a+\rho>a$ cores in total. The larger core allows any single original head to fit both numerical budgets and gives $r_h\ge1$ at every original height. Hence $e_0\ge1$. Also $L\ge\nn[\mu+1:s]$, so $e_0\le\mu$. Lemma \ref{lem:empty-embedding} supplies the emptiness criterion.

Given any coloring of $B$ and output corner $p+q=r+a$, mark $u=\max(0,a-q+1)$ roots red and the others blue. Give red priority to edges between unlike roots and color the cores red. The original corner $(p-u,q-a+u)$ is legal. From a red copy discard its at most $u$ components meeting red roots and trim to $F_p$. Its remaining components lie wholly in the cores or in $B$; deleting the former leaves $L_p$ in $B$. A blue copy, after removing components meeting blue roots, leaves the entire $H_q$ in $B$. This proves the first source.

Given any coloring of $Z$, restore the core edges red. From a red copy first discard wholly internal components, then shorten each remaining star using a core edge to one actual root edge. There are at most $a$ such stars. Length is preserved, and the same resulting list dominates $C_a(L)$. At most $\mu$ original components used any core edge, so it also dominates the original suffix shifted by $\mu$, and hence their componentwise maximum $J$. A blue copy is already wholly retained in $Z$.

The rank identities follow from $e_0=\mu-\delta$. If $r+a>|L|$ or $r+a>t$, a monochromatic endpoint forces a matching in $B$ larger than $k$, a contradiction. For a general fixed internal coloring use the identical root assignment: at most $K$ red components are damaged, blue damage obeys the two stated budgets, and at most $a$ cross-root stars are shortened on the edge-erased graph. Copies arising in different responses are never combined.
\end{proof}

\subsection{Two bounds from the same edge account}
Let $\zeta,\varepsilon,\beta$ have their meanings in the text, and let the original highest-column width be $h$. The actual fees are
\[
c_Z=D\mu+\frac{D\rho+\zeta+2\varepsilon+\beta}{2},\qquad
c_B=D(a+\mu)+\frac{D\rho+\beta-\zeta}{2}.
\]
Write $\lambda_j^B,\lambda_j^Z$ for the two basic source columns, $L_j$ for the original columns, and put
\[
B_j=\mu+\frac a2+\frac{\lambda_j^B+\lambda_j^Z}{2}-L_j,
\qquad C_0=\frac{D\rho+\beta+\varepsilon}{2}-h.
\]
On each graph separately, add any further actual fixed-palette sources and maximize at each height. Let $G_j\ge0$ be the average increase. The canonical envelope and the exact edge fees give, respectively,
\begin{align}
e(G)-E_r&\ge C_0-d+\sum_{j=2}^D(B_j+G_j),\label{eq:two-ledger-potential}\\
\frac{c_B+E_B+c_Z+E_Z}{2}-E_r
&=C_0+\delta+\sum_{j=2}^DB_j+\frac{\lambda_{D+1}^B+\lambda_{D+1}^Z}{2}.
\label{eq:two-ledger-area}
\end{align}
The first layers in the first formula are actual matching numbers $k,a+k$; in the second they are the new ranks $N-a-e_0,N-e_0$. Substitution and columnwise subtraction prove both identities, so their terms $-d$ and $+\delta$ cannot be interchanged. If the second expression is nonnegative, choose the actual graph with larger $c_X+E_X$ to obtain a paid strict descent.

\subsection{The embedded side has a head equal to the degree cap}
\begin{lemma}\label{lem:reverse-owncap}
An original minimal strict source with $n_1=D$ and $s\le\mu<t$ is impossible.
\end{lemma}
\begin{proof}
Put $P=f_D\ge1$, $x=x(\nn)$, and $\omega=(D\rho+\zeta+2\varepsilon+\beta)/2$. The reserve gives $h\le\rho\le d-1$, hence $d+2x+2P\ge h$. Lemma \ref{lem:empty-embedding} gives $\mu\ge s+d+2x-1$, while the blue shifted source gives $t\ge\mu+r$. The top rectangle $(D,2)$ gives $P+g_2-r+1\le h$, so $t+P\ge s+g_2$. Consequently the maximum original column at every $1\le j\le D$ is $t+f_j-r+1$: all nonunit-blue candidates are at most $s+g_2-r+1$, and the remaining maximum is attained by this unit axis. Therefore
\[
E_r=D(t-r+1)+e(\nn)+h.
\]
The same edge-erased graph $Z$ arrows the complete $F_1$ versus $(t-\mu-r+2)K_2$. The endpoint edge bound gives $e(Z)\ge D(t-\mu-r+1)+e(\nn)$. Add the unique actual core fee $D\mu+\omega$ to obtain $e(G)-E_r\ge\omega-h>0$, since $D\ge3$ and $\rho\ge h\ge1$. This is a contradiction.
\end{proof}

\subsection{Clearing the exterior of a persistent region}
\begin{lemma}\label{lem:persistent-region}
Suppose $P_r(G)$, $n_1=D=\ell_r-1\ge3$, $\Delta(G)\le D$, and there is a full-degree factor-critical core region with $k$ actually assignable roots and an exterior. Let its matching capacity be $\mu$, put $M=\mu+k$, and let the number of cores minus roots be $\eta\ge1$. If $M\le t-r+1$, then
\[
 e(G)\ge E_r+\eta-h,\qquad h=L_{D+1}.
\]
\end{lemma}
\begin{proof}
Put $q_0=t-r+1-M$. Delete current degree-$D$ vertices only in the exterior. After $q\le q_0$ deletions color the exterior red and everything else blue. The blue matching number is at most $M$, whereas the blue suffix at the current corner $(1,r+q-1)$ has at least $M+1$ components. The entire red $F_1$ must therefore lie in the exterior, providing a next degree-$D$ center. At $q=q_0$ the exterior still contains the complete $F_1$, so its matching number is at least $s$ and the whole matching number at least $M+s=N-q_0$. Throughout, the full-degree core region and root assignment persist; the final full-degree exposure is at least $\eta$. The matching-rich canonical reserve gives $e(G')\ge E_{r+q_0}+\eta$. The actual deletion fee is $Dq_0$, and the original profile loss is at most $Dq_0+h$. Adding them proves the bound.
\end{proof}

\section{Complete sources of maximum degree at most two}\label{app:binary}
\begin{lemma}\label{lem:binary-cap}A graph of maximum degree at most two cannot be an all-data minimal strict complete source.
\end{lemma}\begin{proof}
First consider an ordinary binary pair. Write
\[
 F=pS_2\sq aK_2,\quad H=qS_2\sq bK_2,\quad p,q\ge1,\quad a,b\ge0,
\]
\[
 R=p+q-1,\quad M=\max(a,b),\quad m=\min(a,b),\quad
 L=3R+2M+m,\quad N=R+M+m.
\]
Suppose $X\to(F,H)$, $\Delta(X)\le2$, and $e(X)<L$. Remove isolated vertices. The graph is a union of paths and cycles. If $\nu(X)\le A+B-1$ with $A,B\ge1$, color entire components red in order until the next component would exceed red matching capacity $A$. In that component, use $2u-1$ consecutive red edges to fill the remaining red capacity $u$, and color the rest blue. Its blue matching capacity is at most its original capacity minus $u$ plus one, so the total blue capacity is at most $B$. Apply this with $A=p+a-1$, $B=q+b-1$ to obtain $\nu(X)\ge N-1$; a zero-budget endpoint is handled by a monochromatic coloring. The potential inequality in the text now forces
\[
 e(X)=L-1,\qquad \nu(X)=N-1,\qquad e(X)-\nu(X)=2R+M.
\]
Let $c$ be the number of odd cycles and $h=\kappa_1(X)$. On each odd cycle reserve one repair edge and alternate the remaining edges red and blue. If $c\le p+q-2$, assign the repair edges so that at most $p-1$ red and $q-1$ blue $S_2$ components can occur. Hence $c\ge R$. Weakening one side to units and retaining only the two-heads on the other gives $h\ge R+M$ through the uniform envelope.

For a path $P_v$, $\kappa_1=\floor{v/3}$; for a cycle $C_v$, $\kappa_1=\ceil{v/3}$. Componentwise comparison gives $e-\nu\ge h+c$. The preceding equality forces equality throughout, leaving exactly the components $P_2,P_3,P_4,P_6,C_3,C_4,C_5,C_7$. If $x_i,z_i$ count these paths and cycles, then
\begin{align*}
 z_3+z_5+z_7&=R,\\
 M&=z_5+2z_7+2z_4+x_3+x_4+2x_6,\\
 m&=1+x_2+x_4+x_6.
\end{align*}
Paths and even cycles permit arbitrary division of their total matching capacity between the two colors without increasing it. If their total capacity is at least two, the monochromatic odd-cycle capacities $1,2,3$ have prefix gaps at most two; the divisible capacity fills such a gap. Red matching capacity can therefore be made $p+a-1$ and blue at most $q+b-1$, a contradiction. Thus paths and even cycles have total capacity at most one, and the number of paths satisfies $c_P\le1\le m$.

Use the following direct cover bound: if a host has a spanning path--cycle cover with $c_P$ paths, then
\[
 e(X)\ge |F|+|H|-c_P-3.
\]
Take $|F|-1$ consecutive vertices in the cover order as $U$, color its internal edges red, and all other edges blue. Red has too few vertices for $F$. At most two cover edges cross $U$, and there are $e(X)-|X|+c_P$ edges outside the cover, so the blue graph has at most $e(X)-|U|+c_P+2$ nonisolated vertices. If the displayed bound failed, blue would also have too few vertices for $H$. Here $|F|+|H|-3=L+m$ and $c_P\le m$, giving $e(X)\ge L$, the required contradiction.

Return to a complete source. If one list consists only of units, the full forest--matching endpoint price in the text applies directly. Otherwise, under the normalized cap $D\le2$, both original lists contain only $1$ and $2$. The case $D=1$ is impossible: monochromatic endpoints force both original heads to be at most one, whereas $\ell_r=2$. If $D=2$, take an original corner $(i,j)$ attaining $\ell_r=3$, with both heads two. For original two-head counts $p,q$ and unit counts $a,b$, its ordinary suffix pair has top-platform length $p+q-r+1$, middle-platform length $M$, and bottom-platform length $m$. Its ordinary profile area is exactly the original $E_r$. The ordinary binary inequality just proved excludes it. Degree zero leaves only unit targets and is handled directly by partitioning the edges between the colors.
\end{proof}
\section{Single-empty exclusion below both original head caps}\label{app:single-below}
\begin{proposition}\label{prop:single-below}
An original all-data minimal strict source cannot satisfy $D\ge6$, $n_1,m_1\le D-1$, and $t\le\mu<s$.
\end{proposition}
\begin{proof}
Use the two successors and their accounts from Appendix \ref{app:sources}. Fix the following notation throughout this proof; $X,Y$ are integers, not graphs:
\begin{gather*}
\alpha=n_1,\ b=m_1,\quad X=\sum_i\floor{(n_i-1)/2},\
Y=\sum_j\floor{(m_j-1)/2},\\
S=s-\mu,\quad T=\mu-t,\quad u=s-f_2,\quad
v_2=c_2(L),\quad V=\mu-f_2,\quad A_*=V+v_2,\\
V_0=\sum_C\floor{(z_C+1)/3},\quad |L|=S+\delta,\quad u=S+V.
\end{gather*}
The common constraints used in the column computations are listed next. Here $q_+$ counts nontrivial cores, $h=L_{D+1}$, and $C_0,B_j,G_j$ are defined in \eqref{eq:two-ledger-potential}:
\begin{align}
C_0&\ge(D/2-1)\rho,\quad T\ge2Y+d-1,\quad t+g_2+Y\le\mu+q_+-d,\label{eq:se-guards1}\\
2X&\le t-a-2d-r+2,\quad S\ge a+r+2d+\delta+2Y-2,\label{eq:se-guards2}\\
\mu&\ge3\rho,\quad V_0\ge\rho,\quad
q_+\le(\mu-\rho)/2,\quad q_+\le a/2+V_0,\label{eq:se-guards3}\\
8V_0+a-2\mu&\ge0,\quad2V_0\le\mu,\quad
3\delta\ge\mu-u-a-\rho,\quad f_3\le S+2\delta,\quad A_*\ge V_0.\label{eq:se-guards4}
\end{align}
The first two lines follow from robust embedding on the empty side and the complete endpoint lemma. The root term $a$ is retained because $A$ together with an optimal barrier of $B$ is an optimal barrier of $G$. A maximum-cardinality optimal barrier therefore has at least $a$ vertices; deleting it preserves the full source and its matching deficit, so Lemma \ref{lem:strong-endpoint} applies. Apply the same lemma to the actual source on $B$ to obtain the bound on $S$.

A core of order at most $D$ has cut at least $D$. Thus at least $\rho$ cores have capacity at least three, proving $\mu\ge3\rho,V_0\ge\rho$. For a nontrivial core,
\[
6z_C+\partial_C\ge3D,\qquad
2Dq_+\le6\mu-D\rho-\zeta-2\varepsilon-\beta.
\]
For $2z_C\le D$, simplicity gives $\partial_C\ge(2z_C+1)(D-2z_C)$ and proves the first inequality; otherwise it is immediate. A capacity-one core consumes at least $2D$ root-cut edges, so there are at most $a/2$ such cores. For every other core, $v(z)=\floor{(z+1)/3}\ge1$ and $8v(z)\ge2z$, proving the remaining structural inequalities. The cheapest $\mu-\delta$ original entries satisfy the window budgets, giving the bounds on $3\delta$ and $f_3$; the height-two $U$ count gives $A_*\ge V_0$. If $L$ is a unit list, we also have
\[
f_2+X\le\mu,\qquad V\ge V_0,\qquad V\ge X.
\]
Since $\alpha,b\le D-1$ and $\alpha+b\ge D+2$, both heads are at least three and $Y\ge1$. Every genuine high rectangle $x+y\ge D+2$ satisfies $f_x+g_y-r+1\le h\le\rho$; in particular $g_b\le\rho+r-2$.

Fix an original attaining pair at each $2\le j\le D$. Classify it, in order, as a unit row $(1,y)$, matching row $(x,1)$, two-head row $(2,y)$, or mixed row $(x,y)$ with $x\ge3,y\ge2$. When needed, height two's $(2,1)$ row is moved wholly into the two-head group and never charged twice. Unit rows have $B_j\ge\delta$. Mixed rows with both heads at least three are positive by $f_x\le X$, $g_y\le Y$, and $\mu\ge2(X+Y)+a+3d+r-3$. For $x\ge3,y=2$,
\[
B_j\ge\mu+a/2-X-g_2+r-1
\ge\frac{\mu+a-3\rho+7d}{4}+\frac{3r}{2}-2>0.
\]
\paragraph{One internal coloring repairs every matching row.}
Put $q_0=T-2Y+1\ge d$ and $w=T-Y+1$. If a matching row of height at least three has loss $L_j=-B_j>q_0$, choose $2w$ vertices in the cores, an even number in each, color incident internal edges red and the others blue. Its blue window residual $Q$ is nonempty because the original blue half-order cost is $t+Y=\mu-w+1$. The trigger gives $f_j>2w+a/2+r-2$, and \eqref{eq:se-guards2} makes the new red suffix longer than $r+a$. The same trigger at height two is used only when $f_2\le\mu$; then $\mu>2w$ again guarantees legality. Both new complete sources have positive rank.

For $m_Q=|Q|\ge1$, the average matching-row width increases by $L_j-q_0+m_Q-1$. Maximizing old and new columns separately on each actual graph leaves matching-row loss at most $q_0$. The old nonnegative columns also give $L_j\le(f_j-q_0)_+$. Hence each remaining loss is at most $f_j/2$, and the total over heights at least three is at most $X$. If no trigger occurs, the same bound already holds for the old columns.
\paragraph{The case $b\le D-2$.}
At height $D$ the winner is either HH (both heads at least three), or B2, namely $(D-1,2)$. Put $c_l=g_l-g_{l+1}$. The common top-column payments are
\begin{align}
\text{HH:}\quad B_D&\ge X+q_0+Y+3a/2+2d+2r-4,\label{eq:se-HH}\\
\text{B2:}\quad B_D&\ge T+c_1+a/2-\rho+g_b,\label{eq:se-B2}\\
C_0+B_D-X-d&\ge X+a/2+g_b+Y+2d+r-2+(D/2-3)\rho>0.\label{eq:se-B2pay}
\end{align}
These follow, respectively, from $B_D\ge\mu+a/2-X-Y+r-1$, the top-rectangle bound, and substitution of \eqref{eq:se-guards1}--\eqref{eq:se-guards2}. If $L$ is nonunit and either $f_2>\mu$ or $v_2>a$, every two-head row has, respectively, $B_j\ge A_*/2>0$ or $B_j\ge\mu-f_2+v_2-a/2>0$. Thus \eqref{eq:se-HH} or \eqref{eq:se-B2pay} pays the only remaining loss $X+d$.

If $L$ is all units, then $B_{l+1}=u+\delta-c_l$ exactly for $j=l+1\le b$. For $l\ge3$, $c_l\le Y\le u$; the boundary two-head row at $b+1$ is at least $1+a/2$. Only the two lowest heights remain. If both are two-head rows,
\[
B_2+B_3=2u+2\delta-t+g_3
\ge r+2d+2Y-2+T+g_3>0.
\]
If height two is a unit row and height three a negative two-head row, put $D_3=-B_3$. Then
\[
D_3+X\le q_0-A',\qquad A'=3d-\rho+Y+r+2\delta+g_3-1.
\]
This follows by combining $D_3+X\le g_2-g_3-S-\delta$ with the common constraints; $A'$ need not be positive. The remaining low-row combinations and their positive remainders are listed below. The losses $X,d$ are already included, and unlisted nonnegative rows are discarded.
\begin{longtable}{P{.25\textwidth}P{.65\textwidth}}
\toprule
 & \\ \midrule\endhead
$B_2<0$, HH & \eqref{eq:se-HH} pays $X+q_0+d$.\\
$B_2<0$, B2 &
$(D/2-2)\rho+T+S+\delta+a/2+g_b-d>0$,
using $-B_2+X\le c_1-S-\delta$.\\
$B_2=\delta,\ B_3<0$, HH &
$(D/2-2)\rho+2Y+a/2+3d+r+2\delta+g_3-1>0$.\\
$B_2=\delta,\ B_3<0$, B2 &
$(D/2-3)\rho+3Y+c_1+a/2+g_b+2d+r+3\delta+g_3-2>0$.\\
\bottomrule\end{longtable}
Finally suppose $L$ is nonunit, $f_2\le\mu$, and $1\le v_2\le a$. The two actual source columns give internal and boundary two-head bounds
\begin{gather*}
2B_{l+1}\ge A_*+\delta+u-c_l\quad(1\le l<b),\\
B_{b+1}\ge(V+A_*-\rho+1)/2>0,\\
2(B_2+B_3)\ge2A_*-\rho+V+r+2d+2Y-2+T+g_3>0
\quad\text{if both rows are two-head},\\
2(B_3+q_0)\ge2\delta+V+Y+4d+r+g_3-1>0,\\
2(-B_2)\le f_2-A_*-\delta-S-T-g_2\quad\text{if }B_2<0.
\end{gather*}
Internal two-head rows of height at least four are nonnegative. If both low rows are two-head, use the third bound jointly. If height two is negative and height three is not two-head, HH uses the same coloring to pay $q_0$; B2 retains the old height-two row and has remainder at least
\[
\frac{T+V+A_*+a+3d+2r+2\delta+3Y-4}{2}
+a/2+g_b+(D/2-5/2)\rho>0.
\]
If height two is a unit row and height three is a negative two-head row, HH uses $B_3>-q_0$. For B2 add the height-three bound directly to the top column; twice the remainder is at least
\[
a+2r+4d+2\delta+4Y+V+g_3+2g_b-4+(D-5)\rho>0.
\]
The substitutions are $g_2+2X\le\mu+\rho-3d-r+2-Y$ and $2g_2\le\mu+q_+-d-Y$; the coefficient is therefore $4d$. All remaining low rows are nonnegative or belong to the repaired matching group. This exhausts $b\le D-2$ and contradicts strictness through \eqref{eq:two-ledger-potential}.
\paragraph{The boundary $b=D-1$.}
Use the complete-area account \eqref{eq:two-ledger-area}, with no recoloring and no $-d$ term. Put height two into the two-head group. At heights $3,\ldots,D-1$, each negative matching part is at most both $\delta$ and $(\rho-2\delta)_+$, hence at most $\rho/3$. If their total is $L_m$, then
\[
L_m\le(D-3)\rho/3,\qquad C_0-L_m\ge D\rho/6.
\]
For nonunit $L$ with $f_2>\mu$ or $v_2>a$, all two-head rows are nonnegative and the proof is complete. For a nonunit short prefix $f_2\le\mu$, $1\le v_2\le a$, let $K$ be the negative internal two-head rows. If $K\ne\varnothing$, sum their disjoint exact-head classes to obtain
\[
\delta+\sum_{j\in K}B_j
\ge\frac{A_*+3\delta+u-t+g_b}{2}
\ge\frac{A_*+T-a-\rho+g_b}{2}\ge-a/2.
\]
If the top row is also two-head, it is at least $V+a/2+r-1-g_b$. Combined with the preceding bound this is at least $1-\rho+(T+g_b)/2$, or at least $1-\rho$ if $K$ is empty. The reserve $D\rho/6$ pays either with strict remainder. If the top row is HH or B2, respectively,
\begin{align*}
B_D&\ge(\mu+2a+3d+3r-5)/2,\\
B_D&\ge3\mu/8+3a/4-5q_+/8+9d/8+5r/4-3/2.
\end{align*}
The first is greater than $a/2$. In the second use $q_+\le(\mu-\rho)/2$; after subtracting $a/2$ the remainder is at least $\mu/16+a/4+5\rho/16+9d/8+5r/4-3/2>0$.

For unit $L$ and nonempty $K$, the exact rows give
\[
\delta+\sum_KB_j\ge u+2\delta-t+g_b.
\]
If the top row is two-head, adding it and using the budget on $3\delta$ gives at least $(8V_0+a-2\mu)/6+4r/3-1-2\rho/3$, strictly positive after adding $D\rho/6$. If the top row is HH or B2, first bound the internal loss by
\[
-\mu/3-a/3+V_0/3-2\rho/3+r/3+g_b.
\]
Adding the corresponding top bound and $D\rho/6$, the HH remainder is at least
\[
\mu/6+2a/3+V_0/3+(D-4)\rho/6+3d/2+11r/6-5/2+g_b>0.
\]
For B2 use $q_+\le(2a+2V_0+\rho)/3$, the weighted average of the two stated bounds; its remainder is at least
\[
(\mu-2V_0)/24+9d/8+(D/6-7/8)\rho+19r/12-3/2+g_b>0.
\]
If $K$ is empty, retain $\delta\ge0$ and the positive top row directly. All original winners and residual types have now been exhausted, and the complete-area account is strictly positive. Choose the actual graph with larger $c_X+E_X$ to obtain a positive smaller-rank strict counterexample, contradicting the same all-data minimality.
\end{proof}
\section{Block selection retaining ports and internal degree suppression}\label{app:blocks}
Throughout this section use the intrinsic region $A,R,B$ of the original minimum.
Put $\zeta=\sum_{v\in A}(D-d_G(v))$, $\epsilon=e(G[A])$, $\beta=e(A,B)$,
$L_A=\zeta+2\epsilon+\beta$, $M=a+\mu$,
$R_0=(D\rho+\beta-\zeta)/2$, and $Q=R_0-h$.
Each root meets at least two actual cores, so $L_A\le(D-2)a$. The degree identities are
\begin{equation}\label{eq:block-account}
 e(G)=DM+R_0+e(B),\qquad
 \omega:=e(R)-D\mu=\frac{D\rho+L_A}{2},\qquad
 R_0+\omega=D\rho+\epsilon+\beta.
\end{equation}
\begin{lemma}[An actual block prefix]\label{lem:block-prefix}
Put $p=\floor{D/2}$. Order the positive-weight nontrivial blocks of $R$ by increasing matching capacity, and take the shortest prefix whose weight reaches $h$. Let $z,Z,l,u,w$ denote its total capacity, largest capacity, number of selected blocks, number of unselected positive blocks, and selected weight. Then
\[
 l\le h,\quad h\le w\le h+p-1,\quad
 \mu\ge z+uZ\ge(u+1)Z.
\]

If there are $b_0=(Z-r+1)_+$ independent original degree-$D$ vertices outside all vertices of the selected blocks, a paid strict descent exists. Otherwise $u\le p$ and
\begin{equation}\label{eq:block-failure}
 D\rho+L_A+2W_-\le2h+2p(p+1)-2,
 \qquad W_-=\sum_{w_L<0}(-w_L).
\end{equation}
If $a>0$, then also
\begin{equation}\label{eq:block-parity}
 Q\ge
 \begin{cases}
 p(\rho-u),&D=2p+1,\\
 (p-1)(\rho-u)+1,&D=2p.
 \end{cases}
\end{equation}
\end{lemma}
\begin{proof}
The block decomposition in Appendix \ref{app:matching} makes each block factor-critical and its capacity additive.
After any subset of blocks is retained, each nontrivial component is still a tree-like gluing of odd factor-critical blocks at cut vertices. Such a gluing is factor-critical: after deleting a vertex on one side, use its perfect matching there and a perfect matching avoiding the shared vertex on the other side. In particular, within each original core the retained maximum matching can avoid any prescribed port.
Keep the original matching assigning roots to distinct original cores, and avoid its port in each assigned core. This gives
\[
 \nu(G-E(\mathcal S))=\nu(G)-\sum_{L\in\mathcal S}z_L.
\]
The reverse inequality follows because a matching uses at most $a$ edges meeting roots. Shared cut vertices between blocks do not invalidate this identity.

Each block is induced and its vertices have original degree $D$, so
$w_L=e(L)-Dz_L=(D-\partial L)/2\le p$; positive weights are positive integers.
The total block weight is $\omega\ge h$. The prefix therefore exists and has the stated three bounds.

Take the independent set $I$ in the statement. In the erasure graph its vertices still have degree $D$; degree counting and Hall's theorem give a matching saturating $I$. Thus
$1\le K=z+b_0\le\nu(G)<N$.
For $G'=(G-E(\mathcal S))-I$, Lemma \ref{lem:menu} supplies the complete $P_{r+K}$ relation since $r+b_0>Z$. The actual fee is
\[
 e(G)-e(G')=DK+w\ge DK+h\ge E_r-E_{r+K},
\]
contrary to the same all-data minimum.

If $u\ge p+1$, outside the selected blocks there are at least
$2uZ+a+\rho-h\ge(D+1)Z$ original degree-$D$ vertices. A greedy argument supplies $Z$ independent vertices, so $u\le p$.
The total positive weight is at most $h+p-1+pu$. Comparing it with $\omega+W_-$ proves \eqref{eq:block-failure}.
Let $q_0$ count portless original cores. If $D$ is odd, $q_0=0$; if $D$ is even and $a>0$, strict root Hall expansion gives $q_0\le\rho-1$.
At most $h+u$ original cores contain positive blocks, and every other core has cut at least $D$. A ported core has cut at least $\varepsilon_D$, equal to one in odd degree and two in even degree. Summing these disjoint actual cuts gives
\[
 L_A\le(D-\varepsilon_D)(h+u)-D\rho+\varepsilon_Dq_0.
\]
Substitute into $Q\ge D\rho/2-L_A/2-h$, and use $h\le\rho$ to obtain \eqref{eq:block-parity}.
\end{proof}
\begin{lemma}[One common internal choice]\label{lem:internal-saving}
In each original core of capacity $z_C$, fix an internal edge set $E_C$ minimizing its matching number subject to
$\Delta(C-E_C)\le D-1$. Put
$s_C=z_C-\nu(E_C)$, $v_C=\floor{(z_C+1)/3}$,
$S=\sum s_C$, $V=\sum v_C$, $c_+=|\,\{C:z_C>0\}\,|$, and $n_+=|R^+|$.
Then
\[
3V\ge\mu-a-\rho,\qquad V\ge\rho+L_A/D,
\]
\[
 S\ge
 \begin{cases}
 \mu/5,&D=5,6,\\
 \mu/4,&D=7,\\
 \mu/3,&D=8,\\
 3\mu/8,&D\ge9,\\
 \mu/2,&D\ge13.
 \end{cases}
\]

For $D\ge9$, the same choice also satisfies
\begin{equation}\label{eq:internal-strong}
 S\ge V,\quad 5S+V-n_+\ge\mu/14,\quad
 S+V-n_+/4\ge\mu/8,\quad S+V\ge\omega+c_+.
\end{equation}
\end{lemma}
\begin{proof}
The first inequality follows from $3v_C\ge z_C-1$.
A core of capacity zero has cut $D$; one of capacity one has cut at least $3(D-2)\ge D$; and $v_C\ge1$ for capacity at least two.
Thus $\partial C+Dv_C\ge D$ for every core, and summation proves the second inequality.

It suffices to construct a vertex set $T$ suppressing internal full degree. Its incident edge set has matching number at most $|T|$, so the fixed minimum $E_C$ has at least the resulting saving.
Select vertices independently with probability $\theta$ and add each full-degree vertex whose closed neighborhood is missed. The expected size is at most
\begin{equation}\label{eq:suppress-expectation}
 (2z+1)\bigl[\theta+(1-\theta)^{D+1}\bigr].
\end{equation}
For $D=5$, every odd core has a port. Take a near-perfect matching exposing a port. On each matched pair, choose either endpoint with probability $3/10$ each, or neither with probability $2/5$; choose the exposed port with probability $3/10$.
After adding missed full-degree vertices, the expectation is at most
$(3/10+2401/25000)(2z+1)-2401/25000$.
Thus the saving is at least $\ceil{2599z/12500-3/10}$. Writing $z=5k+r$ gives $\ceil{z/5}$ except for $z=1,6,11$. The first has no internal full-degree vertex. For the other two, uniform sets of $3,6$ vertices followed by augmentation have expectations below $5,9$, giving savings at least $2,3$.

For the other small degrees use $\theta=1/4$ in \eqref{eq:suppress-expectation}.
The table lists the initial capacities requiring separate treatment and the expected size for a uniform $k$-set followed by augmentation. In general this expectation is
$k+(2z+1)\binom{2z-D}{k}/\binom{2z+1}{k}$.
\begin{center}\small
\begin{tabular}{ccP{.72\textwidth}}
\toprule
$D$&$z$& $k;\quad \mathbb E|T|$\\\midrule
$6$&$3,\ldots,11$&
$1,2,2,3,3,4,4,4,5;\quad
1,\frac94,\frac{16}5,\frac{43}{11},\frac{441}{91},\frac{11}2,
\frac{437}{68},\frac{2141}{285},\frac{11683}{1463}$\\
$7$&$4,5,6,7$&$1,2,2,3;\quad 2,\frac{13}5,\frac{11}3,\frac{54}{13}$\\
\bottomrule
\end{tabular}\end{center}

For $D=6$, the probability coefficient is $6283/16384$, giving saving $z/5$ for $z\ge12$; no deletion is needed for $z\le2$.
For $D=7$, the coefficient is $22945/65536$, giving $z/4$ for $z\ge8$; no deletion is needed for $z\le3$.
For $D=8$, the coefficient $85219/262144<1/3$ directly gives saving $\ceil{z/3}$.

For $D\ge9$, the coefficient is below $307/1000$.
At $z=5,6,7,8,9,10,12$, the expectation bounds are respectively
$3.377,3.991,4.605,5.219,5.833,6.447,7.675$, giving $s_C\ge\ceil{3z/8}$.
Increasing $z$ by eight increases the bound by $4.912<5$; the remaining residue $z\equiv3\pmod8$ starts at $z=19$ with $11.973<12$.
At $z=11$, a uniform three-set has augmented expectation $47/7<7$; for $z\le4$ no deletion is needed.
This proves the saving bound at every capacity, as well as $s_C\ge v_C$.
For $D\ge13$, $\theta=1/5$ gives a coefficient below $1/4$, hence $s_C\ge\ceil{z/2}$ by integrality.

For $z\ge10$, the bounds $s_C\ge3z/8$ and $v_C\ge(z-1)/3$ give
$5s_C+v_C-(2z+1)\ge5z/24-4/3\ge z/14$ and
$s_C+v_C-(2z+1)/4\ge5z/24-7/12\ge z/8$.
For $z=5,\ldots,9$, substitute the savings $2,3,3,3,4$; for $z\le4$, use $s_C=z$.

Finally we prove $s_C+v_C\ge w_C+1$ for every nontrivial core, where $w_C=e(C)-Dz\le p$.
If $2z+1\le D$, then $w_C\le0$. For $2z+1>D$ and $z\le p+2$, the possibilities
$(D,z)=(2p,p),(2p+1,p+1),(2p,p+1),(2p+1,p+2),(2p,p+2)$
use uniform sets of $1,1,2,2,2$ vertices, respectively. Their integer augmented sizes are at most $1,2,2,2,2$, giving $s_C+v_C\ge p+1$.
For $z\ge p+3$, put $c_D=1/5+(4/5)^{D+1}$. Then
\[
 s_C+v_C-p\ge(4/3-2c_D)z-c_D-1/3-p.
\]
The right side increases with $z$, and is positive at $z=p+3$ if $c_D<(p+11)/(6p+21)$.
For $p=4,5$ use $c_9<5/16<16/51$; for $6\le p\le10$ use $c_D<7/27$; for $11\le p\le15$ use $c_D<7/30<26/111$.
For $D\ge32$, instead take $\theta=1/10$, with coefficient below $2/15$; the resulting bound exceeds $(D-14)/30>0$.
Integrality proves the per-core inequality, and summation gives \eqref{eq:internal-strong}.
Every estimate applies to the fixed minimum-matching-cost choice; no later recoloring requires changing $E_C$.
\end{proof}
\begin{lemma}[Internal covers near the top degree]\label{lem:internal-cover}
Put $V=\sum_C\floor{(z_C+1)/3}$. An actual set $T$ with $\Delta(R-T)\le D-c-1$ can be chosen with $|T|\le\lambda V$ as follows:
\[
\begin{array}{c|c|c}
c&D&\lambda\\\hline
1&9,10;\quad11,\ldots,14;\quad D\ge15&3;\quad5/2;\quad2\\
2&D\ge9&4\\
3&16,\ldots,19;\quad20,\ldots,24;\quad D\ge25&3;\quad5/2;\quad2
\end{array}
\]
\end{lemma}
\begin{proof}
Choose vertices with probability $\theta$, then add every unchosen vertex with remaining degree at least $D-c$. The expectation in a core is at most $n\kappa_{D,c}(\theta)$, where
\[
\kappa_{D,c}(\theta)=\theta+
(1-\theta)\sum_{j=0}^{c}\binom Dj\theta^j(1-\theta)^{D-j}.
\]
The coefficient decreases with $D$, and $n=2z+1\le6v+3$.
For $c=1$, at $(D,\theta)=(9,1/3),(11,3/10),(15,1/4)$ the coefficients are
$25315/59049<4/9$,
$379093069720/10^{12}<19/50$, and
$1332022150/4294967296<1/3$.
The first gives integer size at most $3v$.
For the second, no deletion is needed when $v\le1$; $v=2$ gives size at most five, and for $v\ge3$ use
$57v/25+57/50<5v/2+1/2$.
The third gives size at most $2v$.
For $c=2$, use $\theta=1/2$: the coefficient at $D=9$ is $279/512$, which with integer rounding gives $4v$; when $v\le1$, at most four vertices suffice directly.
For $c=3$, the coefficients at $(16,1/3),(20,3/10),(25,1/4)$ are respectively
$57333569/129140163<4/9$,
$374960763152611698817/10^{21}<19/50$, and
$1450882160153428/4503599627370496<1/3$.
The same rounding applies. Take the union of the actual sets in the cores.
\end{proof}
\section{Unit common residuals when an original head equals the degree cap}\label{app:unit-side}
\begin{proposition}\label{prop:unit-single}
An original all-data minimum cannot satisfy
$D\ge3$, $n_1=D$, $t\le\mu<s$, with $L(\nn)$ a nonempty unit list.
\end{proposition}
\begin{proof}
If $m_1=D$, apply Lemma \ref{lem:reverse-owncap}; $m_1=1$ is incompatible with $\ell_r=D+1$.
Thus $2\le b=m_1<D$. Use $M,R_0,Q,L_A$ from Appendix \ref{app:blocks}, and put
\[
 f_j=c_j(\nn),\quad g_j=c_j(\mm),\quad P=f_D,\quad
 S=s-\mu,\quad w=t-r+2,
\]
\[
 X_F=\sum_i\floor{(n_i-1)/2},\qquad X_H=\sum_j\floor{(m_j-1)/2},
\]
\[
 Y=\sum_{j=3}^{D}f_j,\qquad V_e=\mu-f_2,\qquad
 V=\sum_C\floor{(z_C+1)/3}.
\]

\paragraph{One exterior and the full columns.}
Lemma \ref{lem:two-successors}, weakened to $S$ red unit stars, gives
$P_{r+a}(B;1^S,\mm)$ of rank $k+d=N-M>0$.
If $r+a>S$, the last legal red corner, colored entirely blue, forces more than $k$ blue components, a contradiction; the other side is identical.
Thus $r+a\le\min(S,t)$.
Apply Lemma \ref{lem:strong-endpoint} to the original endpoint, writing $\mathfrak b$ for the size of a maximum-cardinality optimal Tutte barrier and using $\mathfrak b\ge a$, and then to this endpoint in $B$. We obtain
\begin{equation}\label{eq:unit-guards}
 w\ge a+2(d+X_F),\qquad
 S\ge a+r+2d+2X_H-2\ge a+r+2\rho+2X_H.
\end{equation}

The union of $A$ and an optimal Tutte barrier of $B$ is an optimal barrier of $G$, by the exact matching decomposition. Its size is at least $a$, so a maximum-cardinality optimal barrier has size $\mathfrak b\ge a$; containment of $A$ in that maximum barrier is not required.
The unit residual means that the original nonunit prefix satisfies both actual budgets and that $U$ has no head two. Consequently
\begin{equation}\label{eq:unit-budgets}
\begin{gathered}
 f_2+X_F\le\mu,\quad 3f_2+Y\le2\mu+a+\rho,\quad
 V_e\ge\max(V,X_F),\\
 3V_e\ge\mu+Y-a-\rho,\qquad Y\ge(D-2)P.
\end{gathered}
\end{equation}

The original top column gives $P+g_2-r+1\le h\le\rho$, whence
$W_2=t-g_2\ge2X_F+a+\rho+P+1>X_F$.
For $b\ge3$, also $P+g_3\le r-1$, so $r\ge P+2\ge3$.
At original heights $2\le j\le D$, there are only three possible winners:

\begin{equation}\label{eq:owncap-columns}
 C_j=\max\bigl(f_j+w-1,\ s+g_j-r+1,\ f_2+g_{j-1}-r+1\bigr),
\end{equation}
Omit terms with nonexistent positive prefixes. Every other mixed row has red threshold at least three, and its excess over the first term is at most
$f_x-f_j-(t-g_y)\le X_F-W_2<0$.
At an internal height $3\le j\le b$, a strict excess of the third term over the second would imply
$g_{j-1}-g_j>S+V_e$, contrary to $g_2\le r+\rho-2$ and \eqref{eq:unit-guards}.
Thus strict mixed debt can occur only at $b+1$.

If no such debt occurs, the following single matching-join account pays the profile.
A nontrivial factor-critical core of capacity $z$ has vertex-cover number at least $z+1$: an independent set of size $z+1$ contradicts a perfect matching after deletion of a vertex outside it.
Singleton cores have a matching into the roots, so $\tau(G[A\cup R])\ge M+\rho$.
Color the exterior of $B$ red at $(1,r-1)$. Since $M<s$, this gives
$B\to((s-M)K_2,H_{r-1})$, and hence $\tau(B)\ge N-M$.
Indeed, a smaller cover could be split into two sets below the respective component demands, and their incident edges colored in priority order to avoid both forests.
Thus $\tau(G)\ge N+\rho$.
Weaken each original side separately to units. On one maximum-degree deletion history of $N$ steps, these two complete sources force degrees
\[
 J_q=\max\bigl(n_{\max(1,q-t)},m_{\max(1,q-s)}\bigr),
 \qquad r\le q\le s+t.
\]
The final remainder $Z$ gives $e(G)\ge\sum J_q+e(Z)$ and $\tau(G)\le N+e(Z)$.
With no mixed debt, $E_r=\sum J_q+h$, so $e(G)\ge E_r$.

Now assume strict boundary debt. Put
$c=D-b$, $\gamma=r-1-g_b$, and $\Delta_g=g_2-g_b$.
Its exact integer condition is $w\le f_2-f_{b+1}-\gamma$.
For $b\ge3$, $\gamma\ge P$; for $b=2$, $\gamma=P-h$.
The complete columns on the same $B$ are $Q_j=S+g_j-r-a+1$ for $1\le j\le b$, all positive.
Subtract column by column: the first difference is $M$, the second is $M+Z_2$, and the remaining differences through $b$ are $M$, where
$Z_2=(t-g_2-S-V_e)_+$.
At height $b+1$ the excess over $M$ is $-V_e-a-\gamma$, and at each higher height it is at most one less.
Including the original top width $h$, we obtain
\begin{equation}\label{eq:unit-payment}
 g:=e(G)-E_r\ge Q+c(V_e+a+\gamma)+(c-1)-Z_2.
\end{equation}
\paragraph{Binary and small-degree boundaries.}
If $b=2,Z_2=0$, use $V\ge\rho+L_A/D$ to obtain
\[
g\ge(D/2-1)\rho+\frac{D-4}{2D}L_A+(D-2)(a+P)+D-3>0.
\]
For $D=3$, use $L_A\le a$ in the negative term; the result is still positive.
If $b=2,Z_2>0$, substitute the exact
$Z_2=w+\gamma-1-S-V_e$ and use \eqref{eq:unit-budgets}:
\[
g\ge(D/2-1)\rho+(D-2)a/2+r+(2D-3)P+D-2>0.
\]
Henceforth $b\ge3$.
For $Z_2=0$, \eqref{eq:unit-payment} gives $g\ge R_0+a+P$.
For $Z_2>0$, its weaker $c=1$ version gives
\begin{equation}\label{eq:unit-positive-Z}
 g\ge Q-w+S+2V_e+a+\Delta_g+1
 \ge Q+S+Y+f_{b+1}+\gamma+\Delta_g+1-\rho.
\end{equation}
Here $Y+f_{b+1}+\gamma\ge DP$.
For $D=4$, use $R_0\ge2\rho-a$ and $S\ge a+r+2\rho$ directly.
If $Z_2=0$, then $g\ge R_0+a+P\ge2\rho+P>0$.
If $Z_2>0$, then $g\ge2\rho+r+4P+1>0$.
Henceforth $D\ge5$. Equation~\eqref{eq:block-failure} gives
$2\rho\le D+1$ in odd degree and $2\rho\le D+4$ in even degree $D\ge6$.
Thus $R_0\ge-a$ pays in the remaining range.

Assume $R_0<-a$. Then
$\zeta\ge D\rho+\beta+2a+2$, and in particular $a\ge\rho+1$.
Apply Lemma \ref{lem:block-prefix}; if its operation succeeds, we are done.
On failure, the untouched original cores give at least $a$ independent full-degree vertices.
Except for original portless $K_{D+1}$ cores, the components outside the selected blocks are $D$-colorable by Brooks' theorem or greedily. Thus they also give an independent set of size at least $(2uZ+a-q_0)/D$.
Failure forces
\begin{equation}\label{eq:unit-buffer-fail}
 Z\ge a+r,\qquad 2uZ+a-q_0\le D(Z-r),\qquad \mu\ge(u+1)Z.
\end{equation}
The parity cut bound gives $L_A\le(D-\varepsilon_D)p-\varepsilon_D$.
For $D=5,6$, this is at most seven and ten, whereas $\zeta\ge D\rho+2a+2$ is at least eleven and twelve.
For $D=7$, only $\rho=h=1$, $3\le a\le4$, $u=3$ remain.
Equation \eqref{eq:unit-buffer-fail} gives $Z\ge a+7r$, $\mu\ge96$, hence $V_e\ge32$ and $R_0\ge-5$.
Equations \eqref{eq:unit-payment} and \eqref{eq:unit-positive-Z} both give positive bounds.
For $D=8$, first $\rho=h=1,a\ge3,q_0=0$. The choice $u=4$ contradicts \eqref{eq:unit-buffer-fail}, so $u\le3$.
Then $L_A\le16$, forcing $a=3,\zeta=L_A=16,\beta=\epsilon=0,R_0=-4,u=3$.
Now $Z\ge14,\mu\ge56,V_e\ge20$, and the same two inequalities are positive.

\paragraph{Paid buffers in higher degree.}
Let $D\ge9$. Each core satisfies $\partial C+(D/2)v_C\ge D$, so
\begin{equation}\label{eq:unit-large-V}
 V\ge2\rho+2L_A/D,\qquad
 R_0+3V+a\ge(D-1)\rho+\epsilon+\beta.
\end{equation}
The probability-$1/4$ suppression construction, whose coefficient is below $1/3$, gives an actual set of size
$t_R\le\mu-V\le2V+a+\rho$.
At the original corner $(P,r-P)$, color root edges, suppression edges, and edges at previously chosen exterior centers blue. The forced red degree-$D$ centers supply
$z_B=(w+P-1-a-t_R)_+$ independent exterior vertices.
The same unit source on $B$ has profile cap $b<D$; deleting these actual full-degree vertices gives extra saving at least $z_B$.
Include this saving in the same $B$ account in \eqref{eq:unit-payment}.
Equation \eqref{eq:unit-large-V} pays $Z_2=0,c\ge3$ and $Z_2>0,c\ge4$.
For the latter use $-Z_2+z_B\ge V_e+L_2+P-t_R$, where $L_2=S+g_2-r-a+1$.
Only $Z_2=0,c=1,2$ and $Z_2>0,c=1,2,3$ remain.

The original forest $H$ forbids actual block-capacity subset sums in
\begin{equation}\label{eq:unit-forbidden}
 [t-d,\ t+2X_H+d-2].
\end{equation}
For a sum $x$ in the first half, color those blocks blue and every other edge red. The red matching is at most $s-r+1$, while the blue odd groups have slot capacity below $t+X_H$.
In the second half, put $k_0=x-t-X_H+1$, choose $2k_0$ vertices distributed evenly in the blue odd groups, and recolor their incident blue edges red.
The red matching is still too small, while the blue slot capacity becomes $t+X_H-1$.
This contradicts the same original corner $(r-1,1)$.
If $t_b$ is the minimum number of vertices reducing $R$ to degree $b-1$, the same coloring test also gives
\[
 t\ge\mu+d+1-t_b.
\]
On the other hand, $G\to(F_1,wK_2)$ supplies $w-1$ successive full-degree independent centers, then the $P$ full-degree centers in one complete copy of $F_1$.
There are $t-r+P+1$ actual independent vertices.
The selected factor-critical block union has independence number at most $z$, so
\begin{equation}\label{eq:unit-buffer-success}
 t+P\ge z+Z
\end{equation}
suffices for Lemma \ref{lem:block-prefix}.

Put $g_H=2(d+X_H)-1$.
If all selected capacities and one unselected capacity are at most $g_H$, their sum is below $t-d$: otherwise the first sum reaching $t-d$ lies in \eqref{eq:unit-forbidden}.
This proves \eqref{eq:unit-buffer-success}.
In all other cases the least unselected capacity $Y_*$ satisfies
$Y_*\ge Z$, $Y_*\ge2(d+X_H)$, and $\mu\ge z+uY_*$.
Consequently
\begin{equation}\label{eq:unit-last-buffer}
 t+P-z-Z\ge(u-1)Y_*+d+P+1-t_b.
\end{equation}
For $Z_2=0,c=1$, strict underpayment forces $V\le-Q-1$ and $\rho+1\le u\le p$.
Lemma \ref{lem:internal-cover} gives $t_b\le\lambda_DV$, and $X_H\ge p-1$.
Substitute \eqref{eq:block-parity} into \eqref{eq:unit-last-buffer}. The lower bound increases with $\rho$, so set $\rho=1$.
The bounds over all intervals are as follows.
\[
\begin{array}{c|c}
D& t+P-z-Z\text{ lower bound}\\\hline
9&9-2u\ge1\\
10&10\\
11,13&(2-p/2)u+p/2+9/2\ge3/2\\
12,14&(9/2-p/2)u+p/2+9/2\ge10\\
2p+1\ge15&2u+4\\
2p\ge16&4u+4
\end{array}
\]

For $Z_2=0,c=2$, use $t_b\le4V$ and
$V\le-Q/2-a-\gamma-1$.
The lower bounds in \eqref{eq:unit-last-buffer} are
\[
\begin{cases}
2\rho(u-1)+\rho+7+2p(\rho-1)+4a+4\gamma,&D=2p+1,\\
2\rho(u-1)+\rho+9+2(p-1)(\rho-1)+4a+4\gamma,&D=2p,
\end{cases}
\]
which are positive.
Finally suppose $Z_2>0,c\le3$. We have $X_H\ge p-2$; put
$C=2a+2r+2d+2X_H+\Delta_g-2$.
Underpayment and \eqref{eq:unit-positive-Z} give
\[
t\ge Q+2V_e+C,\qquad
-Q\ge a+r+\rho+2X_H+DP+\Delta_g+2.
\]
With $\delta_u=u-\rho$, this implies $\delta_u\ge5$ and $\rho\le p-5$.
If $u\ge2\rho+3$, substitution of $3V_e\ge\mu+Y-a-\rho$ gives
\[
3(t+P-z-Z)\ge(u-2\rho-3)Z+3S+3g_2+3+3P+Y-a-\rho>0,
\]
using $S\ge a+5$ and $a\ge\rho+1$.
Otherwise $5\le\delta_u\le\rho+2$ and $3\le\rho\le p-5$, hence $p\ge8$.
By Lemma \ref{lem:internal-cover}, take
$\lambda=3$ for $16\le D\le19$, $5/2$ for $20\le D\le24$, and $2$ for $D\ge25$.
Then $t_b\le\lambda V\le\lambda V_e$.
Add $\lambda$ times $t\ge Q+2V_e+C$ to twice $t\ge\mu+d+1-\lambda V_e$:

\[
(\lambda+2)(t+P-z-Z)
\ge\lambda Q+\lambda C+
\bigl[2\delta_u+(2-\lambda)\rho-\lambda-2\bigr]Z
 +2d+2+(\lambda+2)P.
\]
The bracket is at least $1,2,6$ in the three intervals, respectively.
Use $Z\ge\ceil{D/2}$, $C\ge4\rho+2p+4$, and \eqref{eq:block-parity}.
The right side is at least
\[
\begin{aligned}
&[(2-\lambda)p+2]\delta_u+
 [(2-\lambda)p+3\lambda+4]\rho+(\lambda-2)p+4\lambda+4,
 &&D=2p+1,\\
&[(2-\lambda)p+\lambda]\delta_u+
 [(2-\lambda)p+4\lambda+2]\rho+(\lambda-2)p+6\lambda+6,
 &&D=2p.
\end{aligned}
\]
For $\lambda=3$, use $\delta_u\le\min(\rho+2,p-\rho)$.
The odd-degree bound is at least $-p^2+17p/2+5\ge1/2$ for $p=8,9$, and the even-degree bound is at least $(17-2p)\rho+30-p\ge17$.
For $\lambda=5/2$, the respective bounds are
$(27/2-p)\rho+18-p/2\ge20$ and
$(29/2-p)\rho+26-p/2\ge55/2$.
For $\lambda=2$, they are $2\delta_u+10\rho+12$ and $2\delta_u+10\rho+18$.
Thus \eqref{eq:unit-buffer-success} always holds, and the actual paid successor gives the contradiction.
\end{proof}
\section{Nonunit common residuals when an original head equals the degree cap}\label{app:nonunit-side}
\begin{proposition}\label{prop:nonunit-single}
An original all-data minimum cannot satisfy
$D\ge3$, $n_1=D$, $t\le\mu<s$, with a nonunit star in $L(\nn)$.
\end{proposition}
\begin{proof}
First suppose $3\le b=m_1<D$.
Use Lemma \ref{lem:reverse-owncap} if $b=D$; the case $b=2$ is treated at the end of this section and in Appendix \ref{app:binary-empty}.
Keep $f_j,g_j,P,w,X_F,X_H,M,R_0,Q$ from Appendix \ref{app:unit-side};
here $S$ denotes the internal saving in Lemma \ref{lem:internal-saving}.
Put $L=L(\nn)$, $\alpha=\max L\ge2$, and $\delta=\mu-(s-|L|)$.
The actual exterior source is $P_{r+a}(B;L,\mm)$ of positive rank $N_B=k+d+\delta<N$.
The endpoint and empty-side relations give

\begin{equation}\label{eq:nonunit-guards}
\begin{gathered}
 w\ge a+2(d+X_F),\qquad
 \mu\ge t+2X_H+d-1,\\
 s-\mu\ge a+r+2d+\delta+2X_H-2>a,\\
 \mu\ge a+r+3d+2X_F+2X_H-3
 \ge a+r+3\rho+2X_F+2X_H.
\end{gathered}
\end{equation}

Again $W_2=t-g_2>X_F$, so the complete original columns are exactly \eqref{eq:owncap-columns}.
Fix the same sets $E_C$ from Lemma \ref{lem:internal-saving}.
At $(P,r-P)$, color root edges, $E_C$, and edges at previously chosen exterior centers blue.
Until the number chosen reaches $(w+P-1-M+S)_+$, the blue matching is too small and a complete original red copy forces a new degree-$D$ center.
These centers are independent and have no root neighbors.
Fix a maximum independent set of full-degree vertices in $B$, of size $l$; then $l\ge(w+P-1-M+S)_+$.
Let $Q_j$ be the complete columns of the actual source on $B$.
A maximum matching of $B$ covers its original full-degree vertices, so $Q_{D+1}=0$.
Delete the independent set. If $l<N_B$, use positive-smaller-rank minimality; if $l\ge N_B$, use its $Dl$ actual incident edges alone.
Both cases give
\[
 e(B)\ge\sum_{j=1}^{D}\max(Q_j,l).
\]
\[
 B_j=M+\max(Q_j,l)-C_j,\qquad
 g:=e(G)-E_r\ge Q+\sum_{j=1}^{D}B_j.
\]
The first column has $B_1\ge\delta$.
An $MH$ winner gives $B_j\ge\delta$.
At a mixed $A$ winner, the actual head-two prefix gives $B_j\ge\mu+c_2(L)-f_2\ge V$.
An $MF$ winner gives $B_j\ge S+P-f_j$, and if $j\le\alpha$ it also gives $B_j\ge\mu-r_j\ge0$.
At height $D$, an $MF$ winner gives $S$, and an $A$ winner gives $V$.
If there is no strict mixed debt, the matching-join proof in Appendix \ref{app:unit-side} pays.
Thus retain a strict $A$ winner.
Put $J=\{j:B_j<0\}$ and $m=|J|$. Then

\begin{equation}\label{eq:nonunit-negative}
 J\subseteq\{\max(3,\alpha+1),\ldots,D-1\},\qquad
 g\ge Q+V+(m+\varepsilon)S+mP-\sum_{j\in J}f_j,
\end{equation}
Here $\varepsilon=1$ for a top $MF$ winner and zero otherwise.

\paragraph{Higher degrees with negative columns.}
Suppose $D\ge9$.
Equation \eqref{eq:nonunit-guards} gives $2X_F<\mu$.
Each negative column has $f_j\ge S+P+1$, while $\sum_{j\ge3}f_j\le2X_F$.
Since $S\ge3\mu/8$, there are at most two negative columns.
For one such column, $f_j\le r_j\le|R^+|/4$, so $g\ge Q+\mu/8+P$.
For two columns $j<l$, the original prefix above $\alpha$ satisfies the vertex budget
$(j+1)f_j+(l-j)f_l+(D-l)P\le|R^+|$.
Since $f_j,f_l\ge S+P+1$, this gives
$f_j+f_l\le|R^+|-3S-(D-1)P-3$, whence
\begin{equation}\label{eq:nonunit-high}
 \begin{cases}
 g\ge Q+\mu/8+P,&m=1,\\
 g\ge Q+\mu/14+(D+1)P+3,&m=2.
 \end{cases}
\end{equation}
Apply the actual block-prefix lemma.
Untouched original cores supply at least $a+\rho-h$ independent full-degree vertices, nonadjacent to the fixed exterior independent set.
A negative column gives $X_F\ge S+P+1$, so the latter set also satisfies
\[
l\ge3S-\mu+3P+2d+1\ge\mu/8+3P+2d+1.
\]
If the actual buffer is still too small, integrality yields
$Z\ge\mu/8+C$, where $C=a+\rho-h+3P+2d+r+1\ge12$.
Together with $\mu\ge(u+1)Z$, this forces
\[
u\le6,\qquad \mu\ge\frac{96(u+1)}{7-u}.
\]
For $D\ge13$, $S\ge\mu/2$. If $u\ge1$, the exterior set is already larger than $Z$; if $u=0$, then $Q\ge D\rho-2h-p+1>0$, contradicting \eqref{eq:nonunit-high}.
For $9\le D\le12$, $u\le\rho$ gives $Q\ge0$, and otherwise $-Q\le6(u-1)$.
For $m=1,2$, strict underpayment combined with the preceding bound would require, respectively,
$-6u^2+37u-61>0$ and $-42u^2+379u-979>0$.
Their discriminants are $-95$ and $-20831$.
Thus every negative-column case has an actual paid block descent.

If there is no negative column and the top is $MF$, the strict mixed column and the distinct top column give
$g\ge Q+S+V\ge D\rho+c_++\epsilon+\beta-h>0$.
Only a strict top $A$ winner remains. Then $b=D-1$, $V\le-Q-1$, and $\rho+1\le u\le p$.
The forbidden interval \eqref{eq:unit-forbidden}, the actual independent centers, and \eqref{eq:unit-last-buffer} in Appendix \ref{app:unit-side} use only the complete original $H$; they do not require a unit residual.
The same bound $t_b\le\lambda_DV$ and \eqref{eq:block-parity} give the six positive bounds for $c=1$ in that appendix.
Hence $t+P\ge z+Z$, and actual block descent also handles this final case.

\paragraph{The small-degree failure range.}
Let $3\le D\le8$.
Assume the block operation fails and let $q_0$ count portless cores.
Root Hall expansion and the actual cuts give

\[
 L_A\le
 \begin{cases}
 (D-1)(h+u)-D\rho,&D\text{ odd},\\
 (D-2)(h+u-\rho)-2,&D\text{ even},\ a>0,\\
 0,&a=0.
 \end{cases}
\]
\[
 2uZ+a-q_0\le D(Z-r).
\]
The last inequality follows from the Brooks independent set outside the selected blocks.
For even $D=4,6,8$, the case $u=p$ is impossible:
\eqref{eq:block-failure} gives $\rho\le5,5,6$, respectively, whereas the last inequality requires
$\rho\ge a+Dr+1$ with roots, or $\rho\ge Dr$ without roots.
Thus $u\le1,2,3$, respectively.
The resulting lower bounds are
\begin{equation}\label{eq:nonunit-Q-small}
\begin{array}{c|c}
D&Q\text{ lower bound}\\\hline
3&\rho-u\ge0\\
4,\ a>0&\rho-u+1\ge\rho\\
5&5\rho-3h-4\ge-2\\
6,\ a>0&2\rho-3\ge-1\\
7&7\rho-4h-9\ge-6;\quad u\le2:\ Q\ge3\rho-6\ge-3\\
8,\ a>0&3\rho-8\ge-5\\
a=0&D\rho/2-h>0
\end{array}
\end{equation}
\[
D=7,\ u=3:\quad \mu\ge4a+28r;\qquad
D=8,\ a>0,\ u=3:\quad\mu\ge2a+16r-2\rho+2.
\]

If $J$ is empty, it suffices to pay $Q+V$, which is plainly positive for $D=3,4$.
The other complete lower bounds are listed below.
They use the same inequalities $3V\ge\mu-a-\rho$ and $V\ge\rho+L_A/D$, not an additional edge count.
\[
\begin{array}{c|l}
5&5\rho-11h/5-12/5>0\\
6,\ a>0&16\rho/3-7h/3-2\ge1\\
7,\ u\le2&7\rho-22h/7-30/7\ge-3/7\\
7,\ u=3&8\rho/3+a+28r/3-9>0\\
8,\ a>0,\ u\le2&29\rho/4-13h/4-15/4>0\\
8,\ a>0,\ u=3&2\rho+(a+16r-22)/3>0
\end{array}
\]

The bound $-3/7$ also gives $g\ge0$ by integrality; when $a=0$, use $Q>0$.
For $D=4,b=3$, the only possible negative column is three.
A graph of maximum degree four splits into two edge graphs of degree at most two:
complete to a four-regular multigraph, orient Euler tours, split the two-regular bipartite incidence graph into two perfect matchings, and restrict back.
Thus $f_3+g_3\le r-1$ and $f_3\le r-2$.
Equation \eqref{eq:nonunit-guards} gives
$V-f_3\ge(r-f_3+2\rho+2X_H)/3>0$, paying that one negative column.

For $D=5,\ldots,8$, the bounds for all negative patterns follow.
Each row uses \eqref{eq:nonunit-negative}, \eqref{eq:nonunit-guards}, and Lemma \ref{lem:internal-saving}.
Write $X=X_F$; the final column is the bound required to pay $-Q$.

\begin{center}\small
\begin{longtable}{P{.23\textwidth}P{.60\textwidth}c}
\toprule
$D,J$& $\text{available saving}-\sum_{j\in J}f_j$&\\\midrule
$5,6;\ \{3,4\}\nsubseteq J$&
$V+S+P-X\ge\mu/30+a/6+r/2+7\rho/6+X_H+P$&$>2$\\
$5,6;\ \{3,4\}\subseteq J$&
$V+2S+2P-X-f_4\ge\mu/15+a/3+(5\rho-2h)/3+4X_H/3+4/3+2P$&$>2$\\
$7;\ m=1$&
$V+S+P-f_j\ge\mu/12+a/6+r/2+7\rho/6+X_H+P$&$>3$\\
$7;\ m\ge3\text{ or }m=2,\varepsilon=1$&
$V+3S+mP-2X\ge\mu/12+2a/3+r+8\rho/3+2X_H+mP$&$>6$\\
$7;\ m=2,\varepsilon=0,J\ne\{3,4\}$&
$V+2S+2P-X\ge\mu/3+a/6+r/2+7\rho/6+X_H+2P$&$>6$\\
$7;\ m=2,\varepsilon=0,J=\{3,4\}$&
$V+2S+2P-f_3-f_4\ge\mu/6+a/3+(5\rho-2h)/3+4X_H/3+4/3+2P$&$>6$\\
$8;\ m\ge2$&
$V+2S+2P-2X\ge2a/3+r+8\rho/3+2X_H+2P$&$>5$\\
$8;\ m=1$&
$V+S+P-f_j\ge\mu/6+a/6+r/2+7\rho/6+X_H+P$&$>5$\\
\bottomrule
\end{longtable}\end{center}

For $D=7,m=1,u=3$, use $\mu\ge4a+28r$ to strengthen the bound to $>6$.
For $D=8,m=1$, the bound $>5$ with roots uses
$\mu\ge a+r+3\rho+6P+2X_H$; without roots, positivity suffices.
The rows containing both three and four use the top-cap bound $f_4\le r+h-2$, hence
$2X+f_4\le\mu-a-3\rho+h-2X_H-2$ and
$X+f_4\le\frac23(2X+f_4)$.
For $D=7,m=2,\varepsilon=0$, a top $A$ winner means $b=6$, so $X_H\ge2$.
These patterns exhaust $J\subseteq\{3,\ldots,D-1\}$ and pay \eqref{eq:nonunit-Q-small}.

\paragraph{A binary opposite side in small degree.}
Let $b=2$, with blue list $2^C1^W$.
Take the original attaining pair $i=\max(1,r-C)$, $j=r-i$.
Remove $i-1$ maximum red heads and $j-1$ blue two-heads, obtaining $\nn^*,\mm^*$.
The ordinary profile columns at heights $q=2,\ldots,D$ are
$\max(f_q+t-1,f_{q-1}+C-1)$, and the top column is $P+C-1$.
Both terms decrease by exactly $r-2$, as does the unit column; the original top plateau has at least $r-1$ entries.
Thus, positionwise,

\begin{equation}\label{eq:binary-suffix-area}
 \ell^*_{2+l}=\ell_{r+l}\quad(0\le l<N),\qquad
 E^*_{2+K}=E_{r+K},\quad N^*=N,\quad h^*=h.
\end{equation}
The same $G$ with $P_2(\nn^*,\mm^*)$ remains an all-data minimum.
The complete blue forest still embeds in the same $R$, and the red side remains longer than $\mu$ by \eqref{eq:nonunit-guards}.
If its common residual is unit, use Proposition \ref{prop:unit-single}; otherwise relabel with $r=2$.
Reapply the monochromatic test and even-vertex deletion in the odd cores to obtain
$\mu\ge t+d-1$; the endpoint gives $t\ge a+2d+2X$.
Thus \eqref{eq:nonunit-guards} remains valid with $X_H=0$.
Also $h=P+g_2-1$ and $1\le P\le h\le\rho$.
The block operation and its deficits are applied to this current source, not transported from the former index.

At height two only $MF,MH$ remain; at height three only $MF,A$; above three only $MF$.
If column three has no strict $A$ debt, the matching join pays. Otherwise it gives $V$.
For $D=3$, $g\ge Q+V>0$ finishes.
For $D\ge4$, put $\ell=t-M+S$ and $T=M-t+1$.
If $\ell\ge1$, first choose $\ell-1$ exterior full-degree centers at the ordinary corner, then take all $P$ centers from one final complete red forest:

\[
 J\subseteq\{\max(4,\alpha+1),\ldots,D-1\},\qquad
 g\ge Q+V+(m+1)S+mP-\sum_Jf_j.
\]
If $\ell\le0$, direct column comparison gives
\[
 g\ge Q+V+mT-\sum_Jf_j,\qquad T\ge S+1.
\]
The top cannot be negative because $T\ge a+d\ge\rho+1>P$.
Hence again $J\subseteq\{4,\ldots,D-1\}$.
For $D=4$ there are no negative columns.
For $D=5,6$, their sum is at most $X$; use the first table row with $r=2,X_H=0,P=1$.
For $D=7$, one negative column uses the third row; two have sum at most $X$ since three is absent, so use the fifth row with $2P$ replaced by two; three or more use the fourth row with $mP$ replaced by $m$.
For $D=8$, use the last two rows with $2P,P$ replaced by two and one.
In the one-column case with roots, $\mu\ge a+2+9P\ge12$ gives margin at least $16/3>5$; without roots positivity suffices.
Equation \eqref{eq:binary-suffix-area} carries every fee back to the full original profile, with positive smaller rank.

For $D\ge9,b=2$, the next appendix supplies complete payment for the binary empty side.
All cases are therefore exhausted.
\end{proof}
\section{A binary empty side and unrestricted opposite heads}\label{app:binary-empty}
\begin{proposition}\label{prop:binary-empty}
An original strict minimum with $D\ge9$ cannot have
$\nn=(2^z,1^u)$, $z\ge1$, and $s=z+u\le\mu<t$.
\end{proposition}
\begin{proof}
Use the weaker two-budget transform $T$, with vertex budget $C=|R|=2\mu+a+\rho$ and slot budget $\mu$.
Write $L=|T(\mm)|$ and $c=\max T(\mm)$.
The complete original source and two-budget restoration give
$P_{r+a}(B;\nn,T(\mm))$.
As in Lemma \ref{lem:two-successors}, fix the root-edge colors, remove from the same complete copy the whole-star labels accommodated in the original cores, and control the remaining list by the first overflowing windows.
The empty-side embedding and endpoint tests give
\begin{equation}\label{eq:binary-empty-guards}
 k\ge s\ge r+a,\qquad \mu\ge s+d-1,\qquad
 t\ge a+\mu+d+r-1,\qquad \max m_i=D.
\end{equation}
\[
 g_j=c_j(\mm),\quad g=g_D\ge1,\quad
 h=z+g-r+1,\quad 1\le h\le\rho,\quad p=\floor{D/2}.
\]
\[
 2\mu\ge D\rho+\zeta+2\epsilon+\beta,\qquad
 (D-2)\rho+\zeta\le2p(p+1)-2.
\]
The last inequality follows from \eqref{eq:block-failure}.
Choose a nonempty actual set $U\subseteq R$, of size $\eta$, with
$\Delta(R-U)\le D-1$ and $1\le\eta\le\theta C$, where $\theta=4/13$.
The probability-$1/4$ construction with missed closed neighborhoods gives it because
$1/4+(3/4)^{D+1}<4/13$.
If the resulting set is empty, add one vertex; $C\ge D+1$ preserves the bound.
Color root edges and internal edges meeting $U$ red, and all other core edges blue.
At the original corner $(r-1,1)$, a blue head $D$ must occur in $B$.
A red original forest leaves at least $(s-r+2-a-\eta)_+$ matching edges in $B$.
Hence on the same $B$,

\begin{equation}\label{eq:binary-empty-repair}
 \kappa_j(B)\ge P_0:=(s-r+2-a-\eta)_+,\qquad1\le j\le D-1.
\end{equation}
When the positive part is zero, use nonnegativity and no empty-target arrow.
All subsequent estimates take columnwise maxima within the single bound
$e(B)\ge k+\sum_j\kappa_j(B)$.

\paragraph{A unit surviving list.}
Suppose $c=1$. Put $e=t-L\le\mu$, $b=\sum_i(m_i-1)$, and $c_0=\mu-s\ge d-1$.
The actual surviving source also gives
$\kappa_1(B)\ge z+L-r-a+1$.
Thus
$e(B)\ge k+z+L-r-a+1+(D-2)P_0$.
Write $b_j=g_j-g_{j+1}$. The complete original higher columns are
\begin{equation}\label{eq:binary-empty-profile}
 \Lambda_j=\max(s+g_{j+1}-r+1,z+g_j-r+1)
 =s+g_{j+1}-r+1+(b_j-u)_+\quad(1\le j<D),
 \qquad\Lambda_D=h.
\end{equation}
Consequently
$E_r=N+(D-1)(s-r+1)+b+\sum_{j<D}(b_j-u)_++h$.
Using $\sum(x_j-u)_+\le(\sum x_j-u)_+$, put
$Q_0=(g+u-t)_+\le(\rho-5)_+$ to obtain
\[
 E_r\le N+(D-2)s+t+b+2z-D(r-1)+Q_0.
\]
The entire nonunit blue prefix fits the slot budget, so
$b\le(2-1/p)\mu-\delta g$, where $\delta=0$ for even $D$ and $\delta=1-1/p$ for odd $D$.
This follows termwise from $j-1\le(2-1/p)\ceil{j/2}$, retaining the gap at a maximum head.
Put $v=s-a$ and $w=r-2$.
Then $P_0\ge(v-w-\theta(2\mu+a+\rho))_+$ and $v\le\mu-a$.
The same-graph degree identity gives
\[
\begin{aligned}
 e(G)-E_r\ge{}&(D-5+1/p)\mu+
 a-(D-3)v+(D-2)(v-w-\theta(2\mu+a+\rho))_+\\
 &+(D-1)(r-1)+\delta g-z-1
 +D\rho/2+(\beta-\zeta)/2-Q_0.
\end{aligned}
\]
The three middle terms involving $a,v$ are at least both
$(D-2)a-(D-3)\mu$ and
$[1-(D-3)\theta]a-(D-3)(2\theta\mu+\theta\rho+w)$.
The former increases and the latter decreases with $a$.
Their intersection
$a_*=[(1-2\theta)\mu-\theta\rho-w]/(1+\theta)$
gives the common lower bound $-(D-3)\mu+(D-2)a_*$, also valid if the intersection is negative.
Use $-z\ge-\rho-(r-1)+g$ to obtain
\begin{equation}\label{eq:binary-unit-pay}
 e(G)-E_r\ge A_D\mu+B_D\rho+(\beta-\zeta)/2+D-2+\delta-Q_0,
 \quad
 A_D=\frac{5(D-2)}{17}-2+\frac1p,\quad
 B_D=\frac{9(D-2)}{34}.
\end{equation}
For $D=9,10,11$, substitute
$2\mu\ge D\rho+\zeta$ and
$\zeta\le2p(p+1)-2-(D-2)\rho$.
The lower bounds are respectively
\[
 \frac{770\rho-732}{136}>0,\qquad
 \frac{567\rho-422}{85}-(\rho-5)_+>0,\qquad
 \frac{1314\rho+912}{170}-(\rho-5)_+>0.
\]
For $D=9$, $\rho\le5$, so $Q_0=0$.
For $D\ge12$, $A_D>1$; using $Q_0\le\rho$ gives the positive lower bound
$(A_D-1)\mu+(B_D+D/2-1)\rho+D-2+\delta$.

\paragraph{An arbitrary nonunit surviving list.}
Now let $c\ge2$.
Let $e_j$ be the largest number of cheapest entries in the original blue head-$j$ prefix fitting both budgets.
The first-overflow definition gives $g'_j=g_j-e_j$: the last $e_j$ entries of that prefix fit, and adding the next one fails.
Thus $e_j$ is nonincreasing and
\[
 e_1\le\mu,\quad e_2\le C/3,\quad e_j\le\mu/\ceil{j/2}\ (j\ge3).
\]
\[
 \sum_{y=c+1}^Dg_y\le\alpha_c\mu,\quad
 \alpha_c=\max_{c<j\le D}\frac{j-c}{\ceil{j/2}},\quad
 \alpha_c\le\alpha:=2-\frac2p,\quad
 \alpha_c\le\sum_{j=c+1}^{D}\frac1{\ceil{j/2}}.
\]

Put $v=s-r+1-a$, $X=a+v=s-r+1$, and $\tau=v-P_0=\min(v,\eta-1)$. Then
\[
0\le a\le X\le\mu,\qquad
0\le\tau\le X-a,\qquad
\tau\le\theta(2\mu+a+\rho)-1.
\]
Compare \eqref{eq:binary-empty-profile} columnwise with the same actual surviving source.
For $j<c$, both surviving prefixes are nonempty, so
$\Lambda_j-\kappa_j(B)\le a+e_j$.
For $j=c<D$, maximize the surviving-prefix bound and \eqref{eq:binary-empty-repair} to get
$\Lambda_c-\kappa_c(B)\le a+g_{c+1}+\max(e_c,\tau)$.
For $c<j<D$, obtain
$\Lambda_j-\kappa_j(B)\le a+\tau+g_{j+1}+(b_j-u)_+$.
The sum of the final positive parts is at most
$\Delta_0=(\mu/2-X+\rho-1)_+$, since $g_{c+1}\le\mu/2$ and $u=X+g-h$.
Thus, with $\mathcal L=\sum_{j<D}(\Lambda_j-\kappa_j(B))$,
\begin{equation}\label{eq:binary-affine}
\begin{gathered}
\mathcal L\le(D-1)a+\mu+C/3+
\max\{\alpha\mu+(D-3)\tau,\ F_0\mu+B_0\tau\}+\Delta_0,\\
F_0=(D-1)/3,\qquad B_0=(D-7)/3+2\alpha .
\end{gathered}
\end{equation}
To verify this common bound, for $c=2$ use $e_2,\tau<C/3$ directly.
For $3\le c<D$, put $x=\tau/\mu$; the remaining normalized expression is
\[
 f_c(x)=\sum_{j=3}^{c-1}\frac1{\ceil{j/2}}+\alpha_c+
 (D-c-1)x+\max\{1/\ceil{c/2},x\}.
\]
This function is convex, with $f_c(0)\le F_0$ and $f_c(1/2)\le\alpha+(D-3)/2$.
Use the endpoint chord on $[0,1/2]$; for $x\ge1/2$, bound each reciprocal by $x$.
These give the two affine bounds.
If $c=D$, there is no tail term and every column uses $a+e_j$, again satisfying \eqref{eq:binary-affine}.

Finally maximize only over the compatible actual parameter range. Put
\[
b_*=\mu/2+\rho-1,\qquad
T_*=\frac{12\mu+4\rho-13}{17},\qquad
U_*=(D-10/3)T_*+\rho/3.
\]
We have $b_*\le\mu$ and $T_*\ge\mu/2$.
The required optimization is
\begin{equation}\label{eq:binary-opt}
 C/3+\max\{\alpha\mu+(D-3)\tau,F_0\mu+B_0\tau\}
 +(b_*-X)_+-X\le\alpha\mu+U_*.
\end{equation}
Take either $(L_0,K)=(\alpha,D-3)$ or $(F_0,B_0)$.
If $X\ge b_*$, fix $X$ and maximize
$(2\mu+a+\rho)/3+K\tau-X$.
Put $X_0=\theta(2\mu+\rho)-1$.
For $X\le X_0$, the maximum is at most $(2\mu+\rho)/3+(K-1)X$.
For $X>X_0$, the two bounds on $\tau$ meet at $a=(X-X_0)/(1+\theta)$, giving maximum
\[
 (2\mu+\rho)/3+(K-1/3)\frac{4(X+2\mu+\rho)-13}{17}-2X/3.
\]
Both pieces increase with $X$; the second has slope $(12K-38)/51>0$.
Thus set $X=\mu$.
After adding $L_0\mu$, the second parameter pair minus the first is
$(F_0-\alpha)(\mu-2T_*)\le0$, proving the required bound.
If $X<b_*$, use $\tau\le X-a$ to bound the expression by
$(2\mu+\rho)/3+L_0\mu+(K-1)b_*$.
The second pair is again no larger, since the difference is
$-2(F_0-\alpha)(\rho-1)\le0$.
For the first pair, the remaining gap to the desired bound is
\[
\begin{aligned}
&(D-10/3)T_*-\bigl[2\mu/3+(D-4)(\mu/2+\rho-1)\bigr]\\
&\quad\ge
\frac{(21D^2-260D+656)\rho}{204}+\frac{12D-74}{51}>0.
\end{aligned}
\]
Here $\mu\ge D\rho/2$ was used.
The quadratic is seventeen at $D=9$, and its forward difference is $42D-239>0$.
This proves \eqref{eq:binary-opt}.
Substitute it and \eqref{eq:binary-affine} into the same complete edge account, using
$d\le\mu-s+1=\mu-X-r+2$:

\[
 e(G)-E_r\ge
 \left[\frac{5(D-9)}{17}+\frac2p\right]\mu+
 \frac{27D-56}{102}\rho+\frac{\beta-\zeta}{2}
 +\frac{13(3D-10)}{51}+r-2.
\]
For $D=9,10,11$, substitute $2\mu\ge D\rho+\zeta$ and the same upper bound on $\zeta$.
The results are respectively
$(35\rho-31)/6$,
$(1732\rho-962)/255$, and
$(2002\rho+1408)/255$, all positive.
For $D\ge12$, the coefficient of $\mu$ exceeds one, so it absorbs $-\zeta/2$ and all remaining terms are positive.
Thus every surviving height strictly pays the complete original profile, a contradiction.
\end{proof}
\section{Two completely embedded sides and dominating-set colorings}\label{app:guarded}
\begin{proposition}\label{prop:guarded}
For $D\ge6$, an original strict minimum cannot satisfy $\mu+d\ge\max(s,t)$.
\end{proposition}
\begin{proof}
Lemma \ref{lem:empty-embedding} puts both complete original forests in the same actual $R$ and gives

\[
 e(F),e(H)\le M+\rho-d,\qquad
 \mu\ge s+2X_F+d-1,\qquad
 \mu\ge t+2X_H+d-1.
\]
\[
 M=a+\mu,\quad k=\nu(B),\quad z=e(B)-k,\quad
 N=M+k+d,\quad
 e(G)=DM+\frac{D\rho+\beta-\zeta}{2}+k+z.
\]
Exchange colors so that $s\le t$.
If $n_1\le D-1$, group the diagonals so that each nonfirst entry is counted at most once:

\[
 E_r\le N+(D-3)t+e(F)+e(H)+h.
\]
Substitute the preceding bounds together with
$M\ge a+D\rho/2+(\zeta+\beta)/2+\epsilon$ and
$t\le M-a-d+1$ to obtain

\[
 e(G)-E_r\ge\frac{D-6}{2}(a+\rho)+(D-2)d-(D-3)>0.
\]
Henceforth $n_1=D$. The sharper profile bound needed is

\begin{equation}\label{eq:guarded-profile}
 E_r-h\le
 N+(D-2)t-s+e(F)+e(H)
 -(D+m_1-2)-(D-1)(r-2).
\end{equation}
For \eqref{eq:guarded-profile}, split the ordinary diagonals into
$2,\ldots,s+1$, $s+2,\ldots,t+1$, and $t+2,\ldots,s+t$.
The first group omits the first red entry, the middle group uses the red cap $D$, and the last omits the first blue entry.
Each nonfirst entry is assigned once.
After removal of the excess top plateau, each of the $r-2$ discarded initial diagonals still contributes at least $D-1$.
Separating the top width $h$ gives \eqref{eq:guarded-profile}.

We construct an actual set dominating every original full-degree vertex.
First take a maximum independent set $S_B$ of full-degree vertices in $B$, of size $b$.
If its number of root edges is $\beta_B$, deleting it and losing at most $b$ matching edges gives
$(D-1)b\le z+\beta_B$.
In $Y=A\cup R$, choose vertices with probability $p_0$ and add every original full-degree vertex still undominated.
Put $\theta=p_0+(1-p_0)^{D+1}$.
For a vertex with $b_v$ exterior neighbors, convexity bounds its miss probability by
$(1-p_0)^{D+1}+b_v/(D-1)$.
The total exterior incidence at vertices not dominated by $S_B$ is at most $\beta-\beta_B$.
Thus some actual set $S_0$ dominates all full-degree vertices and satisfies

\[
 |S_0|\le\theta(2M+\rho)+\frac{z+\beta}{D-1}.
\]
Put $w=\min(f_D,r-1)\ge1$.
Color all edges meeting $S_0$ blue and all other edges red.
The red degree is at most $D-1$.
The original legal corner $(w,r-w)$ forces a blue forest with $t-r+w+1$ components, so

\begin{equation}\label{eq:guard-dom}
 z\ge(D-1)(t-r+w+1)-(D-1)\theta(2M+\rho)-\beta.
\end{equation}
Put $\lambda=(D-1)(1-2\theta)$ and $K=(\lambda+1)/2$.
Substitute \eqref{eq:guard-dom} and \eqref{eq:guarded-profile} into the exact edge account:

\[
\begin{aligned}
 e(G)-E_r\ge{}&(\lambda-1)M+(K-2)\rho-(\zeta+\beta)/2\\
 &+k+2d+(D-1)w-h+r+m_1-2.
\end{aligned}
\]
For $D\ge8$, take $p_0=1/4$, so $\theta\le1/3$.
Substituting the lower bound for $M$ gives at least

\[
\begin{aligned}
&\frac{D-4}{3}a+\frac{D^2-3D-10}{6}\rho
+\frac{D-7}{6}(\zeta+\beta)+\frac{D-4}{3}\epsilon\\
&\hspace{20mm}+k+2d+(D-1)w-h+r+m_1-2>0.
\end{aligned}
\]
For $D=7$, use $\theta\le5/14$.
The bound is
$5a/7+13\rho/7-(\zeta+\beta)/7+5\epsilon/7
+k+2d+6w-h+r+m_1-2$.
Since $L_A\le5a$, it is at least
$13\rho/7+\epsilon+k+2d+6w-h+r+m_1-2>0$.

For $D=6$, $\theta\le2/5$, and the same account gives

\[
 e(G)-E_r\ge
 k+2d-\rho+5w-h+r+m_1-2-(\zeta+\beta)/2
 \ge9-L_A/2.
\]
The failed actual block-prefix bound gives $L_A\le22-4\rho$.
The lower bound is therefore at least $2\rho-2\ge0$.
A successful block prefix already gives paid descent.
Every case contradicts strict underpayment.
\end{proof}

\begin{corollary}[Local full-degree root--core interface]
\label{cor:local-root-core-interface}
Let $(G;\mathbf n,\mathbf m;r)$ be the original globally minimal strict
datum, after normalization.  Write $N=s+t-r+1$ and assume explicitly
\[
\begin{gathered}
 e(G)=E_r-1,\qquad \Delta(G)=D=\ell_r-1\ge6,\\
 d=N-\nu(G)\ge\rho_D(G)+1,\qquad
 H=\max(n_1,m_1)\le D-1.
\end{gathered}
\]
Put $h=L_{D+1}$.  Suppose
$G-A=\mathcal R\sqcup B$ is an actual disjoint union of induced
subgraphs, every vertex of $\mathcal R$ has degree $D$ in $G$, and
$\mathcal R$ has $a+\eta$ odd factor-critical components, where
$a=|A|$ and $\eta\ge h\ge1$.
Assume an actual matching assigns every vertex of $A$ to a distinct
component of $\mathcal R$.  Put $\mu=\nu(\mathcal R)$ and $k=\nu(B)$.
Then
\[
 \mu<\min(s,t),\qquad
 s,t\ge a+\mu+r+d-1,\qquad \max(s,t)\le k.
\]
No covering condition on the degree-$D$ vertices of $B$ is required,
and $\eta$ need not equal $\rho_D(G)$.
\end{corollary}
\begin{proof}
The root assignment, internal perfect matchings avoiding its ports,
and a maximum matching of $B$ give
\[
 \nu(G)=a+\mu+k,\qquad
 \nu(G-E(\mathcal R))=a+k.
\]
For the upper bounds, any matching uses at most $a$ edges meeting $A$;
the remaining edges contribute at most $\mu+k$, or at most $k$ after
the internal core edges are erased.  Any matching covers at most
$2\mu+a$ of the $2\mu+a+\eta$ full-degree core vertices.  Hence
$\eta\le\rho_D(G)$ and $d\ge\eta+1$.

Deleting all roots reduces both the rank and matching number by $a$,
so the positive deficit remains $d$.  The monochromatic endpoint
tests give $r+a\le\min(s,t)$.  The emptiness and source-transfer
arguments of Appendix~\ref{app:sources} apply to this decomposition:
internal edges are restored on the actual cores, components meeting
roots are discarded or shortened on actual root edges, and no edge
joins $\mathcal R$ to $B$.  Their matching capacities and half-order
and vertex budgets are exactly those above.  In particular the
robust embedding alternative for each whole original forest holds;
no step requires a matching covering the full-degree vertices of $B$.

\paragraph{The single-empty direction.}
Exchange the colors so that $t\le\mu<s$.
Use the fixed residual $L$ constructed in Appendix~\ref{app:sources},
and put
\[
\begin{gathered}
 X=\sum_i\lfloor(n_i-1)/2\rfloor,\qquad
 Y=\sum_j\lfloor(m_j-1)/2\rfloor,\\
 S=s-\mu,\quad T=\mu-t,\quad u=s-f_2,\qquad
 |L|=S+\delta,\quad e_0=\mu-\delta,\\
 v_2=c_2(L),\quad V=\mu-f_2,\quad A_*=V+v_2,\qquad
 V_0=\sum_C\lfloor(z_C+1)/3\rfloor.
\end{gathered}
\]
Here $z_C=\nu(C)$ and $q_+$ denotes the number of nontrivial cores.
Write
\[
 \zeta=Da-\sum_{x\in A}d_G(x),\quad
 \epsilon=e(G[A]),\quad\beta=e(A,B),\quad
 L_A=\zeta+2\epsilon+\beta.
\]
The actual root cut is $Da-L_A$, so the two actual edge fees are
\[
 c_Z=D\mu+\frac{D\eta+\zeta+2\epsilon+\beta}{2},\qquad
 c_B=D(a+\mu)+\frac{D\eta+\beta-\zeta}{2},
 \qquad Z=G-E(\mathcal R).
\]
Thus $C_0=(D\eta+\beta+\epsilon)/2-h\ge(D/2-1)\eta$.

A full-degree core of order $n\le D$ has cut at least
$n(D-n+1)\ge D$.  All root degrees sum to at most $Da$, so at least
$\eta$ cores have order greater than $D$.
There is therefore a core large enough for each single original head
to satisfy both numerical budgets.  The residual construction gives
$1\le e_0\le\mu$ and $\delta\ge0$, as well as the two fixed complete
sources
\[
 P_{r+a}(B;L,\mathbf m),\qquad P_r(Z;J,\mathbf m).
\]
Their actual ranks are $k+d+\delta$ and $a+k+d+\delta$,
positive and strictly smaller than $N$.
The matching deficit of the first source is $d+\delta$;
its monochromatic tests give $r+a\le\min(|L|,t)$.

The robust embedding of the blue forest gives
\[
 T\ge2Y+d-1,\qquad t+g_2+Y\le e(H)\le\mu+q_+-d.
\]
Take any optimal Tutte barrier $S_B$ of $B$.  Oddness of the cores and
the actual matching identity imply
\[
 o\bigl(G-(A\cup S_B)\bigr)-|A\cup S_B|
 =|G|-2\nu(G).
\]
Thus $A\cup S_B$ is an optimal barrier.  A globally maximum-cardinality
optimal barrier consequently has size at least $a$; no containment
property of that global maximum is used.
The original complete endpoint and the endpoint of the source on $B$
therefore give, respectively,
\[
 2X\le t-a-2d-r+2,\qquad
 S\ge a+r+2d+\delta+2Y-2.
\]

Since $D\ge6$, the $\eta$ large cores identified above have capacity
at least three.  Hence $\mu\ge3\eta$ and $V_0\ge\eta$.
Every nontrivial core of capacity $z$ and cut $\partial C$ satisfies
\[
 6z+\partial C\ge3D.
\]
For $2z\le D$, use $\partial C\ge(2z+1)(D-2z)$;
the resulting excess over $3D$ is $2(z-1)(D-2z)\ge0$.
For $2z>D$, use $6z>3D$.
Subtracting the $D(a+\eta-q_+)$ cut edges of the singleton cores
from the total root cut and summing gives
\[
 2Dq_+\le6\mu-D\eta-L_A\le D(\mu-\eta),
 \qquad q_+\le(\mu-\eta)/2.
\]
A capacity-one core has cut at least $3(D-2)\ge2D$, so there are
at most $a/2$ such cores.  Every other nontrivial core satisfies
$v(z)=\lfloor(z+1)/3\rfloor\ge1$ and $8v(z)\ge2z$.
Also $2v(z)\le z$ for all $z\ge0$.  Therefore
\[
 q_+\le a/2+V_0,\qquad 8V_0+a-2\mu\ge0,
 \qquad 2V_0\le\mu.
\]

The cheapest $e_0=\mu-\delta$ red entries satisfy both window budgets.
Their vertex cost is at least $3e_0-u$, and their half-order cost is
at least $e_0+(f_3-S-\delta)_+$.
The vertex budget is $2\mu+q_+\le2\mu+a+\eta$, so
\[
 3\delta\ge\mu-u-a-\eta,
 \qquad f_3\le S+2\delta.
\]
The height-two count of $U$ gives $v_2\ge(f_2-r_2)_+$, where
$r_2=\sum_C\lfloor(2z_C+1)/3\rfloor$.
The identity
$z-\lfloor(2z+1)/3\rfloor=\lfloor(z+1)/3\rfloor$
then gives
\[
 A_*\ge\mu-r_2=V_0.
\]
If $L$ consists of units, the nonunit prefix of the original red list
satisfies the half-order budget, and the height-two count of $U$ is zero.  Hence
$f_2+X\le\mu$, $V\ge V_0$, and $V\ge X$.
Every highest rectangle is bounded by $h\le\eta$.
Moreover, $H\le D-1$ and the original highest head sum at least $D+2$
make both original largest heads at least three, so $Y\ge1$.

These are all the structural, endpoint, and source inputs of
\eqref{eq:se-guards1}--\eqref{eq:se-guards4} and the subsequent column
comparisons in Appendix~\ref{app:single-below}, with $\eta$ replacing
$\rho$.  Its internal colorings use an even number of vertices in
each of these actual odd cores and the same actual graphs $B,Z$.
The resulting single-empty contradiction therefore holds here.

\paragraph{The two short sides.}
Suppose $\max(s,t)\le\mu+d$.
The robust embedding alternative embeds both whole forests in
$\mathcal R$.  Exchange colors so that $s\le t$, and write
$M=a+\mu$ and $z=e(B)-k\ge0$.
The embedding bounds give
\[
 e(F),e(H)\le M+\eta-d,\qquad t\le\mu-d+1.
\]
The short-head profile bound and the actual edge account are
\[
 E_r\le N+(D-3)t+e(F)+e(H)+h,
\]
\[
 e(G)=DM+\frac{D\eta+\beta-\zeta}{2}+k+z,
 \qquad N=M+k+d.
\]
Substitution yields
\[
e(G)-E_r\ge
(D-3)a+\frac{D-4}{2}\eta+(D-2)d-(D-3)-h
 +\frac{\beta-\zeta}{2}+z.
\]
Using $h\le\eta$, $\zeta\le Da$, and $\beta,z\ge0$, this is at least
\[
 \frac{D-6}{2}(a+\eta)+(D-2)d-(D-3)>0.
\]
This proves the short-head guarded exclusion without a covering
assumption on $B$.  If $\min(s,t)\le\mu$, either one side is empty
and the other nonempty, or both are at most $\mu$ and fall in this
guarded case.  Thus $\mu<\min(s,t)$.

\paragraph{Strong lengths.}
Apply Lemma~\ref{lem:length} to
$G-A=\mathcal R\sqcup B$ at row $r+a$; its matching deficit is
exactly $d$.  This gives $s,t\ge a+\mu+r+d-1$.
Finally $s+t=a+\mu+k+d+r-1$ gives $s,t\le k$.
\end{proof}

\section{Double-empty terminals in degrees three through five}\label{app:low-double}
\begin{proposition}\label{prop:low-double}
An original strict minimum cannot satisfy $D\in\{3,4,5\}$,
$\mu\ge\max(s,t)$, and $\max(n_1,m_1)=D$.
\end{proposition}
\begin{proof}
Keep $M=a+\mu$, $z=e(B)-k$, and $L_A=\zeta+2\epsilon+\beta$.
Appendix \ref{app:blocks} gives

\[
\begin{array}{c|c}
D&\text{necessary bounds}\\\hline
3&\rho\le2,\quad L_A\le2h+2-3\rho\\
4&\rho\le5,\quad L_A\le2h+10-4\rho\\
5&\rho\le3,\quad L_A\le10-3\rho\le7
\end{array}
\]
\[
L_A\le(D-2)a,\quad
e(F),e(H)\le M+\rho-d,\quad
\mu\ge s+2X_F+d-1,\quad\mu\ge t+2X_H+d-1.
\]
\begin{equation}\label{eq:low-double-guards}
\begin{gathered}
s-k\ge a+2d+2X_H+r-2,\qquad
t-k\ge a+2d+2X_F+r-2,\\
M-k\ge2a+3d+2X_F+2X_H+r-3.
\end{gathered}
\end{equation}

\paragraph{Actual unused vertices and a dominating set.}
For either complete forest $Q$ of length $l$, put $b=\mu-l+d$.
The original source forces $Q$ after deletion of any at most $b$ vertices of $R$.
Greedily choose $j=\floor{b/D}$ independent vertices, delete their at most $Dj$ neighbors, and extend to exactly $b$ deleted vertices while avoiding the chosen set; here $b\le\mu-1$.
The resulting $j$ isolated vertices cannot belong to a positive-edge forest.
Thus $M+\rho-d-e(Q)\ge j$, giving

\begin{equation}\label{eq:robust-mass}
 D e(Q)\le(D-1)\mu+l+D(a+\rho)-(D+1)d+D-1.
\end{equation}
If $Q$ contains a head $D$, the strong endpoint gives $b\ge D$, and we may also use
$e(Q)\le M+\rho-d-1$.

The following dominating set handles the original maximum-head branches in all three degrees.
Take a maximum full-degree independent set $S_B$ in $B$; as in Appendix \ref{app:guarded},
$(D-1)|S_B|\le z+\beta_B$.
The graph $Y=A\cup R$ has minimum degree at least two.
Take two copies of $Y$ and join corresponding internal degree-two vertices.
The resulting graph has minimum degree at least three.
Reed's domination theorem \cite{reed} gives a dominating set of size at most $3(2|Y|)/8$.
On restricting to the two copies, only degree-two vertices dominated through their twin may need to be added.
Let $c_2$ count those that have original full degree and are not dominated by $S_B$.
Then $(D-2)c_2\le\beta-\beta_B$, and across both copies at most $c_2$ additions are needed.
One copy therefore has an actual set of size at most $3|Y|/8+c_2/2$ dominating the required vertices.
Since $1/[2(D-2)]\le1/(D-1)$, adjoining $S_B$ gives an actual set $T$ with

\[
 |T|\le\frac38(2M+\rho)+\frac{z+\beta}{D-1}.
\]
If $n_1=D$, color every edge meeting $T$ blue.
The red degree falls to $D-1$, and the original corner $(1,r-1)$ gives

\begin{equation}\label{eq:low-dom}
 z\ge(D-1)(t-r+2)-\frac{3(D-1)}8(2M+\rho)-\beta.
\end{equation}
\paragraph{Degree three.}
Put $A=f_2,P=f_3,B=g_2,Q_3=g_3$.
Then $P+Q_3\ge1$ and $A,B>0$.
The exact columns are

\[
\begin{gathered}
E_r=N+L_2+L_3+h,\quad
L_2=\max(A+t,s+B)-r+1,\\
L_3=\max(P+t,A+B,s+Q_3)-r+1,\\
h=\max(0,P+B-r+1,A+Q_3-r+1).
\end{gathered}
\]
If the winner at $L_3$ is mixed, the complete profile is at most
$N+e(F)+e(H)-P-Q_3-2r+2+h$, hence

\[
 e(G)-E_r\ge d-\rho+P+Q_3+2r-3\ge3.
\]
If $L_3$ is a pure-side winner and $L_2$ is won by the opposite side, then $E_r\le2N+2h$.
Equation \eqref{eq:low-double-guards} gives

\[
e(G)-E_r\ge2a+d+2P+2Q_3+r-4+3\rho-3h\ge2.
\]
Otherwise both columns are won by the same side; exchange colors so that $P\ge1$.
Then $E_r=e(F)+3(t-r+1)+h$.
Put $H_0=A+P$.
Equations \eqref{eq:low-dom} and \eqref{eq:robust-mass} yield

\[
\begin{aligned}
e(G)-E_r
&\ge M/2-H_0-d+3\rho/4+2-h-(\zeta+\beta)/2\\
&\ge(s-P-a)/4+3/2-h-(\zeta+\beta)/2\\
&\ge k/4+Q_3/2+B/4+3/4-\rho/4+\epsilon>0.
\end{aligned}
\]
The last step uses $P+B-r+1\le h$, $\rho\le2$, and
$L_A\le2h+2-3\rho$.

\paragraph{Degree four.}
Exchange colors so that $n_1=4$.
Put $A_2=f_2,A_3=f_3,P=f_4$ and $B_2=g_2,B_3=g_3,Q_4=g_4$.
The degree-four two-factor coloring excludes a legal $3/3$ corner, so
$A_3+B_3\le r-1$ whenever $B_3>0$.
The exact columns are

\[
\begin{aligned}
L_2&=\max(A_2+t,s+B_2)-r+1,\\
L_3&=\max(A_3+t,A_2+B_2,s+B_3)-r+1,\\
L_4&=\max(P+t,A_3+B_2,A_2+B_3,s+Q_4)-r+1.
\end{aligned}
\]
By \eqref{eq:low-double-guards}, $A_3+B_2<P+t$.
If $Q_4>0$, the reversed argument also excludes $A_2+B_3$.
Thus the fourth column is a pure four-head term, except possibly $A_2+B_3$ when the blue head is three.
First suppose it is $P+t-r+1$.
If both lower columns are red winners, then $E_r=e(F)+4(t-r+1)+h$.
Put $V_0=A_2+A_3+P$.
Use the $3/4,1/4$ weighted sum of
$V_0\le M+\rho-d-1-s$ and $V_0\le2s+P$.
Equation \eqref{eq:low-dom} gives

\[
\begin{aligned}
e(G)-E_r
&\ge3M/4-V_0-d+7\rho/8+3-h-(\zeta+\beta)/2\\
&\ge(a+k+d+B_2+2B_3)/4+\rho/8+7/2-5h/4-(\zeta+\beta)/2\\
&\ge(a+k+B_2+2B_3)/4+\rho/8-5/4+\epsilon\ge-7/8.
\end{aligned}
\]
Integrality gives nonnegativity.
The complete profile bounds and actual remaining margins for the other five winner pairs follow; denote red, blue, and mixed winners by $R,B,A$.
Each row also uses
$e(G)\ge4M+2\rho+k-a$,
$M\ge s+A_2+A_3+P+d-\rho$, and
$M-s\ge a+2A_3+d-1$.

\begin{center}\small
\begin{longtable}{cP{.35\textwidth}P{.48\textwidth}}
\toprule
$(L_2,L_3)$&$E_r\le$&$e(G)-E_r\ge$\\\midrule
$(B,R)$&$3N+A_3-s+2h$&
$a+2d+4B_3+2r-4+A_2+P+\rho-2h$\\
$(R,A)$&$3N+2A_2-2s+2h$&
$a+2d+4B_3+2r-4+(s-A_2)+A_3+P+\rho-2h$\\
$(B,A)$&$2N+A_2-P+3h$&
$M+d+2B_3+r-2+A_3+2P+\rho-3h$\\
$(R,B)$&$3N+A_2-s+h$&
$a+2d+4B_3+2r-4+A_3+P+\rho-h$\\
$(B,B)$&$2N+s-A_3+2h$&
$2a+2d+5A_3+2B_3+r-4+2\rho-2h$\\
\bottomrule
\end{longtable}\end{center}

Every row is positive by $d\ge\rho+1$ and $h\le\rho$.
A blue four-head winner is handled by exchanging colors.
In the last branch the fourth column is $A_2+B_3-r+1$ and the blue maximum head is three.
If the third column is red, $A_3+B_3\le r-1$ bounds the sum of the third and fourth columns by $A_2+t-r+1$, reducing to the second or third row above.
If the third column is mixed or blue, either second-column winner gives

\[
E_r\le N+e(F)+e(H)+s-A_3-P-3r+3+h,
\]
\[
e(G)-E_r\ge M+d-s+A_3+P+3r-3-h-a
\ge2d+3A_3+P+3r-4-h>0.
\]
\paragraph{Degree five.}
Exchange colors so that $s\le t$.
If $n_1\le4$, put $C_0=n_1+m_1-2+4(r-2)\ge0$.
The same diagonal grouping gives
$E_r\le N+2t+e(F)+e(H)-C_0+h$, and therefore

\[
\begin{aligned}
e(G)-E_r
&\ge2a+3d-2+\rho/2+C_0-h+(\beta-\zeta)/2\\
&\ge2a+5\rho/2+1+C_0-L_A/2\\
&\ge2a+4\rho-4+C_0\ge0.
\end{aligned}
\]
If $n_1=5$, use \eqref{eq:guarded-profile} and \eqref{eq:low-dom}.
The coefficient of $M$ cancels exactly, giving

\[
e(G)-E_r
\ge k+2d-\rho+4+r+m_1-2-h-(\zeta+\beta)/2
\ge8-L_A/2\ge9/2>0.
\]
Every winner configuration in all three degrees is contradictory.
All deletions, dominating sets, and residual graphs belong to the same original host, and exterior surplus is charged once.
\end{proof}
\section{Completion of the intrinsic-region reduction}\label{app:base-gluing}
\begin{proof}[Proof of Proposition \ref{prop:base}]
All data use the same minimum order $(N,e,|V|)$.
For either original list $Q$ of length $l$, Lemma \ref{lem:empty-embedding} gives

\[
l\le\mu
\ \Longleftrightarrow\ Q\subseteq R
\ \Longleftrightarrow\ T(Q)=U(Q)=L(Q)=\varnothing .
\]
An actual complete copy satisfies both budgets and every height-capacity bound; conversely, emptiness of $U$ already gives $l\le\mu$ at height one.
No equivalence of unit/nonunit types for different transforms is needed.

Suppose one side is empty and the other nonempty; exchange colors so that $t\le\mu<s$.
If $m_1=D$, Lemma \ref{lem:reverse-owncap} excludes it.
If $n_1=D>m_1$, Propositions \ref{prop:unit-single} and \ref{prop:nonunit-single} treat the unit and nonunit types of the same current $L(\nn)$.
If both original heads are below $D$ and $D\ge6$, use Proposition \ref{prop:single-below}.

The following three low-degree arguments do not require a single-empty hypothesis.
For $D=3$, two heads at most two cannot attain $\ell_r=4$.
For $D=4$, an attaining pair must be $3/3$, but the degree-four two-factor edge partition avoids both three-heads.
For $D=5$, choose an attaining pair on the last height-six diagonal, oriented as $4/3$.
Its original prefixes satisfy $f_4+g_3-r+1=h$: a larger sum would create a later height-six diagonal.
Choose a legal pair $i+j=r$ in these actual prefixes and weaken to
$P=f_4-i+1$ four-heads and $Q=g_3-j+1$ three-heads, with $P+Q-1=h$.
Lemma \ref{lem:signed-exposure} on the same $G$ gives

\[
\rho\ge h+2\ceil{h/2}\ge h+2.
\]
However, \eqref{eq:block-failure} gives
$5\rho+L_A+2W_-\le2h+10$, hence
$3\rho+L_A+2W_-\le6$, contradicting $\rho\ge h+2\ge3$.
Thus every single-empty case is excluded.
At the same time, every $D=3,4,5$ case with both original heads below $D$ has been excluded.

Now suppose both sides are nonempty and $n_1=D$.
Reverse the colors in Lemma \ref{lem:two-successors} to obtain
$P_{r+a}(B;\nn,L(\mm))$ on the same $B$.
Write $|L(\mm)|=t-\mu+\delta_m$; its rank is $k+d+\delta_m$.
Apply the strong endpoint inequality to the endpoint retaining the entire original $\nn$:

\[
 |L(\mm)|-r-a+2\ge2(d+\delta_m+X_F).
\]
Consequently $a+\mu\le t-r+1$.
Lemma \ref{lem:persistent-region}, on the persistent original region, gives
$e(G)\ge E_r+\rho-h\ge E_r$, a contradiction.
The reversed argument excludes $m_1=D$.

For $D\ge6$, Proposition \ref{prop:guarded} excludes two empty sides.
After excluding one empty side, we therefore have $\mu<\min(s,t)$, and both original heads are at most $D-1$.
For $3\le D\le5$, the low-degree argument forces one original head to equal $D$, whereas two nonempty sides cannot have such a head.
Thus only two empty sides remain, and Proposition \ref{prop:low-double} excludes these three final degrees.
Together with the $D\le2$ argument in Appendix \ref{app:binary}, this proves \eqref{eq:base}.
\end{proof}
\section{Stronger prices for successive degree reductions}\label{app:stage-prices}
Throughout, $N_0$ is the rank of the same original minimum counterexample.

\begin{lemma}\label{lem:strong-odd-price}
Suppose \(j\ge2\), \(P_j(X)\), \(1\le N_X<N_0\),
\(\Delta(X)\le k=2p-1\), \(p\ge4\), \(\ell_j\le k+1\), and
\(0\le v=e(X)-E_j\le2pj-2\). Then

\[
\rho_k(X)\le
\left\lfloor\frac{2p(p-1)+2v}{2p-3}\right\rfloor.
\]
\end{lemma}\begin{proof}
Put \(\rho=\rho_k(X)\), and suppose
\((2p-3)\rho>2p(p-1)+2v\).
In the actual ordinary matching region of \(X\), write \(a\) for the number of roots and \(\mu\) for the matching number of the core union, and decompose the factor-critical cores into their blocks. Each positive block has capacity \(z_Q\ge p\) and weight
\(1\le w_Q\le p-1\), and their total weight satisfies
\(W_+\ge(p-\tfrac12)\rho\).
The assumed inequality implies
\[
W_+-(\rho+v+1)>
p(p-1)-1>0.
\]
Order the positive blocks by nondecreasing capacity and take the first prefix whose weight reaches \(\rho+v+1\). Write its total capacity, largest capacity, number of blocks, and weight as \(z,Z,l,w\), respectively. Then
\[
\rho+v+1\le w\le\rho+v+p-1,\qquad l\le\rho+v+1,
\]
and
\[
W_+-w\ge(p-\tfrac32)\rho-v-p+1>(p-1)^2.
\]
There are therefore at least \(p\) unselected positive blocks, all of capacity at least \(Z\).

Exclude every vertex of the selected blocks. The remaining original degree-\(k\) vertices form a set \(V_0\) with
\[
|V_0|\ge2(\mu-z)+a+\rho-l\ge2pZ-v-1.
\]
Set \(b=(Z-j+1)_+\). If \(b>0\), the small-surplus hypothesis gives
\[
|V_0|\ge2p(b-1)+1.
\]
Since a closed neighborhood has at most \(2p\) vertices, the greedy independent-set procedure supplies \(b\) mutually nonadjacent vertices \(U\subseteq V_0\). If \(b=0\), take \(U=\varnothing\).
Erase the selected block edges and delete \(U\). These edge sets are disjoint, and the exact fee is \(kK+w\), where \(K=z+b\ge1\).

Let \(h\) be the width at height \(k+1\) in the current source, or zero if that height is absent. Its uniform consequence gives \(h\le\rho\). If \(K\ge N_X\), then
\[
e(X)=E_j+v\le kN_X+h+v
\le kK+h+v<kK+w\le e(X),
\]
a contradiction. Hence \(1\le K<N_X\).
The capacity menu has maximum gap at most \(Z<j+b\). After deleting \(U\), the complete menu lemma therefore gives every new legal pair of \(P_{j+K}(Y)\) on the same actual residual graph and the same original lists. The complete profile loss is at most \(kK+h\), so
\[
e(Y)-E_{j+K}\le v+h-w\le-1.
\]
The new rank is positive and strictly smaller than \(N_X<N_0\), a contradiction. No sign assumption on \(N_X-\nu(X)\) was used.
\end{proof}
\begin{lemma}\label{lem:small-even-price}
For an actual positive smaller-rank complete source $P_j(X)$ with
$\Delta(X)\le2s$, $\ell_j\le2s+1$, $s\ge2$, and
$0\le v=e(X)-E_j\le3s$, one has

\[
\rho_{2s}(X)\le
\max\left\{v,\floor{\frac{s^2+v}{s-1}}\right\}.
\]
\end{lemma}
\begin{proof}
The reserve handles nonpositive matching deficit.
For positive deficit, let $h\le\rho$ be the highest-column width.
If $W_+\ge h+v+s^2+1$, take the capacity-ordered prefix first reaching
weight $h+v+1$. It has $w\le h+v+s$ and $l\le h+v+1$;
at least $s$ positive blocks remain. Thus $|V_0|\ge2sZ-v-1$.
We need $b=(Z-j+1)_+\le Z-1$ independent full-degree vertices.

If $Z=s$, then $2s^2-3s-1>(2s+1)(s-2)$, and greedy selection gives
$s-1$ vertices. If $Z>s$, the induced graph $X[V_0]$ has no
$K_{2s+1}$ component: every vertex of such a component already has
its full degree internally, making it an isolated original positive block
of capacity $s$, selected earlier in the order.
Brooks' theorem gives $2s$ colors; since $|V_0|>2s(Z-2)$,
one color class contains $Z-1$ vertices.
Take $b$ of them and use the preceding physical operation.
Its exact cost proves positive rank and gives
$e(Y)-E_{j+K}\le v+h-w\le-1$.
Therefore $W_+\le h+v+s^2$. Combining $s\rho\le W_+$ with $h\le\rho$
proves the assertion.

\end{proof}
\section{A unified descent for actual isolated regular components}\label{app:regular-components}
\begin{proposition}\label{prop:regular-components}
Let the original minimum strict source have $D=2p\ge8$, $H\le D-1$,
$R=r-1$, and highest-column width $\lambda$. Neither of the following is possible:
\begin{enumerate}
\item $\lambda=1$, and the original graph has an isolated connected
$D$-regular component without a perfect matching;
\item $1\le\lambda\le p$, and it has two isolated connected odd
$D$-regular components, where $|C_i|=2P_i+1$,
$J=\sum_i(P_i-\nu(C_i))$, $P=P_1+P_2$, and
$s,t\ge P+R+d-J$. Near-perfect matchings are not required.
\end{enumerate}
In the second case the length condition is automatic if an actual isolated
union $U$ of regular components contains both components and satisfies
$|U|-2\nu(U)\ge\lambda$. The root--core argument is given below;
an order-based capacity is not asserted to be an existing matching.
\end{proposition}

\begin{proof}
\subsection*{The common physical graph and strong lengths}
In the first case choose the actual component $C$ and a maximum-cardinality
optimal Tutte barrier $A$ inside it. Put
$a=|A|$, $\mu_C=\nu(C-A)$, and
$\rho_C=|C|-2\nu(C)\ge1=\lambda$.
This is an actual full-degree root--core decomposition:
the components of $C-A$ are factor-critical,
$\nu(C)=a+\mu_C$, and the exterior is $O=G-C$.
Here the hypotheses of Corollary \ref{cor:local-root-core-interface}
hold with $\mathcal R=C-A$, exterior $B=O$, and $\eta=\rho_C\ge\lambda$.
The normalization is that of the original globally minimal strict source;
all core vertices have original degree $D$, and the chosen maximum matching
assigns the roots to distinct cores. The corollary therefore gives
$\mu_C<\min(s,t)$ and $s,t\ge a+\mu_C+R+d=\nu(C)+R+d$.
No covering condition on the full-degree vertices of $O$ is needed.

Use the following common notation for the two cases:
\[
\begin{array}{c|c|c|c|c|c}
 &W& |W|&\nu(W)&d_*&\tau\\\hline
\mathrm I&C&2P+\varepsilon&P-j&d-j&(P-R)_+\\
\mathrm {II}&C_1\sqcup C_2&2P+2&P-J&d-J&(\max(P_1,P_2)-R)_+
\end{array}
\]

In case I, take $\varepsilon=1$ for odd order and $\varepsilon=2$ for even order;
$P\ge p$ and $j\ge0$. In case II, $P=P_1+P_2$, $P_i\ge p$, and $\varepsilon=2$.
In both cases,

\begin{equation}\label{eq:regular-account}
 e(W)=DP+p\varepsilon,\quad
 \nu(O)=N-P-d_*,\quad s,t\ge P+R+d_*,\quad d_*\ge2 .
\end{equation}
In I, the last inequality follows from
$d\ge\rho_D(G)+1\ge2j+\varepsilon+1$.
In II, $\rho_D(G)\ge\operatorname{def}(W)=2J+2$, so $d_*\ge J+3$.
For the automatic entry in the proposition, take a maximum-cardinality
optimal Tutte barrier inside $U$. Its odd components are factor-critical,
the roots have an actual matching to distinct components, and the
component--root excess is $\operatorname{def}(U)\ge\lambda$.
Corollary \ref{cor:local-root-core-interface}, now with
$\mathcal R=U-A$, $B=G-U$, and $\eta=\operatorname{def}(U)\ge\lambda$,
applies to this same original normalized source and gives
$s,t\ge\nu(U)+R+d\ge P+R+d-J$.
Each $P_i$ is only a conservative upper bound on the number of forest
components consumed in that monochromatic component. If fewer are
consumed, discard additional earlier components to obtain the required full suffix.

\subsection*{One exterior deletion history and common column taxes}
Fix one maximum-degree deletion history of $O$.
Let $C_h$ count deletions made at current degree at least $h$, and put $b=C_D$.
Define

\[
M_h=\begin{cases}\floor{(2P+\varepsilon)/(h+1)},&\mathrm I,\\
\sum_{i=1}^2\floor{(2P_i+1)/(h+1)},&\mathrm {II},\end{cases}
\quad
A_h=\max\{f_h+t-R:f_h>0;\ g_h+s-R:g_h>0\}.
\]
Coloring $W$ blue and using the original corner $(1,R)$ gives
$O\to(F_1,q_FK_2)$, where $q_F=t-R+1-P\ge d_*+1$;
interchange colors for the other endpoint. Their matching deficit is $d_*$, so

\begin{equation}\label{eq:regular-inventory}
q_F\ge2(d_*+\omega(F_1)),\qquad
C_h\ge q_F+f_h-1,\qquad C_h\ge A_h-P
\end{equation}
whenever the indicated head is present.
For the second inequality, use the first $q_F-1$ vertices as a cover
for the matching color; the remaining complete $F_1$ contains $f_h$
vertex-disjoint large stars.

Every column $2\le h\le D-1$ has tax at most $P$ on the same $O$.
A unit axis uses $\nu(W)\le P$.
For $u,v\ge2$, survival after monochromatic restoration costs at most $M_2\le P$.
If both counts fail to survive, the actual vertex-cover tax is at most
$\min(f_u+g_v-R,\,2P+\varepsilon-uf_u,\,2P+\varepsilon-vg_v)$.
If this were at least $P+1$, then $f_u+g_v\ge P+2$, whereas
$uf_u,vg_v\le P+\varepsilon-1\le P+1$, a contradiction.

The heads in column $D$ have opposite parity.
For even red head $u$, color $u/2$ two-factors in each regular component red
and the others blue; the blue degree is below its head.
If the red count survives, the tax is at most $M_u$.
If it does not and there is a larger original head on that side,
the corresponding two odd heads admit two even degree caps,
forcing the other count to be at most $R-1$.
If $u$ is the original maximum, the highest width gives $f_u\le R+\lambda-1$.
Monochromatic restoration on the other side then gives tax $M_v$
if it survives, and width at most $M_v+\lambda-1$ otherwise.
Thus column $D$ has tax at most $M_z+\lambda-1$ when both heads are at least $z$.

In I the highest-column tax is at most $\lambda=1$.
In II use the highest-column tax $\lambda$: its original width is
$\lambda$, so the required exterior lower bound is zero and always holds.
No single-defect internal coloring is needed.
With $\theta=1$ in I and $\theta=\lambda$ in II, the common exterior account gives

\begin{equation}\label{eq:regular-tax}
d_*\ge P-M_z+p\varepsilon-\lambda+2-\theta,
\quad
z=\begin{cases}2,&H=D-1,\\3,&H\le D-2.\end{cases}
\end{equation}
If another low column also has tax $M_2$, add $P-M_2$ to the right side.

\subsection*{The complete maximum-degree stage}
Every odd regular component of order $2P_i+1$ has a maximal independent
set of size at most $P_i-1$.
If its independence number is less than $P_i$, take a maximum independent set.
Otherwise the complementary $P_i+1$ vertices induce exactly $p$ edges.
Choose an inclusion-minimal vertex cover of size in $[2,p]$:
if a minimum cover is a singleton, use all $p$ leaves of that star instead.
The complement of this cover within those $P_i+1$ vertices is an
independent dominating set of size at most $P_i-1$.
In the even-order case I, $\alpha(C)=P+1$ would make $C$ regular bipartite
and hence perfectly matchable, so $\alpha(C)\le P$.
Every independent set $I$ also satisfies $|I|\le|C|-D$.
Combine these internal sets with all $b$ full-degree vertices from the exterior history.
The resulting original full-degree independent set is saturated by a maximum matching,
so its size $K_1$ is less than $N$.
Deletion leaves degree cap $D-1$ and complete-source surplus at most $\lambda-1$.
Lemma \ref{lem:strong-odd-price} yields

\[
B_p(\lambda)=\max\left\{\lambda-1,
\floor{\frac{2p(p-1)+2\lambda-2}{2p-3}}\right\},
\quad B_p(1)=p,\quad B_p(\lambda)+\lambda\le2p+2\ (1\le\lambda\le p).
\]
For the last bound, the difference after clearing the denominator
is at least $3p-7>0$. Consequently,

\begin{equation}\label{eq:regular-first-stage}
L_D-b\le
\begin{cases}P+p+\varepsilon-2,&\mathrm I,\\
P-2+B_p(\lambda),&\mathrm {II}.
\end{cases}
\end{equation}
If $H\le D-2$, the even head in column $D$ is at least four.
Signed exposure charges two more units; combining this with $|I|\le|C|-D$ gives

\begin{equation}\label{eq:regular-first-strong}
L_D-b\le
\begin{cases}
\min(P+p+\varepsilon-4,\ 2P+\varepsilon-p-2),&\mathrm I,\\
\min(P+p-2,\ 2P-3p+2),&\mathrm {II}.
\end{cases}
\end{equation}
\subsection*{The fixed successor and the largest-head case}
If $\tau=0$, the monochromatic component-capacity menu has gaps less than $r$,
and its $DP+p\varepsilon$ edges directly pay the bound $DP+\lambda$.
Suppose $\tau>0$. Once their existence is proved, take the first $\tau$
vertices $T$ of the same exterior history and fix
$Y=G-(W\cup T)$, $K=P+\tau$.
The menu is $\{0,P\}$ in I, where a matching-capacity upper bound $P$ suffices,
and $\{0,P_1,P_2,P\}$ in II. Its largest gap is less than $r+\tau$,
giving every corner of $P_{r+K}(Y)$.
Equation \eqref{eq:regular-account} gives $1\le N-K<N$.

First let $H=D-1$. If $C_{D-1}\ge\tau$, the exterior vertices exist
and their income is at least $(D-1)K+P+p\varepsilon+b$
(use $DK+p\varepsilon$ directly if $b\ge\tau$).
By \eqref{eq:regular-first-stage}, this pays
$(D-1)K+L_D+\lambda$: the remaining margin is
$(\varepsilon-1)(p-1)$ in I and $2p+2-B_p(\lambda)-\lambda$ in II.

If $C_{D-1}<\tau$, \eqref{eq:regular-inventory} forces $P>2d_*$.
Suppose a low column is attained by two nonunit heads.
If both counts are at most $M_2$, \eqref{eq:regular-tax} and
$3M_2\le2P+2$ give $L_h\le2M_2-R<P+d_*+1\le A_h$, a contradiction.
Thus one side survives and this column's tax drops to $M_2$.
We obtain $d_*\ge2(P-M_2)+p\varepsilon>P/2$ in I and
$d_*\ge2(P-M_2)+2p-2\lambda+2>P/2$ in II, again a contradiction;
the last step uses $M_2\le(2P+2)/3$ and $2p-2\lambda+2\ge2$.
All low columns are therefore unit axes, and the actual edge history gives

\[
e(G)\ge DP+p\varepsilon+\sum_{h=1}^{D-1}(A_h-P)+b
=\sum_{h=1}^{D-1}A_h+P+p\varepsilon+b
\ge\sum_{h=1}^{D-1}A_h+L_D+\lambda=E_r.
\]
\subsection*{Existence of inventory and the penultimate cap}
Now $H\le D-2$. Use $z=3$ in \eqref{eq:regular-tax},
with $M_3\le P/2+\varepsilon/4$ in I and $M_3\le P/2$ in II.
Then $C_H\ge2d_*>\tau$, so the required exterior vertices exist.
Let $a$ be the degree at the last of these deletions; $H\le a\le D$.
The preceding cost argument handles $a\ge D-1$.
Suppose $a=D-2$ and put $c=C_{D-1}<\tau$.
Continue the complete $D-1$ stage in the same first-stage graph.
Equation \eqref{eq:regular-first-strong} bounds the resulting surplus
by $p-2$ in I and $2p-1$ in II.
Lemma \ref{lem:small-even-price} bounds the exposure of the same exterior
$O_2=O-\{\text{first }c\text{ vertices}\}$ by $p+1$ and $2p$, respectively.

Rank exhaustion is included. In I the remaining edge count is at most
$(D-1)(N-K_2)+p-2$, so $K_2>N$ is impossible and $K_2=N$ leaves no degree-$D-2$ vertex.
In II it is at most $(D-1)(N-K_2)+2p-1$.
For $K_2=N+1$ there are no edges; for $K_2=N$ there are at most $(D-2)+1$ edges,
hence at most two full-degree vertices, which can be covered simultaneously
by a matching. Augmentation preserves that coverage. Both boundary exposures are zero.

If both sides of column $D-1$ fail to survive, then
$u,v\ge2$, $u+v=D$, and $M_u+M_v\le P$ give $L_{D-1}\le P-R$.
Substitution of \eqref{eq:regular-first-strong} pays immediately.
Otherwise $L_{D-1}-c\le M_2+p+1$ in I, and at most $M_2+2p$ in II.
Column $D$ on the same first-stage exterior also gives

\[
L_D-b\le
\begin{cases}
\max(M_3,M_4+p-2),&\mathrm I,\\
M_3+p,&\mathrm {II}.
\end{cases}
\]
In I a positive platform with even head at least $p+2$ has signed
exposure greater than $p$ and cannot survive.
The other live platforms lose at most $M_4$ and have width at most $p-2$.
In II a live platform has width at most $p$; a dead platform has
original width at most $M_3+\lambda-1$. These alternatives are maximized, not added.
The exact cost gives

\[
q:=e(Y)-E_{r+K}\le
L_{D-1}+L_D-2P-p\varepsilon-b-c+\lambda-1.
\]
In I the right side is at most

\[
M_2+p+1-2P-p\varepsilon+
\min\{2P+\varepsilon-p-2,\ \max(M_3,M_4+p-2)\}.
\]
For $\varepsilon=1$ and $M_2<p$, the first term in the minimum gives $M_2-p<0$.
If $M_2\ge p$, then $P\ge(3p-1)/2$; the other two terms give
$M_2+M_3+1-2P<0$ and
$M_2+M_4+p-1-2P\le-2p/5<0$.
For $\varepsilon=2$, the two terms are at most
$(13-5P-6p)/6<0$ and $(1-14P)/15<0$.
In II the common upper bound is

\[
M_2+p-1-2P+\min(2P-3p+2,M_3+p).
\]
If $M_2\le2p-2$, the first term gives at most $-1$.
Otherwise $P\ge3p-2$, and the second gives
$2p-1/3-5P/6\le-p/2+4/3<0$.

\subsection*{Common elimination for every lower exterior cap}
Let $\gamma=D-a\ge3$.
A positive highest corner must have two even heads, so $a\ge H\ge6$;
if $\gamma=3$, necessarily $p\ge5$.
If the actual first height on $Y$ exceeds $a+1$, a degree-cap avoiding
coloring already contradicts that complete source.
Otherwise assume $q\ge0$ and use the arbitrary-surplus price
$B(a)=\floor{a/2}+1\le p-1$.
At the successive cuts of the same exterior history,
$L_h\le P+C_{h-1}$ for $a+2\le h\le D-1$, and
$L_D\le M_{\gamma+1}+\lambda-1+C_{D-1}$.
Failure of one of these inequalities would, on restoring the corresponding
internal coloring, give a positive height-$h$ source on a graph of cap $h-2$.
These cuts supply column bounds for one final cost only.
Cancellation, with $L_0=L_{a+1}$, gives

\begin{equation}\label{eq:regular-cancel}
q\le L_0-b-2P+M_{\gamma+1}-p\varepsilon+2\lambda-2.
\end{equation}
If both counts attaining $L_0$ fail to survive, their head sum is
$a+2\ge8$, so $M_u+M_v\le P$ and $L_0\le P-R$;
the upper bound is negative.
Otherwise actual monochromatic restoration gives a platform
of width at least $L_0-M_2-\tau$ on the same $Y$.
Use the exposure price if this width is positive and $q\ge0$ otherwise:

\begin{equation}\label{eq:regular-live}
q\ge L_0-M_2-\tau-B(a).
\end{equation}
For $\gamma\ge4$, case I has $M_2+M_{\gamma+1}\le P+\varepsilon-1$,
and case II has this sum at most $P$.
Comparing \eqref{eq:regular-cancel} and \eqref{eq:regular-live},
the purported nonnegative difference is at most $-R-1$ for I with $\varepsilon=1$,
at most $-R-p$ for I with $\varepsilon=2$,
and at most $-R-3$ for II, using $\tau\le P-p-R$ and $\lambda\le p$.

For $\gamma=3$, the only middle column is $L_1=L_{D-1}$,
whose heads satisfy $u+v=D\ge10$ and $u,v\ge3$.
If one side survives, its tax drops to $M_3$ and cancellation gives
$q\le L_0-b+M_3+M_4-3P-p\varepsilon+2\lambda-2$.
If neither side survives, $L_1\le M_u+M_v-R$, and the bound becomes
$q\le L_0-b+M_u+M_v+M_4-3P-p\varepsilon+2\lambda-2-R-C_{D-2}$.
The required capacity inequalities are, respectively,

\[
M_2+M_3+M_4\le2P,\qquad M_2+M_4+M_u+M_v\le2P.
\]
The reciprocal sums are $47/60$ and at most $109/120$.
For I, integer rounding gives total capacity at most $n-1$ for
$n=2P+1$ and at most $n-2$ for $n=2P+2\ge12$.
For II, round separately in each component of order $2P_i+1$ and add.
Comparison with \eqref{eq:regular-live} is therefore again strictly negative.
Every exterior cap is covered, and the fixed successor has positive smaller rank
and $e(Y)<E_{r+K}$, a contradiction.
\end{proof}

\section{Complete platform splits at head gaps six and seven}\label{app:gaps67}
Use the all-root graph $Y=(G-A_0)-E(C)$, and write
$R=r+a-1$, $n=2Z+c$, and $M_j=\floor{n/(j+1)}$.
The root debt and cost are \eqref{eq:rootdebt}--\eqref{eq:Ycost};
no further port income is added.
The low-deficit estimate proved in Proposition \ref{prop:largegap-low}
does not use a platform upper bound when $\gamma\ge7$.
Together with Lemma \ref{lem:gap6}, it gives
$d>Z/2$ and $C_B(Y)>Z$, unless the original source is already avoided.

\begin{lemma}\label{lem:clip}
For $Z,k,L,t\ge0$,

\[
\min(Z+k,L)-\min(k,(L-t)_+)\le\max(Z,t).
\]
\end{lemma}
\begin{proof}
Split into $L\le t$, $t<L\le k+t$, and $L>k+t$.
The left side is at most $t,t,Z$, respectively.
\end{proof}
\begin{proposition}\label{prop:gaps67}
The original minimum strict source cannot have $\gamma=6$ or $7$.
\end{proposition}
\begin{proof}
First let $\gamma=7$. The original head-sum guard gives $B\ge9$.
If $B=9$, both original maximum heads are nine and caps $8/8$ give an avoiding coloring.
Thus $B\ge10$, $D\ge17$, $n/Z\le36/17$, and $Z\ge8c$.

\subsection*{The low and middle platforms at gap seven}
If $L_{B+3}^0\le Z$, use the actual source widths
$w_h=(L_h^0-\tau_h)_+$ for $B+3\le h\le D+1$ and put

\[
k=\max\left((Z-R)_+,\max_{B+3\le h\le D+1}w_h\right)\le Z.
\]
The same history has the first $k$ vertices of current degree at least $B$.
After deletion, $P_{r+a+Z+k}(X;\nn,\mm)$ has rank $\nu(Y)+d-k>0$.
The difference between the full original tax and actual high-stage income is

\[
\begin{split}
e(X)-E_{\rm new}&\le W-(D-1)\rho-1+T,\\
T&=-7Z+\min(Z+k,L_{B+1}^0)\\
&\quad+\sum_{h=B+2}^{D+1}
\bigl[\min(Z+k,L_h^0)-\min(k,C_{h-1}(Y))\bigr].
\end{split}
\]
The omitted column $B+2$ has
$\tau_{B+2}\le n(1/4+1/(B+1))\le135Z/187<Z$,
so Lemma \ref{lem:clip} charges $Z$.
Every other column pays $\eta_{h-B}$ because $k\ge w_h$,
and the first high column pays at most $2Z$. Therefore

\[
T\le-4Z+n\left(\sum_{j=5}^{10}\frac1j+\frac6{B+1}\right)
\le-\left(1+\frac{709}{13090}\right)Z.
\]
Using $W\le Z$ now gives a strictly underpriced positive smaller-rank successor.

\subsection*{The high platform at gap seven}
If $v=L_{B+3}^0>Z$, both attaining heads are at least four.
Their point capacities sum to at most $n(1/5+1/(B+1))<Z$.
At least one side survives monochromatic restoration, giving width
$w\ge v-M_4>Z-M_4$ on the same $Y$.
Take $\ell=\floor{n/8}$ and $h=(\ell-R)_+$.
Since $M_4+\ell\le13n/40<Z$, the first $h$ vertices of degree at least $B+2$
actually exist. Delete them and use the single allowance $\ell$
in Lemma \ref{lem:fixed-restoration}.
All heads at least seven are retained; the new rank satisfies
$N'\ge2(\mu-Z)+R+2d+\ell-h>0$ and $N'<N$.

The required floor inequalities are

\[
\begin{array}{lll}
2M_2-M_3\le Z,&M_2+M_3-M_4\le Z,&M_2+M_4-M_5\le Z,\\
2M_3-M_5\le Z,&M_2+M_5-M_6\le Z,&M_3+M_4-M_6\le Z,\\
M_2+M_6-\ell\le Z,&M_3+M_5-\ell\le Z,&2M_4-\ell\le Z.
\end{array}
\]
After clearing the positive denominators, the differences from the continuous bounds are
$2Z-5c-9$, $14Z-23c-48$, $8Z-11c-25$,
$2Z-2c-5$, $4Z-5c-12$, $54Z-43c-120$,
$50Z-59c-147$, $10Z-7c-21$, and $18Z-11c-35$.
All are nonnegative from $Z\ge8c$.
The unit-axis tax $Z+M_j-\ell$ therefore dominates the low columns.
The single low-head losses on high columns are $M_2,\ldots,M_6$.
The boundary with two small heads satisfies
$M_5+M_6\le M_2$ and $2M_6\le M_2$;
all remaining high columns have a fully retained head.
The full tax is at most

\[
(B+1)Z+2\sum_{j=2}^{6}M_j-3\ell.
\]

Hence

\[
e(X)-E'\le W-(D-1)\rho-1-6Z+2\sum_{j=2}^{6}M_j-3\ell+5h.
\]
For $h=0$, use $2\sum M_j\le153n/70$ and $W\le Z$;
the bound is $-22Z/35+153c/70-(D-1)\rho-1<0$.
For $h>0$, use $h=\ell-R$ and $2\ell\le n/4$ to obtain

\[
-\frac9{70}Z+\frac{341}{140}c-5R-(D-1)\rho-1
\le-\frac9{70}Z-\frac{359}{140}a-5(r-1)
 -\left(D-\frac{481}{140}\right)\rho-1<0.
\]
\subsection*{The low and middle platforms at gap six}
Now $D=B+6\ge14$ and $B\ge8$.
If $L_{B+3}^0\le Z$, choose the same $k\le Z$ using $B+3\le h\le D$.
Combine the original row shift directly and charge the highest column once:

\[
E_r-E_{r+a+Z+k}\le Da+\sum_{h=1}^{D}\min(Z+k,L_h^0)+\lambda .
\]
The tax at $B+2$ is at most $65Z/84<Z$ and uses the clipping identity.
The other four taxes sum to at most $\beta n+4t$, where
$\beta=533/840$ and $t=(n/(B+1)-R)_+$. Thus

\[
e(X)-E_{\rm new}\le W-(D-1)\rho-1-3Z+\beta(2Z+c)+4t.
\]
For $t=0$, use $W\le Dc/2$ and $Z\ge Dc/2$ to obtain
$[-(1-\beta)D+\beta]c-(D-1)\rho-1<0$.
For $t>0$, the coefficient $3-2\beta-8/(B+1)$ is positive.
The same substitution gives $A_Bc-4R-(D-1)\rho-1$, where

\[
A_B=-(1-\beta)(B+6)+\beta+4+\frac{24}{B+1}
\le A_8=\frac{367}{168}<4.
\]
If $A_B\le0$ this is negative immediately.
Otherwise use $c\le a+\rho$ and $R=a+r-1$.

\subsection*{The high platform at gap six}
Let $v=L_{B+3}^0>Z$, $\ell=M_6$, and $h=(\ell-R)_+$.
The attaining heads are at least four; one side survives and gives
$w>Z-M_4>\ell$, using $M_4+\ell\le12n/35<Z$.
The first $h$ vertices of the same history have current degree at least $B+2$.
Use one allowance $\ell$, retaining every head at least six.
Positive rank and all corners follow as in the preceding high-platform proof.
The total tax through column $D$ is at most

\[
(B+1)Z+2\sum_{j=2}^{5}M_j-3\ell .
\]
Here $Z\ge7c$ gives the low-column floor inequalities.
The only double loss on a high column at $B=8$ is $5/5$,
covered by $2M_5-\ell\le M_2$, since
$2\floor{n/6}-\floor{n/7}\le\floor{n/3}$.
If all first $h$ vertices have degree at least $B+3$, the actual cost gives

\[
e(X)-E'\le
\begin{cases}
-Z/5+19c/10-3R-(D-1)\rho-1,&h>0,\\
-37Z/35+103c/70+18/7-(D-1)\rho-1,&h=0 .
\end{cases}
\]
Both are strictly negative using $Z\ge Dc/2$ and $c\le R-1+\rho$.
Otherwise $C_{B+3}<h$ and the final graph has cap $q=B+2=D-4$.
Even restoring a conservative extra highest-column tax $\ell$, one has

\[
e(X)-E'\le U_2
 :=W-\frac{22}{35}Z+\frac{153}{70}c-4R-(D-1)\rho-1.
\]
If $c\le5$, substitute the weight and length bounds:
$U_2\le(13D-127)c/70-(D-5)\rho-5\le-(5D+635)/70<0$.
If $c\ge6$, the highest-column width of the fixed successor is at least

\[
H_2=\frac{11}{35}Z-\frac{12}{35}c+R+1\ge q/2.
\]
This follows from the actual attaining pair: one head is at least six,
and the other loses at most $M_4-\ell$.
If a higher column remains positive, the degree caps avoid it.
Otherwise Proposition \ref{prop:general-price} gives
$e(X)-E'\ge(q/3-1)H_2-q^2/6$.
The simplified bound $T_q(x)\ge(q/3)(x-q/2)$ for $x\ge q/2$
follows directly from its definition.
For even $q=2m$, the two terms are $mx-m^2$ and
$\floor{((q+1)x-q(q-1)/2+3)/3}$.
For odd $q=2m+1$, the second term has slope larger by
$(5m+1)/(3(3m+2))$; at $x=q/2$, even after subtracting one for rounding,
its excess is $(8m^2+2m-1)/(2(3m+2))>0$.
The first term satisfies the same bound.

Subtract $U_2$, use $Z\ge Dc/2$, $W\le Dc/2$, and then $R\ge c-\rho+1$.
The difference is at least

\[
\frac{(D-1)(11D-59)}{210}c+
\frac{2(D-4)}3\rho+
\frac{D+5}{3}+\frac{D-4}{3}-\frac{(D-4)^2}{6}
\ge\frac{31D^2+140D-696}{210}>0.
\]
This is a contradiction. The low and high platforms exhaust both gaps.
\end{proof}
\section{Payment on one package at head gaps four and five}\label{app:gaps45}
\begin{proposition}\label{prop:gaps45}
The original minimum strict source cannot have $\gamma=4$ or $5$.
\end{proposition}
\begin{proof}
Choose the actual fractional package of Section \ref{sec:fractional},
unless its construction already supplies a complete avoiding coloring.
All of $a,\rho,\zeta,\varepsilon,\beta$ in this appendix belong to this package.
Take all its positive-weight whole cores $C$, and put

\[
n=2Z+c,\quad e(C)=DZ+W,\quad
Y=(G-A)-E(C),\quad R=r+a-1,\quad
\Gamma=\zeta+\varepsilon+\lambda-W.
\]
\[
c\le R-1+\rho,\quad 2Z\ge Dc,\quad
\nu(Y)\ge Z+R+d>Z,\quad L_h^0=(L_h-a)_+.
\]
The failure inequality for unequal even-sum caps gives
$\zeta\le(D-2)(\rho-\lambda)/(p-1)$;
the odd-sum case gives
$\zeta\le[(D-p)\rho-(D-1)\lambda]/(p-1)$.
Here the odd cap has $p\ge3$ and $\rho\ge\lambda$.
Combining with $\Gamma\le\zeta/2-D\rho/2-\beta/2+\lambda$ gives

\begin{equation}\label{eq:gap45-debt}
\begin{array}{ll}
\gamma=5:&\Gamma\le-(D+1)\rho/4,\\
\gamma=4:&\Gamma\le-g\rho-h_0\lambda,\quad
g=(D+2)/4,\ h_0=(D-6)/4.
\end{array}
\end{equation}
The balanced case has $\zeta=0$ and is stronger.
At gap five, the smallest $B=7$ gives $D=12,A=7$ and is avoided by even caps $6/6$;
hence $B\ge8,D\ge13$ below.
For $B=8,D=13$, the actual highest pair $7/8$ permits odd cap seven
and even cap six, giving $\zeta\le\rho-2\lambda$ and $\Gamma\le-6\rho$.
At gap four, $B\ge6,D\ge10$.

\subsection*{Actual roots supplied by the complete source}
The point-cover tax on columns $2,\ldots,B+1$ of the same $Y$
is at most $Z$, except that the $2/2$ column may cost an additional
$\tau_{22}=(\floor{(c-R)/2})_+\le\rho/2$.
For the higher column $B+j$, put

\[
\theta_j=\frac{B+j+1}{B(j+2)+j+1},\qquad
S_5=\sum_{j=2}^{5}\theta_j\le\frac{3426274}{3125265},
\quad
S_4=\sum_{j=2}^{4}\theta_j\le\frac{1873}{2091}.
\]
All three terms of Lemma \ref{lem:point-cover-tax} are at most $n\theta_j$.
The heads are at least three and have fixed sum; the reciprocal expression
is largest at the endpoint pair $(j+1,B)$.
Comparing the common envelope with the one value $e(Y)-\nu(Y)$ gives

\[
d\ge
\begin{cases}
4Z-nS_5-\tau_{22}-\lambda+1-\Gamma>Z/2,&\gamma=5,\\
3Z-nS_4-\tau_{22}+1-\Gamma>Z,&\gamma=4.
\end{cases}
\]
Use $n/Z\le13/6$ and $11/5$, respectively:
$4-(13/6)S_5\ge4-(13/6)(3426274/3125265)=1+450508/721215>1/2$ and
$(11/5)S_4=20603/10455<2$; the root-debt bounds absorb the additional terms.
The same complete forest--matching endpoint therefore supplies
$C_B(Y)>Z$ or $C_B(Y)>2Z$.

\subsection*{Gap five: the low platform}
If $L_{B+4}^0\le Z$, set
$k=\max((Z-R)_+,w_{B+4},w_D)\le Z$.
Delete the first $k$ actual vertices of degree at least $B$
and use the whole-$C$ menu.
This gives every corner of $P_{r+a+Z+k}$ on the original lists,
with positive rank $\nu(Y)+d-k$.
Charge the original highest column once.
Clipping charges $Z$ for each of $B+2,B+3$;
the two followed columns use their actual point-cover taxes $J_h$.
Thus

\[
e(X)-E_{\rm new}\le\Gamma-Z+J_{B+4}+J_D,\qquad
J_{B+4}+J_D\le nS_B,\quad
S_B=\frac{B+5}{6B+5}+\frac{B+6}{7B+6}.
\]
For $B\ge9$, $(15/7)S_9=9255/9499<1$, giving strict negativity.
For $B=8$, $S_8=774/1643<1/2$.
Put $s_0=\min(R,c)$. In each head pair $u\le v$ in these two columns,
$u\ge5$ and $v\le8$. The margin between the point cap and the cover term is

\[
\frac{n\theta-n/(u+1)}{1-\theta}=\frac{nu}{v(u+1)}
\ge\frac{5n}{48}>c.
\]
Hence $J_{B+4}+J_D\le nS_8-(2-S_8)s_0$.
For $R\ge c$ the bound is negative directly.
For $R<c$, use $R\ge c-\rho+1$; the bound is strictly less than
$\Gamma+(2-S_8)\rho<0$.

\subsection*{Gap five: the high platform}
If $v=L_{B+4}^0>Z$, both attaining heads are at least five.
Take $\ell=M_5$ and $h=(\ell-R)_+$.
One side survives monochromatic restoration, giving
$w\ge v-M_5>Z-M_5>\ell$.
The same history therefore supplies $h$ vertices of current degree at least $B+3$.
Use the single allowance $\ell$. Every head at least five is retained;
$N'\ge2(\mu-Z)+R+2d+\ell-h>0$ and $N'<N$.
The tax through column $D$ is

\[
(B+1)Z+2(M_2+M_3+M_4)-2\ell.
\]
Two-nonunit candidates on the low columns are dominated by
$2M_2-M_3,M_2+M_3-M_4,M_2+M_4-M_5,2M_3-M_5\le Z$,
from the floor bounds with $Z\ge7c$.
Every high column has a head at least five because $B\ge8$.
Charging the top once, the actual cost gives

\[
e(X)-E'\le\Gamma-4Z+2(M_2+M_3+M_4)-2\ell+2h
\le\Gamma-\frac{13}{15}Z+\frac{47}{30}c-2\min(R,\ell)<0.
\]
\subsection*{Gap four: the low platform}
If $L_{B+2}^0\le Z$, follow every actual width at $B+2,\ldots,D$
and choose $k\le Z$. This gives a positive smaller-rank successor on the same original lists.
The full original tax minus income on that history is

\[
e(X)-E_{\rm new}\le\Gamma-2Z+\sum_{j=2}^{4}J_{B+j}-1
\le\Gamma-2Z+nS_4-1<0.
\]
\subsection*{Gap four: the high platform}
If $v=L_{B+2}^0>Z$, both actual heads are at least three.
One side survives monochromatic restoration and gives
$w\ge v-M_3>Z-M_3$.
Take $\ell=M_5$ and $h=(\ell-R)_+$.
Since $M_3+\ell\le5n/12<Z$, the first $h$ vertices have current degree at least $B+1$.
One fixed allowance $\ell$ gives the complete new lists and positive smaller rank,
retaining all heads at least five.
The low-column floor inequalities are as above.
At $B=6$, the only double loss on a high column is $4/4$ and requires
$2M_4-\ell\le M_2$.
For $n\ge20$, the left side is at most $7n/30+1\le n/3-1$.
For $n=11$, $12\le n\le14$, $15\le n\le17$, and $18\le n\le19$,
it equals $3,2,4,3$, respectively, each at most $M_2$.
Consequently,

\[
e(X)-E'\le\Gamma-3Z+2(M_2+M_3+M_4)-3\ell+3h-1.
\]
For $h=0$, the upper bound is
$\Gamma-13Z/15+16c/15+2<0$.
If $h>0$ and all first $h$ vertices have degree at least $B+2$,
the extra income $h$ gives $\Gamma-Z/5+7c/5-2R$,
again negative by $Z\ge Dc/2$, $c\le R-1+\rho$, and \eqref{eq:gap45-debt}.

The only remaining case is $h>0,C_{B+2}<h$, with cap $q=B+1=D-3$ on the same $X$.
One actual attaining head is at least five; the other loses at most $M_3-\ell$.
The new highest width is therefore at least

\[
H_0=Z/6-5c/12+R+1,\qquad
e(X)-E'\le U=2Z/15+47c/30-3R-g\rho-h_0-1.
\]
If $U<0$ the descent is paid.
If $U\ge0$, solve for a lower bound on $Z$, substitute it in $H_0$,
and use $c\le R-1+\rho$:

\[
H_0\ge\frac{19}8R+\frac{5D-28}{16}\rho+
\frac{5D-30}{16}+\frac{37}8
\ge\frac{5D+27}{8}>q/2.
\]
The price valid for either deficit sign on the same $X$ gives
$e(X)-E'\ge(q/3-1)H_0-q^2/6$.
The difference from $U$ has slope $(5D-42)/90>0$ as a function of $Z$.
At the real value $Z_0$ where $U=0$, the same calculation still gives
$H_0\ge(5D+27)/8$.
The difference is therefore at least

\[
\frac{D-6}{3}\,\frac{5D+27}{8}-\frac{(D-3)^2}{6}
=\frac{D^2+21D-198}{24}>0.
\]
This contradicts the upper bound $U$ on that same graph.
All platform branches have been covered.
\end{proof}
\section{Gap three: low platforms and actual attaining pairs}\label{app:gap3}
\begin{proposition}\label{prop:gap3}
The original minimum strict source cannot have $\gamma=3$.
\end{proposition}
\begin{proof}
Write $D=B+3$ and use the same actual fractional package from
Section \ref{sec:fractional}, with $Y,Z,c,n,R,\rho,\Gamma,M_j$
as in Appendix \ref{app:gaps45}.
If $B=5$, then $D=8,A=5$ and even caps $4/4$ give an avoiding coloring.
Thus $B\ge6,D\ge9$.
The actual positive cores give $Z\ge\ceil{D/2}c\ge5c$ and $n/Z\le11/5$.
The failure inequalities also give

\begin{equation}\label{eq:gap3-debt}
\Gamma\le-g\rho-h_0\lambda,\qquad
(g,h_0)=
\begin{cases}
((3D+5)/8,(D-9)/8),&D\text{ odd},\\
((3D+2)/8,(D-10)/8),&D\text{ even}.
\end{cases}
\end{equation}
\[
\Gamma\le-(D+1)\rho/4,\qquad c\le R-1+\rho.
\]
For odd $D$, $\rho\ge\lambda+1\ge2$.
For even $D$, Lemma \ref{lem:rho1} excludes $\rho=1$ using only the
conditional theorem in Appendix \ref{app:regular-components}.
Thus $\rho\ge2$ is available below.
Point-cover taxes on low columns are at most $Z$, with an additional
$\tau\le\rho/2$ at $2/2$.
The two high-column capacity coefficients sum to at most
$1/3+5/17<2/3$. Applying these taxes to the complete source on the same $Y$ gives

\[
d\ge2Z-\frac23n-\Gamma+1-\tau
\ge\frac8{15}Z-\Gamma+1-\tau>Z/2.
\]
The complete forest/matching endpoint therefore supplies $C_B(Y)>Z$.
All subsequent vertices are taken from this one maximum-degree deletion history.

\subsection*{High platform: $v=L_{B+2}^0>Z$}
The heads of an actual attaining pair sum to $B+3$.
It is therefore $3/B$ or has both heads at least four.
Put $\ell=M_4$ and $h=(\ell-R)_+$.
The surviving width exceeds $Z-M_3>\ell$ in the first case and
$Z-\ell>\ell$ in the second, since $M_3+\ell\le9n/20<Z$.
Hence this same history supplies $h$ vertices of degree at least $B+1$.
Restore the original lists with the single parameter $\ell$
of Lemma \ref{lem:fixed-restoration}.
Every head at least four is retained; losses at heights one, two, and three
are at most $Z-\ell,M_2-\ell,M_3-\ell$. The new rank satisfies

\[
N'\ge2(\mu_{\mathcal J}-Z)+R+2d+\ell-h>0,\qquad N'<N.
\]
The complete column tax, with the highest column charged once in $\Gamma$,
is at most $(B+1)Z+2M_2+2M_3-2\ell$.
The needed floor inequalities are
$2M_2-M_3\le Z$, $M_2+M_3-M_4\le Z$, and $2M_3-M_4\le Z$.
They follow by summing the corresponding core inequalities using $Z\ge5c$;
direct floors cover $c=1,Z=5,6$.
Thus the fixed successor has upper cost

\[
e(X)-E'\le\Gamma-2Z+2M_2+2M_3-2\ell+2h-1.
\]
If $h=0$, this is at most $\Gamma-7Z/15+23c/30+1<0$.
If $C_{B+2}(Y)\ge h$, the additional actual saving $h$ bounds it by
$\Gamma-Z/15+29c/30-R<0$.
It remains that $h>0,C_{B+2}(Y)<h$, so $X$ has cap $q=B+1=D-2$.

For an actual $3/B$ winner, the same fixed successor retains uniform width

\[
L\ge Z/10-9c/20+R+1=:L_0>0.
\]
If $B$ is odd, even caps $2/(B-1)$ already give an avoiding coloring.
For even $B$, the signed exposure is at least
$L+(B-2)\ceil{L/2}\ge BL/2$.
The complete original cost bound and the cap price on this same graph are

\[
U=Z/3+7c/6-2R-g\rho-h_0-1,\qquad
e(X)-E'\ge \frac{B(B-2)}6L_0-\frac{(B+1)^2}6.
\]
It suffices to consider $U\ge0$. At its smallest permitted $Z=Z_0$,

\[
L_0(Z_0)\ge\frac45R+
\left(\frac{3g}{10}-\frac45\right)\rho+\frac{3h_0}{10}+\frac{21}{10}
\ge\frac{3B+48}{20}>q/B.
\]
\[
\left[\frac{B(B-2)}6L_0-\frac{(B+1)^2}6\right]-U
\ge\frac{3B^3+22B^2-136B-20}{120}>0.
\]
The slope of this difference in $Z$ is $(B^2-2B-20)/60>0$;
the numerator equals $604$ at $B=6$ and increases thereafter.
This contradicts the whole permitted domain.

If both winning heads are at least four, the highest width of the fixed
successor is at least $H_0=Z/5-2c/5+R+1$, with the same upper bound $U$.
At the boundary $U=0$,

\[
H_0\ge
\begin{cases}
(21D+67)/40,&D\text{ odd},\\
(21D+46)/40,&D\text{ even}.
\end{cases}
>q/2.
\]
The general cap price gives $(q/3-1)H_0-q^2/6$.
After subtracting $U$, the slope is $(D-10)/15$.
Its boundary lower bounds for odd $D\ge11$ and even $D\ge12$ are, respectively,

\[
(D^2+42D-415)/120>0,\qquad
(D^2+21D-310)/120>0.
\]
For $D=9$, actual $4/5$ exposure gives $2H_0-49/6$.
Here $H_0\ge32/5$, and the difference at $U=0$ is at least $139/30$,
with slope $1/15$.
For $D=10$, the integer width is at least $\ceil{32/5}=7$;
the exact price $T_8(x)=3x-9$ gives $2H_0-9$.
The difference is at least $19/5$, again with slope $1/15$.
This excludes the entire high platform.

\subsection*{Low platform: $L_{D-1}^0\le Z$}
Use the actual near-top cap colorings of Lemma \ref{lem:near-top}.
Whole cores with ports satisfy the strict internal caps.
Only actual isolated odd regular cores may need one exception;
their number is $c_0\le\rho$, and $c_0=0$ for odd $D$.
An actual height-$D$ winner has both heads at least four and an earlier
larger original head on at least one side, so its tax is $\tau_D\le M_4+c_0$.
At height $D-1$, the winner satisfies $\alpha+\beta=D$ with both heads at least three.
If $A-\alpha\ge2$ or $B-\beta\ge2$, its tax is $\tau_{D-1}\le M_3+c_0$.
The sole remaining case is
$A=\alpha+1,B=\beta+1,A+B=D+2$, called the diagonal case.
For even $D$, this gives $A=5,B=D-3$ and even caps $4/(B-1)$ avoid the pair.
In the nondiagonal case,

\[
\tau_{D-1}+\tau_D\le M_3+M_4+2c_0
\le9n/20+2c_0\le99Z/100+2c_0.
\]
Follow the two actual widths and choose
$k=\max(\kappa,(Z-R)_+)\le Z$.
The established $C_B(Y)>Z$ supplies the same roots.
Original-list restoration gives every new pair, rank $\nu(Y)+d-k>0$, and

\[
e(X)-E'\le\Gamma-Z+\tau_{D-1}+\tau_D-1
\le\Gamma+2c_0-Z/100-1<0.
\]
Consider the diagonal case with odd $D\ge11$.
Point covering permits
$\tau_{D-1}=\floor{11n/39}$ and $\tau_D=M_4$,
since $\alpha,\beta\ge4$ and $nD/(5D-16)\le11n/39$.
Both taxes are at least $M_4$, so the followed widths satisfy $\kappa\le Z-M_4$.
Put $m=\min(R,M_4)$; then $k\le Z-m$, and the first high column
actually refunds $m$. As $Z\ge6c$ and $M_4\ge2c$,

\[
e(X)-E'\le\Gamma-Z-m+\frac{94n}{195}-1
\le\Gamma+\frac4{15}c-m-1<0.
\]
For the last strict inequality, if $R\ge c/2$ then $m\ge c/2$.
Otherwise $\rho>c/2$, whereas $\Gamma\le-3\rho$.

\subsection*{The diagonal boundary at $D=9$}
Only original largest heads $5,6$ and a height-eight $4/5$ winner remain.
Write $f=f_4,f'=f_5,g=g_5,g'=g_6$.
A zero ninth column needs no tax.
If it is attained by $5/5$, its tax $M_5$ and the preceding tax sum to
at most $83n/174$, with actual refund $\min(R,M_5)$; the preceding argument pays.
It remains that height nine is attained by $4/6$. Thus

\[
L_8^0=f+g-R,\qquad
L_9^0=\max(f+g',f'+g)-R
=L_8^0-\min(f-f',g-g').
\]
If the $4/5$ point-cover tax is at most $M_3$, the case is already excluded.
If $f>M_4$, an actual $5/4$ cap coloring preserves a red size-four star in
the exterior, making each high-column tax $M_4$, with total below $Z$.
If $g'>M_6$, an actual $3/6$ cap coloring preserves a blue size-six star;
the taxes $M_5,M_6$ again sum to less than $Z$.
Thus $f\le M_4,g'\le M_6$.
If $g>M_5$, the same $3/6$ coloring gives taxes $M_5,M_4$, already excluded.
We may therefore assume $f\le M_4,g\le M_5$.

Choose $n-5f+1$ actual vertices in the package; color their incident edges
blue and the remaining internal edges red.
The $f$ stars of head at least four in a complete original red forest
cannot all stay in the unchosen vertices. A real size-four star survives outside;
blue component loss is at most the number of chosen vertices.
The complete source at the original extreme pair $(1,r+a-1)$
therefore gives height-eight tax $n-4f$.
Reversing colors gives $n-5g$. Together with the zero-column option,

\[
\tau_8\le\min(L_8^0,n-4f,n-5g).
\]
The original highest $5/6$ column gives $f'+g'-R=\lambda-a\le\rho$.
Since $4/6$ attains height nine,
$L_9^0\le f+\rho-f'\le f+\rho-1$.
We may set this column to zero with $\tau_9=L_9^0$.
Eliminating $g$ gives

\[
\tau_8\le\min\left(n-4f,\frac{n+5f-5R}{6}\right),\qquad
\tau_8+\tau_9\le\frac{14n-15R}{29}+\rho-1.
\]
The latter is the intersection bound for the two affine expressions
in $f$, attained at $f=5(n+R)/29$.
Follow only the eighth-column width.
The original failure inequality at $D=9$ gives $\Gamma\le-4\rho$.
Consequently the same fixed successor satisfies

\[
e(X)-E'\le\Gamma-Z+\tau_8+\tau_9-1
\le-\frac Z{29}-\frac R{29}-\frac{73\rho}{29}-\frac{72}{29}<0.
\]
\end{proof}
\section{Gap two: actual root supply and completed stages}\label{app:gap2-stages}
Throughout, $D=B+2\ge9$. Use the actual fractional package
$a,\rho,C,Z,c,n,W,R,Y,\Gamma$, and original columns $L_h^0=(L_h-a)_+$.
Let $c_0$ count zero-boundary positive cores;
$c_0\le\rho$, and $c_0=0$ for odd $D$.
The exact package inequalities give

\begin{equation}\label{eq:gap2-debt}
\begin{array}{ll}
D=2m:&\Gamma\le-(m+1)\rho/2-(m-3)\lambda/2,\\
D=2m+1:&\Gamma\le-(m+1)\rho/2-(m-2)\lambda/2.
\end{array}
\end{equation}
\[
\Gamma\le-(D+1)\rho/4,\quad \Gamma\le-B/2,\quad
c\le R-1+\rho,\quad Z\ge\ceil{D/2}c\ge5c,\quad n=2Z+c.
\]
\subsection{Actual inventory without a platform assumption}
Columns $2,\ldots,B+1$ have point-cover tax at most $Z$,
except for an additional $\tau_{22}\le(\rho-1)/2$ at $2/2$.
An actual height-$D$ attaining pair has both heads at least three and
an earlier larger original head on at least one side.
The two directions of Lemma \ref{lem:near-top} give tax $M_3+c_0$;
that argument only needs target heads at least three.
Put $Q_0=\sum_{h=2}^D L_h^0$. The actual edge account and
the canonical coordinates on this same graph give

\[
\begin{split}
e(Y)-\nu(Y)&\le Q_0+d-1-(D-1)Z+\Gamma,\\
e(Y)-\nu(Y)&\ge Q_0-BZ-\tau_{22}-M_3-c_0.
\end{split}
\]
\[
d\ge Z-M_3-c_0-\Gamma+1-\tau_{22},\qquad
2d+R-Z\ge R/2+(D-6)\rho/2+7/2>0.
\]
The complete forest/matching endpoint on this same $Y$ has deficit exactly
$d$ and matching target at least $2d$.
Hence one maximum-degree deletion history satisfies

\begin{equation}\label{eq:gap2-stock}
C_B(Y)>Z-R,\qquad L_D^0\le M_3+c_0+C_{D-1}(Y).
\end{equation}
The second inequality uses that same near-top source:
after all vertices of current degree at least $D-1$ are deleted,
cap $D-2$ forbids a positive height-$D$ uniform source.

\subsection{The low platform and one original-list successor}
If $R\ge Z$, the whole-core menu uses the original lists directly,
with upper cost $\Gamma-1<0$.
Otherwise put $t=Z-R>0$ and delete the first $t$ vertices of the same history,
obtaining $X_t$.
The menu gives $P_{r+a+Z+t}(X_t)$ with rank $\nu(Y)+d-t>0$.
If the last vertex is deleted at degree $q\in\{B,D-1,D\}$, actual cost gives

\begin{equation}\label{eq:gap2-raw}
g_t:=e(X_t)-E_t
\le\Gamma-(D-q)Z+\sum_{j=q+1}^{D}(L_j^0-C_j)-1.
\end{equation}
This is one deletion, and all $C_j$ belong to that one history.
Its rank is strictly smaller than the original rank.
If $L_{D-1}^0\le Z$, then also $L_D^0\le Z$.
The full shift loses at most $Da+\lambda+B(Z+t)+2Z$, whereas the
actual saving is at least $Da-\zeta-e(A)+DZ+W+Bt$.
Their difference gives $\Gamma-1<0$.
Hence assume $L_0:=L_{D-1}^0>Z$.
The case $q=D$ in \eqref{eq:gap2-raw} always exits;
so does $q=D-1$ when $L_D^0\le Z$.

\subsection{Final cap $B$: surviving heads and three stage lengths}
Suppose $q=B$ and write
$u=C_{D-1}(Y)<t$, $b=C_D(Y)\le u$, and $\ell_0=R+u<Z$.
Equation \eqref{eq:gap2-stock} shows that failure to pay forces

\begin{equation}\label{eq:gap2-L0}
L_0\ge2Z-M_3-c_0+b+1-\Gamma,\qquad
g_t\le\Gamma-2Z+L_0+M_3+c_0-b-1.
\end{equation}
If $L_0$ has an actual winner with both heads at least three and both
counts are bounded by their corresponding point capacities,
then $L_0\le2M_3-R$ and
$g_t\le\Gamma-Z/2+3c/4+c_0-R-b-1<0$.
Otherwise one side actually survives. Monochromatic restoration gives
uniform width at least $L_0-M_3-t$ on this same $X_t$.
Use $g_t\ge0$ if this displayed width is nonpositive, and
Lemma \ref{lem:surplus-price} if it is positive. In either case,

\[
g_t\ge L_0-M_3-t-\floor{D/2}.
\]
\[
0\le\Gamma+\floor{D/2}+\rho/2+c_0-a/2-r.
\]
For odd $D=2m+1$, the right side is at most $1-a/2-r<0$.
For even $D=2m$, Lemma \ref{lem:rho1} excludes $\rho=1$;
for $\rho\ge2$ the right side is at most $(7-m)/2-a/2-r<0$.
If the winner is $2/B$ and its high-head count exceeds $M_B$,
the monochromatic loss is smaller and the same comparison applies.
Only an actual $2/B$ winner and the three ranges of $\ell_0$ remain.

First, suppose $L_D^0\le Z$ and $\ell_0>M_2$.
Failure to pay also gives $L_0\ge Z+u+b+1-\Gamma$.
If neither side survived, this would force $M_B>Z$.
Thus one side survives with loss at most $M_2$.
The same-graph surplus price and \eqref{eq:gap2-raw} then require
$0\le\Gamma+M_2-\ell_0-b+\floor{D/2}-1<0$.
The middle and short stages below impose no condition $L_D^0\le Z$.

If $M_3\le\ell_0\le M_2$, delete all $u$ high-stage vertices and
use the fixed intersection lists with allowance $\ell=\ell_0$ on $X_u$.
All heads at least three are retained; the size-two count loses at most
$M_2-\ell$, and units lose at most $Z-\ell$.
The complete column tax and actual highest width give

\[
\begin{split}
\mathrm{tax}&\le(B+1)Z+2M_2-\ell,\\
g_u&\le U=\Gamma-Z+2M_2-R-1,\\
H&\ge2Z-2M_2-M_3+R-c_0+b+1-\Gamma.
\end{split}
\]
Here $N'\ge\nu(Y)+d-Z+R>0$, and the rank strictly decreases.
Any higher positive column is already excluded by edge coloring under cap $B$.
Proposition \ref{prop:selector}, valid for either deficit sign, gives
$g_u\ge\alpha H-K$, where
$\alpha=B/2-1$ and $K=\floor{B/2}(\floor{B/2}+1)$.
If the displayed $H$ is nonpositive, $g_u\ge0$ still implies that lower bound.
Put

\[
\Omega=\frac B2(\Gamma+\rho)+K-B+\left(\frac B2-1\right)c_0.
\]
\[
U-(\alpha H-K)
\le\frac{6-B}{12}Z-\frac{B+6}{24}c+\Omega-\alpha b.
\]
If $B=2m$, the package inequality gives
$\Omega\le-m^2(\rho-1)/2+(m-1)c_0\le0$;
$\rho=1$ is excluded, while $\rho\ge2$ gives
$\Omega\le-(m-2)^2/2<0$.
For $B=2m+1$,
$\Omega\le-m-3/4-m(2m+1)(\rho-1)/4<0$.
Thus the middle stage is impossible.

If $\ell_0<M_3$, do not assume the high stage can be prolonged.
Instead put $\ell=M_3$ and $h=M_3-R>u$.
Since $h<t$, \eqref{eq:gap2-stock} actually supplies enough degree-$B$ vertices.
The saving is at least $Bh+u+b$ and the final cap is $B$.
The same fixed lists have the positive smaller rank and complete tax just proved.
The actual bounds are

\[
\begin{split}
g&\le U=\Gamma-Z+2M_2+M_3-2R-u-b-1,\\
H&\ge H_0=2Z-M_2-2M_3+R-c_0+b+1-\Gamma,\\
U-(\alpha H_0-K)
&\le\frac{7-B}{6}Z-\frac{B-1}{12}c-R-u-\frac B2b+\Omega<0 .
\end{split}
\]
The last inequality uses $B\ge7$, including the zero leading coefficient at $B=7$.
Each operation retains the same package's root edges, exterior, and complete
original tails. Failure of the earlier $X_t$ only supplies a necessary bound
on an original column; its income is not added to that of the new graph.
This excludes the low and middle platforms entirely.
For arbitrary high platforms with final cap $B$, only
$\ell_0>M_2$ and $L_D^0>Z$ remain.

\section{Gap two: fixed complete lists and whole-core gaps}\label{app:gap2-high}
Continue with Appendix \ref{app:gap2-stages}. Now $L=L_D^0>Z$.
An actual winner is either $3/B$ or has both heads at least four.
All inventories still belong to one maximum-degree deletion history of $Y$,
with $b=C_D(Y)$.

\subsection{One common fixed-list operation}
Put $h_3=(M_3-R)_+$.
If $b\ge h_3$, choose allowance $\ell=M_3$ and delete the first $h_3$
actual degree-$D$ vertices of the history.
The tax on the first $D$ columns is $(B+1)Z+2M_2-M_3$, so

\[
e(X)-E'\le\Gamma-Z+2M_2-M_3-1
\le\Gamma-Z/6+5c/12<0.
\]
Henceforth $b<h_3$. Choose once
$\ell=\max(M_4,R+b)\le M_3$ and $h=\ell-R\ge b$.
For a $3/B$ winner, $M_3+M_B<Z$ forces one side to survive,
and $C_{D-1}(Y)\ge L-M_3>Z-M_3$.
Since $M_3+M_4<Z$, there are enough vertices of current degree at least $D-1$
to delete $h$ of them.
For two winning heads at least four,
$C_{D-1}(Y)\ge L-M_4>Z-M_4$ gives the same supply.
The first $b$ vertices have degree $D$, so the saving is at least $(D-1)h+b$,
and the resulting $X$ has cap $D-1$.

The fixed intersection lists retain all heads at least four.
Losses at heights two and three are at most
$(M_2-\ell)_+,(M_3-\ell)_+$, and the unit loss is at most $Z-\ell$.
Every pair belongs to $P_{r+a+h+\ell}$, and

\[
N'\ge\nu(Y)+d-Z+R>0,\qquad N'<N.
\]
\begin{equation}\label{eq:gap2-fixed}
\begin{split}
\mathrm{tax}_{1,\ldots,D}&\le(B+1)Z+2M_2+2M_3-3\ell,\\
g=e(X)-E'&\le U=\Gamma-Z+2M_2+2M_3-2\ell-R-b-1.
\end{split}
\end{equation}
The low-column bounds are
$2M_2-M_3\le Z$, $M_2+M_3-M_4\le Z$, and $2M_3-M_4\le Z$.
They follow from $n/Z\le9/4$ and the respective bounds
$\ceil{5n/12},\ceil{23n/60},\ceil{3n/10}$.
The original highest column is charged once, in $\Gamma$.
If the new lists have a positive column above height $D$,
edge coloring under cap $D-1$ already gives complete avoidance.
Otherwise every subsequent price is applied to this one $X$.

\paragraph{The original degree-nine and degree-ten exposure bounds.}
Here the bound is obtained by an actual operation on the original strict
source, rather than by assuming a nonnegative-surplus baseline at the
original minimum rank.  To distinguish the ordinary region from a selected
package, put
\[
 \rho_0=\rho_D(G),\qquad \lambda=L_{D+1},\qquad
 N=s+t-r+1.
\]
The normalization gives
\[
 e(G)=E_r-1,\qquad \nu(G)<N,\qquad
 1\le\lambda\le\rho_0,\qquad \ell_r=D+1.
\]
The ordinary decomposition of \(G\) has positive whole-core weights
\[
 1\le w_Q\le\lfloor D/2\rfloor,\qquad
 W_+\ge\lceil D\rho_0/2\rceil .
\]
Order these actual cores by nondecreasing matching capacity \(z_Q\).
Take the first prefix whose total weight \(w\) reaches \(\lambda\);
write \(z\) for its total capacity and \(Z\) for its largest capacity.
Then
\[
 \lambda\le w\le\lambda+\lfloor D/2\rfloor-1.
\]
If \(D=9\) and \(\rho_0\ge6\), the unselected weight is at least
\[
 \lceil9\rho_0/2\rceil-\rho_0-3\ge18.
\]
Since each weight is at most four, at least five positive cores remain.
Each has at least \(2Z+1\) vertices, all of original degree nine.
Their union therefore contains more than \(10Z\) such vertices.
Greedy coloring with ten colors supplies at least \(Z+1\) independent
vertices in this union.

If \(D=10\) and \(\rho_0\ge7\), the unselected weight is at least
\[
 5\rho_0-\rho_0-4=4\rho_0-4\ge24.
\]
At least five positive cores again remain.  Every positive core has
capacity at least five.  If \(Z=5\), greedy eleven-coloring of these
cores gives at least five independent vertices, which is already more
than the required \(Z-1=4\).  If \(Z>5\), each remaining core has order
at least \(2Z+1>11\).  None of these connected core graphs is
\(K_{11}\), so Brooks' theorem gives a ten-coloring of their union.
There are at least \(5(2Z+1)>10Z\) vertices, giving at least \(Z+1\)
independent vertices.  This argument does not require the later
exclusion of an actual \(K_{11}\) component.

In either case choose
\[
 b=(Z-r+1)_+\le Z-1
\]
of these independent original full-degree vertices, and call the chosen
set \(I\).  It lies in unselected cores, and has no neighbors in a
selected core.  Thus it still has degree \(D\) in the graph obtained by
removing the selected core vertices.  For every \(T\subseteq I\),
\[
 D|T|=e(T,N(T))\le D|N(T)|,
\]
so Hall's theorem gives a matching saturating \(I\), disjoint from the
selected cores.  Adjoining their internal near-perfect matchings gives
\[
 K:=z+b\le\nu(G)<N.
\]
In particular, the successor rank \(N-K\) is positive and strictly
smaller than \(N\).  Erase the internal edges of the selected cores and
delete \(I\), obtaining one graph \(Y\).  Independence and the separation
of these edge sets give the exact fee
\[
 e(G)-e(Y)=D(z+b)+w=DK+w.
\]
Since \(r+b>Z\), Lemma~\ref{lem:menu} gives the complete source
\(P_{r+K}(Y;\nn,\mm)\) at every new legal corner, with the same two
original lists.  The full profile obeys
\[
 E_r-E_{r+K}\le DK+\lambda.
\]
Consequently
\[
 e(Y)-E_{r+K}
 \le -1+\lambda-w\le-1,
\]
contradicting minimality at the positive smaller rank.  Hence
\[
 D=9\ \Longrightarrow\ \rho_D(G)\le5,\qquad
 D=10\ \Longrightarrow\ \rho_D(G)\le6.
\]
The package deficiency used in Appendix~\ref{app:gap2-high} is at most
the original ordinary exposure, so the same upper bounds apply there.
Only the constructed smaller source uses the inductive lower bound;
no zero-surplus assertion has been applied to the original strict source.

\subsection{An actual $3/B$ attaining pair}
The fixed-list width is at least $H=L-\ell-M_3+R$.
For $\ell=M_4$,

\[
U\le U_1=\Gamma+8Z/15+23c/30-R+1,\qquad
H\ge H_1=Z/10-9c/20+R+1.
\]
Only $U_1\ge0$ needs consideration. At its smallest real boundary $Z_0$,

\[
H_1(Z_0)=-3\Gamma/16-19c/32+19R/16+13/16
\ge2-3\Gamma/16-19\rho/32.
\]
Substitution of \eqref{eq:gap2-debt} gives at least two for even $D\ge10$
and strictly more than two for odd $D\ge9$.
At the lowest boundaries $D=10,9$, use the original exposure bounds
$\rho\le6,5$, respectively; all remaining coefficients are nonnegative at higher $D$.
If $\ell=R+b$, put
$U_{\rm base}=\Gamma-Z+2M_2+2M_3-1$ and
$x=U_{\rm base}/3-\ell\ge0$. Then

\[
U=3x,\quad H\ge H_b+x,\quad
H_b\ge Z/18-23c/36-\Gamma/3+R+4/3.
\]
\[
0\le\Gamma+2Z/15+17c/30+2
\ \Longrightarrow\
H_b\ge-3\Gamma/4-7c/8+R+1/2>2.
\]
For odd $D$, the final cap $D-1$ splits into even caps $2/(B-1)$,
contradicting the positive $3/B$ source.
For even $D\ge10$, signed exposure gives $\rho_{D-1}(X)\ge BH/2$.
The general prefix price gives

\[
g\ge A H-K,\qquad
A=(D-4)(D-2)/6,\quad K=(D-1)^2/6.
\]
Here $A\ge8$ and $2A-K=(D^2-10D+15)/6>0$.
For $\ell=M_4$, the lower bound minus $U_1$ is already positive at $Z_0$
and has slope $A/10-8/15>0$.
For $\ell=R+b$, the difference is at least $AH_b-K+(A-3)x>0$.
This excludes the entire class.

\subsection{Two heads at least four with extra original width}
Now suppose

\begin{equation}\label{eq:gap2-extra-width}
L\ge Z+b+1-\Gamma.
\end{equation}
When the original-list successor $X_t$ ends at degree $D-1$,
failure to pay in \eqref{eq:gap2-raw} gives exactly this inequality.
Both winners are fully retained, so $H\ge L-2\ell+R$.
Proposition \ref{prop:selector} gives

\[
g\ge\alpha H-K,\quad
\alpha=(B-1)/2,\quad
K=\floor{(B+1)/2}\bigl(\floor{(B+1)/2}+1\bigr).
\]
For $\ell=M_4$, the upper and lower bounds and $R\ge c+1-\rho$ give

\[
U-(\alpha H-K)\le
\frac{19-3B}{30}Z+\frac{2-9B}{30}c-\frac{B+1}{2}b+\Omega',
\quad
\Omega'=\frac{B+1}{2}(\Gamma+\rho)+K-B+1.
\]
\[
\Omega'\le
\begin{cases}
(3-m)/2-m(2m+1)(\rho-1)/4<0,&B=2m\ge8,\\
(5-m^2)/2-m(m+1)(\rho-2)/2<0,&B=2m+1\ge7.
\end{cases}
\]
The second line uses $\rho\ge2$ for odd-degree failure packages.
For $\ell=R+b$, we have $U=U_{\rm base}-3\ell$ and
$H\ge Z-\ell+1-\Gamma$.
Since $\alpha\ge3$ and $\ell\le n/4$,

\[
U-(\alpha H-K)\le
\frac{B+1}{2}\Gamma+\frac{1-3B}{12}Z+
\frac{3B+7}{24}c+K-\frac{B+1}{2}
\le\frac{-3B^2+4B+15}{24}<0.
\]
Thus \eqref{eq:gap2-extra-width} always gives a strict exit.

\subsection{Stopping early supplies the width, or an actual core dominates}
Only final cap $B$ remains, with
$\ell_0=R+u>M_2$ and $u=C_{D-1}(Y)<Z-R$.
Let $g(C)$ be the maximum gap of the subset sums of the actual whole-core capacities.
If $g(C)\le R+u$, stop after deleting all $u$ high-stage vertices.
For each new pair, the interior interval in \eqref{eq:menuinterval}
contains $R+u$ consecutive integers, so it contains an actual subset sum.
This gives the complete original-list source $P_{r+a+u+Z}$ on the same $X_u$,
with positive rank $\nu(Y)+d-u$.
Its cost is

\[
e(X_u)-E_u\le
\Gamma-DZ-(D-1)u-b-1+\sum_{j=1}^D\min(Z+u,L_j^0)
\le\Gamma-Z+L-b-1.
\]
Failure to pay gives \eqref{eq:gap2-extra-width}.
The fixed-list operation needs
$h=\max(M_4-R,b)\le u$, since $R+u>M_2$.
Its actual roots are therefore available on the same history, and the
proved comparison applies. This excludes the whole domain $g(C)\le R+u$.

Order capacities as $z_1\le\cdots\le z_c$.
The maximum subset-sum gap is
$\max_i(z_i-\sum_{j<i}z_j)$:
successively adjoining translated subset-sum intervals proves the upper bound,
and an item beyond the preceding sum creates the corresponding gap.
For $i<c$, $z_i\le Z/2$, whereas $\ell_0>M_2\ge Z/2$.
Thus only the last item can realize the gap in an unexcluded case, and

\begin{equation}\label{eq:gap2-dominant}
2z_{\max}-Z>R+u>M_2,\qquad
z_{\max}>\frac{5Z+c}{6}.
\end{equation}
This is an actual whole core, not a formal budget,
and the remaining core capacities are genuinely additive.
The next appendix excludes this exact complement.

\section{Gap two: a complete exit from the dominant-core complement}\label{app:gap2-dominant}
\begin{proposition}\label{prop:gap2}
The original minimum strict source cannot have $D\ge9$ and $\gamma=2$.
\end{proposition}
\begin{proof}
Appendices \ref{app:gap2-stages}--\ref{app:gap2-high} reduce the complement to
$L=L_D^0>Z$, an actual winner with both heads at least four,
and \eqref{eq:gap2-dominant}.
Here $\rho\ge2$: use the failure inequality for odd degree
and Lemma \ref{lem:rho1} for even degree.
If only two positive cores remain, in even degree both weights equal $D/2$
and both boundaries are zero; Proposition \ref{prop:regular-components}
excludes these two actual components.
In odd degree their total weight is at most $D-1<D\rho/2$.
Thus $c\ge3$.
Every other core has capacity at least $\ceil{D/2}$, so

\[
Z>6\ceil{D/2}(c-1)+c>3D(c-1).
\]
Use the single fixed-list operation of Appendix \ref{app:gap2-high}.
A sufficiently long degree-$D$ stage already exits.
Otherwise choose $\ell=\max(M_4,R+b)\le M_3$ and $h=\ell-R$;
the same complete source supplies actual degree-$(D-1)$ roots.
The fixed lists have positive smaller rank and every new pair, with

\[
g\le U=\Gamma-Z+2M_2+2M_3-2\ell-R-b-1,\qquad
H=(L-2\ell+R)_+.
\]
\[
g\ge\alpha H-K,\quad
\alpha=(D-3)/2,\quad
K=\floor{(D-1)/2}\bigl(\floor{(D-1)/2}+1\bigr)\le(D^2-1)/4.
\]
A higher positive column is already excluded by edge coloring.
If the displayed $H$ is nonpositive, $g\ge0$ still permits the same linear bound.

\subsection*{Unbounded actual mass for $D\ge10$}
For $\ell=M_4$,
$U\le U_1=\Gamma+8Z/15+23c/30-R+1$ and
$H\ge H_1=Z/5-2c/5+R+1$.
Using $R\ge c+1-\rho$ gives

\[
U_1-(\alpha H_1-K)
\le\frac{25-3D}{30}Z-
\left(\frac7{30}+\frac{3\alpha}{5}\right)c
+\Gamma+(\alpha+1)\rho+K-2\alpha.
\]
For $D\ge12$, use $\Gamma\le-(D+1)\rho/4$, $\rho\le c$,
$Z>3D(c-1)$, and $c\ge3$.
The result is $(-7D^2+77D+50)/20<0$.
At $D=10$, the exact package inequality gives strictly less than $-10-5/2+12=-1/2$.
At $D=11$, use $\rho\le c-1$ and the odd-degree inequality
to obtain strictly less than $-88/5-17/5+19=-2$.
For $\ell=R+b$, put
$x=(\Gamma-Z+2M_2+2M_3-1)/3-\ell\ge0$. Then

\[
U=3x,\quad H\ge H_b+2x,\quad
H_b\ge Z/9+2R/9-7\rho/9-2\Gamma/3+22/9
\ge Z/9+D/6+37/18.
\]
\[
\alpha H_b-K>\frac{6D^2-8D-102}{36}>0,\qquad
2\alpha-3=D-6>0.
\]
Here $Z>6D$, so this position is also impossible.

\subsection*{Original highest pairs and the isolated boundary at $D=9$}
Here $B=7$, and the original highest pair is $4/7$ or $5/6$.
The actual prefix construction at original rank gives

\[
T_9(x)=\max\left(0,\min\left\{\ceil{9x/2}-20,
\floor{(49x-139)/14}\right\}\right),\qquad
T_9(\rho_{\rm all})\le\lambda-1
\quad\text{if }T_9(\rho_{\rm all})>0.
\]
Only the actual descent branch is used, with no zero baseline at original rank.
The $5/6$ signed bound gives $\rho_{\rm all}\ge\lambda+4$;
both price terms at that value are at least $\lambda$, a contradiction.
For $4/7$, $\rho_{\rm all}\ge\lambda+2$ gives the same contradiction when
$\lambda\ge3$. Since $T_9(5)=3$, the remainder satisfies

\[
\lambda\le2,\quad \rho_{\rm all}\le4,\quad
\lambda+1\le\rho\le4,\quad
\Gamma\le-3\rho-\lambda,\quad Z\ge5c,\quad c\ge\rho+1.
\]
Package deficiency is at most original exposure: exterior roots are
assigned only to cores outside the package.
Positive whole-core weights are at most four and total at least $9\rho/2$,
giving the final core-count bound.
Dominance gives $Z>31c-30$, hence $Z>63,c\ge3$.
A high $3/7$ pair is excluded already; $5/5$ is avoided under final cap eight
by even caps $4/4$. Only $4/6$ remains.
The fixed lists have actual highest width

\[
H=L-2\ell+R\ge R+1-c/2
\ge c/2-\rho+2\ge(5-\rho)/2>0.
\]
If $\ell=R+b$, the same representation $U=3x$ gives

\[
H\ge H_b+2x,\qquad
H_b\ge Z/9-7c/9-2\Gamma/3+R+5/3
\ge(7R+16\rho+6\lambda+17)/9\ge62/9.
\]
The cap-eight price $g\ge3H-20$ exceeds $U$, so $\ell=M_4$ below.
First suppose the size-six side originally also has a size-seven head.
Write $L=f_4+g_6-R$ and $\lambda_0=(\lambda-a)_+$.
The original highest candidate $f_4+g_7-R\le\lambda_0$ gives
$g_6-g_7\ge L-\lambda_0$, in particular $g_6>M_6$.
Monochromatic core restoration followed by deletion of the same $h$ vertices
gives actual uniform width $L-M_6-h$ on this same $X$.
Use its whole-core price once:

\[
g\ge4(L-M_6-h)-20-H.
\]
\[
\begin{split}
U-\bigl[4(L-M_6-h)-20-H\bigr]
&\le\Gamma-\frac{11}{21}Z+\frac{73}{42}c-4R-b+16\\
&\le\Gamma-\frac{37}{42}c-4R-b+16<0.
\end{split}
\]
The last step uses $\Gamma\le-7,c\ge3,R\ge2$.
If six is largest on that side, seven is largest on the other.
This forbids $r=2$, which would give highest height twelve rather than ten.
Thus $r\ge3$ and $R\ge c-\rho+2$.
The ordinary cap-eight price and upper bound give

\[
\begin{split}
U-(3H-20)&\le\Gamma-Z/15+59c/30-4R-b+18\\
&\le-Z/15-61c/30+\rho-\lambda+10.
\end{split}
\]
For $c\ge4$, $Z>31c-30$ and $\rho\le c-1$ make the right side
smaller than $10-31c/10<0$.
If $c=3$, then $\rho=2,\lambda=1$.
The case $R\ge4$ pays directly.
For $R=3,\Gamma\le-8$, the difference is at most
$39/10-Z/15\le-11/30$.
The only unpaid form is

\[
c=3,\quad\rho=2,\quad\lambda=1,\quad R=r=3,\quad a=1,\quad\Gamma=-7.
\]
The package has exactly three positive whole cores, no zero or negative cores,
and $e(A)=0$.
The handshake and failure inequalities give
$W=(18+\zeta+\beta)/2$, $\Gamma=-8+(\zeta-\beta)/2$, and $\zeta\le2$.
Thus $\zeta=2,\beta=0,W=10$.
The package together with its root is an actually isolated original region $V_0$.
A root-saturating matching and core near-perfect matchings expose two original
degree-nine vertices. Conversely, one root can match into at most one of the
three odd cores, so every matching exposes at least two full-degree vertices.
Hence $\rho_9(G[V_0])=2$.
Exposure is additive across actual isolated components, making exterior exposure
at most $\rho_{\rm all}-2\le2$.
Lemma \ref{lem:signed-exposure} for single stars $S_4,S_7$ implies that a cap-nine
graph without a strict $3/6$ coloring has exposure at least three.
Interior and exterior each admit that coloring. Exchange its two colors, if necessary, so that cap three is assigned to the original largest head six and cap six to the original largest head seven. Their union then gives complete avoidance in the original graph.
All smaller stars and unit tails are handled because an original largest star
is forbidden in each target.

The low and middle platforms, all final degrees and actual winners,
successful subset gaps, and their dominant complement are now covered.
Every minimality invocation has positive smaller rank, or a complete avoiding
coloring is constructed on the original graph.
Income from different residual graphs is never added.
\end{proof}

\section{An automatic actual entry at gap one}\label{app:gap1-entry}
\begin{lemma}\label{lem:gap1-entry}
Let $D\ge8$ and $B=D-1$. An unexcluded original minimum strict source satisfies
$R+b<M_2$, $C_B(Y)>M_2-R$, and $L_D^0\ge Z+b+1-\Gamma>Z$
on the same actual fractional package and one maximum-degree deletion history,
where $b=C_D(Y)$.
\end{lemma}
\begin{proof}
Take the package from Section \ref{sec:fractional}, with
$R=r+a-1$, $n=2Z+c$, $M_j=\floor{n/(j+1)}$,
$Y=(G-A)-E(C)$, and $c_0$ its number of zero-boundary positive cores.
All quantities refer to this actual package:

\[
c\le R-1+\rho,\quad Z\ge\ceil{D/2}c\ge4c,\quad
\nu(Y)=\nu(G)-a-Z>Z,\quad s,t\ge Z+R+d,
\]
\[
\Gamma\le-\frac{D+2}{4}\rho-\frac{D-6}{4}\lambda<0,\qquad c_0\le\rho.
\]
For odd $D$, $c_0=0$. At $D=8$, original exposure is at most five:
the original-rank prefix price $T_8(x)=\min(4x-16,3x-9)$ exceeds
$x-1\ge\lambda-1$ for $x\ge6$, contradicting strictness.

\subsection*{Canonical coordinates and actual monochromatic refunds}
Point-cover taxes on columns $2,\ldots,D-1$ of the same $Y$ are at most $Z$,
except for an additional $\tau\le(\rho-1)/2$ at $2/2$.
The $2/3$ extra tax for $D=8$ is
$\pos{\floor{(5c-6R-Z)/11}}=0$,
by $Z\ge4c$, $c\le R-1+\rho$, $\rho\le5$, and $R\ge1$.
A height-$D$ winner may have smaller head two;
the two directions of Lemma \ref{lem:near-top} still give tax $M_2+c_0$.
If the unsafe count exceeds capacity, use its monochromatic survival.
Otherwise place exceptions in the opposite direction and retain its earlier
larger original head. No assertion that every suffix contains that head is needed.
The actual edge account and same-graph canonical coordinates give

\begin{equation}\label{eq:gap1-dgate}
d\ge Z-M_2-c_0-\Gamma+1-\tau.
\end{equation}
Let $A_h$ be the largest original unit candidate at height $h$.
For $2\le h\le D-1$, strong lengths give $A_h-a\ge Z+d+1$.
If both counts of a mixed winner were at most $M_2$,
its shifted width would be at most $2M_2-R$. The unit candidate exceeds it by at least

\[
Z+d+1-(2M_2-R)\ge R-c+2-c_0-\Gamma-\tau
\ge\frac72+\frac{D-8}{4}\rho+\frac{D-6}{4}\lambda>0.
\]
Every actual mixed winner thus has a count greater than $M_2$.
Monochromatic core restoration reduces this canonical coordinate's tax from $Z$ to $M_2$.
\subsection*{Insufficient roots force all low columns onto unit axes}
Put $t=(Z-R)_+$. The complete largest-head forest/matching endpoint gives
$C_B(Y)\ge2d+D-3$.
If $C_B(Y)<t$ while a mixed candidate attains some $3\le h\le D-1$,
the refund improves \eqref{eq:gap1-dgate} to
$d\ge2(Z-M_2)-c_0-\Gamma+1-\tau$.
Insufficient supply would then require

\[
0>3Z-4M_2+R+D-1-2c_0-2\Gamma-2\tau>0.
\]
Indeed, $3Z-4M_2\ge(Z-4c)/3\ge0$; the other terms are at least
$R+D+(D-4)\rho/2+(D-6)\lambda/2>0$.
Thus insufficient supply gives
$E=\sum_{h=1}^{D-1}A_h+L_D+\lambda$, where $A_1=N$.
For the complete unit axes, one deletion history satisfies

\begin{equation}\label{eq:unit-history}
C_h(Y)\ge A_h-a-Z\qquad(1\le h\le D-1).
\end{equation}
To prove this without crossing the original row, consider $A_h=s+g_h-r+1$.
Let $T$ consist of the package roots and all vertices deleted at current degree
at least $h$, with $t_T=a+C_h$.
If $A_h>Z+t_T$, strong length $s>Z$ permits

\[
k=\max(0,Z+t_T-s+1)\le\min(t_T,g_h-1).
\]
Choose $k$ vertices of $T$ as $T_B$, and put $T_R=T\setminus T_B$.
Color core interiors and edges incident with $T_R$ red, and everything else blue.
Red matching size is at most $Z+t_T-k$; every blue size-$h$ star meets $T_B$,
so there are at most $k$ disjoint such stars.
Set $I=s-Z-t_T+k\ge1$, $J=g_h-k\ge1$, with $I+J\ge r$.
There are legal positive $i\le I,j\le J,i+j=r$.
Complete red $F_i$ needs more matching edges, and complete blue $H_j$
needs more size-$h$ stars, a contradiction.
The reverse direction is identical; $h=1$ is the same argument for two matchings.
Use the actual identity $e(Y)=\sum_{h=1}^D C_h(Y)$, retain $b=C_D(Y)$,
and add the original root and core saving once:

\begin{equation}\label{eq:unit-area}
e(G)-E\ge a+Z-L_D+b-\Gamma\ge Z-L_D^0+b-\Gamma.
\end{equation}
\subsection*{The single automatic entry}
If $R+b\ge M_2$, choose $\ell=M_2$ and $h=(M_2-R)_+\le b$.
All added vertices have current degree $D$.
The fixed lists retain every nonunit head,
and the first $D$ columns cost $(2Z-\ell)+(B-1)Z+\ell=DZ$.
Actual saving gives $e(X)-E'\le\Gamma-1<0$, with rank at least
$\nu(Y)+d-Z+\ell-h>0$. Hence $R+b<M_2$ in an unexcluded input.
Equation \eqref{eq:gap1-dgate} gives

\[
2d+R-M_2\ge R-c-2c_0-2\Gamma+2-2\tau
\ge4+(D-6)(\rho+\lambda)/2>0.
\]
The complete endpoint supplies $C_B(Y)>M_2-R$.
If the history also supplies $t=Z-R$ vertices, the actual original-list menu successor
has upper cost $\Gamma-Z+L_D^0-b-1$.
Nonnegativity forces $L_D^0\ge Z+b+1-\Gamma$.
If the history is shorter than $t$, \eqref{eq:unit-area} and $e(G)-E=-1$
force the same inequality.
The cases are complementary: their output is a necessary constraint on one
original column, not income that can be combined later.
\end{proof}

\section{A size-three attaining head at gap one}\label{app:gap1-three}
\begin{proposition}\label{prop:gap1-three}
For $D\ge8$ and $\gamma=1$, an actual height-$D$ attaining pair $3/(D-2)$
excludes the original minimum strict source.
\end{proposition}
\begin{proof}
Use the same package and deletion history as Appendix \ref{app:gap1-entry}.
Near-top factor coloring gives uniform width at least $(L-M_3-c_0)_+$
for $3/(D-2)$ on this same $Y$, where $L=L_D^0$.
Earlier larger heads are retained by the two-direction restoration and cumulative closure.
Choose

\[
\ell=\max(M_5,R+b)<M_2,\qquad h=\ell-R\ge b,\qquad
X=Y-\{\text{the first }h\text{ vertices}\}.
\]
The automatic entry actually supplies these vertices:
the first $b$ have degree $D$ and the others degree at least $B=D-1$.
The saving is at least $Bh+b$, and the final cap is $B$.
The auxiliary uniform width on this same $X$ is

\[
H=L-M_3-c_0-h
\ge Z-M_3-c_0+1-\Gamma-(h-b)>0.
\]
For $\ell=M_5$, positivity follows from
$M_3+M_5\le5n/12$, $Z\ge4c$, and $-\Gamma>c_0$;
for $\ell=R+b$ it is immediate.
The fixed allowance lists retain every head at least five.
Other cumulative losses are at most $(M_j-\ell)_+$, and unit loss is at most $Z-\ell$.
They give every pair of $P_{r+a+h+\ell}$,
with $N'\ge\nu(Y)+d-Z+R>0$ and $N'<N$.
The complete tax on the first $D$ columns is

\[
\mathrm{tax}\le
\begin{cases}
DZ+2M_2+M_3+M_4-4\ell,&M_5\le\ell\le M_4,\\
DZ+2M_2+M_3-3\ell,&M_4\le\ell\le M_3,\\
DZ+2M_2-2\ell,&M_3\le\ell\le M_2.
\end{cases}
\]
\[
g=e(X)-E'\le U=\Gamma+\mathrm{tax}-DZ+h-b-1.
\]
The low mixed-column bounds are
$2M_2-M_3\le\ceil{5n/12}\le Z$,
$M_2+M_3-M_4\le\ceil{23n/60}\le Z$,
$M_2+M_4-M_5\le\ceil{11n/30}\le Z$, and
$2M_3-M_5\le\ceil{n/3}\le Z$.
The highest pre-top column has head sum at least nine, so one head is at least five;
its tax is at most $M_2$.
The original top column is charged only in $\Gamma$.
A positive new column above height $D$ is avoided under cap $D-1$;
otherwise the new highest height is at most $D$.

For odd $D$, the two even target caps $2,D-3$ sum to $D-1$,
contradicting the positive auxiliary source.
Let $D=2m$, $m\ge4$. Signed exposure for that same source and the whole-core price give

\[
\rho_{D-1}(X)\ge(m-1)H,\qquad
g\ge AH-K,\quad A=(2m-3)(m-1)/2,\quad K=m(m-1).
\]
The actual highest width of the fixed lists is at most exposure.
The auxiliary $H$ only bounds exposure; it is not substituted for that highest width.
We have $A\ge15/2$ and $A\ge m+1$.
For $\ell=M_5$, put $Q=37/60+5A/12$. Subtracting column by column gives

\[
\begin{split}
U-AH+K\le{}&
(37/30-A/6)Z+Qc+(A+1)\Gamma\\
&-(A+1)(R+b)+Ac_0+K-A-1\\
\le{}&[Q+A-(A+1)(m+1)/2]\rho+K-2A-2<0.
\end{split}
\]
The first line's $Z$ coefficient is at most $-1/60$, and
$Qc-(A+1)R\le Q\rho-(A+1)$.
The bracket in the second line is at most $-113/60-13A/12<0$;
the constant is $K-2A-2=-m^2+4m-5\le-5$.
For $\ell=R+b$, the piecewise tax decreases with $\ell$, so

\[
U\le\Gamma+9Z/10+9c/20+7/3,\qquad
H\ge Z/2-c/4-c_0+1-\Gamma.
\]
\[
\begin{split}
U-AH+K
&\le(9/10-A/2)Z+(9/20+A/4)c\\
&\qquad +(A+1)\Gamma+Ac_0+K-A+7/3\\
&<K-(m-1)A-m+10/3
\le13/3-2m<0.
\end{split}
\]
Here $Z\ge mc$ makes the first two terms negative, and the package inequality gives
$(A+1)\Gamma+Ac_0\le-(m-2)A-m+1$.
Both positions contradict the lower bound on the same $X$, completing the exit
for all original parameters.
\end{proof}

\section{A size-two attaining head at gap one: an actual two-neighbor cover}\label{app:gap1-two}
\begin{proposition}\label{prop:gap1-two}
Let $D\ge13$, $\gamma=1$, and let $2/(D-1)$ attain height $D$.
Then the original minimum strict source does not exist.
\end{proposition}
\begin{proof}
Reverse colors for the other orientation.
On the same package put $B=D-1$,
$n=2Z+c$, $c\le R-1+\rho$, and $Z\ge\ceil{D/2}c\ge7c$.
A weaker form of the exact package inequality is

\[
\Gamma\le-\bigl((D+1)\rho+(D-5)\lambda\bigr)/4.
\]
\[
\ell=R+b<M_2,\quad C_B(Y)>M_2-R,\quad
L=f_2+g_B-R\ge Z+b+1-\Gamma.
\]
The original $3/B$ candidate has zero width:
for even $D$ use even caps $2/(B-1)$;
for odd $D$ its signed exposure is at least $D-2$, exceeding the
original bound $(D+1)/2$.
Thus $f_3+g_B\le r-1$, in particular $f_3\le R-1$ and $g_B\le R$.
Odd-degree packages already have $\rho\ge2$;
Lemma \ref{lem:rho1} excludes $\rho=1$ for even degree.
Hence $\rho\ge2,c\ge\rho\ge2$.
Delete precisely the complete degree-$D$ stage of length $b$,
obtaining the cap-$B$ graph $X$; the current row is $\ell+1$.

\subsection*{Longer rows: $\ell\ge M_6$}
Ordinary allowance intersection gives every new pair and
$N'\ge\nu(Y)+d-Z+R>0$, $N'<N$.
Low mixed columns satisfy
$M_u+M_v-M_{u+v-1}\le5n/12+1\le Z$.
Their taxes are $2Z-\ell$ for the first column and
$Z+(M_j-\ell)_+$ thereafter.
At height $D$ the head sum is at least fourteen, so both heads cannot be damaged;
its tax is at most $M_2$. Consequently

\[
g=e(X)-E'\le U=\Gamma+2M_2+M_3+M_4+M_5-5M_6-1,
\qquad H\ge Z-M_2+1-\Gamma>0.
\]
A higher positive column is directly avoided by edge coloring.
Otherwise the whole-core price on the same graph gives $g\ge\alpha H-K$,
where $\alpha=(D-3)/2\ge5$ and
$K=\floor{B/2}(\floor{B/2}+1)$.
Using $M_6\ge n/7-1$, upper minus lower is at most

\[
(239/210-\alpha/3)Z+(239/420+\alpha/3)c
+(\alpha+1)\Gamma-(\alpha+1)+K+5<0.
\]
The first two terms are at most $-43c/420$ by $Z\ge6c$ and $\alpha\ge5$.
The remainder is at most $(-D^2+2D+39)/8<0$ by $\rho\ge2,\lambda\ge1$.

\subsection*{One cover for a short row $\ell<M_6$}
Independently select each vertex into $S_0$ with probability $2/7$.
Add an unselected internal degree-$D$ vertex if it has at most one selected neighbor;
add an internal degree-$(D-1)$ vertex if it has none.
No other vertex needs repair.
The actual set $S$ leaves internal degree at most $D-2$,
and each inclusion probability is at most

\[
\alpha_D=\frac27+\left(\frac57\right)^D
\left(1+\frac{2(D-1)}7\right)
\le\alpha_{12}=\frac27+\frac{29\cdot5^{12}}{7^{13}}<\frac9{25}.
\]
Only the independent neighbor choices for each vertex are used;
different neighborhoods need not be independent.
Choose such a set in the actual isolated core union $C$, and enlarge it to

\[
k=|S|=\floor{(9n+c)/25},\qquad k_0=(9n+c)/25=18Z/25+2c/5.
\]
\[
k\ge M_2,\qquad k<Z,\qquad k-2\ell\ge c.
\]
The last inequality follows from
$k_0-2n/7-c=(13n-168c)/175>0$ and integrality of $2\ell+c$;
also $k_0\le59Z/75<Z$.
Exterior ports are neither covered nor precolored.

\subsection*{Two actual internal colorings and one intersection}
In the first state, color $C$ entirely blue, obtaining $(F,H_{\rm mono})$.
In the second, color internal edges incident with $S$ red and the others blue,
obtaining $(F_{\rm cov},H_{\rm rest})$.
The second red matching is at most $k$, with vertex budget $n$.
Blue matching is at most $\floor{(n-k)/2}$, with vertex budget $n-k$;
its cap $D-2$ retains all original size-$B$ heads.
Take upward cumulative closures for every residual list.
Define the fixed cumulative lists

\[
K_j=\min(f_j,(F_{\rm cov})_j+\ell),\qquad
J_j=\min(g_j,(H_{\rm mono})_j+\ell,(H_{\rm rest})_j).
\]
For a new pair $I+J'=2\ell+1$, use the first state's old pair
$(I,J'-\ell)$ if $J'>\ell$.
Otherwise $I>\ell$ and use the second state's pair $(I-\ell,J')$.
Each old pair is positive and legal; suffix trimming according to the fixed lists
gives every pair of $P_{2\ell+1}(X;K,J)$.
Since $f_3\le R-1<\ell$, the red truncation retains all original heads at least three.
The blue size-$B$ head is retained by $g_B\le R\le\ell$ and the second state's strict cap.
Red and blue length losses are at most $k-\ell,Z-\ell$, respectively, so

\[
N'\ge\nu(Y)+d-k+R>0,\qquad N'<N.
\]
In the corresponding allowance convention, the first blue residual is
$H_{\rm rest}$. Its internal blue degree cap is $D-2$, whereas
$B=D-1$, so $(H_{\rm rest})_B=g_B>0$ and every lower original blue
threshold survives. Moreover $g_B\le R\le\ell$ gives $J_B=g_B>0$
in the final fixed list. Thus no absent blue threshold is used as a
candidate in the following comparison.

\subsection*{Complete column taxes, including the $2/2$ candidate}
Write the two losses as $x_j,y_j$. Then

\[
\begin{gathered}
x_1\le k-\ell,\quad x_2\le M_2-\ell,\quad x_j=0\ (j\ge3),\\
y_1\le Z-\ell,\quad y_B=0,\quad
y_j\le\max((M_j-\ell)_+,\floor{(n-k)/(j+1)})\ (2\le j<B).
\end{gathered}
\]
The first two column taxes are $Z+k-\ell$ and $Z+M_2-\ell$.
For $3\le h\le B$, the four candidate types give
$k+y_h$, $M_2+y_{h-1}$, $Z$, and $\ell+y_2$.
The last is at most $Z$, and the unit/size-two candidate in the second column
also satisfies its stated bound.
Use these two independent bounds, without assuming $\ell\le\floor{k/3}$:

\[
2k/3+\ell+1\le13Z/15+1\le Z,\qquad
\ell+(n-k)/3\le23n/63+1/3\le299Z/378+1/3<Z.
\]
These use $k<4Z/5$, $\ell<Z/3$, $k\ge n/3-1$,
$n\le13Z/6$, and $Z\ge14$.
The actual height-$D$ pair $2/B$ gives tax $M_2$ and
$H\ge Z-M_2+1-\Gamma>0$.
Thus the same fixed lists satisfy

\[
g\le U=\Gamma-1+k-Z+2M_2-2\ell+
\sum_{h=3}^{B}\max(0,k+y_h-Z,M_2+y_{h-1}-Z).
\]
\subsection*{Analytic endpoints of one cost function}
First replace all point capacities by real quotients.
The resulting continuous upper bound is nondecreasing in $k$:
the branches of $k+y_h$ have slopes $1,h/(h+1)$,
and those of $n/3+y_{h-1}$ have slopes $0,-1/h$.
For $h\ge6$, the sole negative branch satisfies
$n/3+(n-k)/h-Z\le4n/9+1/6-Z<0$ and cannot be active.
The remaining possible decreases sum to at most
$1/3+1/4+1/5=47/60$; the exterior $+k$ makes the total slope at least $13/60$.
It is therefore safe to replace $k$ by $k_0$.
This continuous upper bound decreases in $\ell$, with slope magnitude at most $B$.
Since $\ell\ge R\ge c-\rho+1$, replacing $\ell$ by $c$ costs at most $B(\rho-1)$.
Define

\[
\bar y_j(t)=\max(n/(j+1)-t,(n-k_0)/(j+1)),
\]
\[
\begin{split}
\Phi_\alpha(Z,c)={}&k_0-Z+2n/3-2c-\alpha(Z-n/3)\\
&+\sum_{h=3}^{B}\max(0,k_0+\bar y_h(c)-Z,n/3+\bar y_{h-1}(c)-Z).
\end{split}
\]
For fixed $Z$, this is convex in $c$.
For $0\le c\le Z/6$ and $\alpha\ge5$, its two analytic endpoints are

\[
\Phi_5(Z,0)=-43Z/1050,\qquad
\Phi_5(Z,Z/6)=-179Z/1800.
\]
Expand the finite sum directly: at the first endpoint the contributions
for $h=3,4,5,6$ are $Z/3,Z/6,Z/15,Z/175$;
at the second, those for $h=3,4,5$ are $5Z/18,7Z/72,Z/60$.
All other terms vanish. Hence $\Phi_\alpha<0$.
Compare with the same $X$'s whole-core price $g\ge\alpha H-K$.
The remaining constants are at most

\[
(\alpha+1)\Gamma+(D-1)\rho+K-\frac{3(D-1)}2
\le-\frac{(D-1)(D-9)}8<0.
\]
As before, first avoid any new column exceeding the cap; otherwise apply the price.
The two row lengths exhaust the cases.
Every new pair, unit, and actual port is retained within the same fixed intersection.
\end{proof}

\section{Weighted whole-core descent at the even near-top layer}\label{app:even-weighted}
This appendix proves the exit for $D=2p\ge12$ and $H=D-1$; the final coverage uses only $D=12$. All objects belong to one original minimal source. Put $R=r-1$ and $\lambda=L_{D+1}$. The strong odd-cap price in Appendix \ref{app:stage-prices} applies to the actual successor without requiring a positive successor matching deficit.
\subsection{A whole-core package from one capacitated flow}
The highest column has a pair of even heads $\alpha+1,\beta+1$, where $\alpha+\beta=D$ and $3\le\alpha\le\beta$ are odd; two odd heads are excluded by an even-cap coloring. The minimum-excess factor formula and the actual switch construction in Section \ref{sec:fractional} give constant-cap deficiency at least $\lambda$. Truncating the caps to $\min(\alpha,d_x),\min(\beta,d_x)$ changes no edge set's total excess, so the maximum interval-cap certificate deficiency is also at least $\lambda$.

For $\alpha<\beta$, assign root capacities
$c_x=\min(d_x,1+\kappa(D-d_x))$, $\kappa=(\alpha-1)/(\beta-\alpha)$.
In the central case assign capacity $d_x$ to nonfull roots and one to full roots. Fix a factor certificate. Besides the matching-cut inequality of Section \ref{sec:fractional}, every zero-boundary odd component must contain an exposed vertex: otherwise it has a perfect matching. Thus the $c_*$ term in the half-cut inequality can be replaced by the number of vertices of this same exposed set in these components. Weight the inequalities by $\theta=(\beta-\alpha)/(D-2)$ and $1-\theta$, and average over the same integral flow decomposition. The certificate deficiency is at most the unassigned flow mass $\delta$. In the central case use only the half-cut inequality. Hence $\delta\ge\lambda$.

Choose a smallest core family $I$ attaining $\delta=|I|-c(N(I))$, and put $T=N(I)$. If some root had $c_x=d_x$, removing all its neighboring cores would release at least their number in root capacity, preserving deficiency and decreasing the family, a contradiction. Thus $c_x=1+\kappa(D-d_x)<d_x$ in the noncentral case, and all roots of $T$ are full in the central case. Write
\[
t_0=|T|,\quad \sigma=|I|-t_0,\quad
\zeta=\sum_{x\in T}(D-d_x),\quad \beta_T=e(T,V\setminus(T\cup V(I))).
\]
\[
|Q|=2z_Q+1,\quad \mu=\sum_Iz_Q,\quad M=\mu+t_0,\quad
W_I=\sum_I(e(Q)-Dz_Q).
\]
Since $\kappa\ge1/(p-3)$, we have $\zeta\le(p-3)(\sigma-\lambda)$. Let $J$ consist of all internal core edges and all original edges incident with $T$, and set $O=G-E(J)$. The ordinary root assignment and near-perfect core matchings give
\begin{align}
\nu(J)&=M,\qquad \nu(O)=\nu(G)-M,\qquad e(J)=DM+w,\\
w&=p\sigma-\zeta/2+\beta_T/2,\qquad W_I=w+\zeta+e(T),\\
2w&\ge(p+3)\sigma+(p-3)\lambda+\beta_T,\quad
w\ge p\lambda,\quad \sigma\ge\lambda. \label{eq:weighted-pack}\\
s,t&\ge M+R+d,\qquad N\ge2M+R+2d. \label{eq:weighted-length}
\end{align}
For the matching identity, every matching of $J$ has at most $t_0$ root edges and $\mu$ core edges. Retain the original root assignment and match each entered core after deleting its attachment vertex to attain equality. After actually deleting $T$, the cores in $I$ are isolated, giving the stated matching number of $O$. The complete exterior sources in Section \ref{sec:regions} give the strong lengths.
\subsection{One weighted vertex selection}
For a positive-weight core $Q$, put $n=2z+1$, $\delta_Q=D-2w_Q$, and $\alpha=D-3$. Select vertices independently with probability $1/5$ to obtain $A$, then take a maximal independent set $B$ among the internal degree-$D$ vertices not dominated by $A$. Put $S=A\cup B$. Then $\Delta(Q-S)\le D-1$, while $B$ has no deficit, internal edge, or edge to $A$. For the debt of this same set,
\[
a_Q=\sum_{v\in S}(D-d_Q(v))+e(Q[S]),\qquad k_Q=|S|,
\]
\[
\mathbb E[\alpha k_Q+2a_Q]\le
n\{\alpha/5+D/25+\alpha(4/5)^{D+1}\}+9\delta_Q/25.
\]
Simplicity and positive weight give $n\ge D+1$. For even $D\ge12$,
$(D-3)(4/5)^{D+1}<1/2$: at $D=12$ this is
$18\cdot4^{13}<5^{13}$, and the ratio after increasing $D$ by two is
$16(D-1)/(25(D-3))<1$.
The difference between the target $\alpha z-2w_Q$ and the expectation exceeds
$(13D^2-132D+5)/50\ge293/50>5$.
Integrality supplies an actual choice with
\[
1\le k_Q\le z_Q-1,\qquad
(D-3)(z_Q-k_Q)-2a_Q\ge2w_Q+6,\qquad
e(Q)-e(Q-S)=Dk_Q-a_Q.
\]
In a nonpositive-weight core take a maximal independent set of internal full-degree vertices; its debt is zero and its size at most $z_Q$. Use the empty set in a singleton. Put $s_0=\sum_I(z_Q-k_Q)$ and $A_c=\sum_{w_Q>0}a_Q$. Then
$(D-3)s_0\ge2(W_I+A_c)$.
Fix a highest-degree history of $O$ and put $b=C_D(O)$. Actually delete $T$, the sets $S_Q$, and the first full-degree stage of $O$ to obtain
\[
\begin{aligned}
K_1&=M-s_0+b, & X_1&=G-S_1,\\
\Delta(X_1)&\le D-1, & &\p_{r+K_1}(X_1),\\
e(G)-e(X_1)&=DK_1-a, & a&=W_I+A_c-w\ge0.
\end{aligned}
\]
The selected roots, core sets satisfying the capacity bounds, and the full-degree independent set of $O$ can be charged using the same root assignment and near-perfect matchings, giving $1\le K_1\le\nu(G)<N$. The original positive-block selector gives $\rho_D\le p+2$, hence $\lambda\le p+2\le Dr-1$. A strict descent finishes the proof. Otherwise the actual surplus
$v_1=e(X_1)-E_{r+K_1}$ satisfies
$0\le v_1\le a+\lambda-1\le D(r+K_1)-2$.
Apply Lemma \ref{lem:strong-odd-price} with $j=r+K_1$, $k=D-1=2p-1$, and $N_X=N-K_1$. The actual source is fixed and complete, its rank is positive and smaller, and $v_1\le a+\lambda-1\le2pj-2$. Apply Lemma \ref{lem:local} to the single-star consequence of the current highest column: $\ell_j-1\le\Delta(X_1)\le D-1$, so the current profile height is at most $D=k+1$. The highest-column consequence gives $(L_D-K_1)_+\le\rho_{D-1}(X_1)$; use the single-term price, which does not assume positive matching deficit, to obtain the next bound. If $L_D\le K_1$, that bound holds directly.
\[
L_D-b\le M-s_0+
\floor{\frac{2p(p-1)+2a+2\lambda-2}{2p-3}}.
\]
Put $B=s_0+w-\lambda$ and $U_0=(D-3)s_0-2(W_I+A_c)\ge0$. The denominator times $B+1$, minus the numerator, is exactly
\[
U_0+(2p-1)(w-\lambda+1)-2p(p-1)
\ge U_0+p+(2p-1)(p-1)(\lambda-1)>0.
\]
\begin{equation}\label{eq:weighted-top}
L_D+\lambda\le M+w+b.
\end{equation}
\subsection{Full column taxes, inventory, and payment}
Put $M_u=\sum_I\floor{(2z_Q+1)/(u+1)}$ and $F=\mu-M_2$. Each core has the full integral damage menu $[0,z_Q]$ for two nonunit heads: choose $x$ vertices and color their incident edges red. When $2x\le z_Q+1$, the remaining blue vertex count is below $3(z_Q-x+1)$; interchange colors for the other half. Root slots add $[0,t_0]$. Thus every low-column tax is at most $M$. If a nonunit attaining pair has a count exceeding its $M_u$, color all cores that color and allocate root slots: its tax is at most $t_0+M_2$. Otherwise its width is at most $2M_2-R$.

At height $D$, take $u+v=D+1$ with $u$ even. An even-cap $(u,v-1)$ coloring and root slots give tax $t_0+M_{\min(u,v)}$, unless both counts are below their corresponding taxes. In that dead case, if an original $(u+1)$ head exists, the even-cap coloring forces $f_{u+1}+g_v\le R$. Otherwise the opposite list has a $(v+1)$ head and the highest column forces $f_u+g_{v+1}-R\le\lambda$. Both alternatives give the safe tax $t_0+M_2+\lambda-1$. The complete matching reserve on this same $O$ yields
\[
F\ge0,\quad 3F\ge\mu-t_0-\sigma,\quad
d\ge F+w-2\lambda+2.
\]
\[
d\ge2F+w-2\lambda+2
\quad\hbox{if a live nonunit pair attains a low column.}
\]
For $1\le h\le H$, define the actual unit axes by
\[
 A_h=\max\bigl(\{f_h+t-R:f_h>0\}\cup
                    \{g_h+s-R:g_h>0\}\bigr).
\]
The set of candidates is nonempty in this range. Each included candidate
has width at least $M+d+1$, by the strong lengths, and its whole-piece
restoration gives $C_h(O)\ge A_h-M$.
 The complete forest/matching endpoint also gives
$C_H(O)\ge2d+H-2$. Let $g\le\mu$ be the largest gap in whole-core subset sums and put $\tau=(g-t_0-R)_+$. If $C_H<\tau$, then $\tau\le\mu-t_0-R$. A dead nonunit pair cannot attain a low column, since
\[
A_h-(2M_2-R)\ge w-2\lambda-\sigma+R+3
\ge(p-3)\lambda+\beta_T/2+R+3>0.
\]
If a live nonunit winner exists, the strengthened reserve gives
\[
C_H-\tau\ge
\frac{\mu-t_0}{3}+\left(2p-\frac{16}{3}\right)\lambda+\beta_T+D+R+1>0.
\]
If no such winner exists, all low columns are unit axes. Coarea on the same history and \eqref{eq:weighted-top} give
$e(G)\ge\sum_{h\le D-1}L_h+M+w+b\ge E_r$. Thus $C_H\ge\tau$ in all cases.

Delete the first $\tau$ vertices of the same history of $O$, obtaining $Y=O-U$ and $K=M+\tau$. The menu
$\operatorname{Sums}(z_Q)+[0,t_0]$ has gap at most $R+\tau$. Each new pair can therefore be restored using monochromatic whole cores and actual root slots, giving $\p_{r+K}(Y)$ with the fixed original lists. Equation \eqref{eq:weighted-length} gives $1\le N-K<N$. Charge only the disjoint edge sets $J$ and $O$. If $\tau\le b$, the saving is $DK+w\ge DK+\lambda$. If $\tau>b$, it is
\[
(D-1)K+M+w+b\ge(D-1)K+L_D+\lambda
\ge\sum_{h=1}^{D+1}\min(K,L_h)=E_r-E_{r+K}.
\]
Both give a strict positive smaller-rank successor, a contradiction. The auxiliary cost of $X_1$ was used only to prove the high-column constraint and was not added to the final saving.
\section{The complete degree-eleven near-top exit}\label{app:D11}
Take $D=11,H=10$. Section \ref{sec:fractional}, at a highest attaining pair, supplies one root package. Let $C$ consist of its positive-weight whole cores and write
$Z=\nu(C)$, $c=\#C$, $W=e(C)-11Z$, $\delta=\zeta+\varepsilon$,
$Y=(G-A)-E(C)$, $R=r+|A|-1$, and $\Gamma=\delta+\lambda-W$.
The odd-cap failure inequality $(\alpha-1)\zeta\le\beta\rho-10\lambda$, with $\alpha\ge3,\beta\le8$, gives
\begin{gather}
\rho\ge\lambda+1,\quad \zeta\le4\rho-5\lambda,\quad
\Gamma\le-7\rho/2-3\lambda/2\le-17/2, \label{eq:D11debt}\\
c\ge\rho+1\ge\lambda+2,\quad Z\ge6c,\quad c\le R-1+\rho,\\
s,t\ge Z+R+d,\quad \nu(Y)=\nu(G)-|A|-Z>Z,\\
e(G)-e(Y)=11(|A|+Z)+\lambda-\Gamma.
\end{gather}
\subsection{The necessary high column on one exterior history}
Fix a highest-degree history of $Y$, and put $b=C_{11}(Y)$, $L=L_{11}-|A|$, and $M_2=\floor{(2Z+c)/3}$. The integral core menu gives tax $Z$ at every low column. A pair $x+y=12$ at height eleven has an earlier higher head on at least one side, say $f_{x+1}>0$. Every core has a port. A near-perfect matching exposing a port vertex covers all its internal degree-eleven vertices. Give this matching the appropriate color and split the remaining degree-ten graph into even caps, obtaining red cap $x$ and blue cap $y-1$. Retaining the earlier heads through
$\max(f_x-M_x,f_{x+1})>0$, one complete copy gives height-eleven tax at most $M_2$. The reserve forces
\[
d\ge Z-M_2-\Gamma+1.
\]
A unit-axis width is at least $Z+d+1$, strictly exceeding the doubly dead nonunit width $2M_2-R$: the difference is at least $3-\rho-\Gamma>0$. If a nonunit pair attains a low column, its tax drops to $M_2$, so
$d\ge2(Z-M_2)-\Gamma+1$.
The complete size-ten forest/matching endpoint gives $C_{10}(Y)\ge2d+8$.
Put $\tau=(Z-R)_+$. If $C_{10}<\tau$ and a nonunit low-column winner exists, these inequalities force
$0>3Z-4M_2+R-2\Gamma+10>0$, a contradiction.

If $b\ge\tau$, the two-color whole-core menu and the first $\tau$ full-degree vertices immediately give paid strict descent. If $b<\tau\le C_{10}$, the same menu saves at least
$11(|A|+Z)+\lambda-\Gamma+10\tau+b$; failure to exit forces
$L\ge Z+b+1-\Gamma$. If inventory is insufficient, every low column is a unit axis. The actual original-row coloring in Appendix \ref{app:gap1-entry} gives $C_h(Y)\ge L_h-|A|-Z$ for $1\le h\le10$. Coarea on the same history again gives that conclusion. Every input not already excluded therefore satisfies
\begin{equation}\label{eq:D11lower}
L_{11}-|A|-Z-b\ge1-\Gamma.
\end{equation}
\subsection{Internal full-degree sets and a weighted one-port set}
\begin{lemma}\label{lem:odd-deficit-independent}
Let \(D\ge3\) be odd. Let \(Q\) be a connected simple graph of order \(2m+1\), with
\(\Delta(Q)\le D\) and \(0<c_Q=D|Q|-2e(Q)<D\).
Put \(F=\{v:d_Q(v)=D\}\).
Then \(Q[F]\) has a maximal independent set of size at most \(m-1\).

\end{lemma}\begin{proof}
The one-vertex graph has \(c_Q=D\) and is excluded, so \(m\ge1\).
For any independent \(I\subseteq F\),
\[
D|I|\le D(2m+1-|I|)-c_Q,
\]
so \(|I|\le m\).
Choose a minimum-cardinality maximal independent set \(I\) of \(Q[F]\).
Suppose \(|I|=m\); otherwise the conclusion holds.
Put \(B=V(Q)\setminus I\) and
\(h=e(Q[B])=(D-c_Q)/2\le(D-1)/2\).

For \(b\in B\cap F\), set \(X=N_I(b)\), \(d_b=|X|\), and
\(C_b=\{v\in B\cap F:N_I(v)\subseteq X\}\).
Then \(D-h\le d_b\le D\), and
\[
D|C_b|\le D|X|+2h<D(|X|+1),
\]
hence \(|C_b|\le d_b\).
Take a maximal independent set \(K\) of \(Q[C_b]\) containing \(b\).
The set \((I\setminus X)\cup K\) is maximal in \(Q[F]\):
the removed vertices of \(X\) are adjacent to \(b\), \(K\) dominates \(C_b\), and every other full-degree vertex of \(B\) has a neighbor in \(I\setminus X\).
Minimality of \(|I|\) forces \(|K|\ge|X|\).
Thus \(K=C_b\), \(|C_b|=d_b\), and \(C_b\) is independent.
It follows that
\[
d_b(D-d_b)\le e(C_b,B\setminus C_b)\le h.
\]
If \(d_b<D\), put \(z=D-d_b\in[1,h]\).
Then
\[
z(D-z)\ge D-1>h,
\]
because \(z(D-z)-(D-1)=(z-1)(D-z-1)\ge0\), a contradiction.
Consequently every \(b\in B\cap F\) has \(d_b=D\).
The sets \(X,C_b\) then form an isolated \(K_{D,D}\), contradicting connectedness and odd order.

If \(B\cap F=\varnothing\), every vertex of \(B\) has a positive deficiency, so
\(c_Q\ge|B|=m+1\). Since \(I\ne\varnothing\), one of its vertices has \(D\) distinct neighbors in \(B\), giving \(m+1\ge D\), again a contradiction.
The claimed maximal independent set therefore exists.
\end{proof}
For $D=11$, a positive core with cut at least three has weight at most four. The independent set $S_Q$ just constructed therefore satisfies
$4(z_Q-|S_Q|)\ge w_Q$, with zero debt.
A one-port core has internal degree ten at its port vertex $x$, eleven elsewhere, and weight five. We claim that it has a set $S$ dominating all full-degree vertices such that
\[
4|S|+e(Q[S])+\mathbf1_{x\in S}\le4z_Q-5.
\]
For order $n\ge21$, select $S_0$ independently with probability $1/4$ and add all undominated full-degree vertices $U$. There are no edges between $S_0$ and $U$. With $t=(3/4)^{12}$,
$\mathbb E|S|=n/4+(n-1)t$,
$\mathbb Ee(S)\le e(Q)(1/16+t)$, and
$\mathbb E\mathbf1_{x\in S}=1/4$.
The target minus the expected cost is at least
$(11922717n-237438087)/33554432>0$.
For orders $13,15,17,19$, a uniform three-set leaves an expected number
$0,14/455,4/17,210/323$ of full-degree vertices undominated. Each is below one, so a dominating three-set exists. Its cost is at most sixteen, while the target is at least nineteen.

Let $k$ be the total number of selected core vertices and let
$c_{\rm del}=\sum(e(Q[S_Q])+\mathbf1_{\text{one-port vertex selected}})$.
Then
\[
c\le k\le Z-c,\qquad 4(Z-k)-c_{\rm del}\ge W.
\]
\subsection{Contradiction from the same first stage}
Delete $A$, then the sets $S_Q$, then all $b$ degree-eleven vertices on the same history of $Y$, obtaining $X_1$.
\paragraph{The degree-ten price is Proposition~\ref{prop:selector}.}
Retain the notation of Appendix~\ref{app:D11}.  Let
\[
 K_1=|A|+k+b,\qquad r_1=r+K_1,\qquad
 N_1=N-K_1,\qquad h_1=(L_{11}-K_1)_+.
\]
Here \(h_1\) denotes the height-eleven column of the unchanged two lists
on the new row.  The previously constructed sets dominate all internal
degree-eleven vertices, the roots have been deleted, and the complete
degree-eleven stage of the same exterior has been deleted.  Hence
\(\Delta(X_1)\le10\).

The exterior stage is an independent set of original degree-eleven
vertices in \(Y\), so Hall counting gives \(b\le\nu(Y)\).  Using
\(k\le Z-c\) and \(\nu(G)=|A|+Z+\nu(Y)\), we obtain
\[
 K_1\le\nu(G)-c<N,\qquad 0<N_1<N.
\]
Moreover \(K_1\ge c>\lambda\), so the original height-twelve column
has disappeared.  The highest value of the new profile is therefore
at most eleven.  The actual vertex-deletion operation gives
\(P_{r_1}(X_1;\nn,\mm)\) for every new legal corner; its two lists
are fixed.  Minimality at this positive smaller rank gives
\[
 v_1:=e(X_1)-E_{r_1}\ge0.
\]
The exact fee \(11K_1-\delta-c_{\rm del}\), together with the full-profile
loss at most \(11K_1+\lambda\), also gives
\[
 v_1\le\delta+c_{\rm del}+\lambda-1.
\]

All hypotheses of Proposition~\ref{prop:selector} now hold for this
actual \(X_1\), with \(D=10\) and \(p=5\).  It gives
\[
 v_1\ge\max\{0,5\rho_{10}(X_1)-30-h_1\}.
\]
The constants come from leaving the six largest positive whole cores
of \(X_1\) unselected.  If \(Z'\) is the largest selected capacity,
those six actual cores supply at least \(6(2Z'+1)>11Z'\) original
degree-ten vertices, hence the required independent buffer.
No estimate involving a number of selected cores is needed.
The proof of that proposition handles a further deletion parameter
reaching the current rank by its actual edge fee, rather than by
invoking a zero-rank source.  It applies for either sign of
\(N_1-\nu(X_1)\).

Finally \(h_1\le\kappa_{10}(X_1)\le\rho_{10}(X_1)\), so
\[
 4\rho_{10}(X_1)\le v_1+30.
\]
Since \(L_{11}-K_1\le h_1\) and
\(4(Z-k)-c_{\rm del}\ge W\), the whole comparison is
\[
\begin{aligned}
 L_{11}-|A|-Z-b
 &\le k+\rho_{10}(X_1)-Z\\
 &\le \frac{\delta+\lambda+c_{\rm del}+29+4k-4Z}{4}\\
 &\le \frac{\delta+\lambda-W+29}{4}
  =\frac{\Gamma+29}{4}.
\end{aligned}
\]
Together with \eqref{eq:D11lower}, this forces \(\Gamma\ge-5\),
contradicting \eqref{eq:D11debt}.  No structural conclusion about the
original minimum source has been assigned to \(X_1\).

\section{Degree nine: the original factor certificate and two one-port regions}\label{app:D9}
\subsection{A global coloring when the highest width is two}
The original positive-block selector gives $W_+\le\lambda+19$ and $9\rho\le2\lambda+38$.
Write $a$ for the odd cap. Since the highest head is at most eight, $a\in\{3,5,7\}$.
On one certificate,
$\Psi_x=q+e(S,T)-x|S|-(9-x)|T|$,
$\Psi_1\le\rho$, and $\Psi_9\le0$; the latter follows because every counted component has a positive odd number of $S$ ports.
Thus $\lambda\le\lfloor3\rho/4\rfloor$, and the edge bound gives $\rho\le4$.
If $a\ge5$ and $\Psi_a>0$, then $|S|-|T|\ge1$ and $\Psi_1\ge a\ge5$, impossible. Hence the highest heads are $4/7$, and a maximum canonical certificate, with a largest root set, satisfies
\[
|S|-|T|=1,\quad \eta=\max\Psi_3\in\{1,2\},\quad
\lambda\le\eta,\quad \eta+2\le\rho\le4.
\]
First suppose $\lambda=2$, so the original highest multiplicities are $(u,v)=(1,2)$ or $(2,1)$. The canonical certificate has three slots per $S$ root and six per $T$ root. Its other side consists of counted-component tokens and tokens for actual $S$--$T$ edges. A matching covers all slots and all edge tokens, leaving exactly two component tokens free. The canonical factor construction gives, in every counted component $Q$, the internal degree function
$f_Q(x)=3-d_G(x,T)$ and factors with degrees $f_Q\pm\mathbf1_x$ for any $x\in Q$; uncounted components admit strict colors compatible with the default ports. This is the canonical-certificate construction in the minimum-deficiency formula: maximality of the root set excludes $f_Q=0,d_Q$, and an obstruction to a one-vertex deviation factor could be merged into the roots without decreasing deficiency, enlarging the root set.

An alternating step may release an edge token. Reconstruct all original edges each time: edges within $S$ are blue and within $T$ red; edge tokens assigned to $S$ slots are red and those assigned to $T$ slots blue; free edges take their labels. Component ports to $S$ are blue by default and those to $T$ red. For a component assigned to an $S$ root, turn its actual assigned port red and use $f_Q-\mathbf1_x$ internally; for a $T$ root, turn the port blue and use $f_Q+\mathbf1_x$. A free red component uses the plus factor, a free blue component the minus factor. Uncounted components retain their compatible colors. All caps are strict except at the free resources.

If the two free tokens share an $S$ root, then for $(u,v)=(2,1)$ merge their red deviations at that root, using the actual compensating ports. Only one red high center remains, excluding two disjoint red four-stars. For $(1,2)$ label both resources blue and choose their high centers so that their unique seven-stars must use the common $S$ leaf. A common $T$ root is symmetric. Mixed edge/component tokens follow the same actual-edge rules.

Starting from each free token, follow arcs $w\to m(l)$ whenever $w$ meets slot $l$. If the reachable sets intersect, flip the two disjoint prefixes before their first intersection; this leaves two free tokens adjacent to one slot. If they are disjoint, every root's slots belong wholly to one region, and
$q_i+e(S_i,T_i)=3|S_i|+6|T_i|+1$.
Each counted component has at least one $S_i$ port, so $|S_i|-|T_i|\ge1$. The union of the two actual root regions is one certificate of the original graph, with
$\Psi_1\ge\sum_i[1+2(|S_i|-|T_i|)]\ge6>\rho$, a contradiction.
The alternating procedure therefore succeeds. Both complete original targets are excluded by their missing highest-star packing.
\subsection{Actual regions and strong lengths when the highest width is one}
Let $\lambda=1$, let $U$ be the union of counted components, and put $O=G-(S\cup T\cup U)$. Handshaking on the two sides of the same certificate gives
\[
q-e(T,U)=\eta+3+\zeta_T+2e(T)+e(T,O),\qquad
\zeta_S+2e(S)+e(S,O)+\sum_Q(e(Q,S)-1)=6-\eta.
\]
At least $\eta+3$ components have no $T$ port, and at most two have more than one $S$ port. Thus at least two actual components have total cut one.
They must also be shown near-perfect. Fix one maximum matching $M$ minimizing uncovered degree-nine vertices, write $m(v)$ for its coverage indicator, and put
$L_Q=e_M(Q,S)+e_{G-M}(Q,T)+u_M(Q)$.
Every $L_Q$ is positive and odd. Expanding the matching-cut count gives
\begin{align*}
\rho-\Psi_1={}&U_S+U_O+\sum_{S}(1-m(s))+\sum_T
\{8+\mathbf1_{d(t)=9,m(t)=0}-(d(t)-m(t))\}\\
&+2e_M(S)+2e_{G-M}(T)+e_M(S,O)+e_{G-M}(T,O)
+\sum_Q(L_Q-1).
\end{align*}
All terms on the right are nonnegative, while the left is at most one. Hence every $L_Q=1$. The selected components without $T$ ports have internal near-perfect matchings:
$|Q_i|=2P_i+1$, $P_i\ge5$. Put $P=P_1+P_2$, let $A$ be their external port endpoints, $a=|A|\in\{1,2\}$, and put
$O_0=G-(A\cup Q_1\cup Q_2)$, $k=\nu(O_0)$.
The actual root debt satisfies $\zeta_A+e(A)\le5$.
Define, without assuming matching equality,
$\Delta_0=a+P+k-\nu(G)\in[0,a]$ and $d_0=d-\Delta_0\ge2$.

Let \(Q_1,Q_2\) be the two actual odd components selected above. They are isolated in \(G-A\), have orders \(2P_i+1\), and satisfy
\(\nu(Q_i)=P_i\ge5\). Factor-criticality is not assumed.
Put
\[
\mathcal Q=Q_1\sqcup Q_2,\quad P=P_1+P_2,\quad
O_0=G-(A\cup V(\mathcal Q)),\quad k=\nu(O_0).
\]
Here \(a=|A|\in\{1,2\}\),
\(\Delta_0=a+P+k-\nu(G)\in[0,a]\),
\(d_*=d_0=d-\Delta_0\ge2\), and
\(\xi=\zeta_A+e(A)\le5\).
The original cap is nine, the largest heads are at most eight, and the width at height ten is \(h=1\).
In the numerical Appendix \ref{app:single-below} comparison use
\[
D=9,\qquad \mu=P,\qquad q_+=2,\qquad \eta=h=1.
\]
The parameter \(\eta\) is a comparison parameter; it is neither the actual exposure nor \(q_+-a\).
In this paragraph write \(R=r-1\), \(b=m_1\), and \(v_2=c_2(L)\); \(f_h,g_h\) denote the cumulative head counts of the two original lists.

Vertex deletion supplies the complete source \(P_{r+a}\) on
\(G-A=\mathcal Q\sqcup O_0\). Its rank is \(P+k+d_*\), its matching number is \(P+k\), and its matching deficit is \(d_*>0\).
Consider a complete original forest \(F\) of length \(q\le P\), and put \(b_F=P-q+d_*\).
The positive deficit makes the corresponding extreme pair in row \(r+a\) legal. Coloring \(O_0\) entirely with the matching color shows that
\[
\mathcal Q\longrightarrow(F,(b_F+1)K_2).
\]
Indeed, the other complete target in that extreme pair has exactly \(k+b_F+1\) star components, and hence requires a matching of that size. An avoiding coloring of \(\mathcal Q\) with matching number at most \(b_F\) would therefore avoid the entire original pair. The monochromatic test also gives \(b_F+1\le P\).

For every fixed deletion set of at most \(b_F\) points in the two boxes, color the incident edges with the matching color and all other internal edges with the forest color. This proves the existence of a complete \(F\) after that particular deletion. It does not assert a common copy for different deletion sets.
Choose \(2\lfloor b_F/2\rfloor\) deleted points, an even number in each box. Since each remaining box has odd order, its forest half-order budget gives
\[
q+\omega(F)\le P-\lfloor b_F/2\rfloor,
\qquad P\ge q+2\omega(F)+d_*-1.
\]
Deleting any \(b_F\) points and counting vertices gives
\[
q+e(F)\le2P+2-b_F,\qquad e(F)+d_*\le P+2.
\]
Only the actual orders, isolation, and known matching capacities have been used.

For a single-empty direction, interchange colors so that \(t\le P<s\).
Take the same two-budget residual \(L=L(\mathbf n)\) as in Appendix \ref{app:sources} and put
\(e_0=s-|L|\), \(\delta=P-e_0\).
At least one entry fits because every box has at least eleven vertices and every head is at most eight; thus \(1\le e_0\le P\).
The two fixed complete sources, constructed by the actual Appendix \ref{app:sources} root coloring and trimming of a single response, are
\[
P_{r+a}(O_0;L,\mathbf m),\qquad
P_r(Z;C_a(L)\vee\mathbf n[P+1:],\mathbf m),
\quad Z=G-E(\mathcal Q).
\]
Their ranks are \(k+d_*+\delta\) and \(a+k+d_*+\delta\), both positive and strictly smaller than \(N\).

If their height-\(j\) columns are \(\lambda_j^B,\lambda_j^Z\), set
\[
B_j=P+a/2+(\lambda_j^B+\lambda_j^Z)/2-L_j.
\]
On each of these two fixed graphs, maximize each column over its old source and the additional actual palettes described below. Let \(G_j\ge0\) be the average increase, exactly as in the two-ledger definition of Appendix \ref{app:sources}.
The exact fees are
\[
c_B=9(P+a)+8-\xi,\qquad c_Z=9P+8.
\]
Using the actual matching lower bound \(\nu(Z)\ge a+k-\Delta_0\), the potential comparison has constant
\[
C_{\mathrm{pot}}=7-\frac{\xi+\Delta_0}{2}\ge\frac72,
\]
and gives
\[
e(G)-E_r\ge C_{\mathrm{pot}}-d_*+
\sum_{j=2}^{9}(B_j+G_j).
\]
The separate full-area comparison is
\[
\frac{c_B+E_B+c_Z+E_Z}{2}-E_r
=\left(7-\frac{\xi}{2}\right)+\delta
+\sum_{j=2}^{9}B_j+
\frac{\lambda_{10}^B+\lambda_{10}^Z}{2}.
\]
Its constant is at least \(9/2\). The terms \(-d_*\) and \(+\delta\) belong to different ledgers and are not interchanged.

Write
\[
X=\sum_i\lfloor(n_i-1)/2\rfloor,\quad
Y=\sum_j\lfloor(m_j-1)/2\rfloor,\quad
S=s-P,\quad T=P-t,\quad u=s-f_2,
\]
\[
V=P-f_2,\quad A_*=V+v_2,\quad
V_0=\sum_{i=1}^2\left\lfloor\frac{P_i+1}{3}\right\rfloor\ge4.
\]
The complete embedding bounds proved above give
\[
T\ge2Y+d_*-1,\qquad t+g_2+Y\le P+2-d_*.
\]
The complete forest/matching endpoint gives
\[
2X\le t-a-2d_*-r+2.
\]
For its root term, if \(\Delta_0=0\), an optimal barrier of \(G-A\) together with \(A\) is an optimal barrier of \(G\). If \(\Delta_0\ge1\), use \(2d\ge2d_*+a\). The endpoint in \(O_0\), whose actual deficit is \(d_*+\delta\), also gives
\[
S\ge a+r+2d_*+\delta+2Y-2.
\]
The two capacities \(P_i\ge5\) yield every remaining structural guard:
\[
P\ge3\eta,\quad V_0\ge\eta,\quad
q_+=2\le(P-\eta)/2,\quad q_+\le a/2+V_0,
\]
\[
8V_0+a-2P\ge0,\qquad 2V_0\le P.
\]
Since \(2P+2\le2P+a+\eta\), the point budget for the lowest \(P-\delta\) entries gives
\(3\delta\ge P-u-a-\eta\).
The half-order budget gives \(f_3\le S+2\delta\); the height-two count of the fixed residual gives \(A_*\ge V_0\).
When \(L\) is all-unit, additionally
\[
f_2+X\le P,\qquad V\ge\max(V_0,X).
\]

The matching-row repair in Appendix \ref{app:single-below} has an actual realization in these two boxes.
Set \(q_0=T-2Y+1\ge d_*\) and \(w=T-Y+1\).
Then \(0<w\le P-1\). Select \(2w\) vertices, an even number in each box. Color internal edges meeting these points red and all other internal edges blue.
The red matching number is at most \(2w\). The blue half-order and vertex budgets are \(P-w\) and \(2P+2-2w\), respectively.
The original blue forest has half-order cost
\(t+Y=P-w+1\), so its window residual \(Q\) is nonempty.
No perfect matching after an arbitrary deletion is used.

If this repair is triggered at height \(j\ge3\), the trigger
\(f_j>2w+a/2+r-2\), together with \(f_j\le X\) and the displayed endpoint bounds, gives \(s-2w>r+a\).
At height two, use \(f_2\le P\) with the same trigger and the lower bound on \(S\) to get the same conclusion.
Thus both repaired sources have positive rank; \(2w>0\) and the residual length bound make the ranks strictly smaller.
The repaired average column has \(B_j\ge-q_0+|Q|-1\ge-q_0\).
Combining this with the old column bounds leaves loss at most
\[
\min(q_0,(f_j-q_0)_+)\le f_j/2.
\]
The sum over heights at least three is at most \(X\).

Finally \(h=\eta=1\) means that every rectangle with \(x+y\ge11\) satisfies
\(f_x+g_y-r+1\le1\); in particular \(g_b\le r-1\).
All the Appendix \ref{app:single-below} guards, both actual complete sources, both ledgers, and the only additional internal palette have now been established with
\((D,\rho,d,\mu,q_+)=(9,\eta,d_*,P,2)\).
The Appendix \ref{app:single-below} column comparisons apply verbatim to these numerical inputs. Their positive average cost supplies an actual paid positive-rank successor and excludes the single-empty direction. Interchange colors for the other direction.

If \(\max(s,t)\le P+d_*\), the two original guarded bounds are
\[
E_r\le9P+a+k-7d_*+11,\qquad
e(G)\ge9P+9a+3+k.
\]
Their difference is at least \(8a+7d_*-8>0\), so this case is impossible.
If that double guard fails and one list still has length at most \(P\), the other is longer than \(P\), which is the single-empty case just excluded.
Both lists therefore have length greater than \(P\).
The two monochromatic restorations in the same source \(G-A\) give complete endpoints in \(O_0\), and their monochromatic tests imply \(s,t\le k\).
Using \(s+t=a+P+k+R+d_*\), obtain
\begin{equation}\label{eq:D9length}
s,t\ge a+P+R+d_*\ge P+R+d,
\end{equation}
where the last step uses \(\Delta_0\le a\). These are the strong lengths required by the subsequent port and cap operations.
\subsection{One edge-erased graph with adapted ports}
If the unique internal port vertex $x_i$ satisfies $\nu(Q_i-x_i)=P_i-1$, include both the internal edges and the port in $X_i$ (type E). Otherwise include only the internal edges and retain the port (type L). Let $e$ be the number of E types, $k_0=2-e$, and put
$X=X_1\cup X_2$, $Y_0=G-E(X)$.
If a matching uses an E port, its internal capacity drops by one. For an L component fix a near-perfect matching exposing the internal port; it combines with every matching of $Y_0$. Hence
\[
\nu X=P,\quad \nu Y_0=N-P-d,\quad e(X)=9P+w,\quad w=8+e.
\]
\[
M_z=e+\floor{\frac{2P+4-2e}{z+1}}\quad(z\ge2).
\]
The number $M_z$ is an actual damage bound: at most $e$ affected stars use E ports, each occupying a core vertex; at most $k_0$ use L ports while meeting an erased edge, each occupying $z$ core vertices; every other affected star occupies $z+1$ vertices. Every low-column tax is at most $P$, except height three may cost one more. Indeed, a surviving monochromatic count gives $M_2\le P$. If both counts are dead, choose a vertex cover set containing all E internal port vertices, color its incident erased edges red and the others blue. A valid tax is
$\max(2P+4-e-vB,e+B-1)$ or its color reverse.
If both exceed $P$, then $uA,vB\le P+3-e$.
At height three use
$\min(2A-R,2P+4-e-2A)\le P+1$.
At higher low columns at least one head is three, so
$A+B\le5(P+3-e)/6\le P+1$,
contrary to $A+B\ge P+2$ required by width at least $P+1$.
At height nine retain the earlier head to obtain tax at most $M_z$, where $z=3$ for $H=7$ and $z=2$ for $H=8$. The reserve on the same graph gives
\[
d\ge P-M_z+7+e,\qquad
d\ge2(P-M_2)+7+e
\ \text{if a low nonunit column is live}.
\]
\subsection{Two stages on one history and strict payment}
Fix a highest-degree history of $Y_0$. Put $b=C_9$, $c=C_8$, and
$\tau=(\max(P_1,P_2)-R)_+$.
The menu $\{0,P_1,P_2,P\}$ has gap at most $\max(P_1,P_2)$. Whenever the first $\tau$ vertices exist, the single graph
$Y=Y_0-T$ has $\p_{r+P+\tau}$ with the fixed original lists. Equation \eqref{eq:D9length} ensures positive smaller rank.

Take maximal independent sets inside the two cores, requiring each E internal port vertex to belong to its set; let their total size be $i$. Degree-counting Hall shows that each such independent set can be saturated by an internal matching. Neighborhoods of full-degree vertices also give
\[
2\le i\le\min(P,2P-16).
\]
Delete these vertices, then the first $b$ vertices of the same exterior history. The actual cap-eight graph $X^{(1)}$ has $K_1=i+b<N$ deleted vertices and cost $9K_1-\delta_1$, where $0\le\delta_1\le k_0$. The cap-eight small-surplus price of Appendix \ref{app:stage-prices} gives
$\rho_8(X^{(1)})\le R_8(e)$,
where $R_8(0)=6$ and $R_8(1)=R_8(2)=5$.
Transfer exposure to $Y_1=Y_0-\{\text{first }b\text{ vertices}\}$ as follows: first remove the L-type pendant leaves already deleted by the internal operation, restrict an optimal matching, and then restore these leaves one at a time. Adding a pendant leaf cannot increase final full-degree exposure: if its newly full root is uncovered, cover it by the new leaf and augment further. Hence
\[
L_9-b\le i+R_8(e),\qquad L_9-b\le M_z+R_8(e).
\]
The complete ordinary endpoints give $C_H\ge2d$, and every unit axis $A_h$ gives $C_h\ge A_h-P$.
For $H=7$, $d>P/2$, so inventory suffices. For $H=8$, if $C_8<\tau$, a nonunit winner would again force $d>P/2$ through the strengthened reserve, a contradiction. Thus all low columns are unit axes, and coarea together with the preceding bounds and $w\ge R_8+1$ pays the full area. Inventory therefore suffices unless the source has already exited. If the final deletion degree is at least eight, the saving is
$9P+w+8\tau+\min(b,\tau)$.
For $b\ge\tau$ use the full degree-nine saving; for $b<\tau$ use $L_9-b\le P+R_8$. Both pay.

Only final degree seven remains; here $H=7$ and $b\le c<\tau$.
Inside the cores of $X^{(1)}$, take a maximal independent set of current degree-eight vertices and continue the same exterior history to $c$. The cumulative cost is
$8K_2+K_1-\delta_2$, with $0\le\delta_2\le k_0$: port debt is charged only once across both stages.
For positive rank, the actual surplus is at most $R_8+k_0$; the cap-seven small-surplus price gives $\rho_7\le R_7(e)=8-e$.
For nonpositive rank,
$e(X^{(2)})\le8(N-K_2)+R_8+k_0\le8(N-K_2)+8$.
Thus $K_2>N+1$ is impossible; $K_2=N+1$ leaves no edge; and $K_2=N$ leaves at most eight edges, at most one full degree-seven vertex, and zero exposure.
The same pendant-leaf argument transfers this to actual $Y_2$, giving
\[
L_8-c\le M_2+R_7(e)\quad\text{in the live case},\qquad
L_8\le M_u+M_v-R\quad\text{in the dead case }u+v=9.
\]
\[
e(Y)-E_{r+P+\tau}\le L_8+L_9-2P-w-b-c.
\]
In the live case with $e=0$, the upper bound is both
$M_2-10$ and $M_2+M_3-2P+6$.
Use the former if $M_2\le9$; otherwise $P\ge13$ and the latter is at most $-5P/6+25/3<0$.
For $e=1,2$ the bounds are respectively $-5P/6+37/6$ and $-5P/6+5$, strictly negative since $P\ge10$.
In the dead case, $1/(u+1)+1/(v+1)\le11/24$ gives the upper bound
$-7P/12+R_8(e)-31/6+7e/12-R<0$.
Every branch retains the original strict one-edge difference on the same graph. Restore the fixed menu and actual vertex labels to obtain an avoiding coloring at an original legal pair. Degree nine is completely excluded.
\section{Common preliminaries for degrees eight and ten}\label{app:even-small}
Here $D=2p$, where $p=4,5$. Write $\rho=\rho_D(G)$, $\lambda=L_{D+1}$, and let $c$ count the actual isolated odd-order connected $D$-regular components.
\subsection{Eliminating isolated complete graphs}
If $C=K_{2p+1}$ is an actual component, the strong length bound gives $s,t\ge p+r+1$. Put $Y=G-E(C)$ and $c_h=\lfloor(2p+1)/(h+1)\rfloor$. Removing the stars contained in $C$ from one complete copy gives the two complete sources. Set
\[
 R_h=\max_{v\ge h}(f_v-c_v)_+,\quad S_h=\max_{v\ge h}(g_v-c_v)_+,
\quad K_h=\min(f_h,R_h+u),\quad L_h=\min(g_h,S_h+u),
\]
where $1\le u\le r-1$. Since $u\ge1$, the displayed minima retain every positive original
threshold: $f_h>0$ implies $K_h>0$, and $g_h>0$ implies $L_h>0$.
The candidate comparisons below therefore only use surviving heads. At a new corner $I+J=r+u$, either $I>u$ or $J>u$, so use the old corner $(I-u,J)$ or $(I,J-u)$, respectively. Thus the nested counts define fixed lists satisfying $\p_{r+u}$ on the same $Y$. Their rank is at least $r-u+1\ge2$ and at most $N-u$.
Let $\delta_h=(c_h-u)_+$. The complete cost is at most
\[
 u(D+1)+\sum_{j=1}^{D+1}\max_{\alpha+\beta-1=j}(\delta_\alpha+\delta_\beta).
\]
The following list gives every boundary case; entries not displayed in a column-loss vector are zero. When $r=2$, odd--odd heads are excluded by a strict even-factor coloring, so the listed even--even pairs are exhaustive.
\[
\begin{array}{c|c|c|c}
 D&r&u&\text{upper bound on cost}\\ \hline
8&r\ge3&2&18+(2H+4)\le36\\
8&2,\ (4,6)&1&9+(6+5+4+3+3+3+2+1)=36\\
10&r\ge4&3&33+(2H+2)\le55\\
10&3&2&22+(3H+5)\le55\\
10&2,\ (4,8)&1&11+(8+6+5+5+4+4+4+4+2+1+1)=55\\
10&2,\ (6,6)&1&11+(8+6+5+5+4+4+2+1+1)=47
\end{array}
\]
Here $H\le D-1$. Each cost is at most $e(C)$, giving $e(Y)\le E'-1$, contrary to the original minimality.
\subsection{Interpolation in one factor certificate and maximal exposure}
For one certificate $S,T$, define
\[
 \Psi_x=q+e(S,T)-x|S|-(D-x)|T|.
\]
A counted component consists of vertices of original full degree and satisfies $|Q|+e(Q,T)$ odd. The componentwise token assignment from a maximum matching gives $\Psi_1,\Psi_{D-1}\le\rho$. Every nonisolated counted component has cut at least two, so the half-cut count gives $\Psi_p\le c$. Consequently
\begin{equation}\label{eq:even-interpolation}
 \eta_a\le\left\lfloor c+\frac{|D-2a|}{D-2}(\rho-c)\right\rfloor.
\end{equation}
The actual whole-graph defect construction requires $\lambda\le\eta_a$ at a top corner: otherwise its independently placed defects avoid both targets.
\paragraph{The Brooks argument for \(\rho\le p+1\).}
Assume \(D=2p\), \(p\in\{4,5\}\), and use the original normalized strict
source.  Thus
\[
 e(G)=E_r-1,\qquad \nu(G)<N,\qquad
 1\le\lambda=L_{D+1}\le\rho=\rho_D(G).
\]
Suppose that the preceding complete-list argument has excluded every
actual component \(K_{2p+1}\).  In the ordinary decomposition, the
positive whole-core weights satisfy
\[
 1\le w_Q\le p,\qquad W_+\ge p\rho.
\]
If \(\rho\ge p+2\), order these cores by nondecreasing capacity and take
the first prefix whose weight reaches \(\lambda\).  Write \(w,z,Z\)
for its weight, total capacity, and largest capacity.  Then
\[
 w\le\lambda+p-1\le\rho+p-1,\qquad
 W_+-w\ge(p-1)(\rho-1)\ge p^2-1.
\]
As each unselected weight is at most \(p\), at least \(p\) positive
whole cores remain, each with capacity at least \(Z\).  Their union
contains at least
\[
 p(2Z+1)>2pZ
\]
original full-degree vertices.

Each remaining core is a connected graph of maximum internal degree
at most \(2p\).  A core isomorphic to \(K_{2p+1}\) would already give
every one of its vertices all \(2p\) neighbors internally.  It would
therefore be an actual isolated component of \(G\), excluded above.
Brooks' theorem consequently gives a \(2p\)-coloring of the union
of the remaining cores.  Components of smaller maximum degree,
including odd cycles, also admit \(2p\) colors.  Thus this actual union
has an independent set of size at least \(Z+1\).

Choose \(b=(Z-r+1)_+\le Z-1\) vertices from it.  They avoid every selected
core, retain original degree \(D\), and have no neighbors inside any
selected core.  Hall counting in the complement of the selected core
vertices supplies a matching saturating these \(b\) vertices.  Together
with the selected near-perfect core matchings this gives
\[
 1\le K=z+b\le\nu(G)<N.
\]
Erase the selected internal edges and delete these independent vertices.
The actual fee is \(DK+w\), and \(r+b>Z\) gives the complete same-list
source \(P_{r+K}\) by Lemma~\ref{lem:menu}.  Its positive rank is smaller.
Since \(w\ge\lambda\) and the full profile loses at most \(DK+\lambda\),
\[
 e(Y)-E_{r+K}\le-1+\lambda-w\le-1,
\]
a contradiction.  Therefore \(\rho\le p+1\).
This is an original-source operation with a positive smaller successor;
the nonnegative-surplus baseline at the original rank is not used.

Return to the equality case $\rho=p+1$. Equality forces every nonnegative loss to vanish:
\[
 \lambda=\rho=p+1,\quad W=p(p+1),\quad c=p+1.
\]
The root degree defect, internal root edges, root-to-exterior edges, and negative block weights all vanish. Thus $A\cup R$ is an actual isolated regular region containing all $p+1$ odd regular components, each with a near-perfect matching. All other full-degree vertices are covered by the same matching. Write their capacities as $a_1\le\cdots\le a_{p+1}$; the preceding exclusion of cliques gives $a_i\ge p+1$.
Take the smallest component $C_1$, of capacity $a=a_1$. In each of the other $p$ noncomplete regular components, Brooks theorem supplies $\lceil(2a_i+1)/(2p)\rceil$ independent vertices, at least $a+1$ in total. Extend their union to a maximal independent set $T$ among all exterior full-degree vertices; put $k=|T|\ge a+1$. Hall counting into its neighborhood gives a matching saturating $T$, so $K=a+k\le\nu G<N$. The two monochromatic menus on $C_1$, followed by deletion of $T$, give complete $\p_{r+K}$ on $Y=G-(V(C_1)\cup T)$. The actual fee is $DK+p$, whereas
\[
 E_r-E_{r+K}=DK+\min(K,p+1)-\sum_{j\le D}(K-L_j)_+.
\]
Failure to pay therefore forces $e(Y)=E_{r+K}$, $L_D\ge K$, and $\Delta Y\le D-1$. The odd-cap equality price gives $\rho_{D-1}(Y)\le p$.
For an odd cap $L$, a complete uniform source with even head $e$ and odd head $o$ has width at most $(o-1)\rho_L/(L-1)$. Indeed, interpolate the endpoint token bound in the same certificate with $q+e(S,T)\le L|S|$; each counted component contributes an actual edge to $S$ in the latter inequality. Apply this to column $D$ on the same $Y$. At least one original count exceeds $M_h=\lfloor(2a+1)/(h+1)\rfloor$, since otherwise the two counts cannot sum to $L_D+r-1\ge2a+r$. The surviving side gives
\[
 a-M_{\min(e,o)}\le\left\lfloor\frac{p(o-1)}{2p-2}\right\rfloor.
\]
Thus $p+1\le a\le3p+1$. If $e\ge4$, also $e<o$ and $e^2+3e\le4p+4$, impossible for $p=4,5$. Hence column $D$ is attained only by $(2,D-1)$, and in particular $H=D-1$.
In a noncomplete component of capacity $z\ge p+1$, take $\lfloor z/p\rfloor+1$ vertices from a Brooks independent set. Some exterior vertex has at most one neighbor in this set, since
\[
 (2z+1)-(p+1)(\lfloor z/p\rfloor+1)>0.
\]
Adding it gives an actual stock of $\lfloor z/p\rfloor+2$ vertices inducing at most one edge. Put $b=a_2$. Delete the first two components and use the remaining stocks to fill the subset-menu gap $\max(a,b-a)$. If enough vertices are available, the internal-edge loss is at most $p-1$, and the fee is at least $DK+2p-(p-1)=DK+\lambda$. Otherwise
\[
 b-a\ge \sum_{i=3}^{p+1}(\lfloor a_i/p\rfloor+2)+2.
\]
The gap cannot equal $a$, since for $a\le3p+1$ the stock is at least $a-1$. Writing $A=a_{p+1}$ and $P=\sum a_i$, this inequality gives $A<2b-pa-p^2<2b$. Thus the largest gap in the full actual subset-sum menu is exactly $g=b-a<P/p$.
We must still supply enough exterior high-head vertices. Take the largest $p$ components, of total capacity $P_0$, and put $M_2=\sum\lfloor(2a_i+1)/3\rfloor$. The preceding inequality gives $b>pa+p(p+1)\ge2p(p+1)$, hence $P_0>2p^2(p+1)$. For a nonunit attaining pair, exchange colors so that $f\le g$, and set
\[
 x=\min(f-1,\lfloor P_0/2\rfloor),\quad y=P_0-x,\quad
 k=(2x-P_0+p-2)_+.
\]
If the column width exceeds $P_0$, then $x<f,y<g$. Moreover $k\le p-2$ and $x\ge Dk$. Choose $k$ vertices in different components, extend their neighborhood to $x$ vertices, color incident edges red and the rest blue. There are at most $x$ red stars, while the blue nonisolated support has fewer than $3(y+1)$ vertices. Every low-column tax on this same remainder is therefore at most $P_0$; unit columns use the complete matching capacity. The column-$D$ tax is at most $M_2$: color the selected components with the head-$2$ color. If that side does not survive, the prohibition of the original odd--odd pair $(3,D-1)$ already bounds the old width by $M_2$. The top-column tax is at most $p$. Each actual near-perfect regular component has an odd-cap factor of degrees $\alpha-1\pm1_v$ at any specified vertex, so distribute the $p$ defects according to the top multiplicities. This factor existence follows directly from the strict nonempty-root bound in \eqref{eq:even-interpolation} and the change of parity at the empty root; factor-criticality is unnecessary.
Apply the canonical price on this one actual remainder. The removed edge count is $DP_0+p^2$, its matching number decreases by $P_0$, and full-degree exposure decreases from $p+1$ to $1$. This yields
\[
 d\ge P_0-M_2+p^2-p+2,\qquad P_0\le3d-3p^2+4p-6.
\]
Hence $g=b-a\le P_0/p-a<d$. For the union $U$ of all $p+1$ components, strong length and complete source give two endpoints on $O=G-U$, each with matching target at least $(d+1)K_2$. Therefore deletion of any at most $d$ vertices from $O$ leaves both complete original forests. Take the first $\tau=(g-r+1)_+$ vertices in its highest-degree history. Every step has degree at least $H=D-1$, and $K=P+\tau<N$. The subset menu gives complete $\p_{r+K}$. If payment still fails, writing $b_0$ for the full $D$-stage length and $L=L_D$ gives
\[
 b_0<\tau,\quad L\ge P+b_0+p^2,\quad
 q_Y=e(Y)-E_{r+K}\le L-P-b_0-p^2.
\]
The head-$2$ side in column $D$ must survive. Thus the same $Y$ has a complete uniform source of width $L-M-\tau$, with $M\le(2P+p+1)/3$. The odd-cap deficiency bound and general price give, respectively,
\[
 \rho_{D-1}(Y)\ge L-M-\tau,\qquad \rho_{D-1}(Y)\le q_Y+p.
\]
But $P-M-\tau>0$: indeed $\tau<P/p$, $P\ge(p+1)^2$, and
$P-M-\tau\ge P(1/3-1/p)-(p+1)/3>0$.
The first bound is thus at least $q_Y+p^2+1$, a contradiction. This excludes $\rho=p+1$.
\subsection{One further layer and completion of degree eight}
Let $b_p$ count all positive blocks of weight exactly $p$; each is an actual isolated regular factor-critical component. Order positive blocks by nondecreasing capacity and take the first prefix whose total weight reaches $\lambda$. Write $w$ for its weight, $t$ for its last block's weight, $Z$ for its largest capacity, and $u$ for the number of unselected positive blocks. Each weight belongs to $\{1,\ldots,p\}$, and $w\le\lambda+t-1$.
If $u\le p-1$ and $t=p$, at least one weight-$p$ block is selected. The unselected weight is at most $(p-1)u+b_p-1$, so
\[
W_+\le\lambda+p-1+(p-1)^2+b_p-1
=\lambda+p(p-1)+b_p-1.
\]
If $t\le p-1$, the unselected weight is at most $(p-1)u+b_p$, giving
\[
W_+\le\lambda+p-2+(p-1)^2+b_p
=\lambda+p(p-1)+b_p-1.
\]
A strict violation therefore forces at least $p$ unselected positive blocks, not merely unselected total weight at least $p$. Each has capacity at least $Z$, so the block inventory supplies at least $2pZ$ original full-degree vertices outside the selected blocks. If $Z=p$, greedy coloring supplies the required $Z-1$ independent vertices; if $Z>p$, Brooks' theorem does so, with the sole clique exception already excluded. Take $b=(Z-r+1)_+\le Z-1$ additional roots. The original capacity menu and actual fee give a paid successor of positive smaller rank. Hence
\[
p\rho\le W_+\le\lambda+p(p-1)+b_p-1.\tag{*}
\]
If $\rho=p$, then (*) gives $\lambda+b_p\ge p+1$. When $b_p\ge2$, use the two-component conclusion of Proposition \ref{prop:regular-components}. When $b_p=1$, equality forces $\lambda=p,W=p^2$ and zero root loss. Interpolation then forces $c=p$ near-perfect components; strong length follows from the actual isolated region $A\cup R$, and the same proposition applies. The case $b_p=0$ is impossible. Hence $\rho\le p-1$.
At the central odd cap the actual width is at most $c$. If $p$ is odd, take $a=p$ directly. If $p$ is even, take $a=p-1$. A width larger than $c$ would force $\rho\ge c+p-1$ in \eqref{eq:even-interpolation}, contrary to $\rho\le p-1$, except possibly for equality with $c=0$. Equality makes every loss in the half-cut count zero, so the certificate region is actually isolated and regular. Since $c=0$, its order is even; its positive deficiency has the same parity as that order and cannot equal $1$.
For $D=8$, the sole top pair is $(4,6)$, so $\lambda\le c\le\rho\le3$. If $\lambda=1$, one actual odd regular component without a perfect matching is excluded by Proposition \ref{prop:regular-components}. If $\lambda\ge2$, there are at least two odd regular components. Total exposure is at most three, so each has deficiency one. The Tutte barrier, single-empty exclusion and two guards give strong length on their union, as proved in that proposition. Its two-component conclusion applies. This completes degree eight.
\section{Degree ten: the deficiency-two region and complete payment}\label{app:D10}
\subsection{A nearly full-degree decomposition: hypotheses, actual successors and row costs}
This subsection is used only in the following original-source range: $D=2p\ge8$, $H\le D-1$, $h=L_{D+1}$, and

\begin{equation}\label{eq:D10-original-domain}
 1\le h\le\rho_D(G)\le p-1,\qquad
 d=N-\nu(G)\ge\rho_D(G)+1.
\end{equation}

Suppose this same $G$ has an actual partition $A\sqcup R_0\sqcup B_0$. The region $R_0$ is a union of $q$ induced odd factor-critical components and has no edge to $B_0$. One maximum matching assigns all $a=|A|$ roots through actual edges to distinct components and attains the remaining capacities inside the components and $B_0$. Write

\[
 \mu=\nu(R_0),\quad k=\nu(B_0),\quad
 \eta=q-a\ge1,\quad \nu(G)=a+\mu+k,
 \quad \varepsilon_R=\sum_{v\in R_0}(D-d_G(v))\le2.
\]

The matching identity and absence of crossing edges are hypotheses: total deficiency two in an induced region alone is insufficient. Put $\beta=e(A,B_0)$, $e_A=e(G[A])$ and $\sigma=Da-\sum_{v\in A}d_G(v)$. We first exclude $t\le\mu<s$; the reversed case is symmetric.

\paragraph{A fixed common residual and two physical graphs.}
Use the common window-and-capacity residual $L$ of Appendix \ref{app:sources}: it is the coordinatewise maximum of the window residual for the actual budgets
$\sum\lceil n_i/2\rceil\le\mu$, $\sum(n_i+1)\le|R_0^+|$, and the head-capacity residual. The latter has counts
$\max_{l\ge j}(f_l-\sum_C\lfloor|C|/(l+1)\rfloor)_+$. Put $e_0=s-|L|$, $\delta=\mu-e_0$, and let $C_a(L)$ replace the first $\min(a,|L|)$ entries by units and reorder. The whole-copy deletion argument of that appendix gives the complete sources

\[
 P_{r+a}(B_0;L,\mm),\qquad
 P_r(Z;C_a(L)\vee\nn[\mu+1:s],\mm),\quad Z=G-E(R_0).
\]

Full degree of the components is not needed for the copy-deletion budget: factor-criticality, the one matching and absence of $R_0$--$B_0$ edges suffice. Full degree was used there only to obtain a large component; below we still obtain at least $\eta\ge1$ such components. Hence $e_0\ge1$, both successors have $1\le N_B,N_Z<N$, and $r+a\le\min(|L|,t)$. Each new angle is obtained from one complete original copy, without fixing the copy across different angles.

Let $\lambda_j^B,\lambda_j^Z$ be the column widths at the respective new rows and $C_j=L_j$ the original widths. Define

\[
 B_j=\mu+\frac a2+\frac{\lambda_j^B+\lambda_j^Z}{2}-C_j,
 \quad C_0'=\frac{D\eta+\beta+e_A-\varepsilon_R}{2}-h.
\]

Let $G_j\ge0$ be the average increase obtained by taking the larger old or additional-palette column on each fixed graph, with the same palette as Appendix \ref{app:single-below}. Each actual deletion fee differs from its full-degree formula by exactly $\varepsilon_R/2$. Thus the two distinct ledgers of Appendix \ref{app:sources} become

\begin{align}
 e(G)-E_r&\ge C_0'-d+\sum_{j=2}^D(B_j+G_j),\label{eq:D10-ledger-potential}\\
 \frac{(c_B+E_B)+(c_Z+E_Z)}2-E_r
 &\ge C_0'+\delta+\sum_{j=2}^D B_j,\label{eq:D10-ledger-area}
\end{align}

Here $c_B=e(G)-e(B_0)$ and $c_Z=e(G)-e(Z)$. The second ledger does not subtract an extra $d$, and the two deletion fees are never added.

\paragraph{Effect of the degree deficit on structural bounds.}
For a component $C$, put $z_C=(|C|-1)/2$, $V_0=\sum_C\lfloor(z_C+1)/3\rfloor$, and let $q_+$ count nontrivial components. An odd component of order at most $D$ satisfies
$\partial C+\sum_C(D-d_G)\ge D$. Total boundary plus deficit is at most $Da+2<D(a+1)$, so there are at most $a$ small components and at least $\eta$ components of order at least $D+1$. Thus $\mu\ge3\eta$ and $V_0\ge\eta$. The componentwise simplicity calculation gives

\[
 2Dq_+\le6\mu-D\eta+\varepsilon_R-\sigma-2e_A-\beta,
 \qquad q_+\le\frac{\mu-\eta}{2}.
\]

The second inequality does not discard the deficit: initially its right side has the extra term $\varepsilon_R/(2D)<1/2$. The baseline is integral or half-integral and $q_+$ is integral, giving the asserted bound by rounding. A capacity-one component has boundary plus deficit at least $3(D-2)\ge2D$, so $2Dq_1\le Da+2$ and $2q_1\le a$. For other components, $8\lfloor(z+1)/3\rfloor\ge2z$. Consequently

\[
 q_+\le a/2+V_0,\qquad 8V_0+a-2\mu\ge0,\qquad2V_0\le\mu.
\]

A maximum matching can fully cover at most $a$ odd components, so at least $\eta$ vertices of $R_0$ remain exposed, at most $\varepsilon_R$ of them non-full. Thus $\eta\le\rho_D(G)+\varepsilon_R$, and in particular $d\ge\eta-1$.

\paragraph{Complete endpoints and list constraints.}
Use the notation of Appendix \ref{app:single-below}, explicitly

\begin{gather*}
 \alpha=n_1,\quad b=m_1,\quad X=\sum_i\floor{(n_i-1)/2},\quad
 Y=\sum_j\floor{(m_j-1)/2},\\
 S=s-\mu, T=\mu-t, u=s-f_2=S+V, V=\mu-f_2,\quad
 v_2=c_2(L),\ A_*=V+v_2,\\
 q_0=T-2Y+1\ge d,\quad w=T-Y+1,\quad c_l=g_l-g_{l+1}.
\end{gather*}

Deletion-resistant embeddings, complete endpoints and window budgets use the actual component matching, not full degree, and give

\begin{gather}
 T\ge2Y+d-1,\quad t+g_2+Y\le\mu+q_+-d,\quad
 2X\le t-a-2d-r+2,\label{eq:D10-list-guards}\\
 S\ge a+r+2d+\delta+2Y-2,\quad
 3\delta\ge\mu-u-a-\eta,\quad f_3\le S+2\delta,\quad A_*\ge V_0.\nonumber
\end{gather}

For a unit residual $L$, also $f_2+X\le\mu$ and $V\ge\max(V_0,X)$. Both $\alpha,b\ge3$, so $X,Y\ge1$. The correct original high-rectangle constraint is $f_x+g_y-r+1\le h$ for $x+y\ge D+2$, in particular $g_b\le h+r-2$. Here $h$ must not be replaced by $\eta$. The constant term satisfies $C_0'\ge p\eta-h-1\ge0$.

\paragraph{Top row HH or B2: payment of every low row.}
The unit, matching, two-head and mixed winners of Appendix \ref{app:single-below} exhaust each column. Its one palette bounds the matching group's total loss by $X$; legality uses \eqref{eq:D10-list-guards}, independently of the degree deficit. The two top-row substitutions are

\begin{align}
 B_D&\ge X+q_0+Y+3a/2+2d+2r-4 &&(\mathrm{HH}),\nonumber\\
 C_0'+B_D-X-d
 &\ge p\eta+(d-h)+q_0+Y+3a/2+2r-5,\label{eq:D10-HH}\\
 B_D&\ge T+c_1+a/2-h+g_b &&(\mathrm{B2}),\nonumber\\
 C_0'+B_D-X-d
 &\ge X+a/2+g_b+Y+2(d-h)+r-3+(p-1)\eta.\label{eq:D10-B2}
\end{align}

For $b\le D-2$, the boundary at height $b+1$ must still be paid; it is not automatically nonnegative. Its negative part is at most $g_b\le Y$ for unit $L$, at most $g_b/2\le Y/2$ for short nonunit $L$ ($f_2\le\mu,1\le v_2\le a$), and zero for long nonunit $L$. Internal two-head rows of height at least four are nonnegative.
Here long nonunit means that $L$ is nonunit and $f_2>\mu$ or $v_2>a$. This case is paid directly by \eqref{eq:D10-HH} or \eqref{eq:D10-B2}. For unit $L$, if both low rows are two-head, their sum is $2u+2\delta-t+g_3>0$. If height two is unit and height three is negative two-head, put $D_3=-B_3$. Then

\[
 D_3+X\le q_0-A',\quad
 A'=3d-\eta+Y+r+2\delta+g_3-1>0,
\]

using $\eta\le d+1$. If height two is negative and three is not two-head, the same palette also bounds the extra loss by $q_0$, so \eqref{eq:D10-HH} pays all HH cases. After subtracting the boundary loss $Y$, the two special unit B2 cases have the following remainders. For short nonunit cases subtract $Y/2$; the final display is twice the remainder:

\begin{align*}
 &p\eta+2(d-h)+3Y+3a/2+r+2\delta+g_b-4>0,\\
 &(p-1)\eta+2(d-h)+2Y+c_1+a/2+g_b+r+3\delta+g_3-3>0,\\
 &(p-\tfrac12)\eta+2(d-h)+(V+A_*)/2+a+r+\delta+2Y+g_b-\tfrac72>0,\\
 &(D-1)\eta+4(d-h)+a+2r+2\delta+3Y+V+g_3+2g_b-6>0.
\end{align*}

These follow by adding the top row and \eqref{eq:D10-list-guards} to, respectively, the inequalities $-B_2+X\le c_1-S-\delta$, $D_3+X\le q_0-A'$, the short-case bound $2(-B_2)\le f_2-A_*-\delta-S-T-g_2$, and $2(B_3+q_0)\ge2\delta+V+Y+4d+r+g_3-1$ from that appendix. The conditions $p\ge4,\eta\ge1,d-h\ge1,r\ge2,Y\ge1$ make each displayed remainder positive. In the short case two low two-head rows have positive sum by the same list budgets; a negative height-three row satisfies $B_3>-q_0$; all other cases are repaired matching rows. This exhausts the low-row combinations. Each single expression has already lost $\varepsilon_R/2\le1$ and each doubled expression $\varepsilon_R\le2$; no second subtraction is made.

\paragraph{The top boundary $(2,D-1)$.}
If $b=D-1$ but the highest winner is HH or B2, the preceding payment still applies, now with no additional height-$b+1$ boundary loss. Only the highest winner $(2,D-1)$ remains, and we use \eqref{eq:D10-ledger-area}. For a matching winner the original rectangle bound gives

\[
 B_j\ge S+\delta+T+g_j-f_j
 \ge a+2d+2\delta+2Y+T+g_j+g_b-h-1>0.
\]

Also $\mu\ge2X+2Y+a+3d+r-3$. Combining $8V_0+a\ge2\mu$ with $Y\ge2g_b$ gives $8V_0\ge8g_b+8$, hence $V_0\ge g_b+1$. In the long nonunit case all two-head rows are nonnegative and the top row is positive, so $C_0'\ge0$ suffices. For short nonunit $L$, let $K$ be the negative internal two-head rows. If $K$ is empty the top row is at least $(V+A_*+r-1-g_b)/2>0$. Otherwise summing the disjoint head classes gives

\[
 \delta+\sum_KB_j\ge(A_*+T-a-\eta+g_b)/2,
\]
\[
 \delta+B_D+\sum_KB_j
 \ge V+r-1+(A_*-\eta)/2+(T-g_b)/2>0.
\]

For unit $L$, an empty $K$ gives top row at least $a/2+r>0$. Otherwise use $3\delta\ge\mu-u-a-\eta$, retaining the full term $T$, to obtain

\begin{align*}
 C_0'+\delta+B_D+\sum_KB_j
 &\ge C_0'+(8V_0+a-2\mu)/6+T+4r/3-1-2\eta/3\\
 &\ge(p-2/3)\eta+(T-h)+4r/3-2>0.
\end{align*}

Here $T\ge2Y+d-1\ge h+2Y$. All original winners and residual types have now been covered. The first ledger contradicts $e(G)<E_r$; in the second, choose the actual successor with larger $c_X+E_X$, giving a positive smaller-rank strict descent.

\paragraph{Both guards and strong length.}
If $\max(s,t)\le\mu+d$, the short-head calculation of Appendix \ref{app:guarded}, before using $h\le\eta$, gives

\[
 e(G)-E_r\ge\frac{D-6}{2}a+\frac{D-4}{2}\eta
 +(D-2)d-(D-3)-h-\frac{\varepsilon_R}{2}>0.
\]

The smallest permitted values $D=8,a=0,\eta=1,h=1,d=2,\varepsilon_R=2$ already give seven, and the remaining increments are nonnegative. If the guard fails and one list is no longer than $\mu$, the single-empty proof applies. Otherwise both lists exceed $\mu$. Delete the actual roots and color $R_0$ wholly red or wholly blue. On the same $B_0$ this gives

\[
 P_{r+a}(B_0;\nn[\mu+1:s],\mm),\qquad
 P_{r+a}(B_0;\nn,\mm[\mu+1:t]).
\]

Both ranks equal $k+d>0$ and their matching deficit is $d$. Positive deficit makes the endpoint angles legal; monochromatic tests give $s,t\le k$. Since $s+t=a+\mu+k+d+r-1$, we obtain

\begin{equation}\label{eq:D10-strong}
 s,t\ge a+\mu+r-1+d.
\end{equation}

This conclusion applies only to the verified physical partition. Below it is constructed from the same degree-ten certificate and matching; an arbitrary low-cut graph is not silently identified with such a partition. With $\varepsilon_R=0$, the argument also applies to a maximum-cardinality optimal Tutte barrier inside an actual isolated regular region.

\paragraph{The width-one, exposure-four entrance.}
Suppose $\lambda=1,\rho=4$, and let $b_5$ count weight-five positive ordinary blocks. If $b_5>0$, such a block is an actual isolated regular factor-critical component and the first part of Proposition \ref{prop:regular-components} applies. Otherwise the prefix-selection inequality of Appendix \ref{app:even-small} gives

\[
 20=5\rho\le W_+\le\lambda+5(5-1)+b_5-1=20.
\]

Every inequality is an equality. The exact root ledger $W_+-W_-=5\rho+L_A/2$ forces $W_-=L_A=0$, where $L_A$ is a sum of nonnegative root deficits, internal root edges and root-to-exterior edges. Thus the original root-component region is actually isolated and ten-regular, with matching deficiency four. At least one entire connected component has no perfect matching, so the first part of Proposition \ref{prop:regular-components} again applies. Hence for $c=0,\lambda=1$ only $\rho\le3$ remains.
For $c\ge2,\rho\le4$, write $z=|S|-|T|$ in the same certificate.
Then $\Psi_3=\Psi_5+2z=\Psi_1-2z$. If $z\le0$, it is at most $c$;
if $z\ge1$, it is at most $\rho-2\le2\le c$.
The central $5/5$ cap also has width at most $c$. Thus every highest corner
not avoided by an actual defect allocation satisfies $\lambda\le c$.
Let $U$ be the union of all actual odd regular components, so
$\operatorname{def}(U)\ge c\ge\lambda$.
Choose two, and put $W=C_1\sqcup C_2$, $|C_i|=2P_i+1$,
$P=P_1+P_2$ and $J=P-\nu(W)$.
The case $c=2$ with deficiencies $3,1$ is possible; near-perfect matchings cannot be assumed.
The actual barrier above gives $s,t\ge\nu(U)+R+d\ge P+R+d-J$.
The second part of Proposition \ref{prop:regular-components} uses conservative
capacities $P_i$, the actual deficit $d_*=d-J$, and highest-column tax $\lambda$;
its full proof requires no near-perfect matching and therefore applies here.

\subsection{Exhausting the remaining certificates}
Appendix \ref{app:even-small} gives $\rho\le4$. A top pair is $(4,8)$ or $(6,6)$. Central-cap width is at most $c$; the conservative-capacity argument above excludes $c\ge2$, while the single-component result excludes $\lambda=1,c=1$. With the width-one, exposure-four boundary proved above, interpolation leaves exactly
\[
(\lambda,c,\rho)=(1,0,3),\quad(2,1,4),\quad(2,0,4),
\]
all with top pair $(4,8)$ and $H\in\{8,9\}$.
First consider $(1,0,3)$. A strict whole-graph $3/7$ coloring already avoids the original corner. Otherwise an actual factor certificate has $\Psi_3\ge1$. The bounds $\Psi_1\le3,\Psi_5\le0$ in that same certificate, together with
$\Psi_3=\Psi_5+2(|S|-|T|)$ and $\Psi_1=\Psi_3+2(|S|-|T|)$, force
\[
 |S|=|T|+1,\quad(\Psi_1,\Psi_3,\Psi_5)=(3,1,-1).
\]
Let $C$ be induced by $S,T$ and all counted components, and put $O=G-C$. Write $\zeta$ for its original degree defect, $\ell$ for its external cut, and $\delta_Q=e(Q,S\cup T)$. The exact handshake identity is
\[
 \zeta+2e(S)+2e(T)+\ell+\sum_Q(\delta_Q-2)=2.
\]
Every $\delta_Q$ is positive and even. If the internal even terms summed to two, $C$ would be isolated and regular, and $\Psi_1=3$ would force an odd-order component, contrary to $c=0$. Therefore
\[
 e(S)=e(T)=0,\quad\delta_Q=2,\quad\zeta+\ell=2,\quad e(C)=5|C|-1.
\]
Since $q+e(S,T)=10|T|+4$, we have $q\ge4$. Simplicity gives $|Q|\ge11$ for every $Q$, hence $|C|\ge45$.
Fix the same maximum matching $M$ exposing exactly three full-degree vertices. Its componentwise token inequality is
\[
1\le e_M(Q,S)+e_{G-M}(Q,T)+\#\{\text{uncovered vertices in }Q\}.
\]
In this matching and piece construction, an external $S$-port is an
actual $S$--$O$ edge and an external $T$-port is an actual $T$--$O$ edge.
These differ from edges joining roots to the counted components.
Summation gives $\Psi_1\le\rho$. Equality covers all of $S$, makes every nonfull vertex of $T$ have degree nine and remain uncovered, puts no external $S$-port in $M$, and puts every external $T$-port in $M$. Thus $e_M(S,O)=0$ and $e_{G-M}(T,O)=0$. Write $\xi=e(T,O)$, $\zeta_T=\sum_T(10-d_G)$ and $m=\nu(C)$. The restricted matching attains the same Tutte lower bound, so
\[
 |C|-2m=3+\zeta_T+\xi,\qquad \nu(C-U)=m-|U|\quad(U\subseteq S).
\]
The latter holds for the entire subset: $e(U,T)$ cancels between the root-crossing term and the degrees in $T$, and the deficiency increases by exactly $|U|$.
For $(2,1,4)$, let $J$ be the unique actual odd regular component. Its local exposure is a positive odd number at most three. Interpolation and certificate parity give $\Psi_3^J\le1$. Clearing its local roots gives exactly one and does not decrease the global resource certificate. That same certificate therefore satisfies
\[
 (\Psi_1,\Psi_3,\Psi_5)=(4,2,0),\quad |S|=|T|+1.
\]
Its sole zero-cut counted component is $J$. The remaining region $C$ has the preceding total-defect-two structure; equality in the token count simultaneously forces $J$ to be near-perfect. Fix $W=C\sqcup J$. Then
\[
 |W|-2\nu(W)=4+\zeta_T+\xi,\quad |W|\ge56,
 \quad\nu(W-U)=\nu(W)-|U|\quad(U\subseteq S).
\]
For $(2,0,4)$, the same certificate has no zero-cut component and its midpoint account is zero. It directly gives an actual isolated regular region $W$ with $\operatorname{def}(W)=4$ and $|W|\ge56$. No odd-order component of $W$ is required.
\paragraph{Termwise expansion of endpoint equality.}
Let $m_M(v)$ indicate coverage by the one maximum matching $M$. Let $Q_\Sigma$ be the union of all counted components and let $U_Q,U_T$ count original full-degree uncovered vertices there and in $T$. Summing the component parity inequalities and subtracting the same $S$--$T$ edges gives

\begin{align*}
 q+e(S,T)\le{}&m_M(S)+\sum_{v\in T}(d_G(v)-m_M(v))+U_Q\\
 &-2e_M(S)-2e_{G-M}(T)-e_M(S,O)-e_{G-M}(T,O).
\end{align*}
\[
 m_M(S)\le |S|,\qquad
 d_G(v)-m_M(v)\le9+\mathbf1_{\{d_G(v)=10,\ m_M(v)=0\}}.
\]

Hence $\Psi_1\le U_Q+U_T\le\rho$. Equality of the two endpoints in the current certificate makes every displayed nonnegative loss zero. This proves coverage of $S$, degree nine and noncoverage of every non-full $T$ vertex, exclusion of external $S$-ports from $M$, and inclusion of all external $T$-ports in $M$. In the first branch $M|_C$ exposes exactly $3+\zeta_T+\xi$ vertices. The reverse ordinary-matching cut bound is

\[
 |C|-2\nu(C)\ge q+e(S,T)-|S|-\sum_{v\in T}(d_C(v)-1)
 =3+\zeta_T+\xi.
\]

Thus $M|_C$ is maximum. In $C-U$, use roots $(S-U,T)$. Counted components are unchanged and $e(U,T)$ cancels between the crossing-edge term and the degree sum over $T$, increasing the bound by $|U|$. Therefore
$|C|-|U|-2\nu(C-U)\ge |C|-2\nu(C)+|U|$, or $\nu(C-U)\le\nu(C)-|U|$. The general deletion bound gives the reverse inequality. This proves the identity for whole subsets $U\subseteq S$, rather than inferring it from singleton equalities. The branch adjoining $J$ uses the same calculation with four instead of three; equality of the empty-root token of $J$ forces exactly one exposed vertex there.

\paragraph{A factor-critical decomposition with every port accounted for.}
Take a maximum-cardinality optimal Tutte barrier $A'$ of $W-S$. The whole-subset identity gives $\operatorname{def}(W-S)=\operatorname{def}(W)+|S|$, so $A_0=S\cup A'$ is optimal in $W$, and its residual components are factor-critical. Equality for $M|_W$ assigns $A_0$ to distinct components, with uncovered vertices in distinct unassigned components. Every $T$ port belongs to the original $M$, so its internal endpoint is uncovered by $M|_W$. Add their distinct external endpoints as roots, using those very matching ports to the distinct available components. Thus

\[
 A=A_0\cup V_{\mathrm{out}},\quad R_0=W-A_0,\quad B_0=O-V_{\mathrm{out}},
 \qquad\nu(R_0)+|A|=\nu(W)+\xi=B.
\]

All external $S$-ports meet the old roots, and every other port has its external endpoint in the new roots. Thus $E(R_0,B_0)=\varnothing$. The remaining part of the same $M$ is maximum in $B_0$, since otherwise it would augment in $G$; hence $\nu(G)=B+\nu(B_0)$. The root-component differences are $3+\zeta_T$ and $4+\zeta_T$, with $\varepsilon_R\le\zeta_T\le2$. These verify every hypothesis of the first subsection. Strong length is therefore obtained legitimately, not inferred from an arbitrary low cut.

\subsection{The actual piece, strong lengths, and port column taxes}
For the first two cases write $W$ uniformly ($W=C$ in the first), $X=E(W)\cup E(W,O)$, and $O=G-W$. The same matching using every external $T$-port gives the lower bound. Any matching of the piece using $k_S$ external $S$-ports has at most $m-k_S$ internal edges. Hence
\[
 B:=\nu(X)=m+\xi,\quad |W|=2P+\varepsilon,\quad
 B=P-j,\quad \varepsilon\in\{1,2\},\quad j\ge0.
\]
Here $P\ge22$ in the first case and $P\ge27$ in the second. The exact fee is
\[
 e(X)=10B+w=10P+E_0,\quad E_0=5\varepsilon-1+\ell,\quad w=10j+E_0.
\]
Take a maximum-cardinality optimal barrier $A'$ in $W-S$ and put $A_0=S\cup A'$. This is an optimal barrier of $W$, with factor-critical remaining components. Add the distinct exterior endpoints of the $\xi$ original matching ports to the roots. The resulting actual decomposition has no $R_0$--$B_0$ edge, component-root excess $3+\zeta_T$ or $4+\zeta_T$, total original defect at most two, and $\nu(R_0)+|A|=B$. Thus \eqref{eq:D10-strong} gives $s,t\ge B+R+d$, where $R=r-1$. Put
$d_Y=N-B-\nu(O)$. Then $d-\xi\le d_Y\le d$, with $d_Y\ge2$ in the first case and $d_Y\ge3$ in the second.
For $z\ge1$, put
\[
 M_z=\ell+\left\lfloor\frac{|W|-\ell}{z+1}\right\rfloor.
\]
This includes port-crossing stars centered outside. If there are $t\le\ell$ actual crossing stars, each uses at least one internal vertex, while every other star uses $z+1$ vertices, so $|W|\ge(A-t)(z+1)+t$. Unit stars still use the exact piece matching capacity $B$.
Every low-column tax, at heights $2,\ldots,9$, is at most $P$. For a unit axis, color the whole piece. If a nonunit attaining pair has one count larger than $M_h$, color the whole piece with that color, giving tax at most $M_2\le P-1$. Otherwise let the actual counts be $A,B'$ and heads $u,v\ge2$. Choose a vertex set containing all internal port endpoints $Z$, of size
$\max\{|Z|,|W|-(v+1)B'+1,0\}$, and color its incident edges red and the rest blue. This gives tax
\[
 \max\{|W|-vB',|Z|+B'-1\}.
\]
Reversing colors gives $\max\{|W|-uA,|Z|+A-1\}$; the original width itself is a third bound. The additional endpoint terms are at most $P$. If all three bounds exceeded $P$, then $A+B'\ge P+2$ but $uA,vB'\le P+1$, contradicting $u,v\ge2$.
At column ten, $u+v=11$; take the even head $u$. Only for the existence argument, add a virtual edge or loop to fill the internal total defect two, and decompose into five $2$-factors. Avoid the factor containing the virtual edge and take $u/2$ of the others. This gives an exact $u$-factor consisting entirely of actual edges. Color it and all ports red and the rest blue; no new blue $v$-star appears. A surviving red head loses at most $M_u$. If it does not survive, the forbidden earlier odd head, or the top-head count bound from $\lambda$, gives tax at most $M_z+\lambda-1$, where both heads are at least $z$. Use tax $\lambda$ at the top column.
Fix one highest-degree history of the same $O$, and write $b=C_{10}$ and $c_9=C_9$. The two complete endpoints have matching targets at least $d_Y+1$; the endpoint price raises these to $2d_Y$. The single matching account on this same $O$ gives
\begin{equation}\label{eq:D10-dy}
 d_Y\ge P-M_z+j+E_0-2(\lambda-1).
\end{equation}
If one low-column tax improves to $M_2$, add $P-M_2$ on the right. Put $\tau=(B-R)_+$ and $K=B+\tau$. When $\tau=0$, the whole-piece menu already pays. Otherwise take the first $\tau$ vertices of the same history and set $Y=O-T$. Vertex deletion first gives row $r+\tau=B+1$. Then use the whole-piece budget $B$: at every new corner one index is larger than $B$, and discarding at most $B$ affected components from one complete copy gives all of $\p_{r+K}(Y)$. Strong lengths and the actual endpoints ensure $0<K<N$.
\subsection{Cap reduction and stock on the same history}
Delete all internal endpoints $Z$ first, at actual fee $10|Z|-\delta$, where
$\delta=\sum_Z(10-d_G)+e(Z)\le1$. Next take a maximal independent set $I$ among current full-degree-ten vertices of $W-Z$. The full-degree Hall matching and necessity of the $S$-subsets give
\[
 i_0=|Z|+|I|\le m+\xi=B.
\]
All ports have disappeared, so the exterior is exactly the previously fixed $O$. Delete its complete degree-ten stage. The resulting cap-nine graph $X_1$ has $K_1=i_0+b<N$, actual fee $10K_1-\delta$, and excess at most $\lambda-1+\delta$. If its new top exceeds the cap, pull back an avoiding coloring directly. Otherwise Appendix \ref{app:stage-prices} gives
\[
 \begin{array}{c|c|c}
 \lambda&\rho_9(X_1)& L_{10}-b\\ \hline
 1&\le5+\delta&\le B+5+\delta\\
 2&\le6&\le B+6
 \end{array}
\]
The same $O_1$ is an actual disjoint component of $X_1$, so it has the same exposure bound. In the first case $w\ge6+\delta$: if $\xi=2$, both endpoints lie in the independent $T$, so $\delta=0,w=6$; if $\xi=1$, then $w\ge10$; if $\xi=0$, then $w\ge14$. In the second case $w\ge11$.
If $H=8$, both heads are at least three. Equation \eqref{eq:D10-dy} and $M_3\le P/2+(\varepsilon+3\ell)/4$ give $d_Y>P/2$, so $C_8\ge2d_Y>P\ge\tau$. If $H=9$ and $C_9<\tau$, then $B>2d_Y$. Each unit axis is at least $B+d_Y+1$, while a doubly dead nonunit column is at most $2M_2-R$. Their difference is at least $4\varepsilon-\ell+R-2(\lambda-1)>0$. Hence a nonunit attaining pair must survive, improving one column tax and giving $d_Y\ge2P/3-2(\lambda-1)>P/2$ (using $P\ge22,27$, respectively), a contradiction. If every low column is a unit axis, the same-history edge partition and the displayed table directly give $e(G)\ge E_r$. Thus the required $\tau$ vertices always exist. If $\tau\le b$, all $\tau$ exterior deletions have current degree ten. The disjoint actual edge sets cost
\[
10B+w+10\tau=10K+w.
\]
The complete original profile loss satisfies
\[
E_r-E_{r+K}=\sum_{h=1}^{11}\min(K,L_h)\le10K+\lambda.
\]
Since $w\ge\lambda$, this branch is paid, including $\tau=b$. If $\tau>b$ and the final deletion has degree at least nine, the exterior fee is at least $9\tau+b$, so the total is at least $9K+B+w+b$. The two table rows give
\[
B+w+b-L_{10}-\lambda\ge
\begin{cases}
w-6-\delta\ge0,&\lambda=1,\\
w-8\ge3,&\lambda=2.
\end{cases}
\]
This pays $9K+L_{10}+\lambda$, leaving only final degree eight. Each branch uses the same $K=B+\tau$ and the same actual residual graph. If $\tau>0$, then $K=2B-R$ and strong lengths give $N-K\ge2R+2d>0$; if $\tau=0$, then $K=B<N$. The fixed whole-piece menu already supplies every new pair, and the original strict one-edge margin is retained.
\subsection{Termwise payment for the final two columns}
Now $H=8$, $b\le c_9<\tau$, and the even head at column ten is at least four. The signed odd-cap reserve requires two extra exposed vertices for a positive uniform platform. Therefore
\[
 \begin{array}{c|c|c}
 \lambda&L_{10}-b&L_{10}-K_1\\ \hline
 1&\le M_3+3+\delta&\le3+\delta\\
 2&\le M_3+4&\le4
 \end{array}
\]
Continue the complete degree-nine stage of that same $X_1$, to cumulative $K_2$, at fee $9K_2+K_1-\delta$. The exact profile difference bounds excess by $3+2\delta$ and $5+\delta$, respectively. If $K_2>N$, the edge count is negative; if $K_2=N$, at most six edges remain, so there is no full-degree-eight vertex. At positive rank first handle a top above the cap by direct avoiding color, then use the cap-eight small-excess price:
\[
 \rho_8(O_2)\le6+\delta\quad(\lambda=1),\qquad
 \rho_8(O_2)\le7\quad(\lambda=2).
\]
If column nine survives, one whole-piece restoration gives $L_9-c_9\le M_2+\rho_8(O_2)$. Otherwise $L_9\le M_u+M_v-R$, where $u+v=10$.
The cost formula for the unique final graph $Y$ retains the original strict subtraction of one:
\[
 e(Y)-E_{r+K}\le L_9+L_{10}-2B-w-b-c_9+\lambda-1.
\]
Using
\[
 M_2+M_3\le\frac{7P}{6}+\frac{7\varepsilon}{12}+\frac{17\ell}{12},
\quad
 M_u+M_v+M_3\le\frac{25(2P+\varepsilon)}{36}+\frac{83\ell}{36},
\]
termwise substitution gives
\[
\renewcommand{\arraystretch}{1.6}
\begin{array}{c|c|c}
 &\text{column nine survives}&\text{column nine is dead}\\ \hline
\lambda=1&-\frac{5P}{6}+\frac{101}{12}<0&
-\frac{11P}{18}+\frac{83}{36}<0\\
\lambda=2&-\frac{5P}{6}+\frac{113}{12}<0&
-\frac{11P}{18}+\frac{119}{36}<0
\end{array}
\]
Here $P\ge22$ in the first row and $P\ge27$ in the second. This exhausts zero, one, or two ports and every surviving or dead platform.
\paragraph{Exact excesses of both cap reductions and the final constants.}
Let $L_h$ be the original complete column widths. The first $K_1$ deletions cost exactly $10K_1-\delta$, and lose at most $10K_1+\lambda$ in area. Thus

\[
 e(X_1)-E_{r+K_1}\le\lambda-1+\delta.
\]

When the last exterior deletion has degree eight, $H=8$ and the even head in the actual height-ten column is at least four. The signed cap-nine price contributes two additional exposed vertices, giving $L_{10}-K_1\le3+\delta$ for $\lambda=1$, or $\le4$ for $\lambda=2$. Continue the complete degree-nine stage on the same graph to cumulative length $K_2$, costing $9K_2+K_1-\delta$. Hence

\begin{align*}
 e(X_2)-E_{r+K_2}
 &\le(9K_2+L_{10}+\lambda)-(9K_2+K_1-\delta)-1\\
 &=L_{10}-K_1+\lambda-1+\delta\\
 &\le\begin{cases}3+2\delta,&\lambda=1,\\5+\delta,&\lambda=2.\end{cases}
\end{align*}

This uses the actual profile difference, not a substitute bound by $9N'$. If $K_2=N$ there are at most six edges, hence no full degree-eight vertex; if $K_2>N$ the edge count is negative. At positive rank, a head above the cap gives an avoiding coloring directly; otherwise the small-excess price, with the successor's own deficit sign, gives the cap-eight exposure bounds above.
There is only one final graph, $Y=O-T$. In the present branch $b\le c_9<\tau$, the exterior vertices cost at least $8\tau+b+c_9$, and the whole piece costs $10B+w$. Therefore

\begin{align*}
 e(Y)-E_{r+K}
 &\le(8K+L_9+L_{10}+\lambda)-(10B+w+8\tau+b+c_9)-1\\
 &=L_9+L_{10}-2B-w-b-c_9+\lambda-1.
\end{align*}

Substitute $B=P-j$ and $w=10j+5\varepsilon-1+\ell$ first, retaining the cost of matching slack. For a surviving ninth column and $\lambda=1$ the bound is

\[
 M_2+M_3-2P-8j+10+2\delta-5\varepsilon-\ell
 \le-\frac{5P}{6}+10-\frac{53\varepsilon}{12}
      +\frac{5\ell}{12}+2\delta-8j
 \le-\frac{5P}{6}+\frac{101}{12}.
\]

For $\lambda=2$ it is

\[
 M_2+M_3-2P-8j+13-5\varepsilon-\ell
 \le-\frac{5P}{6}+13-\frac{53\varepsilon}{12}
       +\frac{5\ell}{12}-8j
 \le-\frac{5P}{6}+\frac{113}{12}.
\]

For a non-surviving ninth column, $u,v\ge2,u+v=10$, and convexity gives $1/(u+1)+1/(v+1)\le1/3+1/9$. Adding the $1/4$ from $M_3$ gives $25/36$. The two branches are respectively

\begin{align*}
 &-\frac{11P}{18}+4+\delta-R-\frac{155\varepsilon}{36}
       +\frac{47\ell}{36}-8j
       \le-\frac{11P}{18}+\frac{83}{36},\\
 &-\frac{11P}{18}+6-R-\frac{155\varepsilon}{36}
       +\frac{47\ell}{36}-8j
       \le-\frac{11P}{18}+\frac{119}{36}.
\end{align*}

Only $\varepsilon\ge1,\ell\le2,0\le\delta\le1,j\ge0,R\ge1$ were used. No unlisted port case is hidden in a finite enumeration. Substitution of $P\ge22$ or $P\ge27$, respectively, makes all four bounds strictly negative, without truncating list lengths or unit tails.

It remains to consider the isolated regular $W$ in $(2,0,4)$. Here $|W|=2B+4=2P+2$, $P\ge27$, $B=P-1$, and $e(W)=10B+20$. Take $\ell=\delta=0$. Low-column tax is $P$, column-ten tax is $M_z+1$, and $d_Y\ge P-M_z+9$. The preceding stock proof gives the same final $Y$. The first cap-nine shadow has exposure at most six, giving $L_{10}-b\le B+6$. If $\tau\le b$, the fee $10K+20$ pays $10K+2$. If $\tau>b$ and the final degree is at least nine, the fee is at least $9K+B+20+b$, exceeding $9K+L_{10}+2$ by at least twelve. At final degree eight the cap-eight exposure is at most seven. The final excesses for a surviving and a dead ninth column are, respectively,
\[
 M_2+M_3-2P-6\le-\frac{5P}{6}-\frac{29}{6}<0,
\quad
 M_u+M_v+M_3-2P-13-R
 \le-\frac{11P}{18}+\frac{25}{18}-13-R<0.
\]
Every successor comes from the same actual deletion history, has positive smaller rank, and carries every legal corner. This completes degree ten.
\section{Factor obstructions with fixed ports and a common deleted-edge coloring}\label{app:factor}
We supply the factor interface used in the main text. Let $a,b\ge1$, $D=a+b$, and $\Delta(G)\le D$. For disjoint root sets $S,T$, let $q(S,T)$ count components $Q$ of $G-S-T$ whose vertices all have degree $D$ in $G$ and for which $a|Q|+e(Q,T)$ is odd. Put
\[
 \delta(S,T)=a|S|+b|T|-e(S,T)-q(S,T).
\]
We use the degree-deviation form of the integer interval-factor theorem \cite{lovasz,luyu}: the minimum total deviation is $-\min\delta$, with a zero-deviation decomposition when this expression is nonpositive. Choose a certificate minimizing $\delta$, then maximizing $|S|+|T|$. Write its value as $-k<0$ and set $\Lambda=a|S|+b|T|-e(S,T)$, so $q-\Lambda=k$. Every choice below belongs to this one actual certificate.

\subsection{One matching of root tokens}
Give each root in $S$ exactly $a$ clones and each root in $T$ exactly $b$ clones. Right-side tokens are actual $S$--$T$ edges and counted components. An edge token is adjacent to its endpoint clones; a component token is adjacent to clones of roots having actual ports into it. Removing any set of roots and comparing $\delta$ shows that its number of clones is at most its incident edge and component tokens. Identical clone neighborhoods give Hall's inequalities also for partial clone sets. Thus all clones can be matched.

All $S$--$T$ edge tokens can be matched simultaneously. Their Hall inequalities follow from
\[
 e(S_0,T_0)\le D\min(|S_0|,|T_0|)
 \le a|S_0|+b|T_0|.
\]
Start by matching every edge token and augment along alternating paths to cover the clones without uncovering an edge token. Exactly $\Lambda$ component tokens are then matched to roots, leaving $k$ unmatched components. This fixes the responsibility for all root edges at once; incompatible root matchings cannot be substituted in different restorations.

\subsection{Internal factors of one odd obstruction component}
On a counted component $Q$, set $f(v)=a-d(v,T)$. Since $d_G(v)=D$,
\[
 d_Q(v)-f(v)=b-d(v,S).
\]
For disjoint $X,Y\subseteq V(Q)$ write the internal factor deviation as
\[
 \delta_Q(X,Y)=f(X)+\sum_{v\in Y}(d_Q(v)-f(v))
                  -e_Q(X,Y)-q_f(X,Y).
\]
Adding $X,Y$ to the original root sets and canceling the cross terms gives
\[
 \delta(S\cup X,T\cup Y)-\delta(S,T)=\delta_Q(X,Y)+1.
\]
For $X\cup Y\ne\varnothing$, secondary extremality excludes equality on the right. Since $f(Q)$ is odd, deviation parity gives $\delta_Q(X,Y)\ge1$. Single-vertex root choices also give $1\le f(v)\le d_Q(v)-1$. Changing the demand at any prescribed vertex $v$ to $f(v)-1$ or $f(v)+1$ decreases a nonempty-cut deviation by at most one and makes the empty-cut total even. The interval-factor theorem therefore supplies each of these exact degree sequences. The exposed vertex may be any actual port endpoint.

For an uncounted component $K$, the remaining red and blue capacities are $A_v=a-d(v,T)$ and $B_v=b-d(v,S)$. Single-vertex root moves and secondary extremality give $A_v,B_v\ge1$, while $A_v+B_v\ge d_K(v)$. Every internal-cut inequality follows from the same global deviation comparison; the empty cut has no counted odd obstruction. The interval-factor theorem gives a zero-deviation internal coloring with red degree at most $A_v$ and blue degree at most $B_v$.

\subsection{An arbitrary actual missing edge}
Choose any actual edge $uv$ in an unmatched $Q$. On $Q-uv$ prescribe upper bounds $h=f$ and lower bounds $l=f-\mathbf1_{\{u,v\}}$. They are legal because $1\le f\le d_Q-1$. A nonempty-cut interval deviation is no smaller than the original $\delta_Q$: endpoints left in the remainder have slack and create no new odd obstruction; an endpoint in the second root set leaves $d-l$ unchanged; deleting an edge crossing the two root sets increases deviation by one. Other components are unchanged. The empty cut also satisfies the interval condition. Hence the interval factor exists.

Since $f(Q)$ is odd and a factor's degree sum is even, its total deficit is exactly one. Precisely one endpoint $x$ has red degree $f(x)-1$, and the other endpoint $y$ has red degree $f(y)$. Including default port colors, their red-blue degrees in $G-uv$ are $(a-1,b)$ and $(a,b-1)$; all other vertices satisfy the $(a,b)$ caps. Restoring $uv$ red makes only the red degree of $y$ equal to $a+1$; restoring it blue makes only the blue degree of $x$ equal to $b+1$.

\subsection{Simultaneous global assembly}
Color edges inside $S$ blue and inside $T$ red. An $S$--$T$ edge uses the color assigned by its matched clone, consuming one exceptional capacity of that root. Default $S$ ports are blue and $T$ ports red. If a component token is matched to an $S$ root, switch its chosen port red and use the internal $f-1$ factor at that endpoint. For a token matched to a $T$ root, switch the port blue and use $f+1$. The clone matching exactly pays these root exceptions, and the internal change cancels the port change. Use zero-deviation colorings on uncounted components.

Independently choose an actual edge $e_i$ in each unmatched $Q_i$, and put $M=\{e_1,\ldots,e_k\}$. These edges lie in disjoint actual components and form a matching. The construction gives one common coloring of $G-M$. Each of the two restorations of $e_i$ has just the stated maximum-star exception. All restorations are independent, with the ports and exterior fixed. This proves the common base coloring; it does not additionally assert that some restoration avoids the complete targets. The latter is supplied by the actual payment or whole-forest arguments in the main text.

\subsection{Expanded verification of the standard theorem, every cut and the assembly}
All preceding steps use the same ordinary interval-factor theorem. To fix its precise version, let $X$ be finite, with nonnegative integer capacities $A,B$ satisfying $A(v)+B(v)\ge d_X(v)$. For $L\subseteq E(X)$ put

\[
 \operatorname{dev}(L)=\sum_v[(d_L(v)-A(v))_++(d_{X-L}(v)-B(v))_+].
\]

Let $t(P,Q)$ count components of $X-P-Q$ on which $A+B=d_X$ pointwise and $A(C)+e(C,P)$ is odd. The precise theorem of Lovász \cite[Theorem 7.3]{lovasz} is

\begin{equation}\label{eq:factor-deviation-exact}
 \min_L\operatorname{dev}(L)=\max_{P\cap Q=\varnothing}
 [t(P,Q)+e(P,Q)-B(P)-A(Q)].
\end{equation}

Substitution of the constants $A=a,B=b,P=T,Q=S$ gives the asserted $-\min\delta$, with no minimum-degree assumption. For $0\le l\le u\le d_X$, take $A=u,B=d_X-l,P=Y,Q=Z$. A spanning factor with degrees in $[l,u]$ exists exactly when every disjoint $Z,Y$ satisfies

\begin{equation}\label{eq:factor-interval-exact}
 u(Z)+\sum_{v\in Y}(d_X(v)-l(v))-e_X(Z,Y)-q_{lu}(Z,Y)\ge0.
\end{equation}

Here $q_{lu}$ counts only components on which $l=u$ pointwise and $u(C)+e(C,Y)$ is odd. A component containing a slack vertex is not counted. This is the ordinary interval-factor condition, not a parity-factor condition.

\paragraph{The internal cut identity and parity.}
Adding $X,Y\subseteq Q$ to the global roots increases the root fee by exactly
$a|X|+b|Y|-e_G(X,T)-e_G(Y,S)-e_Q(X,Y)$. The old contribution one from $Q$ is removed. For a new component, the internal expression $f(C)+e_Q(C,Y)$ differs from $a|C|+e_G(C,T\cup Y)$ by $-2e_G(C,T)$. This proves the exact identity $\Delta\delta=\delta_Q+1$ above. Minimality gives $\delta_Q\ge-1$, and equality for nonempty roots would enlarge the global roots; thus then $\delta_Q\ge0$. Modulo two,

\[
 q_f\equiv f(Q-X-Y)+e_Q(Q-X-Y,Y),\qquad
 d_Q(Y)\equiv e_Q(Y,X)+e_Q(Y,Q-X-Y).
\]

Substitution gives $\delta_Q\equiv f(Q)\equiv1\pmod2$, so a nonempty cut is at least one. A singleton first root gives $f(v)-q_f\ge1$, and a singleton second root gives $d_Q(v)-f(v)-q_f\ge1$, proving the degree margins. For $f\pm\mathbf1_v$, only one root fee changes if $v$ is a root, and otherwise only the parity of one residual component changes. Every nonempty cut loses at most one; the empty cut is zero because $Q$ is connected and total demand is even. This checks all cuts for both specified-vertex factors.

\paragraph{Hall's proof of simultaneous saturation.}
For a root set $W$ write $\operatorname{cap}(W)=a|W\cap S|+b|W\cap T|$. Let $\operatorname{inc}(W)$ count $S$--$T$ edges incident with $W$, each once, and let $n_Q(W)$ count adjacent counted components. Removing these roots leaves every old counted component not adjacent to $W$ unchanged. Hence

\[
 0\le\delta(S-W,T-W)-\delta(S,T)
 \le-\operatorname{cap}(W)+\operatorname{inc}(W)+n_Q(W).
\]

Partial clone sets have the same neighborhoods as all clones at the roots they touch, proving every Hall inequality. For simultaneous edge-token saturation, maximize the matched edge-token count among clone-saturating matchings. Follow alternating paths from an unmatched edge token. Reaching a matched component token permits a reversal that improves the objective and keeps all clones saturated. Otherwise all reached tokens are edge tokens and all their neighboring clones are reached; the unmatched initial token makes the token count exceed its neighborhood, contradicting the edge-token Hall condition. No compatibility of two independently chosen matchings is assumed.

\paragraph{Uncounted components.}
If $A_K(v)=a-d_G(v,T)\le0$, adding $v$ to $S$ changes global $\delta$ by $A_K(v)-q_{\rm new}\le0$. A decrease contradicts minimality and equality contradicts maximal root size. Thus $A_K(v)\ge1$. Similarly $B_K(v)=b-d_G(v,S)\ge1$, and $A_K+B_K\ge d_K$. For any local roots $X,Y$, the capacity cut expression equals $\delta(S\cup X,T\cup Y)-\delta(S,T)$: capacity equality means original full degree, parity is unchanged, and the old contribution of $K$ is zero. Every cut is nonnegative, so \eqref{eq:factor-deviation-exact} gives an internal zero-deviation coloring. This includes a singleton component.

\paragraph{Every endpoint position in the missing-edge interval.}
In $Q-uv$ the upper fee is unchanged; at an endpoint in the second root set both degree and lower demand decrease by one; if the edge crosses the root sets, the negative crossing term increases by one. Every potentially changed residual component contains an endpoint and is no longer counted because $l<h$ there. When both endpoints remain, a deleted bridge leaves two components each containing an endpoint, whereas a nonbridge leaves both in one component. With one endpoint remaining, only its component can change. Components containing neither endpoint are unaffected. Hence the new $q_{lh}\le q_f$, and every nonempty cut remains at least one. For empty roots $Q-uv$ has at most two components, each containing an endpoint, and the count is zero. The factor degree sum is even and $f(Q)$ is odd; their difference lies in $\{0,1,2\}$, so it equals one. The obtained factor determines the direction.

\paragraph{Exhaustion of actual edges and full root degrees.}
Every edge lies inside $S$, inside $T$, between $S$ and $T$, at a port to a residual component, or inside a residual component. Each class is treated exactly once above. At an $S$ root the red edges are precisely the edge tokens and component ports assigned to its $a$ clones. They are distinct actual edges, so its red degree is exactly $a$ and its blue degree is $d_G-a\le b$. A $T$ root similarly has blue degree exactly $b$. A switched port in a matched component is compensated by the internal $f\mp\mathbf1_v$ factor. Unmatched components change no ports. Distinct components share neither vertices nor internal edges, so all missing edges form a matching and all restorations hold simultaneously. No external edge is left uncolored, and the root matching is unchanged between restorations.

\section{Exact prices for low-degree successors}\label{app:small-prices}
The successor $X$ has a fixed complete source, positive rank $N_X<N$, degree cap $q$, and profile height at most $q+1$. Let $h$ be its height-$(q+1)$ width, $\rho=\rho_q(X)$, and $d_X=N_X-\nu(X)$. Every price concerns this one actual $X$. If $d_X\le0$, the canonical inequality gives $e(X)-E_X\ge\rho$. If $d_X>0$, then
\begin{align}
q=5:\quad&e(X)-E_X\ge\max(0,2\rho-h-4),\label{eq:price5}\\
q=4:\quad&e(X)-E_X\ge\max(0,\floor{3\rho/2}-h-4),\label{eq:price4}\\
q=6,\ \rho\ge5:\quad&e(X)-E_X\ge\max(0,2\rho-h-4).\label{eq:price6}
\end{align}
These include $h=0$ and do not assign the original minimal object's structure to the successor.

\subsection{Caps five and four}
Only the positive bonuses require proof. Write $\mathcal H=2\rho-4$ for cap five and $\mathcal H=\floor{3\rho/2}-4$ for cap four, and assume $\mathcal H>h$. In the ordinary full-degree region of $X$, order positive blocks by capacity and take the first prefix whose weight reaches $\mathcal H$. Block weights are at most two in both cases. Thus $\mathcal H\le w\le\mathcal H+1$, with at most $\mathcal H$ selected blocks. Let $z$ be their total capacity, $Z$ their largest capacity, and $u$ the number of unselected positive blocks.

For cap five, $W_+\ge5\rho/2$ and $Z\ge3$, giving
\[
 u\ge(W_+-w)/2\ge\rho/4+3/2,\qquad u\ge3.
\]
At least $2uZ+\rho-\mathcal H$ actual full-degree vertices remain outside the selected blocks. This exceeds $6(Z-1)$: at $Z\ge3$ the difference is at least $6u-\rho-8\ge\rho/2+1>0$. Greedily choose $b=(Z-r_X+1)_+\le Z-1$ independent full-five vertices. For cap four, $W_+\ge2\rho$, $Z\ge2$, and again $u\ge\rho/4+3/2$. The exterior full-degree count exceeds $5(Z-1)$ by at least
\[
 4u+\rho-\mathcal H-5
 \ge5\rho-3\mathcal H-7\ge\rho/2+5>0.
\]
The required independent full-four vertices therefore also exist.

Erase the selected block edges and delete these exterior vertices. The actual saving is $q(z+b)+w$, whereas the complete loss is at most $q(z+b)+h$. The block matching identity and the degree-counting Hall matching of the exterior vertices give $z+b\le\nu(X)<N_X$. The capacity menu covers every new pair. The known smaller-rank inequality consequently forces $e(X)-E_X\ge w-h\ge\mathcal H-h$. When the claimed bonus is nonpositive, use the smaller-rank source inequality directly.

\subsection{Cap six and complete deletion stages}
For cap six take $\mathcal H=2\rho-4>h$. Positive block weights are at most three, capacities at least three, and $W_+\ge3\rho$. The shortest prefix has $\mathcal H\le w\le\mathcal H+2$, leaving
$u\ge\ceil{(\rho+2)/3}\ge3$ unselected blocks. There are at least $2uZ-\rho+4$ exterior vertices. If $Z=3$, this is at least fifteen, so greedy seven-coloring supplies the required at most two independent vertices. If $Z\ge4$, the exterior contains no $K_7$: its capacity-three, full-weight block would have preceded the selected prefix. Brooks theorem six-colors the remaining exterior. For $\rho=5,6,7$ use $u\ge3$ directly; for $\rho\ge8$ use the displayed bound. Both give $2uZ-\rho+4>6(Z-1)$. The same actual deletion proves \eqref{eq:price6}.

We need one stage fact on the same graph. Let $q\ge1$, let $\Delta(X)\le q$, and let $I$ be a complete independent degree-$q$ deletion stage: each step removes a current degree-$q$ vertex, until none remains. Every original full-degree vertex belongs to $I$ or is adjacent to it, since otherwise its degree never changes. Take the ordinary full-degree decomposition of $X$, with $a$ roots in $A$ and $a+\rho_q(X)$ odd cores. Every core vertex has original degree $q$, and every edge leaving a core goes to $A$.

If a core $Q$ avoids $I$, each of its vertices is dominated by $I\cap A$. For $n_Q=|Q|\ge q$, its cut has at least $n_Q\ge q$ edges. For $1\le n_Q\le q$, simplicity gives
\[
e_X(Q,A)=qn_Q-2e(Q)\ge n_Q(q-n_Q+1)\ge q,
\]
where the last inequality is equivalent to $(n_Q-1)(q-n_Q)\ge0$. The core cuts use distinct edges, whereas all root degrees sum to at most $qa$. At most $a$ cores can avoid $I$. At least $\rho_q(X)$ mutually disjoint cores therefore meet $I$ at distinct vertices, proving $|I|\ge\rho_q(X)$. This is a total cut budget, not an assumed injection of avoided cores into different roots.
In particular, every complete degree-six stage of an original cap-six object with top width $\lambda=1$ has length $k$ satisfying
\begin{equation}\label{eq:six-shadow}
 k\ge L_6-4.
\end{equation}
Otherwise its deletion leaves a positive smaller-rank cap-five source of height-six width $L_6-k>4$. Comparing the actual saving $6k$ with the complete loss bounds the new excess by zero. But the uniform source gives $\rho_5\ge L_6-k$, and either \eqref{eq:price5} for positive deficit or the matching-rich price gives strictly positive excess, a contradiction. If all original heads are at most four, signed exposure for the actual $3/4$ column strengthens \eqref{eq:six-shadow} to $k\ge L_6$.

\section{Degree six: strong lengths, complete profile taxes, and structural entry}\label{app:D6-structure}
Throughout we use the same original minimum strict source, with $D=6$ and all original heads at most five. Write $R=r-1$, $\lambda=L_7$, let $c$ count actual isolated connected odd-order six-regular components, and put $\rho=\rho_6(G)$. No bound is placed on core orders, list lengths, or unit multiplicities.

\subsection{Strong lengths for an isolated full-degree region}
\begin{lemma}\label{lem:isolated-length}
In any original minimum with $D\ge6$ and $H\le D-1$, let $W$ be a union of actual isolated $D$-regular components and put $P=\nu(W)$. If $\eta=|W|-2P\ge\lambda$, then $O=G-W$ satisfies
\[
 \max(s,t)\le\nu(O),\qquad s,t\ge P+R+d.
\]
\end{lemma}
\begin{proof}
Choose a maximum-cardinality optimal Tutte barrier $A$ in $W$. Every remaining component is odd and factor-critical: adding one vertex of an even component preserves deficiency and increases barrier size; a Tutte failure set in a non-factor-critical odd component does the same. Put $\alpha=|A|$ and $\mu=\nu(W-A)$. There are $\alpha+\eta$ odd cores, $P=\alpha+\mu$, and actual edges match the roots into distinct cores. A root--root edge or two roots entering the same core would give a matching smaller than $P$.

Apply Corollary \ref{cor:local-root-core-interface} with
$\mathcal R=W-A$, $B=O$, $a=\alpha$, and this $\eta\ge\lambda$.
The source is the original normalized global minimum, all core vertices
have original degree $D$, and the matching just constructed saturates
the roots into distinct cores. Thus every hypothesis of that corollary
has been checked on the same graph. It gives
\[
 \mu<\min(s,t),\qquad \max(s,t)\le\nu(O),\qquad
 s,t\ge\alpha+\mu+r+d-1=P+R+d.
\]
In particular, no full-degree covering property of $O$ and no minimality
property of an auxiliary graph have been used.
\end{proof}

\subsection{Excluding an actual isolated seven-clique}
Let $J$ consist of four disjoint stars of sizes $7,6,6,1$, so $e(J)=20<21$. Every coloring of $J$ has one of the following, allowing color exchange:
\[
 (5,2,1)\text{ in one color};\qquad
 (4,1),(5,2)\text{ in opposite colors};\qquad
 (4,2,1),(2,2,1)\text{ in opposite colors}.
\]
To prove this, suppose neither color has $(5,2,1)$ and one large star has red capacity $d\ge5$. Let $x,y$ be the other large stars' red capacities and $z$ the unit's red capacity. Either $x,y\le1$, or at most one of $x,y,z$ is positive. In the first case the other blue capacities are at least five; the blue forbidden pattern forces $d$ to be its original capacity and $z=1$, giving the second outcome. If only $z$ is positive the same conclusion holds. If only $x$ or $y$ is positive, the blue restriction forces two red capacities at least five, again giving the second outcome. If all are zero, blue contains the forbidden pattern. Finally, if all monochromatic capacities are at most four, the seven-star splits $4/3$ and both six-stars have capacity at least two in both colors, giving the third outcome.

Color $K_7$ accordingly: monochromatically in the first case; as a $K_5$ in the $(4,1)$ color and its complement in the $(5,2)$ color in the second; and as a Hamilton seven-cycle in the $(2,2,1)$ color and its complement in the $(4,2,1)$ color in the third. Every forest with heads at most five in either color is coordinatewise dominated by the corresponding list: seven vertices allow at most three positive stars, two stars have at most five edges, and three at most four; two stars in $K_5$ have at most three edges; the complement of the seven-cycle has maximum degree four. Thus complete forests are preserved color by color with any fixed exterior. Replacing an actual isolated $K_7$ by $J$ gives a same-rank source with one fewer edge, contradicting common minimality.

\subsection{The complete positive-block entry}
For blocks of the original ordinary region put $z_Q=\nu(Q)$ and $w_Q=e(Q)-6z_Q$. They satisfy
\[
 1\le w_Q\le3\quad(w_Q>0),\qquad
 W_+-W_-=3\rho+L_A/2,\quad L_A\ge0.
\]
Order by capacity and take the first prefix reaching weight $\lambda$, with largest capacity $Z$, total capacity $z$, weight $w$, and block count $l$. Then $\lambda\le w\le\lambda+2$ and $l\le\lambda$. If three positive blocks remain, there are at least $6Z$ exterior actual full-six vertices. After excluding $K_7$, Brooks six-coloring supplies the required $(Z-r+1)_+\le Z-1$ independent vertices. The actual saving $6K+w$ pays $6K+\lambda$; the capacity menu and Hall matching give $1\le N-K<N$. Hence at most two blocks remain, and
\[
 3\rho+L_A/2+W_-\le W_+\le\lambda+8.
\]
Balanced Euler coloring in even-regular components gives $\lambda\le c\le\rho$, so
\begin{equation}\label{eq:D6-entry}
 1\le\lambda\le c\le\rho\le4.
\end{equation}
If $\rho=4$, equality forces $\lambda=c=4$, $L_A=W_-=0$, and $W_+=12$. The last selected block and both remaining blocks have weight three, giving at least three actual factor-critical components. All four odd regular components have near-perfect matchings.

The exact even-capacity factor decomposition excludes a $3/5$ attaining pair and its reverse. Thus the highest column is attained by $4/4$, and $d\ge\rho+1$. Suppose $c=2,\rho=3$. Every weight-three block is an actual isolated six-regular factor-critical component, so at most two such blocks exist. At most two positive blocks remain after the first prefix reaching weight $\lambda$. If $\lambda=1$, the prefix consists of one block and there are at most three blocks in total, giving $W_+\le3+3+2=8<3\rho$. The same contradiction holds if $\lambda=2$ and the prefix consists of one block. Otherwise its first block has weight one, the prefix has two blocks, and there are at most four blocks in total. Hence
\[
W_+\le1+3+3+2=9.
\]
The identity $W_+-W_-=3\rho+L_A/2$, with $W_-,L_A\ge0$, forces
\[
W_+=9,\qquad W_-=0,\qquad L_A=0.
\]
No unique ordering of the remaining weights is needed. Zero root loss makes every original root full-degree, with no root-root or root-exterior edge. Thus $A\cup R$ is an actual isolated six-regular region. Put $a=|A|$ and $\mu=\nu(R)$. Its $a+\rho=a+3$ cores give
\[
|A\cup R|=a+2\mu+(a+3)=2(\mu+a)+3,
\]
which is odd. The exterior maximum matching covers every original degree-six vertex, so no actual odd six-regular component lies outside this region. The region contains exactly $c=2$ odd-order components, with any others of even order; its total order is therefore even, a contradiction. Hence $c=2$ implies $\rho=2$. Near-perfectness was not assumed in this argument.
Consequently
\[
 c=2\Longrightarrow\rho=2,\qquad
 c=3\Longrightarrow\rho=3.
\]
Each odd component has a near-perfect matching. If $c=1,\rho=3$, then $\lambda=1$ and among at most three positive blocks only one can have weight three, giving $W_+\le7<9$, a contradiction. Thus $c=1$ implies $\rho\le2$, and its positive odd matching deficiency is at most $\rho\le2$ and is also one. We can therefore work with actual whole components and omit the finer ordinary-block weight subdivisions.

\subsection{Column costs and stages on one exterior}
Choose $j$ such actual components, of capacities $z_1,\ldots,z_j\ge4$ and total capacity $P$, with $j\ge\lambda$. Lemma \ref{lem:isolated-length} supplies strong lengths. Put
\[
 M_2=\sum_{\ell=1}^j\floor{(2z_\ell+1)/3},\qquad
 M_3=\sum_{\ell=1}^j\floor{(2z_\ell+1)/4}.
\]
On the same exterior $O$, each column of heights two through five has tax at most $P+\floor{(j-1)/2}$; height six has tax at most $M_2$, or $M_3$ when there are no size-five heads; height seven has tax at most $\lambda$. These follow from the point-cover restoration in Appendix \ref{app:prices}: strong lengths preserve the complete side in a unit direction; disjoint nonunit stars give the cover tax on $2P+j$ vertices, with the original $R\ge1$ deducted. A dead $3/4$ height-six candidate is at most $M_3+\lambda-1$, and a dead $2/5$ candidate at most $M_2-1$. Since every $z_\ell\ge4$ gives $M_2-M_3\ge j$, dead candidates are also bounded by $M_2$. At the highest column, balanced $3/3$ coloring gives one degree-four exception per odd component; orient the exceptions and charge the original $\lambda$ top positions. Every live candidate is obtained from an arbitrary coloring of the same $O$, so canonical pricing may maximize each column.

Fix one maximum-current-degree deletion history on $O$. Let $C_h$ count vertices deleted at degree at least $h$, and put $b=C_6$. A source $O\to(F,qK_2)$ forces
\[
 C_h\ge q+f_h-1\quad(f_h>0).
\]
If $C_h<q$, color edges covered by these roots in the matching color and the rest in the other color to avoid the targets. Otherwise reserve the first $q-1$ roots for the matching color; the $f_h$ stars in the remaining complete $F$ must be hit by distinct later deletions. If this endpoint has matching deficit $d$, Appendix \ref{app:matching} also gives
\[
 q\ge2\bigl(d+\omega(F)\bigr),\qquad
 \omega(F)=\sum_i\floor{(n_i-1)/2}.
\]
All intermediate heads and unit tails are retained. Every later use fixes this exterior and this history.

\section{Degree six: complete descent using one actual odd regular component}\label{app:D6-one}
\begin{proposition}\label{prop:D6-one}
An original degree-six source with $\lambda=1$ and an actual isolated odd six-regular component $C$ having a near-perfect matching does not exist.
\end{proposition}
Only regularity, order $2P+1$, and matching number $P$ are used below; factor-criticality is unnecessary.
Appendix \ref{app:D6-structure} excludes $K_7$, giving $P\ge4$; regularity gives $e(C)=6P+3$.
The required strong lengths and exterior endpoints retaining the complete $F_1,H_1$ are established directly below.

\subsection{Scope of the proposition and reconstruction of its hypotheses}
Here an ``original source'' means precisely the common globally minimal datum of the main text: $P_r(G;\nn,\mm)$, $e(G)=E_r-1$, $D=6$, $H\le5$, $d=N-\nu(G)\ge\rho_6(G)+1$, and $\lambda=L_7=1$. Attainment by $4/4$ and prohibition of $5/3$ and its reverse follow from complete cap colorings of the original graph. The graph $C$ is an entire connected component of $G$, is six-regular, has order $2P+1$ and $\nu(C)=P$. The exterior $O=G-V(C)$ is arbitrary. In particular $E(C,O)=\varnothing$; a factor-critical region with ports cannot be substituted.

We assume neither $c=\rho_6(G)=1$ nor factor-criticality of $C$. All source and cost premises used below are derived from the stated near-perfect-matching hypothesis.

\paragraph{Strong length from the same actual component.}
Choose a maximum-cardinality optimal Tutte barrier $A$ inside $C$. An even residual component permits adding one vertex while preserving deficiency and enlarging the barrier; a non-factor-critical odd component permits adding its Tutte failure set, with the same contradiction. Thus every residual component is odd and factor-critical. Put $a=|A|$ and $\mu=\nu(C-A)$. Since $|C|-2\nu(C)=1$,

\[
 q-a=1,\qquad P=a+\mu,\qquad \nu(G)=a+\mu+\nu(O).
\]

One maximum matching must assign each root to a different odd component and expose just one vertex in each unassigned component; otherwise the displayed equality is unattainable. This is the actual full-degree decomposition of Appendix \ref{app:sources}. Its root-component difference is $\eta=1=\lambda$, while the original $d\ge2$; the global exposure is not reset to one. The structural budgets of Appendix \ref{app:single-below} use this decomposition's $\eta=1$, its high rectangles use $h=\lambda=1$, and $d\ge\eta+1$ holds. At $D=6$, the actual two-guard remainder is $4d-3>0$. Neither the single-empty case nor both guards can occur.
Both lists therefore exceed $\mu$. Deleting $A$ and coloring the residual components wholly red or wholly blue gives two complete sources on the same $O$, each of rank $\nu(O)+d$. Their positive matching deficit makes the endpoint angles legal and gives $s,t\le\nu(O)$. Substitution into $s+t=P+\nu(O)+d+R$ yields $s,t\ge P+R+d$. Factor-criticality is used for $C-A$, not for $C$.

\paragraph{Two endpoints retaining the complete original forests.}
The original angle $(1,R)$ is legal. For any coloring of $O$, color $C$ entirely blue. Any red $F_1$ lies wholly in $O$. Otherwise at most $P$ components of the blue $H_R$ can meet $C$, since one representative edge per component forms a matching. Thus at least $t-R+1-P$ disjoint blue edges remain in $O$. Reversing colors gives the other endpoint. Consequently

\begin{gather}
 O\to(F_1,q_FK_2),\quad O\to(H_1,q_HK_2),\qquad
 q_F=t-R+1-P,\quad q_H=s-R+1-P,\label{eq:D6-full-endpoints}\\
 q_F,q_H\ge d+1,\quad
 \nu(O)=s+q_F-1-d=t+q_H-1-d.\nonumber
\end{gather}

Both endpoints have deficit $d$ because $C$ is an actual disjoint component with $\nu(C)=P$, not by inheritance on an arbitrary successor. Centers, all smaller stars and unit tails remain in $F_1,H_1$.

\paragraph{Complete column taxes on one exterior and their total ledger.}
For one component, $M_z=\lfloor(2P+1)/(z+1)\rfloor$. At heights $2\le h\le5$, a unit-axis monochromatic restoration loses at most $P$ components. For two nonunit heads, a side whose count exceeds $M_z$ survives monochromatic restoration with tax at most $M_2\le P$. If both sides are below their capacities, the vertex-cover restoration gives

\[
 \min\{f_u+g_v-R,\ 2P+1-uf_u,\ 2P+1-vg_v\}\le P.
\]

Otherwise the first term being at least $P+1$ forces $f_u+g_v\ge P+R+1\ge P+2$, while the other two force $uf_u,vg_v\le P$, contradicting $u,v\ge2$. Height six has pairs $3/4$, $2/5$ and their reverses. A union of two-factors gives the even-head side its exact degree and creates no star of the opposite head; a surviving side loses at most the corresponding $M_z$. If the even-head side does not survive, the original $5/3$ prohibition and $\lambda=1$ bound the original width by the same tax. Thus the height-six tax is $M_2$ when a head five occurs and $M_3$ otherwise. At height seven use the original width one as the tax, requiring only a zero bound on the exterior, with no prescribed exceptional vertex inside $C$.
Let $M=M_2$ or $M_3$ accordingly. The canonical price on this one $O$, taking columnwise maxima rather than adding totals from different colors, gives

\[
 e(O)-\nu(O)\ge(E_r-N)-(4P+M+1).
\]
\[
 e(O)-\nu(O)=(E_r-1-6P-3)-(N-P-d)
              =(E_r-N)+d-5P-4.
\]

Subtracting gives $P-M\le d-3$. If one column tax improves from $P$ to $M$, the same ledger gives $2(P-M)\le d-3$. These are exactly the bounds used below, without $c=1$, $\rho_6(G)=1$, or the ability to expose a specified vertex.

\paragraph{Stages, rank and all new angles.}
The root set $I$ of a complete degree-six stage is independent and has original degree six. For $U\subseteq I$, $6|U|=e(U,N(U))\le6|N(U)|$. Hall gives a matching saturating $I$; augmentation to a maximum matching preserves covered vertices. Thus $|I|\le\nu(G)<N$. Stages in $C$ and $O$ combine because there are no edges between them. Their cost is $6|I|$; the cap-five successor uses its own exposure and deficit in Appendix \ref{app:small-prices} and need not itself be minimal.
The whole-component capacity menu is $\{0,P\}$. If $P<r$, every new angle already has a side exceeding $P$, and $6P+3$ pays the full loss, at most $6P+\lambda$. Otherwise first delete $\tau=P-R$ vertices along the same exterior history, then delete $C$. The intermediate row is $r+\tau=P+1$ and every final angle has a side exceeding $P$. Color $C$ wholly on that side and remove at most $P$ components from the one complete copy, obtaining the required entire suffix. Thus the fixed $Y$ satisfies all of $P_{r+K}$, with $K=P+\tau=2P-R$. Strong length gives $N\ge2P+R+2d>K$. Every cost bound concerns this same $Y$ and retains the minus one from $e(G)=E_r-1$.
This reconstructs every load-bearing premise used below. The proposition therefore has its stated scope for an actual near-perfect component; it does not extend to regions with ports, regions without near-perfect matchings, or arbitrary cap-five successors.

\subsection{A short full-degree stage}
An odd six-regular graph has a maximal independent set of size at most $P-1$. Degree counting gives $\alpha(C)\le P$. If equality holds, choose an independent set $J$ of size $P$ and put $B=V(C)-J$. Then $|B|=P+1$ and $e(C[B])=3$. This three-edge graph has an inclusion-minimal vertex cover $T$ with $2\le|T|\le3$: if its cover number is one, take the three leaves of its star; otherwise take a minimum cover. The set $I=B-T$ is independent and dominates $T$. Each vertex of $J$ has six neighbors in $B$, not all in $T$, so $I$ is maximal and $|I|\le P-1$.

Fix the maximum-degree history on $O$ and put $b=C_6(O)$. The set $I$ and the exterior degree-six stage are nonadjacent and form a complete independent full-six stage in the original graph. Since $\lambda=1$, \eqref{eq:six-shadow} gives
\begin{equation}\label{eq:D6-one-shadow}
 b\ge L_6-P-3.
\end{equation}
A maximal independent set of size at most $P-2$ improves this bound by one.

\subsection{An original head of size five}
Put $M=\floor{(2P+1)/3}$. The complete profile tax on the same exterior is $4P+M+1$. Comparison with its actual edge and matching counts gives
\[
 P-M\le d-3,\qquad P\le3d-8.
\]
If $P<r$, deleting $C$ by the binary capacity menu already pays. Otherwise put $\tau=P-R$ and $K=P+\tau=2P-R<N$. If $C_5\ge\tau$, take the first $\tau$ vertices of the fixed history. When $b\ge\tau$ all are degree six and payment is immediate. For $b<\tau$ the complete accounting gives
\[
 e(Y)-E_{r+K}\le L_6-P-b-3\le0.
\]
Failure of strict exit forces $b=L_6-P-3$.

If $C_5<\tau$, the size-five endpoint gives $q_F\ge2(d+\omega(F_1))\ge2d+4$ and $C_5\ge q_F+f_5-1$, hence $P>2d$. If a height-three, -four, or -five column is attained by two nonunit heads whose counts are both at most $M$, its candidate is at most $2M-R$. An actual unit axis is at least $P+d+1$, strictly larger by $P-M\le d-3$, a contradiction. Thus one count exceeds $M$, and monochromatic restoration of $C$ lowers this column's tax from $P$ to $M$. The same exterior's canonical price then forces $2(P-M)\le d-3$, or $P\le(3d-7)/2$, again contradicting $P>2d$. Hence heights one through five are attained by actual unit axes
\[
 A_h=\max\{f_h+t-R\ (f_h>0),\ g_h+s-R\ (g_h>0)\}.
\]
The endpoint stage inequalities $C_h\ge A_h-P$ imply
\[
 e(G)\ge6P+3+\sum_{h=1}^5(A_h-P)+b
       \ge\sum_{h=1}^5A_h+L_6=E_r-1.
\]
Since the original graph has exactly $E_r-1$ edges, equality holds throughout and again $b=L_6-P-3$.

Thus the short stage must have $|I|=P-1$ and $k=|I|+b=L_6-4$. The actual disjoint union
$Z=(C-I)\sqcup(O-A_6)$ has cap five, a positive smaller-rank complete source, new height-six width four, and $e(Z)=E_{r+k}$. The two deficit signs in \eqref{eq:price5} give $\rho_5(Z)\le4$. A $3/4$ attaining pair for the original sixth column would give signed exposure at least eight, a contradiction. Only $2/5$ or its reverse remains. The forbidden original $5/3$ pair, together with size-four heads on both sides, bounds the size-five count by $r-2$ and the size-two count below by $L_6+1$.

For $P\ge5$, restore $C$ entirely in the size-two color and delete the same $A_6$. The graph $O-A_6$ retains positive uniform width
\[
 L=L_6-M-b=P+3-M\ge5,
\]
forcing $\rho_5(O-A_6)\ge5$, contrary to $\rho_5(Z)\le4$ on the actual disjoint union. For $P=4$, the complement of $C$ is two-regular. If it has a cycle of length at least four, two adjacent cycle vertices form a maximal independent set of size two in $C$, improving \eqref{eq:D6-one-shadow}. Otherwise $C=K_{3,3,3}$. Label the parts $a_i,b_i,c_i$ and color the cycles
\[
 a_1b_1a_2c_1a_1,\qquad a_3b_2c_2b_3c_3a_3
\]
in the size-two color, and all remaining edges in the size-five color. The former contains at most two disjoint two-stars; the latter has maximum degree four. Thus the same $O-A_6$ retains width $L_6-2-b=P+1=5$, again a contradiction.

\subsection{Both original largest heads are four}
Put $M_2=\floor{(2P+1)/3}$ and $M_3=\floor{(2P+1)/4}=\floor{P/2}$. The complete profile tax gives $P-M_3\le d-3$. The ordinary endpoint therefore gives $q\ge2d+2\ge P+8$, so $C_4\ge q$ on the fixed history and supplies $\tau=(P-R)_+$ roots.

Every complete degree-six stage has length $k\le\nu(G)<N$, and its cap-five successor has excess at most zero. If the new sixth-column width $h=(L_6-k)_+$ is positive, the actual $3/4$ source gives signed exposure at least $h+2\ceil{h/2}$. Formula \eqref{eq:price5} yields positive excess even for $h=1$, a contradiction. If $k>L_6$, the exact profile loss is at most $5k+L_6+1$, again giving strict payment. Hence every complete stage has length exactly $L_6$, with successor exposure at most two. Fixing the exterior stage of length $b$, all maximal independent sets in $C$ have the same size $\alpha=L_6-b$. Restore $C$ in the size-three color of an actual sixth-column pair. If $L_6-M_3-b>0$, the same exterior cap-five graph has signed exposure at least three, a contradiction. Thus
\[
 \alpha\le M_3,\qquad b\ge L_6-M_3.
\]

If $C_5\ge\tau$, deleting $C$ and the first $\tau$ exterior roots gives excess at most $M_3-P-3<0$. If $C_5<\tau$, retaining the first four columns and the complete fifth and sixth columns gives
\begin{equation}\label{eq:D6-one-no5}
 e(Y)-E_{r+K}\le
 \min(K,L_5)-2P-C_5+M_3-3,\qquad K=2P-R<N.
\end{equation}
For $4\le P\le7$, use $\min(K,L_5)\le K$ to obtain $M_3-R-C_5-3<0$.

Finally let $P\ge8$. In $C$, delete a maximal independent set $I$ of size $\alpha$, then a maximal independent set $J$ among the remaining current-degree-five vertices. The set $J$ is still independent in the original $C$, so $|J|\le\alpha$ and the total stage length is at most $2M_3\le P$. Combine it with the exterior's complete degree-six and degree-five stages. The resulting cap-four union $Z_4$ has $K_0$ deleted roots, exactly $L_6$ of degree six, and actual saving $5K_0+L_6$. If $K_0\ge N$, direct edge counting gives $e(Z_4)\le5(N-K_0)\le0$. If $K_0<N$, there is a positive smaller-rank complete source with excess at most zero. Formula \eqref{eq:price4} and the top-width bound $h\le\rho_4$ give, for both deficit signs,
\[
 \rho_4(O-A_{\ge5})\le\rho_4(Z_4)\le9.
\]

Height five is attained only by $3/3$ or $2/4$. If both $3/3$ counts are at most $M_3$, then $L_5\le2M_3-R$. Otherwise restore $C$ monochromatically and the same $C_5$ roots. A cap-four graph has no positive $3/3$ source, so $L_5\le M_3+C_5$. Both make \eqref{eq:D6-one-no5} negative. For $2/4$, the size-four count is at most $R$. If the size-two count is at most $M_2$, then $L_5\le M_2$ already pays. Otherwise the same cap-four exterior retains width $(L_5-M_2-C_5)_+$, so its exposure bound gives $L_5\le M_2+C_5+9$. Therefore
\[
 e(Y)-E_{r+K}\le M_2+M_3+6-2P
 \le \frac{79-10P}{12}<0\quad(P\ge8).
\]
This covers every capacity and every original lower tail. Only the near-perfect matching, strong lengths, and the same exterior were used, never a prescribed exposed vertex. Thus the proof has exactly the stated scope.

\section{Degree six: two actual components and degree-five compensation}\label{app:D6-two}
Let $\lambda\le2$ and choose two actual isolated odd six-regular near-perfect components of capacities $4\le a\le b$, with $P=a+b$. Strong lengths give $s,t\ge P+R+d$. Write $O$ for the exterior and put
\[
 M_2=\floor{(2a+1)/3}+\floor{(2b+1)/3},\quad
 M_3=\floor{(2a+1)/4}+\floor{(2b+1)/4}.
\]
The complete same-graph tax $4P+M_2+\lambda$ and $e(C_1\sqcup C_2)=6P+6$ give
\begin{equation}\label{eq:D6-two-bound}
 P-M_2\le d+\lambda-7,\qquad P\le3d+3\lambda-19.
\end{equation}
If $b>3a$, Brooks six-coloring of the larger component supplies $(a-r+1)_+$ independent full-six vertices. Deleting the smaller component and these vertices saves $6K+3>6K+\lambda$, with $K\le2a-1<P\le\nu(G)$, giving immediate exit. Hence only
\[
 g=\max(a,b-a)\le P/2,\qquad
 \tau=(g-R)_+\le\floor{P/2}-1
\]
remains. The subset capacities $0,a,b,P$ have largest gap $g$. After deleting $\tau$ exterior vertices, deleting both components therefore gives the complete original-list source $P_{r+P+\tau}$.

Fix the maximum-degree history on $O$ and write $b_6=C_6$, $c_5=C_5-C_6$. If an original size-five head occurs, the complete ordinary endpoint gives $q\ge2d+4>\tau$, hence $C_5\ge q$. Otherwise both largest original heads are four, and an actual $4/4$ pair gives $q\ge2d+2>\tau$ and $C_4\ge q$. These comparisons use \eqref{eq:D6-two-bound}. In either case $\nu(O)\ge q>\tau$, so finally
\[
 K=P+\tau<P+\nu(O)=\nu(G)<N.
\]

Each odd regular component has a short maximal independent set. Together with the same exterior full-six stage, these give an original full-six independent set of size $2\le k_0\le P-2+b_6$. Hall matching covers it, so $k_0<N$. Its deletion leaves a cap-five source of excess at most $\lambda-1$, with sixth-column width $h=(L_6-k_0)_+\le\rho_5$. The positive-deficit price \eqref{eq:price5} gives $h\le\lambda+3$; the matching-rich price gives the stronger $h\le\lambda-1$. Thus the fixed history satisfies
\begin{equation}\label{eq:D6-two-shadow}
 b_6\ge L_6-P-\lambda-1.
\end{equation}

Without size-five heads, we must improve $C_4\ge\tau$ to $C_5\ge\tau$. If $b_6\ge\tau$, this is immediate. Otherwise first take $\tau$ degree-at-least-four vertices from the fixed history. Their exterior saving is at least
$4\tau+2b_6+\min(c_5,\tau-b_6)$. Failure of strict payment requires
\[
 L_6\ge P-\tau+2b_6+\min(c_5,\tau-b_6)+7-\lambda.
\]
Here $M_3\le\floor{(P+1)/2}$ and $M_3+\tau\le P$. If the size-three side of an actual sixth-column $3/4$ candidate dies, its width is at most $M_3+\lambda-1$, contradicting this bound. Hence the same $O-A_6$ has a positive $3/4$ source of width $h_*=L_6-M_3-b_6$. Deleting the same $c_5$ vertices reduces its cap to four. The degree-decomposition definition directly gives
$c_5\ge\kappa_5(O-A_6)\ge h_*$; no reversal of $\kappa\le\rho$ is used. If $c_5<\tau-b_6$, the preceding inequality forces
$c_5\ge b_6+c_5+7-\lambda$, a contradiction. Thus $C_5\ge\tau$ on the same history.

Delete both components and the first $\tau$ exterior roots. If $b_6\ge\tau$, all-degree-six saving pays immediately. Otherwise the actual saving is at least $6P+6+5\tau+b_6$, while the complete loss is at most $5K+L_6+\lambda$. Equation \eqref{eq:D6-two-shadow} gives
\[
 e(Y)-E_{r+K}
 \le L_6-P-b_6+\lambda-7\le2\lambda-6\le-2.
\]
This is an actual positive smaller-rank descent for every new pair and the full original profile. Only the two actual near-perfect components and strong lengths were used. Any further odd regular components in the exterior contribute no extra claimed income.

\section{Degree six: termination with three and four actual components}\label{app:D6-three-four}
By \eqref{eq:D6-entry}, Appendix \ref{app:D6-one} excludes $\lambda=1$, and Appendix \ref{app:D6-two} applies to two actual near-perfect components when $\lambda=2$. Only $(c,\lambda,\rho)=(3,3,3)$ and $(4,4,4)$ remain. Let $W$ be the union of all odd six-regular components, of total capacity $P$, and let $O$ be its exterior. Each has a near-perfect matching, so $e(W)=6P+3c$. Lemma \ref{lem:isolated-length} gives $s,t\ge P+R+d$. We first establish the capacity menu and adequate exterior roots, then pay on one final $Y$.

\subsection{Three components}
Order the capacities as $a\le b\le e$. The largest subset-sum gap is $g=\max(a,b-a,e-a-b)$. If $g>a$, the other two capacities satisfy $b+e>3a$. Brooks six-coloring on them supplies at least $a$ independent full-six vertices. Deleting the smallest component and $(a-R)_+$ vertices pays $6K+3$. Hence failure of exit implies
\[
 g=a\le P/3.
\]
On the same exterior, heights two through five each have tax at most $P+1$, height six at most $M_2$, and the top column at most three. The total tax $4P+M_2+7$ and actual edge-matching accounting give
\[
 P-M_2\le d-3,\quad M_2\le(2P+3)/3,\quad P\le3d-6.
\]
Consequently $g\le d-2$.

\subsection{Weakly independent additions for four components}
Order the capacities as $a\le b\le c_0\le e$, initially selecting the two smallest components. A remaining component of capacity $z\ge4$ has an independent set of size $\floor{z/3}+1$. The count $2z-4\floor{z/3}-3\ge1$ permits one further vertex with at most one induced edge. Thus the other two components supply
\[
 Q=\floor{c_0/3}+\floor{e/3}+4
\]
original full-six vertices inducing at most two edges altogether. If $(\max(a,b-a)-R)_+\le Q$, choose the required prefix, delete the two smallest components and these vertices, and save at least $6K+6-2=6K+4$. This pays the complete original loss. The vertex count is at most the other capacities, so $K\le\nu(G)<N$.

Otherwise put $g=\max(a,b-a)$. The integer failure inequality gives
\[
 c_0+e\le3g-14.
\]
If $g=a$, then $a\ge14$ and $P\ge56$. If $g=b-a$, then $b\ge3a+14$. These bounds also control the other subset gaps $c_0-a-b$ and $e-a-b-c_0$. Hence the full four-component menu has largest gap exactly $g$, with $g\le P/3$.

\subsection{The exact tax of three factor-critical components}
The four-component entry guarantees at least three actual isolated factor-critical components from the original ordinary region. Choose any three, of total capacity $P_3\ge12$ and order $2P_3+3$. They belong to the original $\mu$, so $s,t>P_3$; their exterior still contains the fourth component. Define $M_2,M_3$ as before. The total tax at heights two through five is exactly bounded by $4P_3$. Only a sparse $2/2$ height-three candidate could cost one more. For $r\ge3$, the cover bound is already at most $P_3$. For $r=2$, an extra one requires even $P_3$ and both counts $(P_3+2)/2$. Choose an actual vertex $v$ and a set $T$ of $P_3/2$ other vertices containing its six neighbors. Color edges meeting $T$ red and the rest blue. Red has at most $P_3/2$ disjoint two-stars. In blue, $v$ is isolated and positive support has at most $3(P_3+2)/2-1$ vertices, so the blue demand also fails internally. The same exterior retains a $2/2$ source, refunding the last unit. The nonunit cover bounds at heights four and five are
\[
 \floor{(5(2P_3+3)-6R)/11}\le P_3,\quad
 \floor{(3(2P_3+3)-4R)/7}\le P_3,\quad
 \floor{(2(2P_3+3)-3R)/5}\le P_3.
\]
Height six costs at most $M_2$, since $M_2-M_3\ge3$. The original top width is four. Orient the three Euler degree-four exceptions red or blue, with $c_R+c_B=3$ and $c_R\le f_4-1,c_B\le g_4-1$. The actual pair $(f_4-c_R,g_4-c_B)$ is legal and sums to $r$, retaining $S_4/S_4$ on the same exterior. Thus the top tax is three.

The total tax is at most $4P_3+M_2+3\le14P_3/3+4$. These three components contribute exactly three to the original exposure, leaving exposure one outside. Counting this reserve once and comparing the actual $e-\nu$ accounts gives
\begin{equation}\label{eq:D6-triple-bound}
 P_3\le3d-21.
\end{equation}
If the four-component menu gap is $g=a$, any three capacities sum to at least $3g$. If $g=b-a$, any three sum to at least $a+2b=3a+2g$. Equation \eqref{eq:D6-triple-bound} gives $g<2d+2$ in both cases.

\subsection{Final payment on the same exterior}
For both three and four components we now have
\[
 g\le P/3,\qquad g<2d+2.
\]
Use an actual original $4/4$ pair and restore all of $W$ in the matching color. The same $O$ has a complete ordinary endpoint with matching demand $q\ge2d+2$. Thus $C_4(O)\ge q>g$ and $\nu(O)\ge q$. Take $\tau=(g-R)_+$ and $K=P+\tau<\nu(G)<N$. Deleting the first $\tau$ exterior roots and all of $W$ preserves every new pair by the subset-capacity menu. If $\tau=0$ or $C_6\ge\tau$, payment is immediate.

Otherwise put $b=C_6<\tau$, $c_5=C_5-C_6$, and $M=M_2$. For both three and four components,
\[
 M+\tau\le P.
\]
If the first $\tau$ degree-at-least-four roots have not already paid strictly, actual edge counting requires
\begin{equation}\label{eq:D6-many-failure}
 L_6\ge P-\tau+2b+\min(c_5,\tau-b)+(3c+1-\lambda).
\end{equation}
A dead size-three side in a sixth-column $3/4$ candidate has width at most $M_3+\lambda-1\le M$; a dead $2/5$ candidate has width at most $M-1$. Equation \eqref{eq:D6-many-failure} excludes both. Thus the same $X_5=O-A_6$ has a positive uniform source of width $h_*=L_6-M-b$, so $c_5\ge\kappa_5(X_5)\ge h_*$. If $c_5<\tau-b$, \eqref{eq:D6-many-failure} would give $c_5\ge b+c_5+3c+1-\lambda$, a contradiction. Hence $C_5\ge\tau$ on the fixed history.

The final saving is at least $6P+3c+5\tau+b$, and the complete successor excess satisfies
\[
 e(Y)-E_{r+K}\le L_6-P-b-(3c+1-\lambda).
\]
A negative right side finishes. Otherwise put
\[
 L=L_6-M-\tau,\qquad h=(L_6-P-\tau)_+.
\]
The same $Y$ has the positive uniform source of width $L$, with
$\rho_5(Y)\ge L\ge b+3c+1-\lambda$ and $h\le L$. If matching-rich, its price is at least $L$; if matching-poor, \eqref{eq:price5} gives at least $2L-h-4\ge L-4$. Even the latter exceeds the upper bound by
\[
 P+b+3c-3-\lambda-M-\tau
 \ge b+3c-3-\lambda>0.
\]
Both deficit signs are impossible, excluding the complete three- and four-component domains.

Together with Appendices \ref{app:D6-one}, \ref{app:D6-two}, and entry \eqref{eq:D6-entry}, this exhausts every original strict source at degree six. Original lists, rows, arbitrary lower-head multiplicities, and unit tails were never replaced by finite templates.


\begin{thebibliography}{99}
\bibitem{erdos} T. F. Bloom, Erdős Problems, Problem 561, \url{https://www.erdosproblems.com/561}.
\bibitem{fln} P. Fu, Z. Luo and Z. Ni, Size Ramsey minimal graphs for uniform star forests, arXiv:2606.04439v3 (2026), \url{https://arxiv.org/abs/2606.04439v3}.
\bibitem{you} C. You, Size Ramsey numbers of star forests with two components on one side, arXiv:2610.25958v1 (2026), \url{https://arxiv.org/abs/2610.25958v1}.
\bibitem{lovasz} L. Lovász, Subgraphs with prescribed valencies, Journal of Combinatorial Theory \textbf{8}(4) (1970), 391--416, \url{https://doi.org/10.1016/S0021-9800(70)80033-3}.
\bibitem{luyu} H. Lu and Q. Yu, Constructive proof of deficiency theorem of $(g,f)$-factor, Science China Mathematics \textbf{53}(6) (2010), 1657--1662, \url{https://doi.org/10.1007/s11425-010-3079-7}.
\bibitem{zhao} X. Zhao, H. Zheng, F. Dong, H. Li and Y. Ma, Characterization of equimatchable even-regular graphs, arXiv:2408.15552v3 (2026), \url{https://arxiv.org/abs/2408.15552v3}.
\bibitem{reed} B. Reed, Paths, stars and the number three, Combinatorics, Probability and Computing \textbf{5}(3) (1996), 277--295, \url{https://doi.org/10.1017/S0963548300002042}.
\bibitem{brooks} R. L. Brooks, On colouring the nodes of a network, Proceedings of the Cambridge Philosophical Society \textbf{37}(2) (1941), 194--197, \url{https://doi.org/10.1017/S030500410002168X}.
\end{thebibliography}


\end{document}
